% concatenated from Twisted_de_Rham_complexes.tex
\documentclass[12pt, a4paper]{article}
\usepackage[a4paper,top=2cm,bottom=2cm,left=2.5cm,right=2.5cm,marginparwidth=1.75cm]{geometry}
\usepackage{amsmath, amsthm, amssymb, amscd}
\usepackage{bm, mathrsfs, dsfont}
\usepackage{cases}
\usepackage{bbding}
\usepackage{hyperref}
\usepackage[all,pdf]{xy}
\usepackage{upgreek}
\usepackage{xcolor}
\usepackage{graphicx}
\usepackage{tikz}
\newtheorem{definition}{Definition}[section]
\newtheorem{question}{Question}[section]
\newtheorem{conjecture}{Conjecture}[section]
\newtheorem{lemma}{Lemma}[section]
\newtheorem{proposition}{Proposition}[section]
\newtheorem{theorem}{Theorem}[section]
\newtheorem{corollary}{Corollary}[section]
\newtheorem{remark}{Remark}[section]
\newtheorem{example}{Example}[section]
\newtheorem{construction}{Construction}[section]
\newtheorem{slogan}{Slogan}[section]
\title{De Rham-Higgs comparison for mixed Hodge modules in positive characteristic }
\author{Zebao Zhang}


\begin{document}
%%%%%%%%%%%%%%%%%%%% Text italic %%%%%%%%%%%%%%%%%%%%%%%%%%%%
% Skriptbuchstaben
\theoremstyle{definition}
\newcommand{\Gr}{\mathrm{Gr}}
\newcommand{\dR}{\mathrm{dR}}
\newcommand{\Hig}{\mathrm{Hig}}
\newcommand{\Fil}{\mathrm{Fil}}






\maketitle
\begin{abstract}
Using the nonabelian Hodge theory in positive characteristic established by Ogus, Vologodsky and Schepler \cite{OV,Schepler}, we propose a generalization of the decomposition theorem of Deligne and Illusie \cite{DI} from the perspective of mixed Hodge modules. We confirm this generalization in certain special cases by extending the method in Sheng and the author \cite{SZ2}, thereby obtaining some geometric byproducts.
\end{abstract}
\section{Introduction}
Let $k$ be a perfect field of characteristic $p>0$, and let $X$ be a smooth variety over $k$ endowed with a simple normal crossing divisor $D$. Denote by $X',D'$ the Frobenius twist of $X,D$, respectively. Let $F:X\to X'$ be the relative Frobenius morphism. Suppose that $(X,D)/k$ is $W_2(k)$-liftable. In their famous paper \cite{DI}, Deligne and Illusie proved a remarkable theorem asserting that the truncated de Rham Frobenius pushforward $\tau_{\leq p}F_*\Omega^*_{X/k}(\log D)$ is decomposable. This theorem has inspired numerous non-trivial extensions in subsequent work. In what follows, we focus on some of these developments. The first is the positive characteristic analogue of a result of Kontsevich and Barannikov on twisted de Rham complexes. The second is an observation due to Drinfeld, which implies that the two-sided truncation $\tau_{[a,a+p-1]}F_*\Omega^*_{X/k}$ is decomposable. The third is the recent work of Petrov \cite{P2}, which shows that  the full complex $F_*\Omega^*_{X/k}$ is decomposable for quasi-$F$-split varieties $X/k$, even when $\dim X>p$. Neverthless, as shown in \cite{P1}, the truncation $\tau_{\leq p}$ is generally necessary. The fourth is a broad generalization to twisted coefficients, grounded in the non-abelian Hodge theory in positive characteristic developed by Ogus-Vologodsky \cite{OV} and Schepler \cite{Schepler}. The final is the work of Sheng and the author \cite{SZ2}, which incorporates the theory of mixed Hodge modules into the fourth extension. 



Throughout this paper, denote by $(E,\theta)$ a nilpotent Higgs bundle on $(X',D')/k$ of level $\leq\ell$ $(\ell<p)$ and by $(H,\nabla):=C^{-1}_{(\tilde X,\tilde D)/W_2(k)}(E,\theta)$ its inverse Cartier transform. Motivated by the above achievements, we propose the following.

\begin{question}\label{Question 1}
Let $K^*_\mathrm{dR}\subset\Omega^*(H,\nabla)$ and $K^*_\mathrm{Hig}\subset\Omega^*(E,\theta)$ be subcomplexes inspired by the theory of mixed Hodge modules, and let $f\in\Gamma(X,\mathcal O_X)$ be a suitable global function on $X$. Do we have a de Rhm-Higgs comparison of two-sided truncations
\begin{eqnarray}\label{de-H com for MHM}
\tau_{[a,a+p-\ell-1]}F_*(K^\bullet_\mathrm{dR},\nabla+df\wedge)\cong\tau_{[a,a+p-\ell-1]}(K^\bullet_\mathrm{Hig},\theta-df'\wedge)~\mathrm{in}~D(X')?
\end{eqnarray}
Furthermore, under what geometric conditions on the triple \( (X, D, f) \) does there exist a non-truncated de Rham-Higgs comparison
\begin{eqnarray}\label{de-H com for MHM non-truncated}
F_*(K^\bullet_\mathrm{dR},\nabla+df\wedge)\cong(K^\bullet_\mathrm{Hig},\theta-df'\wedge)~\mathrm{in}~D(X')?
\end{eqnarray}
\end{question}
By default, the trivial examples of $K_{\mathrm{dR}}^*$ and $K^*_\mathrm{Hig}$ are the full complexes $\Omega^*(H,\nabla)$ and $\Omega^*(E,\theta)$, respectively. The non-trivial examples considered in this paper fall into one of the following three types: the weight filtration, the intersection subcomplexes, and the Kontsevich subcomplexes. We answer this question in several special cases below.































\iffalse
Let $k$ be a perfect field of characteristic $p>0$, let $X$ be a smooth variety over $k$ endowed with a simple normal crossing divisor $D$, and let $f\in\Gamma(X,\mathcal O_X)$. Assume that $(X,D,f)/k$ has a $W_2(k)$-lifting $(\tilde X,\tilde D,\tilde f)/W_2(k)$. Denote by $F$ the Frobenius endomorphism of $X$. Let $(E,\theta)$ be a nilpotent Higgs bundle on $(X,D)/k$ of level $\leq\ell$ and let $(H,\nabla):=C^{-1}_{(\tilde X,\tilde D)/W_2(k)}(E,\theta)$. Then we propose the following
\begin{slogan}\label{slogan}
For subcomplexes $K^*_\mathrm{dR}\subset\Omega^*(H,\nabla)$ and $K^*_\mathrm{Hig}\subset\Omega^*(E,\theta)$ induced by the de Rham complex of some kind of mixed Hodge modules, we have a de Rham-Higgs comparison 
\begin{eqnarray}\label{de-H com for MHM}
\tau_{[a,a+p-\ell-1]}F_*(K^\bullet_\mathrm{dR},\nabla+df\wedge)\cong\tau_{[a,a+p-\ell-1]}(K^\bullet_\mathrm{Hig},\theta-df\wedge)~\mathrm{in}~D(X).
\end{eqnarray}
\end{slogan}
Its motivation comes from two parts. The first part is to find an alternative mod $p$ argument for a theorem of Sabbah below.
\begin{theorem}[Sabbah \cite{Sah1}]\label{Sabbah's theorem}
Let $X_\mathbb{C}$ be a smooth projective algebraic variety over $\mathbb C$, and let $f:X_\mathbb{C}\to\mathbb A^1_\mathbb{C}$ be a proper morphism. Then for any mixed Hodge module $(\mathcal M,F_\bullet)$ on $X$, the twisted de Rham complex 
$$
\mathcal M\xrightarrow{\nabla+df\wedge}\mathcal M\otimes\Omega^1_{X_\mathbb{C}/\mathbb C}\xrightarrow{\nabla+df\wedge}\cdots\xrightarrow{\nabla+df\wedge}\mathcal M\otimes\Omega^n_{X_\mathbb{C}/\mathbb C}
$$
and the twisted Higgs complex 
$$
gr^F\mathcal M\xrightarrow[]{gr^F\nabla-df\wedge}gr^F\mathcal M\otimes\Omega^1_{X_\mathbb{C}/\mathbb C}\xrightarrow[]{gr^F\nabla-df\wedge}\cdots\xrightarrow[]{gr^F\nabla-df\wedge}gr^F\mathcal M\otimes\Omega^n_{X_\mathbb{C}/\mathbb C}
$$
have the same finite-dimensional hypercohomology in every degree.
\end{theorem}
 The idea is that by spreading-out technique, we reduce the above theorem to a suitable de Rham-Higgs comparison as in \eqref{de-H com for MHM}. More precisely, for some geometric mixed Hodge module $(\mathcal M,F_\bullet)$, we may expect that there is a geometric logarithmic de Rham bundle $(H_\mathbb{C},\nabla_\mathbb{C},F_\bullet)$ such that $(\mathrm{DR}_X(\mathcal M),F_\bullet)$ can be described by a filtered subcomplex $(K^*_{\mathrm{dR},\mathbb C},F_\bullet)\subset(\Omega^*(H_\mathbb{C},\nabla_\mathbb{C}),F_\bullet)$. Set $(E_\mathbb{C},\theta_\mathbb{C}):=\mathrm{Gr}^F(H_\mathbb{C},\nabla_\mathbb{C})$ and $K^*_\mathrm{Hig}:=\mathrm{Gr}^FK^*_\mathrm{dR}$. Then Sabbah's theorem above says that 
 $$
 \dim_\mathbb C\mathbb H^i(X,(K^\bullet_{\mathrm{dR},\mathbb C},\nabla_\mathbb{C}+df\wedge))=\dim_\mathbb C\mathbb H^i(X,(K^\bullet_{\mathrm{Hig},\mathbb C},\theta_\mathbb{C}-df\wedge))<\infty,~\forall i.
 $$
By spreading-out technique, this equality can be derived by the de Rham-Higgs comparison
\begin{eqnarray}\label{mod p de Rham-Higgs}
F_*(K^\bullet_{\mathrm{dR},p},\nabla_p+df_p\wedge)\cong(K^\bullet_{\mathrm{Hig},p},\theta_p-df_p\wedge)~\mathrm{in}~D(X_p),
\end{eqnarray}
 where the subscript $p$ means sufficiently good mod $p$ reduction. One more words, $(\mathcal M,F_\bullet)$ and $(H_\mathrm{C},\nabla_\mathrm{C},F_\bullet)$ come from geometry ensure that 
 $$
 K^*_{\mathrm{dR},p}\subset\Omega^*(H_p,\nabla_p),~K^*_{\mathrm{Hig},p}\subset\Omega^*(E_p,\theta_p)
 $$
and the inverse Cartier transform of $(H_p,\nabla_p)$ is $(E_p,\theta_p)$.
 
 In what follows, We list some known examples below.
 \begin{example}\label{mod p proof Sabbah}
\upshape
(1) When $\mathcal M=\mathcal O_X$, $F_\bullet$ is the trivial Hodge filtration, and $f$ is constant, Theorem \ref{Sabbah's theorem} specializes to the $E_1$-degeneration of the Hodge to de Rham spectral sequence of $X_\mathbb{C}/\mathbb C$, which can be deduced from the decomposition theorem established by Deligne-Illusie \cite[Th\'eorème 2.1]{DI}.

(2) When $\mathcal M=\mathcal O_X$, $F_\bullet$ is the trivial Hodge filtration, and $f$ is dominant, Theorem \ref{Sabbah's theorem} is due to Kontsevich and Barannikov, which can be deduced from the de Rham-Higgs comparison established by Ogus-Vologodsky \cite[Theorem 4.23, 2]{OV}. If $f$ satisfies some technical conditions, then it also follows from the formality theorem of Arinkin-C\u{a}ld\u{a}raru-Hablicsek \cite[Theorem 4.16]{ACH}.

(3) Suppose that $f$ is constant and let $g:X_\mathbb{C}\to\mathbb P^1_\mathbb{C}$ be another morphism such that $D_\mathbb{C}:=g^*(0)$ is a simple normal crossing divisor on $X_\mathbb{C}$. Set $U_\mathbb{C}:=X_\mathbb{C}-D_\mathbb{C}$ and $j:X_\mathbb{C}-D_\mathbb{C}\hookrightarrow X_\mathbb{C}$ the natural inclusion. Let $\mathbb Q_{h,U_\mathbb{C}}[\dim X_\mathbb{C}]$ be the pure Hodge module on $U$ associated to the constant system $\mathbb Q_{U_\mathbb{C}}$. When $(\mathcal M,F_\bullet)$ is the Beilinson's maximal extension $\Xi_f(j_*\mathbb Q_{h,U_\mathbb{C}}[\dim X_\mathbb{C}])$, Theorem \ref{Sabbah's theorem} is a special case of Kontsevich's double degeneration property ( \cite[Theorem 2.18]{KKP}, \cite[Theorem 1.32]{ESY}), which can be deduced from the decomposition theorem after Kontsevich (see \cite[Theorem D.2.2]{ESY}).

(4) Suppose that $f$ is constant and let $\mathbb V$ be a polarized variations of Hodge structure on $X_\mathbb{C}$ of geometric origin. When $(\mathcal M,F_\bullet)=(\mathcal O_{X_\mathbb{C}}\otimes_{\mathbb C_{X_\mathbb{C}}}\mathbb V,F_\bullet)$, Theorem \ref{Sabbah's theorem} specializes to the formality lemma (see \cite[Lemma 2.2]{Sim}), which can be deduced from the de Rham-Higgs comparison of \cite[Corollary 2.27]{OV}. 

(5) Suppose that $f$ is constant. Let $D_\mathbb{C}\subset X_\mathbb{C}$ be a simple normal crossing divisor, let $j:X_\mathbb{C}-D_\mathbb{C}\hookrightarrow X_\mathbb{C}$ be the natural inclusion, and let $\mathbb V$ be a polarized variations of Hodge structure on $X_\mathbb{C}-D_\mathbb{C}$ of geometric origin. When $(\mathcal M,F_\bullet):=j_*\mathbb V_h[\dim X_\mathbb{C}]$ or $j_{!*}\mathbb V_h[\dim X_\mathbb{C}]$, Theorem \ref{Sabbah's theorem} specializes to \cite[Proposition 8.3.34]{CEGT}, which can be deduced from the de Rham-Higgs comparison of Schepler \cite[Corollary 5.7]{Schepler} or \cite[Corollary 3.6]{SZ2}, respectively.
\end{example}


The second part addresses a refinement of the aforementioned decomposition theorem of \cite[Th\'eorème 2.1]{DI}. Through a detailed study of abstract Koszul complexes, Achinger-Suh \cite[Theorem 1.1]{AS} showed that the two-sided truncation $\tau_{[a,b]}\Omega^*_{X/k}$ is decomposable if $a\leq b<a+p-1$ for $p>2$ and $a\leq b\leq a+1$ for $p=2$. Using the theory of de Rham stack, Drinfeld observed that the mod $p$ de Rham
cohomology of varieties liftable to $W_2$ has the structure of a $\mu_p$-representation. It follows that Achinger-Suh's result can be improved for $p>2$, namely the two-sided truncation $\tau_{[a,b]}\Omega^*_{X/k}$ is decomposable if $a\leq b\leq a+p-1$ (see also \cite{BL, LM, O24}).

Using the nonabelian Hodge theory in positive characteristic established by \cite{OV, Schepler}, it is natural to expect that similar result holds for a general nilpotent Higgs bundle $(E,\theta)$ on $(X,D)/k$ of level $\leq\ell$, namely
\begin{eqnarray}\label{two-sided truncation E}
\tau_{[a,a+p-\ell-1]}F_*\Omega^*(H,\nabla)\cong\tau_{[a,a+p-\ell-1]}\Omega^*(E,\theta)~\mathrm{in}~D(X).
\end{eqnarray}
As shown in \cite[\S3]{SZ2}, if the isomorphism above is constructed following the spirit of \cite{DI}, then it may restrict to an isomorphism on subcomplexes $K^*_\mathrm{dR}$ and $K^*_\mathrm{Hig}$, which are induced by the de Rham complex of some kind of mixed Hodge modules.
Based on \eqref{mod p de Rham-Higgs}, We may further conjecture that the isomorphism \eqref{two-sided truncation E} induces an isomorphism \eqref{de-H com for MHM}.

To the best of my knowledge, for $\theta\neq0$, neither the two-sided truncation $\tau_{[a,a+p-\ell-1]}$ nor the twisting of $\nabla$ (resp. $\theta$) by $df$ (resp. $-df$) has been explored in the literature. This paper will proceed along this direction.\fi







\iffalse
In recent years, the Landau-Ginzburg model in characteristic zero had been extensively studied from various viewpoints, such as mirror symmetry, birational geometry and Hodge theory. Nevertheless, there are little literature about Landau-Ginzburg models in positive characteristic. In this paper, we shall propose a generalization of the decomposition theorem of Deligne and Illusie \cite{DI} to Landau-Ginzburg models in positive characteristic, and then provide it with several evidences. This generalization, coupled with the work of Dodd and Abe-Caro, may shed lights on the existence of the positive characteristic analogue of the theory of mixed Hodge modules.\fi
\subsection{Partial answers}
 This paper builds on two main findings: the first is an analogue of the $\alpha$-transform (see Ogus \cite{O24}); the second is a connection between splittings of the logarithmic cotangent bundle $\Omega^1_{X/k}(\log D)$ and isomorphisms in $\mathrm{Hom}_{D(X')}(F_*K^*_\mathrm{dR},K^*_\mathrm{Hig})$.

\subsubsection{Answers provided by the first finding}
 Let $x\in X$ be a closed point, and let $(\hat H,\hat \nabla),(\hat E,\hat\theta)$ be completions of $(H,\nabla),(E,\theta)$ along $x$ and $x'$, respectively. We observe that the inverse Cartier transform applies to $(\hat E, \hat \theta - df' \wedge)$, but the resulting bundle with integrable connection is not isomorphic to $(\hat H, \hat \nabla + df \wedge)$. However, we find that there exists another Higgs field $\Theta$ on $\hat E$, which may be viewed as a variant of the $\alpha$-transform of $\hat \theta - df' \wedge$, such that its inverse Cartier transform is isomorphic to $(\hat H, \hat \nabla + df \wedge)$, and its associated Higgs complex is isomorphic to that of $(\hat E, \hat \theta - df' \wedge)$. This leads to the following.











\iffalse
A Landau-Ginzburg model over $\mathbb C$ is a pair $(X,f)$, where $X$ is a smooth quasi-projective algebraic variety of dimension $n$ over $\mathbb C$ and $f\in\Gamma(X,\mathcal O_X)$. Our starting point is a theorem of Sabbah stated as follows.
 

Since the aforementioned $E_1$-degeneration of a smooth projective variety $X/\mathbb C$ have a mod $p$ proof provided by Deligne and Illusie \cite{DI}, it is natural to ask 
\begin{question}\label{mod p proof}
 Does the above theorem have a mod $p$ proof for nice $\mathcal M$?   
\end{question}
To address this question, we need a nice description of the de Rham complex of $\mathcal M$ that allows for the mod 
$p$ reduction argument. For a general mixed Hodge module $\mathcal M$ on $X$, it is challenging to find such a description. Nevertheless, the following examples meet our expectations. 
\iffalse we hope that there is a {\itshape geometric} logarithmic de Rham bundle $(H,\nabla,F_\bullet)$ on $X$ with pole along a simple normal crossing divisor $D$ such that 
$$
(DR_X(\mathcal M), F_\bullet)\subset(\Omega^*(H,\nabla),F_\bullet).
$$
Here $\Omega^*(H,\nabla)$ is the de Rham complex of $(H,\nabla)$ and $DR_X(\mathcal M)$ can be \emph{functorially} characterized by $\nabla$.  Let $S$ be any scheme, and let $X'$ be a smooth $S$-scheme endowed with a simple normal crossing divisor $D'$ relative to $S$. Let $(H',\nabla')$ (resp. $(E',\theta')$) be a flat bundle (resp. Higgs bundle) on $(X',D')/S$.
Then we hope that the aforementioned characterization applies to the de Rham complex (resp. Higgs complex) of $(H',\nabla')$ (resp. $(E',\theta')$). The resulting subcomplex of $\Omega^*(H',\nabla')$ (resp. $\Omega^*(E',\theta')$) is called {\itshape mixed Hodge module type, or simply  MHM type. If we want to emphasize the type depends on $\mathcal M$, we call it is $\mathcal M$-type.} This notion is somehow ambiguous, we list some explicit\fi
\begin{example}\label{MHM}
Let $\overline f:\overline X\to\mathbb P^1$ be a compactification of $f:X\to\mathbb A^1$ such that $P:=f^*([\infty])_{red}$ is a SNCD on $\overline X$, and let $\overline D$ be a SNCD on $\overline X$ containing $P$.

$\mathrm{(i)}$ Let $g:( Y, E)\to(\overline X,\overline D)$ be a semistable family (i.e., $g$ is projective, $ E=g^*\overline D$ and $g: Y- E\to \overline X-\overline D$ is smooth). Set
$$
(\overline H,\overline\nabla,\overline F_\bullet):=(R^bg_*\Omega^*_{ Y/\overline X}(\log  E/\overline D),\nabla^{GM},F_\bullet)
$$
to be a Gauss-Manin system associated to $g$. Clearly, $(\overline H,\overline\nabla,\overline F_\bullet)|_{\overline X-\overline D}$ underlies a pure Hodge module $\mathcal M$. Consider the intermediate extension $j_{!*}(\mathcal M)$, for which we have
$$
(DR_X(j_{!*}(\mathcal M)),F_\bullet)^\mathrm{an}\cong(\Omega_{int}^*(H,\nabla),F_\bullet)^\mathrm{an}. 
$$
Here, $j:\overline X-\overline D\hookrightarrow X$ is the natural inclusion, $(H,\nabla,F_\bullet):=(\overline H,\overline\nabla,\overline F_\bullet)|_X$ and $\Omega_{int}^*(H,\nabla)$ is the intersection subcomplex of $\Omega^*(H,\nabla)$ (see \cite[\S2]{SZ2}).

$\mathrm{(ii)}$ Let $g:(Y,E)\to(\overline X,P)$ be a semistable family and let
$$
(H,\nabla,F_\bullet):=(Rg_*\Omega^*_{(Y-E)/X},\nabla^{GM},F_\bullet).
$$
Clearly, the Gauss-Manin system underlies a pure Hodge module $\mathcal M$ on $X$. Consider the mixed Hodge module $j_*\mathcal M|_{\overline X-\overline D}$ with wieght filtration $W_\bullet$, for which we have
$$
(DR_X(W_ij_*\mathcal M|_{\overline X-\overline D}),F_\bullet)^\mathrm{an}\cong((H\otimes W_i\Omega^\bullet_{X/\mathbb C}(\log D),\nabla),F_\bullet)^\mathrm{an}.
$$
Here $D:=\overline D|_X$ and $W_\bullet\Omega^*_{X/\mathbb C}(\log D)$ is the canonical weight filtration on $\Omega^*_{X/\mathbb C}(\log D)$.

$\mathrm{(iii)}$  Notation as in $\mathrm{(ii)}$. Assume $D_0:=f^*([0])\subset D$. Consider the  Beilinson's maximal extension $\Xi_f(\mathcal M|_{X-D_0})$of $\mathcal M|_{X-D_0}$, for which we have
$$
(DR_X(\Xi_f(\mathcal M|_{X-D_0})),F_\bullet)^\mathrm{an}\cong((H\otimes\Omega^\bullet_f,\nabla),F_\bullet)^\mathrm{an}.
$$
Here $\Omega^*_f\subset\Omega_{X/\mathbb C}^*(\log D)$ is the Kontsevich subcomplex (see \cite[Appendix E]{ESY}).
\end{example}
Next, for mixed Hodge modules as in the example above, 
we reduce Theorem \ref{Sabbah's theorem} to a statement in positive characteristic. Let $(H,\nabla,F_\bullet)$ be as in (i) (resp. (ii) or (iii)). Set
$$K_{dR}^*:=\Omega^*_{int}(H,\nabla)~(\mathrm{resp.}~(H\otimes W_i\Omega^\bullet_{X/\mathbb C}(\log D),\nabla)~\mathrm{or}~(H\otimes\Omega^\bullet_f,\nabla))\subset\Omega^*(H,\nabla)
$$
and 
$$(E,\theta):=gr^F(H,\nabla),~K^*_{Hig}:=gr^FK^*_{dR}\subset\Omega^*(E,\theta).$$
Then Sabbah's theorem says that
$$
\mathbb H^i(X,(K_{dR}^\bullet,\nabla+df\wedge))=\mathbb H^i(X,(K_{Hig}^\bullet,\theta-df\wedge)),~\forall i.
$$
By spreading-out technique, to prove the above equality, it suffices to show
$$
\mathbb H^i(X_p,(K_{dR,p}^\bullet,\nabla_p+df_p\wedge))=\mathbb H^i(X_p,(K_{Hig,p}^\bullet,\theta_p-df_p\wedge)),~\forall i,
$$
where the subscript $p$ means sufficiently good $p$ reduction. consequently, it is enough to show
\begin{eqnarray}\label{Sabbah}
F_*(K_{dR,p}^\bullet,\nabla_p+df_p\wedge)\cong(K_{Hig,p}^\bullet,\theta_p-df_p\wedge)~\mathrm{in}~D(X_p).
\end{eqnarray}
Here $F$ is the Frobenius endomorphism of $X_p$. Since $(H,\nabla,F_\bullet)$ comes from geometry, then by Faltings \cite[Theorem 4.1*]{Fa} or Ogus-Vologodsky \cite[Theorem 4.17]{OV}, we know that the inverse Cartier transform of $(E_p,\theta_p)$ is $(H_p,\nabla_p)$. Moreover, $K^*_{dR,p},K^*_{Hig,p}$ are subcomplexes of $\Omega^*(H_p,\nabla_p),\Omega^*(E_p,\theta_p)$, respectively.














To refine the statement of Slogan \ref{slogan}, we specify that, for the reminder of this paper, $K^*_{\mathrm{dR},\theta}$ and $K^*_{\mathrm{Hig},\theta}$ will be taken one of the following.
\begin{example}\label{example MHM}
\upshape
(1) $K^*_{\mathrm{dR},\theta}=\Omega^*(H,\nabla)$ and $K^*_{\mathrm{Hig},\theta}=\Omega^*(E,\theta)$.
   
(2) (Weight filtration \cite[Definition 8.3.19]{CEGT}) Suppose that $(H,\nabla)$ and $(E,\theta)$ have no pole along $D$, namely
    $$
    \nabla(H)\subset H\otimes\Omega^1_{X/k},~\theta(E)\subset E\otimes\Omega^1_{X/k}.
    $$
    Define the weight filtration $W_i$ on $\Omega^q(H,\nabla)=H\otimes\Omega^\bullet_{X/k}(\log D)$ by
 $$
 \begin{array}{c}
 W_i\Omega^q(E,\nabla):=\left\{\begin{matrix}E\otimes\mathrm{Im}(\Omega_{X/k}^i(\log D)\otimes\Omega^{q-i}_{X/k}\to\Omega_{X/k}^q(\log D)),&i\leq q,\\
E\otimes\Omega_{X/k}^q(\log D),&i>q.
 \end{matrix}\right.\\
 \end{array}
 $$
Clearly, we have $\nabla(W_i\Omega^q(H,\nabla))\subset W_i\Omega^{q+1}(H,\nabla)$, from which we obtain a weight $i$ de Rham subcomplex $W_i\Omega^*(H,\nabla)\subset\Omega^*(H,\nabla)$. Similarly, one can define the weight $i$ Higgs subcomplex $W_i\Omega^*(E,\theta)\subset\Omega^*(E,\theta)$.

(3) (Intersection complexes \cite[Definition 8.3.31]{CEGT}, \cite[\S2]{SZ2}) Define the intersection de Rham subcomplex $\Omega^*_\mathrm{int}(H,\nabla)\subset\Omega^*(E,\nabla)$ as follows: let $t_1,\cdots,t_n$ be a system of coordinates on an open subset $U$ of $X$ such that $D|_U:=(t_1\cdots t_n=0)$, then
 $$
 \Omega^\bullet_{int}(E,\nabla)|_U:=\bigoplus_{I\subset\{1,\cdots,n\}}(\sum_{J\subset I}t_J\nabla_{I-J}H)\otimes d_I\log t.
 $$
 Here $t_J:=\prod_{j\in J}t_j$ for $J\neq\emptyset$ and $t_\emptyset:=1$,  $\nabla_I:=\prod_{i\in I}\nabla_i$ for $I\neq\emptyset$ and $\nabla_\emptyset:=id$ ($\nabla=\sum_{i=0}^n\nabla_i\otimes d\log t_i$), and $d_I\log t:=d\log t_{i_1}\cdots d\log t_{i_r}$ for $I=\{i_1,\cdots,i_r\},i_1<\cdots<i_r$ and $d_\emptyset\log t:=1$. Similarly, one can define the intersection Higgs subcomplex $\Omega^*_\mathrm{int}(E,\theta)\subset\Omega^*(E,\theta)$.

(4) (Kontsevich complexes \cite[Definition 2.11]{KKP}, \cite[\S1.3]{ESY}) Let $\tilde g:(\tilde X,\tilde D)\to(\mathbb A_{W_2(k)}^1,\tilde E)$ be a semistble family along a normal crossing divisor $\tilde E$ over $W_2(k)$ such that $0\in \tilde E$. Let $g:(X,D)\to(\mathbb A^1_k,E)$ be the mod $p$ reduction of $\tilde g$. Define the Kontsevich-de Rham subcomplex $\Omega_{g^{-1}}^*(H,\nabla)$ of $\Omega^*(H,\nabla)$ by
$$
\Omega^i_{g^{-1}}(H,\nabla):=H\otimes\Omega^i_{g^{-1}},~\Omega^i_{g^{-1}}:=\mathrm{Ker}(\Omega^i_{X/k}(\log D)\xrightarrow{dg^{-1}\wedge}\Omega^{i+1}_{X/k}(*D)/\Omega^{i+1}_{X/k}(\log D)).
$$
Similarly, one can define the Kontsevich-Higgs subcomplex $\Omega_{g^{-1}}^*(E,\theta)$.
\end{example}

\subsubsection{Local de Rham-Higgs comparison}
As shown in \cite[Remarques 2.2 (ii)]{DI}, \cite[Theorem 2.26, 2]{OV}, \cite[Theorem 5.2 (ii)]{Schepler} and \cite[Proposition 3.8]{SZ2}, \eqref{de-H com for MHM} should be true Zariski locally without truncation. The following theorem provides supporting evidence for this observation.
\fi


\begin{theorem}\label{de Rham-Higgs local}
For any closed point $x\in X$ and any $f\in\Gamma(X,\mathcal O_X)$, we have an isomorphism of complexes
$$
F_*(K^\bullet_{\mathrm{dR}},\nabla+df\wedge)\otimes\hat{\mathcal O}_{X',x'}\cong(K^\bullet_{\mathrm{Hig}},\theta-df'\wedge)\otimes\hat{\mathcal O}_{X',x'}.
$$
\end{theorem}
This theorem gives a satisfactory answer to the above question in a formal neighborhood of a closed point of $X'$. It also sheds light on the global situation discussed below.

\begin{corollary}
The supports of both sides of \eqref{de-H com for MHM} coincide, namely
$$\mathrm{Supp}~\mathcal H^j(K^\bullet_{\mathrm{dR}},\nabla+df\wedge)=\mathrm{Supp}~\mathcal H^j(K^\bullet_{\mathrm{Hig}},\theta-df'\wedge),~\forall j.$$
In particular, if the right hand side support consisting of isolated closed points of $X$ for any $j$, then
$$
F_*(K^\bullet_{\mathrm{dR}},\nabla+df\wedge)\cong(K^\bullet_{\mathrm{Hig}},\theta-df'\wedge)~\mathrm{in}~D(X').
$$
\end{corollary}
When $(E,\theta)=(\mathcal O_X,0)$, $K^*_{\mathrm{dR}}=F_*\Omega^*_{X/k}$, $K^*_{\mathrm{Hig}}=\bigoplus_{i=0}^\infty\Omega^i_{X/k}[-i]$, and $f$ has only isolated singularities, this corollary was mentioned in Ogus \cite{Ogus04}.

Next, let $\tilde{f}: (\tilde{X},\tilde D) \to (\mathbb{A}^1_{W_2(k)},\tilde\Delta)$ be a semistable family over $W_2(k)$. Assume that $\Delta=\sum_{i=1}^m[\tilde \lambda_i]$, where $\tilde\lambda_i\in W_2(k)$. Let $f,\Delta,\lambda_i$ be the mod reductions of $\tilde f,\tilde\Delta,\tilde\lambda_i$, respectively. 

Set $K^*_{\mathrm{dR}} \subset \Omega^*_{X/k}(\log D)$ to be the subcomplex defined as follows: on $X-D$, define $K^*_{\mathrm{dR}} := \Omega^*_{X/k}$; around a fiber $f^{-1}(\lambda_i)$, define $K^*_{\mathrm{dR}} := W_j\Omega^*_{X/k}(\log D)$ or $\Omega^*_{(f - \lambda_i)^{-1}}$ (see \S2.6 for more details). Define $K^*_{\mathrm{Hig}}$ similarly.

\begin{corollary}\label{intro tau<p structure sheaf}
Notation and assumptions as above. Then
$$
\tau_{<p}F_*(K^\bullet_\mathrm{dR},d+df\wedge)\cong \tau_{<p}(K^\bullet_\mathrm{Hig},-df'\wedge)~\mathrm{in}~D(X').
$$
\end{corollary}
As mentioned at the beginning, when \( K^*_{\mathrm{dR}} = \Omega^*_{X/k} \), the above corollary was previously obtained by \cite{OV} and \cite{ACH}, under different technical assumptions on \( f \). 

\subsubsection{Answers provided by the second finding}
 We now proceed to the second finding and its related consequences. Using the stacky approach, Drinfeld observed that there is a $\mu_p$-action on the de~Rham complex $\Omega^\ast_{X/k}$, which induces a decomposition of the two-sided truncation $\tau_{[a,a+p-1]}F_*\Omega^\ast_{X/k}$ for any $a \geq 0$. A complete proof of this result was subsequently given by Bhatt--Lurie~\cite{BL}, Li--Mondal~\cite{LM}, and Ogus~\cite{O24}. Using the theory of abstract Koszul complexes, Achinger-Suh~\cite{AS} established a slightly weaker version of the aforementioned decomposition theorem, proving that the truncated complex $\tau_{[a,a+p-2]} F_*\Omega^\ast_{X/k}$ admits a decomposition for all $a \geq 0$. Moreover, by adapting a variant of the \v{C}ech construction from~\cite{DI}, they also provided an alternative proof of the decomposition of the two-term truncation $\tau_{[a,a+1]} F_*\Omega^*_{X/k}$ in the same paper.


By further investigating the \v{C}ech construction in~\cite{AS}, we find that a much more general result underlies their construction. More precisely, we establish a connection between splittings of the logarithmic cotangent bundle $\Omega^*_{X/k}(\log D)$ and quasi-isomorphisms from $K^*_\mathrm{Hig}$ to $\check{\mathcal C}(\mathcal U',F_*K^*_\mathrm{dR})$, where $\mathcal U$ is an open affine covering of $X$.

\begin{proposition}\label{intro thm quasi-iso}
Let $\Omega^1_{X/k}(\log D)=\bigoplus_{i=1}^\beta\Omega_i$ be a splitting, and set $\underline\Omega:=(\Omega_1,\cdots,\Omega_\beta)$.
Let $\mathcal U$ be an open affine covering of $X$. For suitable subcomplexes 
$K^*_\mathrm{Hig}\subset\Omega^*(E,\theta)$ and $K^*_\mathrm{dR}\subset\Omega^*(H,\nabla)$, there is a quasi-isomorphism
\begin{eqnarray}\label{intro cech construction omega}
\varphi_{(\tilde X,\tilde D),\underline\Omega}:\tau_{<q}K^*_\mathrm{Hig}\to\check{\mathcal C}(\mathcal U',\tau_{<q}F_*K^*_\mathrm{dR}),
\end{eqnarray}
where $q=\infty$ if $\mathrm{rank}(\Omega_i)<p-\ell$ holds for any $1\leq i\leq\beta$ and $p-\ell$ otherwise. In particular, it induces an isomorphism
$$
\Phi_{(\tilde X,\tilde D),\underline\Omega}\in\mathrm{Hom}_{D(X')}(\tau_{<q}F_*K^*_\mathrm{dR},\tau_{<q}K^*_\mathrm{Hig}).
$$
\end{proposition}

Similar to Petrov's result in \cite{P2}, we highlight a special case of the above proposition.

\begin{corollary}
Suppose that the logarithmic cotangent bundle of the $W_2(k)$-liftable pair $(X,D)/k$ splits as a direct sum of subbundles of rank less than $p$. Then the complex $F_*\Omega^*_{X/k}(\log D)$ is decomposable. In particular, if $A$ is an abelian variety over $k$, then $F_*\Omega^*_{A/k}$ is decomposable.
\end{corollary}

According to Yobuko \cite{Y}, an abelian variety $A$ of dimension $g$ over $k$ is quasi-$F$-split if and only if its $p$-rank is $g$ or $g-1$. In particular, if its $p$-rank less than $g-1$, then $A$ is not quasi-$F$-split. Consequently, in that case, the decomposition of $\Omega^*_{A/k}$ cannot be deduced directly from \cite{P2}.

A priori, the construction of $\varphi_{(\tilde X,\tilde D),\underline{\Omega}}$ depends on the \emph{ordered splitting} $\underline{\Omega}$. For example, when $\beta = 2$, we generally have $\varphi_{(\tilde X,\tilde D),(\Omega_1, \Omega_2)} \neq \varphi_{(\tilde X,\tilde D),(\Omega_2, \Omega_1)}$. Nevertheless, the following proposition shows that $\varphi_{(\tilde X,\tilde D),\underline{\Omega}}$ is independent of the choice of $\underline{\Omega}$ up to homotopy. When $\beta = 1$, we simply denote the morphism $\varphi_{(\tilde X,\tilde D),\underline{\Omega}}$ by $\varphi_{(\tilde X,\tilde D)}$, which has already been constructed in \cite{SZ1}.


\begin{proposition}\label{intro thm homotopy}
 Notation and assumptions as in the above proposition. Then $\varphi_{\underline{\Omega}}$ is homotopic to $\varphi_{\mathrm{SZ}}$ for $q = p - \ell$. Let 
\[
\underline{\Omega}' := (\cdots, \Omega_{o-1}, \Omega_o \oplus \Omega_{o+1}, \Omega_{o+2}, \cdots)
\]
for some \(1 \leq o < \beta\). Assume that $\mathrm{rank}(\Omega_i) < p - \ell$ holds for any \(1 \leq i \leq \beta\), and that 
\[
\mathrm{rank}(\Omega_o) + \mathrm{rank}(\Omega_{o+1}) < p - \ell,
\]
then $\varphi_{(\tilde X,\tilde D),\underline{\Omega}}$ is homotopic to $\varphi_{(\tilde X,\tilde D),\underline{\Omega}'}$ for $q=\infty$.
\end{proposition}
Combing Propositions \ref{intro thm quasi-iso} and \ref{intro thm homotopy}, we have the following generalization of Drinfeld's observation. 
\begin{theorem}
$(1)$ Suppose that the nilpotent level of $(E, \theta)$ is strictly less than $p - 2$, namely, $\ell < p - 2$. Then for any $a \geq 0$, we have

$$
\tau_{[a,a+1]}F_*K^*_\mathrm{dR}\cong\tau_{[a,a+1]}K^*_\mathrm{Hig}~\mathrm{in}~D(X').
$$
$(2)$ Under the notation and assumptions above Corollary~\ref{intro tau<p structure sheaf}, we have

$$
\tau_{[a,a+1]}F_*(K^\bullet_\mathrm{dR},d+df\wedge)\cong\tau_{[a,a+1]}(K^\bullet_\mathrm{Hig},-df\wedge)~\mathrm{in}~D(X').
$$
\end{theorem}
As an immediate consequence of the above theorem, we obtain isomorphisms between cohomology sheaves.

\begin{corollary}
Under the notation and assumptions in \textup{(1)} of the above theorem, we have
\[
F_*\mathcal H^a(K^*_{\mathrm{dR}}) \cong \mathcal H^a(K^*_{\mathrm{Hig}}), \quad \forall a.
\]
Under the notation and assumptions above Corollary~\ref{intro tau<p structure sheaf}, we have
\[
F_*\mathcal H^a(K^\bullet_{\mathrm{dR}}, d + df \wedge) \cong \mathcal H^a(K^\bullet_{\mathrm{Hig}}, -df \wedge), \quad \forall a.
\]
\end{corollary}

When \( K^*_{\mathrm{dR}} = \Omega^*(H,\nabla) \) and \( K^*_{\mathrm{Hig}} = \Omega^*(E,\theta) \), the first isomorphism of this corollary has been obtained by \cite{Ogus04}. 

\subsection{Further consequences}
Several interesting consequences arise from the partial answers to Question \ref{Question 1}. In what follows, we examine a few of them. 

\subsubsection{K\"unneth formula}
 Proposition~\ref{intro thm quasi-iso} is somewhat unexpected, yet perfectly reasonable. Suppose that there are splittings
\begin{eqnarray}\label{intro assumption splitting}
(\tilde{X}, \tilde{D}) = \prod_{i=1}^\beta (\tilde{X}_i, \tilde{D}_i), \quad (E, \theta) = (E_1, \theta_1)\boxtimes\cdots\boxtimes(E_\beta,\theta_\beta),\quad
\mathcal{U} = \mathcal{U}_1 \times \cdots \times \mathcal{U}_\beta,
\end{eqnarray}
where each \((E_i, \theta_i)\) is a nilpotent Higgs bundle on \((X_i, D_i)/k\), which is the mod \(p\) reduction of \((\tilde{X}_i, \tilde{D}_i)/W_2(k)\), of level \(\leq\ell\) $(<p-\dim X)$, and each \(\mathcal{U}_i\) is an open affine covering of \(X_i\). Set $$(H_i,\nabla_i):=C^{-1}_{(\tilde X_i,\tilde D_i)/W_2(k)}(E_i,\theta_i),\quad \underline\Omega:=(\pi_1^*\Omega^1_{X_1/k}(\log D_1),\cdots,\pi_\beta^*\Omega^1_{X_\beta/k}(\log D_\beta)),$$
where $1\leq i\leq\beta$ and $\pi_i:X\to X_i$ is the natural projection. We have natural identifications 
$$
\Omega^*(E,\theta)\cong\Omega^*(E_1,\theta_1)\boxtimes\cdots\boxtimes\Omega^*(E_\beta,\theta_\beta),\quad \Omega^*(H,\nabla)\cong\Omega^*(H_1,\nabla_1)\boxtimes\cdots\boxtimes\Omega^*(H_\beta,\nabla_\beta).$$
Let $\Phi_{(\tilde X_i,\tilde D_i)}$ denote the isomorphism from $F_*\Omega^*(H_i,\nabla_i)$ to $\Omega^*(E_i,\theta_i)$ induced by $\varphi_{(\tilde X_i,\tilde D_i)}$.
Define 
$$
\cup:\check{\mathcal C}(\mathcal U_\beta',F_*\Omega^*(H_\beta,\nabla_\beta))\boxtimes\cdots\boxtimes\check{\mathcal C}(\mathcal U_1',F_*\Omega^*(H_1,\nabla_1))\to\check{\mathcal C}(\mathcal U',F_*\Omega^*(H,\nabla))$$
to be the cup product of \v{C}ech complexes. 
\begin{theorem}[K\"unneth formula]
Notation and assumptions as above. Then there are splittings
\begin{eqnarray}\label{intro splitting}
(H,\nabla)\cong(H_1,\nabla_1)\boxtimes\cdots\boxtimes(H_\beta,\nabla_\beta),\quad \varphi_{(\tilde X,\tilde D),\underline\Omega}=\cup\circ(\varphi_{(\tilde X_\beta,\tilde D_\beta)}\boxtimes\cdots\boxtimes\varphi_{(\tilde X_1,\tilde D_1)}).
\end{eqnarray}
In particular, we have
$$
\Phi_{(\tilde X,\tilde D)}=\Phi_{(\tilde X_1,\tilde D_1)}\boxtimes\cdots\boxtimes\Phi_{(\tilde X_\beta,\tilde D_\beta)}\in\mathrm{Hom}_{D(X')}(\Omega^*(E,\theta),F_*\Omega^*(H,\nabla)).
$$
\end{theorem}
The second equality of \eqref{intro splitting} shows that although its right-hand side appears to depend on the splittings in \eqref{intro assumption splitting}, it is in fact independent of them, and depends only on $(\tilde X,\tilde D)/W_2(k)$, $(E,\theta)$, the covering $\mathcal U$, and the ordered splitting $\underline\Omega$.
\iffalse
Denote by $\mathrm{SmProp}_{W_2(k)}^{\dim < p}$ the category of smooth proper varieties over $W_2(k)$ of dimension strictly less than $p$. For any object $\tilde{X}$ in this category and any integer $0\leq i\leq p-2$, we associate the triple
\[
\left( H_{\mathrm{DR}}^i(X/k),\; F^\bullet H_{\mathrm{dR}}^i(X/k),\; \phi^i_{\tilde{X}/W_2(k)} \right),
\]
where $H_{\mathrm{DR}}^i(X/k)$ denotes the de Rham cohomology of $X/k$ (the mod $p$ reduction of $\tilde{X}/W_2(k)$), $F^\bullet$ is the Hodge filtration on the de Rham cohomology, and
\[
\phi^i_{\tilde{X}/W_2(k)} \colon F_k^* \operatorname{Gr}_FH_{\mathrm{dR}}^i(X/k) \longrightarrow H_{\mathrm{DR}}^i(X/k)
\]
is the $k$-linear isomorphism induced by $\Phi_{\tilde{X}/W_2(k)}$. This triple defines a \emph{Fontaine--Laffaille module}, which corresponds, via Fontaine--Laffaille theory, to a crystalline Galois representation of $G_K$ over $\mathbb{F}_p$, where $K$ is the fraction field of $W(k)$ and $G_K := \operatorname{Gal}(\overline{K}/K)$. Consequently, we obtain a functor
\[
T_\mathrm{cris}^i \colon \mathrm{SmProp}_{W_2(k)}^{\dim < p} \longrightarrow \operatorname{Rep}_{\mathbb{F}_p}^{\mathrm{cris}, [0, p-2]}(G_K).
\]
\begin{corollary}
For any $0\leq i\leq p-2$, we have
$$
T_\mathrm{cris}^i(\tilde X_1\times_{W_2(k)}\tilde X_2)\cong\bigoplus_{i_1+i_2=i}T_\mathrm{cris}^{i_1}(\tilde X_1)\otimes_{\mathbb F_p}T_\mathrm{cris}^{i_2}(\tilde X_2),
$$
where $\tilde X_1,\tilde X_2,\tilde X_1\times_{W_2(k)}\tilde X_2\in\mathrm{SmProp}_{W_2(k)}^{\dim<p}$.
\end{corollary}

To state the cohomological consequence of this theorem, we introduce the following notation, which is closely related to the Fontaine-Lafaille theory.
\begin{definition}
A Fontaine-Laffaille module over $k$ is a triple $(V,F^\bullet V,\psi)$, where $V$ is a finite dimensional vector space over $k$, $F^\bullet V$ a decreasing filtration on $V$ such that $F^0V=V$ and $F^mV=0$ for $m\gg0$, and $\psi:\sigma^*_k\mathrm{Gr}_FV\to V$ is an isomorphism of $k$-vector spaces ($\sigma_k$ is the Frobenius endomorphism of $k$).
\end{definition}
Suppose that each $X_i$ is proper. Set 
$$
V^m=\mathbb H^m(X,\Omega^*_{X/k}(\log D)),~F^qV:=\mathrm{Im}(\mathbb H^m(X,\Omega^{\geq q}_{X/k}(\log D))\to\mathbb H^m(X,\Omega^*_{X/k}(\log D))).
$$
One can check that the decomposition $\Phi_{(\tilde X,\tilde D)/W_2(k)}$ induces a $k$-linear isomorphism $$\psi_{(\tilde X,\tilde D)/W_2(k)}:\sigma_k^*\mathrm{Gr}_FV^m\to V^m.$$ Consequently, we obtain a Fontaine-Laffaille module $$\mathrm{FL}^m_{(\tilde X,\tilde D)/W_2(k)}:=(V^m,F^\bullet V^m,\psi_{(\tilde X,\tilde D)/W_2(k)}).$$
Similarly, we have Fontaine-Laffaille modules $\mathrm{FL}^m_{(\tilde X_i,\tilde D_i)/W_2(k)}$ associated to $(\tilde X_i,\tilde D_i)/W_2(k)$  ($i=1,2$).
\begin{corollary}
If $(\tilde X,\tilde D):=(\tilde X_1,\tilde D_1)\times(\tilde X_2,\tilde D_2)$, then for any $m\geq0$ we have
$$
\mathrm{FL}^m_{(\tilde X,\tilde D)/W_2(k)}=\bigoplus_{m_1+m_2=m}\mathrm{FL}^{m_1}_{(\tilde X_1,\tilde D_1)/W_2(k)}\otimes\mathrm{FL}^{m_2}_{(\tilde X_2,\tilde D_2)/W_2(k)}.$$
\end{corollary}
This corollary can be regarded as a positive characteristic analogue of the K\"unneth formula for Hodge structures of smooth quasi-projective varieties over $\mathbb C$.
\fi

\subsubsection{$E_1$-degeneration and vanishing theorem modulo $p^n$}
\iffalse
\begin{definition}\label{cotangent splitting}
Let $X$ be a smooth variety over $k$ endowed with a simple normal crossing divisor $D$.
For some $0\leq\ell<p$, we say $(X,D)/k$ is $\gamma$-$\Omega$-split if it admits a $W_2(k)$-lifting $(\tilde X,\tilde D)/W_2(k)$ and
\begin{eqnarray}\label{ell-omega-split}
\Omega^1_{X/k}(\log D)=\bigoplus_{i=1}^\beta\Omega_i,~\gamma=\max\{\mathrm{rank}(\Omega_i):1\leq i\leq\beta\}
\end{eqnarray}
\end{definition}
It is clear that abelian varieties over $k$ are $(p{-}2)$-$\Omega$-split, owing to the triviality of their cotangent bundles.


\begin{theorem}\label{de Rham-Higgs without truncation}
Let $(X,D)/k$ be an $\ell$-$\Omega$-split pair, and let $(\tilde X,\tilde D)/W_2(k)$ be a $W_2(k)$-lifting of it. Let $(E,\theta)$ be a nilpotent Higgs bundle of level $\leq \ell$ on $(X,D)/k$, and define $(H,\nabla) := C^{-1}_{(\tilde X,\tilde D)/W_2(k)}(E,\theta)$. Let $K^*_{\mathrm{dR}} \subset \Omega^*(H,\nabla)$ and $K^*_{\mathrm{Hig}} \subset \Omega^*(E,\theta)$ be the respective specified subcomplexes.
 Then
\begin{eqnarray}\label{decom without truncation f=0}
F_*K^*_{\mathrm{dR}}\cong K^*_{\mathrm{Hig}}~\mathrm{in}~D(X).
\end{eqnarray}
\end{theorem}
This theorem incorporates the perspective of mixed Hodge modules, suggesting that we may be able to establish an analogue of the theory of mixed Hodge modules in positive characteristic. Inspired by the above theorem, For quasi-$F$-split varieties over $k$, we may expect more than \cite{P2}.
\begin{question}
Let $X$ be a smooth quasi-$F$-split variety over $k$ and let $\tilde X$ be its lifting over $W_2(k)$. Let $(E,\theta)$ be a ninpotent Higgs bundle on $X/k$ of level $\leq p-1$ and set $(H,\nabla):=C^{-1}_{\tilde X/W_2(k)}(E,\theta)$. Then do we have
$$
F_*\Omega^*(H,\nabla)\cong\Omega^*(E,\theta)~\mathrm{in}~D(X)?
$$
\end{question}


\iffalse
Set 
$$
\mathrm{Supp}~(K^\bullet_\mathrm{dR},\nabla+df\wedge):=\bigcup_{j=0}^\infty\mathrm{Supp}~\mathcal H^j(K^\bullet_\mathrm{dR},\nabla+df\wedge)
$$
and define $\mathrm{Supp}~(K^\bullet_\mathrm{Hig},\theta-df\wedge)$ similarly.




The discussion above shows that the solution to Question \ref{mod p proof} depends on the existence of an isomorphism \eqref{Sabbah}. Although proving this in general is quite challenging, its existence remains promising. Before formally summarizing the desired isomorphism \eqref{Sabbah} as a conjecture, we distill Example \ref{MHM} into the following
\begin{definition}\label{three type MHM}
 Let $k$ be a field, let $X$ be a smooth variety of dimension $n$ over $k$, and let $D\subset X$ be a SNCD on $X$. 
 
 $\mathrm{(i)}$ Let $(E,\nabla)$ be a logarithmic flat (resp. Higgs) bundle on $(X,D)/k$.

 $\mathrm{(ii)}$ 

$\mathrm{(iii)}$ 
\end{definition}


\begin{conjecture}\label{conjecture}
Let $k$ be an algebraically closed field of characteristic $p>0$, let $(X,f)$ be a Landau-Ginzburg model over $k$, let $D\subset X$ be a normal crossing divisor, and let $g\in\Gamma(X,\mathcal O_X)$ such that $g^*([0])\subset D$. Assume $(X,D,f,g)/k$ has a $W_2(k)$-lifting $(\tilde X,\tilde D,\tilde f,\tilde g)/W_2(k)$ such that $\tilde g^*([0])\subset\tilde D$. Let $(E,\theta)$ be a nilpotent Higgs bundle on $(X,D)/k$ of level $\leq\ell~(\ell\leq p-2)$. Set $(H,\nabla)=C^{-1}_{(\tilde X,\tilde D)/W_2(k)}(E,\theta)$.
\begin{itemize}
\iffalse \item[$\mathrm{(i)}$] There is an isomorphism
    \begin{eqnarray}\label{Con1}
    \tau_{<p-\ell}F_*\Omega^*(H,\nabla+df\wedge)\cong\tau_{<p-\ell}\Omega^*(E,\theta-df\wedge)~\mathrm{in}~D(X);
    \end{eqnarray}\fi
\item[$\mathrm{(C1)}$] Let $K^*_{dR}=\Omega^*(H,\nabla)$, $\Omega_{int}^*(H,\nabla)$, $W_i\Omega^*_{\log}(H,\nabla)$ or $\Omega^*_{g^{-1}}(H,\nabla)$, $K^*_{Hig}=\Omega^*(E,\theta)$, $\Omega_{int}^*(E,\theta)$, $W_i\Omega^*_{\log}(E,\theta)$ or $\Omega^*_{g^{-1}}(E,\theta)$). Then
\begin{eqnarray}\label{Con2}
\tau_{<p-\ell}F_*(K^\bullet_{dR},\nabla+df\wedge)\cong \tau_{<p-\ell}(K^\bullet_{Hig},\theta-df\wedge)~\mathrm{in}~D(X).
\end{eqnarray}
 \item[$\mathrm{(C2)}$] The statement $\mathrm{(ii)}$ can be strengthened by a two-sided truncation $\tau_{[a,b]}$ for $b=a+p-\ell-1$. i.e., 
\begin{eqnarray}\label{Con3}
\tau_{[a,b]}F_*(K^\bullet_{dR},\nabla+df\wedge)\cong \tau_{[a,b]}(K^\bullet_{Hig},\theta-df\wedge)~\mathrm{in}~D(X).
\end{eqnarray}
\end{itemize}
\end{conjecture}
Note that (i) is a special case of (ii). There is already some evidence supporting this conjecture, which we list below.
\begin{itemize}
    \item For (C1), if we take $K_{Hig}^*=\Omega^*(E,\theta), K_{dR}^*=\Omega^*(H,\nabla)$ and further suppose $f=0$, then it was proved by Ogus-Vologodsky\cite{OV} and Schepler\cite{Schepler}. If  we take $(E,\theta)=(\mathcal O_X,0)$, together with some technical conditions (these conditions do not require $f$ equals $0$), then it was proved by Ogus-Vologodsky\cite{OV} again and Arinkin-C\u{a}ld\u{a}raru-Hablicsek\cite{ACH}. If we take $f=0$, $K^*_{dR}=\Omega^*_{g^{-1}}(H,\nabla)$ (resp. $\Omega^*_{int}(H,\nabla)$ or $W_i\Omega^*_{\log}(\mathcal O_X,d)$), $K^*_{Hig}=\Omega^*_{g^{-1}}(E,\theta)$ (resp. $\Omega^*_{int}(E,\theta)$ or $W_i\Omega^*_{\log}(\mathcal O_X,0)$), then it was proved by Kontsevich and Esnault-Sabbah-Yu\cite{ESY} (resp. Sheng and the authur \cite{SZ2} or the author\cite{Z}).
    \item For (C2), if we take $f=0, K^*_{dR}=\Omega^*_{X/k}, K^*_{Hig}=\bigoplus_{i=0}^{\dim X}\Omega^i_{X/k}[-i]$, then it was proved by Achinger-Suh\cite{AS}, Bhatt-Lurie\cite{BL}, Drinfeld and Li-Mondal\cite{LM}.
\end{itemize}

Due to the potential bad singularities of $f$, proving (C1) is difficult even for $(E,\theta)=(\mathcal O_X,0)$ and $D=\emptyset$.  However, under some technique conditions which differ from those in \cite{OV} or \cite{ACH}, we provide additional evidence for (i) with more general $(E,\theta)$ and $D$.
In fact, $K_{dR}^*,k_{Hig}^*$ can be taken as any other suitable subcomplexes inspired by the theory of mixed Hodge modules. For the mod $p$ proof of Theorem \ref{Sabbah}, (C1) is enough. Nevertheless, (C2) has its own interest. 
\fi

\iffalse
\subsubsection{Global de Rham-Higgs comparison (without truncation)}
To derive various $E_1$-degeneration and vanishing theorem on $(X,D)/k$, what we truly need is a non-truncated version of \eqref{de-H com for MHM}. However, the truncation is necessary in general, as \cite[Corollary 10.7]{P1} demonstrated  the existence of a $W_2(k)$-liftable smooth projective variety of dimension $>p$, yet its de Rham complex $F_*\Omega^*_{X/k}$ is not decomposable. On the other hand, if $X$ is a nice geometric object, then the decomposition of its de Rham complex is possible. Recently, \cite[Corollaries 4.6, 4.10]{P2} showed that if $X$ is a quasi-$F$-split (i.e., the morphism $\mathcal O_X\xrightarrow{f\mapsto[f]^p}F_*W_n(\mathcal O_X)/p$ has an $\mathcal O_X$-linear splitting) variety, then its de Rham complex $F_*\Omega^*_{X/k}$ is decomposable.

In this paper, we identify another type of splitting on $(X,D)$ which induces the decomposition of its de Rham complex $F_*\Omega^*_{X/k}(\log D)$.
\fi
As immediate consequences of the above theorem, we have the following $E_1$-degeneration and vanishing theorem.\fi
Suppose that there exists a smooth formal $W(k)$-scheme $\mathfrak{X}$ endowed with a simple normal crossing divisor $\mathfrak{D}$ relative to $W(k)$, such that their reductions modulo $p$ recover $(X, D)/k$. Denote by $\mathrm{MF}^\nabla_{[0,\ell]}((\mathfrak X,\mathfrak D)/W(k))$ ($0\leq\ell\leq p-2$) the category of Fontaine-Faltings modules of Hodge-Tate weight $\leq\ell$ on $(\mathfrak X,\mathfrak D)/W(k)$.

\begin{theorem}
Let $(M, \nabla, F^\bullet M, \varphi) \in \mathrm{MF}^\nabla_{[0,\ell]}((\mathfrak{X}, \mathfrak{D})/W(k))$, and let $(\tilde{M}, \tilde{\nabla})$ be the $p$-connection corresponding to $(M, \nabla)$. Assume that $\Omega^1_{X/k}(\log D)$ splits into a direct sum of subbundles, each of rank $< p - \ell$. Then the spectral sequence
\[
E_1^{i,j} = \mathbb{H}^{i+j}(\mathfrak{X}, \mathrm{Gr}_F^i \Omega^*(M, \nabla)) \Rightarrow \mathbb{H}^{i+j}(\mathfrak{X}, \Omega^*(M, \nabla))
\]
degenerates at the $E_1$-page. Moreover, if $X$ is projective over $k$ and $L$ is an ample line bundle on $X$ (regarded as a coherent sheaf on $\mathfrak{X}$), then
\[
\mathbb{H}^i(\mathfrak{X}, \Omega^*(\tilde{M}, \tilde{\nabla}) \otimes L) = 0, \quad i > \dim X.
\]
\end{theorem}
The first part of the above theorem slightly strengthens the $E_1$-degeneration of Faltings \cite[Theorem 4.1$^*$]{Fa} (see also the alternative proof by Xu \cite{X}) in our specific setting. In particular, it recovers the fact that the Hodge-to-de Rham spectral sequence of an abelian variety $A/k$ degenerates at $E_1$, even when $A$ is supersingular and $\dim A > p$, a result originally established by Oda \cite{Oda}. The second part of the theorem slightly strengthens the vanishing theorem of Arapura \cite{Arapura} in our specified situation. In particular, it recovers the vanishing of higher cohomology vector spaces of ample line bundles on $A/k$, a fact that was already guaranteed by a classical result of Mumford \cite{Mumford}.

Following the spirit of Deligne–Illusie \cite{DI}, we expect that there exists a comparison between $\Omega^*(M, \nabla)$ and $\Omega^*(\tilde{M}, \tilde{\nabla})$, which furthermore restricts to subcomplexes inspired by the theory of mixed Hodge modules. The following theorem confirms this expectation in a special case.

\begin{theorem}
 Suppose that $(\mathfrak X,\mathfrak D)/W(k)$ is the fiber product of smooth logarithmic formal curves over $W(k)$. Let $(M, \nabla, F^\bullet M, \varphi) \in \mathrm{MF}^\nabla_{[0,p-2]}((\mathfrak{X}, \mathfrak{D})/W(k))$, and let $(\tilde{M}, \tilde{\nabla})$ be the $p$-connection corresponding to $(M, \nabla)$. Then there is an isomorphism of sheaves of abelian groups
 \begin{eqnarray}\label{De-Hig mod p^n}
\Omega^*(M,\nabla)\cong\Omega^*(\tilde M,\tilde\nabla)~\mathrm{in}~D(\mathrm{Ab}(X)),
 \end{eqnarray}
 which restricts to an isomorphism of intersection subcomplexes
 \begin{eqnarray}\label{De-Hig mod p^n int}
\Omega_\mathrm{int}^*(M,\nabla)\cong\Omega_\mathrm{int}^*(\tilde M,\tilde\nabla)~\mathrm{in}~D(\mathrm{Ab}(X)).
 \end{eqnarray}
 In particular, the spectral sequence
\[
E_1^{i,j} = \mathbb{H}^{i+j}(\mathfrak{X}, \mathrm{Gr}_F^i \Omega_\mathrm{int}^*(M, \nabla)) \Rightarrow \mathbb{H}^{i+j}(\mathfrak{X}, \Omega_\mathrm{int}^*(M, \nabla))
\]
degenerates at the $E_1$-page.
\end{theorem}
This theorem sheds light on the issue raised in \cite[Remark 14.19]{X}, and shows that \cite[Theorem 5.6]{Fa} may be further extended from the perspective of the theory of mixed Hodge modules. Note that when $\mathfrak{X}/W(k)$ admits a Frobenius lifting (or more generally, when $X$ can be embedded into a smooth formal scheme $\mathfrak{Y}/W(k)$ endowed with a Frobenius lifting), the corresponding comparison has been studied by Shiho \cite{Shi}, Ogus \cite{O24}, and Wang \cite{W}.

\subsubsection{$E_1$-degeneration and vanishing theorem in characteristic $0$}

Let $X_\mathbb{C}$ be a smooth projective variety over $\mathbb{C}$, and let $f: X_\mathbb{C} \to \mathbb{P}^1_\mathbb{C}$ be a semistable family. Let $D_\mathbb{C}$ be a simple normal crossing divisor on $X_\mathbb{C}$ which is the pullback of an effective divisor on $\mathbb{A}^1_\mathbb{C}$ containing $\infty$, such that the restriction $f|_{X_\mathbb{C} - D_\mathbb{C}}$ is smooth. Let $L$ be an ample line bundle on $X_\mathbb{C}$. Let $K^*_\mathrm{dR}\subset\Omega^*_{X_\mathbb C/\mathbb C}(\log D_\mathbb C)$ be the subcomplex constructed as above Corollary \ref{intro tau<p structure sheaf}.
\begin{theorem}
Notation and assumptions as above. 

$(1)$ For any $i\geq0$, we have
$$
\dim_\mathbb C\mathbb H^i(X_\mathbb C-D_\mathbb C,(K^\bullet_\mathrm{dR},d+df\wedge))\cong \dim_\mathbb C\mathbb H^i(X_\mathbb C-D_\mathbb C,(K^\bullet_\mathrm{dR},-df\wedge)).$$

$(2)$ For any $i>\dim X_\mathbb C$, we have
$$
H^j(X_\mathbb C,K^i_\mathrm{dR})=0,\quad i+j>\dim X_\mathbb C.
$$
\end{theorem}
The first statement is a special case of the main theorem of Sabbah \cite{Sah1}, whose mod~$p$ proof reduces to Corollary~\ref{intro tau<p structure sheaf}. In the case where $D_\mathbb{C} = \emptyset$ and $K^*_{\mathrm{dR}} = \Omega^*_{X_\mathbb{C}/\mathbb{C}}$, the result was proved in \cite{OV}. The second statement is a special case of the Kodaira-Saito vanishing theorem \cite{Saito}, whose mod $p$ proof essentially reduces to the decomposition theorem in \cite[Appendix D]{ESY}.














\iffalse
There have may studies on smooth projective varieties $X/k$ with trivial logarithmic cotangent bundle $\Omega^1_{X/k}(\log D)$, including \cite{MS, Li, J, ACH}. To the best of my knowledge, the positive characteristic analogue of Beauville's conjecture has not yet been explored.

\subsubsection{De Rham-Higgs comparison (with truncation)}
If we do not put any splitting condition on $(X,D)/k$, then as suggested by Slogan \ref{slogan}, there should be a truncated version of Theorem \ref{de Rham-Higgs without truncation}. 
\begin{theorem}\label{two-sided truncation 0}
Notation and assumption as in Slogan \ref{slogan} and let $K^*_{\mathrm{dR},\theta},K^*_{\mathrm{Hig},\theta}$ be as in Example \ref{example MHM}. Then
\begin{eqnarray}\label{[0,p-ell-1]}
\tau_{[0,p-\ell-1]}F_*K^*_{\mathrm{dR},\theta}\cong\tau_{[0,p-\ell-1]}K^*_{\mathrm{Hig},\theta}~\mathrm{in}~D(X).
\end{eqnarray}
If $\ell=p-3$, then for any $a\geq0$ we have
\begin{eqnarray}\label{[a,a+1]}
 \tau_{[a,a+1]}F_*K^*_{\mathrm{dR},\theta}\cong\tau_{[a,a+1]}K^*_{\mathrm{Hig},\theta}~\mathrm{in}~D(X). 
\end{eqnarray}
\end{theorem}

The first statement \eqref{[0,p-ell-1]} can deduced using a similar argument as in the proof of \cite[Corollary 3.6]{SZ2}. To derive \eqref{[a,a+1]}, we need new insights presented in this paper. Note that, as suggested by Slogan \ref{slogan} once again, the truncation in \eqref{[a,a+1]} should ideally be \( \tau_{[a,a+2]} \). However, due to the limitations of our technique, we must restrict it to \( \tau_{[a,a+1]} \).

For $(E,\theta)=(\mathcal O_X,0)$, $K^*_{\mathrm{dR},0}=F_*\Omega^*_{X/k}$ and $K^*_{\mathrm{Hig},0}=\bigoplus_{i=0}^\infty\Omega^i_{X/k}[-i]$, as shown by \cite{AS,BL,LM,O24}, the truncation in \eqref{[a,a+1]} can be ideally strengthened to $\tau_{[a,a+p-1]}$. For $K^*_{\mathrm{dR},0}=F_*\Omega^*_{g^{-1}}$ and $K^*_{\mathrm{Hig},0}=\bigoplus_{i=0}^\infty \Omega^i_{g^{-1}}[-i]$, the methods of \cite{BL,LM,O24} may still applicable. Consequently, for any $a\geq0$, there may exist an isomorphism
$$
\tau_{[a,a+p-1]}F_*\Omega^*_{g^{-1}}\cong \bigoplus_{i=a}^{a+p-1} \Omega^i_{g^{-1}}[-i]~\mathrm{in}~D(X).
$$
However, it is easy to see that \( F_*K^*_{\mathrm{dR},0} \) is no longer an abstract Koszul complex, hence the method of \cite{AS} may not directly apply. Even worse, for \( \theta \neq 0 \), neither of the aforementioned methods may be directly applicable.  
 



Let $g$ be as in Example \ref{example MHM} (4) and let $f\in k[g]$. Set $K^*_{\mathrm{dR},0}:=W_i\Omega^*_{X/k}(\log D)$ (resp. $\Omega^*_{g^{-1}}$) and $K^*_{\mathrm{Hig},0}:=\bigoplus_{i=0}^\infty W_i\Omega^i_{X/k}(\log D)[-i]$ (resp. $\bigoplus_{i=0}^\infty\Omega^i_{g^{-1}}[-i]$).

\begin{corollary}
 With the notation and assumption above,  for any $a\geq0$ we have
\begin{eqnarray}\label{[0,p-1], theta=0}
\tau_{[0,p-1]}(K^\bullet_{\mathrm{dR},0},d+df\wedge)\cong\tau_{[0,p-1]}(K^\bullet_{\mathrm{Hig},0},-df\wedge)~\mathrm{in}~D(X)
\end{eqnarray}
and
\begin{eqnarray}\label{[a,a+1], theta=0, unique}
\tau_{[a,a+1]}(K^\bullet_{\mathrm{dR},0},d+df\wedge)\cong\tau_{[a,a+1]}(K^\bullet_{\mathrm{Hig},0},-df\wedge)~\mathrm{in}~D(X).
\end{eqnarray}
\end{corollary}
As we mentioned in Example \ref{mod p proof Sabbah} (2), \cite{OV} and \cite{ACH} had established isomorphisms between $(\Omega^*_{X/k},d+df\wedge)$ and $\bigoplus_{i=0}^\infty\Omega^i_{X/k}[-i]$ under different technical conditions, both requiring $\dim X<p$. However, our isomorphism \eqref{[a,a+1], theta=0, unique} holds even for $\dim X\geq p$.

Combine Theorem \ref{de Rham-Higgs local}, \eqref{[0,p-1], theta=0} and the spreading-out technique, we deduce a special case of Theorem \ref{Sabbah's theorem}. Let $X_\mathbb{C}$ be a quasi-projective algebraic variety over $\mathbb C$ and let $f:X_\mathbb{C}\to\mathbb A_\mathbb{C}^1$ be a projective morphism of complex algebraic varieties. Set 
$$E:=[x_1]+\cdots+[x_r]+[y_1]+\cdots+[y_s],~x_1,\cdots,x_r,y_1,\cdots,y_s\in\mathbb C$$ 
be a reduced effective divisor on $\mathbb A^1_\mathbb{C}$ such that $D_\mathbb{C}:=f^*E_\mathbb{C}$ is a simple normal crossing divisor on $X_\mathbb{C}$. Suppose that $f$ has only isolated singularities on $X_\mathbb{C}-D_\mathbb{C}$. Let $K^*_\mathrm{dR}\subset\Omega^*_{X_\mathbb{C}/\mathbb C}(\log D_\mathbb{C})$ be a subcomplex defined as   
$$
K^*_\mathrm{dR}|_{X_\mathbb C-f^*(D_\mathbb C-[x_j])}:=\Omega^*_{(g-x_j)^{-1}},~K^*_\mathrm{dR}|_{X_\mathbb C-f^*(D_\mathbb C-[y_j])}:=W_{i_j}\Omega^*_{X_\mathbb C/\mathbb C}(\log D_\mathbb C)
$$
and $K^*_\mathrm{dR}|_{X_\mathbb C-f^*D_\mathbb C}:=\Omega^*_{X_\mathbb C/\mathbb C}$. Similarly, we can define $K^*_\mathrm{Hig}$.

\begin{theorem}
 With the notation and assumption above, for any $i$ we have
 \begin{eqnarray}\label{E_1 deg Sabbah}
 \dim_\mathbb{C}\mathbb H^i(X,(K^\bullet_\mathrm{dR},d+df\wedge))=\dim_\mathbb{C}\mathbb H^i(X,(K^\bullet_\mathrm{Hig},-df\wedge))<\infty.
 \end{eqnarray}
\end{theorem}
\fi
\iffalse
\subsubsection{Gerbes of splittings and liftings}
Let $X$ be a smooth variety over $k$ which is not necessarily $W_2(k)$-liftable. In \cite[\S3]{DI}, Deligne and Illusie constructed two $T_{X/k}$-gerbes: the first one is the gerbe of splittings of $\mathcal  K_1^*:=\tau_{\leq 1}F_*\Omega^*_{X/k}$, denoted by $\mathrm{sc}(\mathcal K_0^*)$; the second one is the gerbe of liftings of $X$ over $W_2(k)$, denoted by $\mathrm{Rel}(X,W_2(k))$. By \cite[Proposition 3.3, Théorème 3.5]{DI}, there is a natural equivalence between these gerbes and their classes are related to the distinguished triangle 
$$
\mathcal O_X\to\mathcal K^*_1\to\Omega^1_{X/k}[-1]\xrightarrow{\delta^1_{\mathcal K^*}[-1]}\mathcal O_X[1]
$$
 as follows:
$$
\mathrm{cl}~\mathrm{sc}(\mathcal K^*_1)=\mathrm{cl}~\mathrm{Rel}(X,W_2(k))=-\delta^1_{\mathcal K^*_0}\in\mathrm{Ext}^2_{\mathcal O_X}(\Omega^1_{X/k},\mathcal O_X)\cong H^2(X,T_{X/k}).
$$
In \cite[\S3]{AS}, Achinger and Suh considered the gerbe of splittings of $\mathcal K^*_a:=\tau_{[a-1,a]}F_*\Omega^*_{X/k}$ for $a>1$. In \cite[Theorem 3.9]{AS}, they showed that there is a natural functor $\wedge^a$ form $\mathrm{sc}(\mathcal K^*_1)$ to  $\mathrm{sc}(\mathcal K^*_a)$ and their classes are related by 
$$
\mathrm{cl}~\mathrm{sc}(\mathcal K^*_a)=\mathrm{ctr}^q(\mathrm{cl}~\mathrm{sc}(\mathcal K^*_1)).
$$

Let $g:X\to\mathbb A^1_k$ be a semistable family along a reduced effective divisor $E$ on $\mathbb A^1_k$. Then the above discussion can be extended to the pair $(X,g)$ (in some literature, this is referred to as a \emph{Landau-Ginzburg model}).  

\begin{theorem}
 There is a suitably defined $\mathcal Hom_{\mathcal O_X}(\Omega^1_{g^{-1}},g\mathcal O_X)$-gerbe of liftings of $(X,g)$ over $W_2(k)$, denoted by $\mathrm{Rel}((X,g),W_2(k))$. For any $a>0$, set $\mathcal G^*_a:=\tau_{[a-1,a]}F_*\Omega^*_{g^{-1}}$. Then
 $$
 \mathrm{cl}~\mathrm{sc}(\mathcal G^*_1)=\mathrm{cl}~\mathrm{Rel}((X,g),W_2(k))=-\delta^1_{\mathcal G^*_0}\in\mathrm{Ext}^2_{\mathcal O_X}(\Omega^1_{g^{-1}},g\mathcal O_X)
 $$
 and 
 $$
 \mathrm{cl}~\mathrm{sc}(\mathcal G^*_a)=\mathrm{ctr}^q(\mathrm{cl}~\mathrm{sc}(\mathcal G^*_1)).
 $$
\end{theorem}
\fi
















\iffalse
Keep the notation and assumption in Conjecture \ref{conjecture}. Although it is generally quite difficult to construct the desired isomorphisms \eqref{Con2}-\eqref{Con3}, the following theorem demonstrates that they always exist in the formal neighborhood of any closed point of $X$.
\begin{theorem}\label{main A}
For any closed point $x\in X$, we have
$$F_*(K^\bullet_{dR},\nabla+df\wedge)\otimes_{\mathcal O_X}\hat{\mathcal O}_{X,x}\cong(K^\bullet_{Hig},\theta-df\wedge)\otimes_{\mathcal O_X}\hat{\mathcal O}_{X,x}~\mathrm{in}~D(\hat{\mathcal O}_{X,x}).$$
\end{theorem}
\iffalse
Assume $X$ admits a system of coordinates $t_1,\cdots,t_n$ such that $D=(t_1\cdots t_n=0)$. Suppose the absolute Frobenius $F$ has a global lifting $\tilde F:\tilde X\to\tilde X$ such that $\tilde F^*(\tilde t_i)=\tilde t_i^p$ for any $1\leq i\leq n$. Here $\tilde t_i$ is a lifting of $t_i$ and $\tilde D=(\tilde t_1\cdots\tilde t_n=0)$. Let 
$$
\{1,\cdots,n\}=\cup_{i=1}^rP_i,~P_i\neq\emptyset,~P_i\cap P_j=\emptyset~\mathrm{for}~i\neq j.
$$
Set
$$
f:=\sum_{i=1}^rf_i,~f_i:=g_i^p\prod_{j\in P_i}t_j^{a_j},~g_i\in\Gamma(X,\mathcal O_X),~a_j\in\mathbb N.
$$

Using a similar argument as in the proof of the theorem above, we generalize the classical Cartier isomorphism (see e.g. \cite[Théorème 1.2]{DI}) and its variant in \cite[Proposition 3.8]{SZ2} to the following
\begin{corollary}[Twisted Cartier isomorphism]
With the above assumptions, there exists an isomorphism
$$F_*(K^\bullet_{dR},\nabla+df\wedge)\cong(K^\bullet_{Hig},\theta-df\wedge)~\mathrm{in}~D(X).$$
\end{corollary}
Combing the twisted Cartier isomorphism above with the $\infty$-homotopy $\mathrm{Ho}$ \cite[Theorem 3.5]{SZ2}, we have
\fi
Using a similar argument as in the proof of the theorem above, coupled with the $\infty$-homotopy $\mathrm{Ho}$ \cite[theorem 3.5]{SZ2}, gives rise to the following
\begin{theorem}\label{Main B}
Suppose 
$$X=X_1\times X_2,~\tilde X=\tilde X_1\times\tilde X_2,~f=\pi^*_1f_1,~D=\pi^*_2D_2,~\tilde D=\tilde\pi^*_2\tilde D_2,$$
where $X_1,X_2$ are  respectively varieties of dimension $n_1,n_2$ over $k$, $\tilde X_1,\tilde X_2$ are respectively $W_2(k)$-liftings of $X_1,X_2$, $f_1\in\Gamma(X_1,\mathcal O_{X_1})$ has only isolated singularities, $D_2\subset X_2$ is a SNCD, $\pi_i:X\to X_i$ (resp. $\tilde X\to\tilde X_i$) is the natural projection for $i=1,2$ and $\tilde D_2$ is a lifting of $D_2$ on $\tilde X_2$. Then:
\begin{itemize}
    \item[$(a)$\ ] we have the desired isomorphism \eqref{Con2};
    \item[$(b)$\ ] for any integer $a$, we have
$$\tau_{[a,a+1]}F_*(K^\bullet_{dR},\nabla+df\wedge)\cong \tau_{[a,a+1]}(K^\bullet_{Hig},\theta-df\wedge)~\mathrm{in}~D(X);$$
\item[$(c)$\ ] if the projective bundle $\mathbb P\Omega^1_{X_2/k}(\log D_2)$ has sections $s_1,\cdots,s_{n_2+1}$ such that for any closed point $x\in X_2$, $s_1(x),\cdots,s_{n_2}(x)$ span $\mathbb P\Omega^1_{X_2/k}(x)$, then
$$
F_*(K^\bullet_{dR},\nabla+df\wedge)\cong (K^\bullet_{Hig},\theta-df\wedge)~\mathrm{in}~D(X).
$$
\end{itemize}
\end{theorem}
Note that $(b)$ (resp. $(c)$) is weaker (resp. stronger) than the one suggested by Conjecture \ref{conjecture} (C2). Though the condition presented in $(c)$ is somehow technical, it is satisfied when $D_2=\emptyset$ and $X_2$ is either a product of curves or an Abelian variety.
As an immediate consequence of this theorem, one has
\begin{corollary}\label{Main C}
Notation and assumption as in Conjecture \ref{conjecture}.
 Then

    $$\tau_{<p-\ell}F_*K^*_{dR}\cong\tau_{<p-\ell}K^*_{Hig}~\mathrm{in}~D(X).$$
If we further assume 
\begin{itemize}
    \item $D=\emptyset$ and $f$ has only isolated singularities, or
    \item $D=\emptyset, f=0$ and $X$ is  either a product of curves or an Abelian variety, 
\end{itemize}
 then
    $$
    F_*\Omega^*(H,\nabla+df\wedge)\cong\Omega^*(E,\theta-df\wedge)~\mathrm{in}~D(X).
    $$
\end{corollary}

\iffalse
On the other hand, if we take $(E,\theta)=(\mathcal O_X,0)$,  suppose that the $W_2(k)$-morphism $\tilde f:\tilde X\to\mathbb A^1_{W_2(k)}$ is a semistable reduction along a relative divisor $\tilde E=\sum_{i=1}^s[\tilde\lambda_i]$ on $\mathbb A^1_{W_2(k)}/W_2(k)$ ($\tilde\lambda_i\in W_2(k)$), and set $\tilde D:=\tilde f^*\tilde E$,   then Conjecture \ref{conjecture} holds. 
\begin{theorem}\label{Main C}
 Assumption as above. Then  for any $0\leq a\leq n-1$ and any polynomial $P(T)\in k[T]$, we have the following isomorphisms in $D(X)$:
 $$
 \begin{array}{c}
 \tau_{<p-\ell}F_*K^*_{dR}\cong\tau_{<p-\ell}K^*_{Hig};\\
\tau_{[a,a+1]}F_*(K^\bullet_{dR},d+dP(f)\wedge)\cong\tau_{[a,a+1]}(K^\bullet_{Hig},-dP(f)\wedge).
\end{array}
$$
Here $K^*_{dR}=W_i\Omega^*_{\log}(\mathcal O_X,d)$ or $\Omega^*_f(\mathcal O_X,d)$ and  $K^*_{Hig}=W_i\Omega^*_{\log}(\mathcal O_X,0)$ or $\Omega^*_f(\mathcal O_X,0)$.
\end{theorem}\fi
\begin{remark}
$\mathrm{(a)}$ Using the fact that $F_*\Omega^*_{X/k}$ is an abstract Koszul complex, Achinger-Suh\cite{AS} established the formality of the two-sided truncations $\tau_{[a,b]}F_*\Omega^*_{X/k}$ for $a\leq b<a+p-1$. Drinfeld observed that there is a $\mu_p$-action on $\Omega^*_{X/k}$, its proof was given by Bhatt-Lurie\cite{BL} and Li-Mondal\cite{LM} by different approaches. As an immediate consequence of Drinfeld's observation, the formality of Achinger-Suh can be strengthened to $a\leq b\leq a+p-1$, which in turn is a generation of the formality of Deligne-Illusie\cite{DI}.

Unfortunately, neither $F_*W_i\Omega^*_{\log}(\mathcal O_X,d)$ nor $F_*\Omega^*_f(\mathcal O_X,d)$ is an abstract Koszul complex, hence the approach of Achinger-Suh may not easy to show the formality of
$$\tau_{[a,b]}F_*W_i\Omega^*_{\log}(\mathcal O_X,d),~\tau_{[a,b]}F_*\Omega^*_{g^{-1}}(\mathcal O_X,d)$$
for $a\leq b<a+p-1$. On the other hand, the approach of Bhatt-Lurie or Li-Mondal may available to the aforementioned formality. 

$\mathrm{(b)}$ As discussed in \S1.1, Corollary \ref{Main C} can be used to provide a mod $p$ proof for Theorem \ref{Sabbah's theorem} in certain special cases. Moreover, it can also be used to provide a mod $p$ proof for the Kodaira-Saito vanishing theorem in certain special cases again. The argument is standard: see, for example, \cite[\S4.7]{OV} or \cite[\S5]{SZ2}.
\iffalse
We prove Theorem \ref{Main B} following the spirit of Deligne-Illusie\cite{DI}. This approach allows us to obtain the second isomorphism in loc.cit., to the knowledge of the author, that result is new. However, our approach can not show
$$
\tau_{[a,a+p-1]}F_*\Omega^*(\mathcal O_X,d+dP(f)\wedge)\cong\tau_{[a,a+p-1]}\Omega^*(\mathcal O_X,-dP(f)\wedge),
$$
which is suggested by Conjecture \ref{conjecture}. It is interesting to use the approach of Bhatt-Lurie or Li-Mondal to achieve the above isomorphism.\fi
\end{remark}
\fi
\subsection{Organization}
In \S2, we fix notation, recall some basic notions, and introduce several complexes that will be considered throughout the paper, motivated by the theory of mixed Hodge modules. In \S3, we prove that Question~\ref{Question 1} holds in the formal neighborhood of any closed point of~$X'$. In \S4, we establish a relation between splittings of the logarithmic cotangent bundle $\Omega^1_{X/k}(\log D)$ and isomorphisms in $D(X')$ between suitable truncations of $\Omega^*(E,\theta)$ and $F_*\Omega^*(H,\nabla)$. We then show that these isomorphisms are independent of the choice of splittings. Furthermore, we verify that all constructions are compatible with the complexes introduced in \S2.6, under suitable technical conditions. Finally, we establish a de Rha–Higgs comparison under two-term truncation. In \S5, we present several applications, including the K\"unneth formula, the lifting of the constructions in \S4 mod $p^n$, and the $E_1$-degeneration and vanishing theorem.







\iffalse
Section 1  is devoted to the proof of Theorems \ref{de Rham-Higgs local}. The key ingredient is described as follows. Take a closed point $x\in X$ and pick a \emph{nice} Frobenius lifting $\tilde F$ around $x$. Suppose that $f(x)=0$. Set $\mathfrak X:=\mathrm{Spf}(\hat{\mathcal O}_{X,x})$ and let $(\hat E,\hat\theta),(\hat H,\hat\nabla)$ be the pull-backs of $(E,\theta),(H,\nabla)$ on $\mathfrak X$, respectively. Note that 
$$
(\hat E\otimes\Omega^\bullet_{\mathfrak X/k}(\log\mathfrak D),\hat\theta-df\wedge)\cong\Omega^*(\hat H,\hat\theta-df\wedge),~(\hat H\otimes\Omega^\bullet_{\mathfrak X/k}(\log\mathfrak D),\hat\nabla+df\wedge)\cong\Omega^*(\hat H,\hat\theta+df\wedge).
$$
In general, the inverse Cartier transform   $C^{-1}_{\tilde F}(\hat E,\hat\theta-df\wedge)$ is NOT isomorphism to $(\hat H,\hat\nabla+df\wedge)$. Fortunately, we find that there is an infinitesimal permutation $\Theta$ of $\hat\theta-df\wedge$ such that (i) $\Omega^*(\hat E,\hat\theta-df\wedge)$ is isomorphic to $\Omega^*(\hat E,\Theta)$, and (ii) $(\hat H,\hat\nabla+df\wedge)\cong C^{-1}_{\tilde F}(\hat E,\Theta)$. From which we get an isomorphism between $\Omega^*(\hat E,\hat\theta-df\wedge)$ and $\Omega^*(\hat H,\hat\nabla+df\wedge)$. Moreover, this isomorphism is compatible with all kinds of complexes considered in Example \ref{example MHM}.

Section 2 is devoted to the proofs of Theorems \ref{de Rham-Higgs without truncation} and \ref{two-sided truncation 0}. The key input in the proof of Theorem \ref{de Rham-Higgs without truncation} is the construction of an $\mathcal L$-indexed $\infty$-homotopy
$$
\mathrm{Ho}_{\underline\Omega}:\Delta_*(\mathcal L)\to\mathcal Hom^*_{\mathcal O_X}(\Omega^*(E,\theta),F_*\Omega^*(H,\nabla)),
$$
which generalizes the construction in \cite{SZ1}. Here, $\mathcal L$ is the sheaf of Frobenius liftings over a given  $(\tilde X,\tilde D)/W_2(k)$, $\Delta_*(\mathcal L)$ is its $\Delta$-complex, and $\underline\Omega:=(\Omega_1,\cdots,\Omega_\beta)$ arises from a splitting $\Omega_{X/k}^1(\log D)=\bigoplus_{i=1}^\beta\Omega_i$ with $\mathrm{rank}(\Omega_i)<p-\ell$. If $\mathrm{Ho}_{\underline\Omega'}$ is another $\mathcal L$-indexed $\infty$-homotopy, it is natural to ask whether they are homotopic? We confirm it in a special case, which is the key ingredient in the proof of Theorem \ref{two-sided truncation 0}.

The final section is devoted to the geometric applications of the preceding two sections, including the gerbes of splittings and liftings, the K\"unneth formula, and the $E_1$-degeneration and vanishing theorem.
\fi
\section{Preliminaries}

\subsection{Convention}
 Throughout this paper, let $p$ denote a positive prime number and let $k$ denote a perfect field of characteristic $p$. Denote by $X$ a smooth variety over $k$, $D$ a simple normal crossing divisor on $X$, and $(\tilde X,\tilde D)/W_2(k)$ a flat lifting of $(X,D)/k$. Let $\sigma:\mathrm{Spec}(k)\to\mathrm{Spec}(k)$ be the morphism induced by the Frobenius endomorphism of $k$. For any open subset \( U \subset X \), denote by \( \widetilde{U} \) the corresponding open subset of \( \widetilde{X} \) whose reduction modulo \( p \) is \( U \). Set
\[
X' := X \times_{\mathrm{Spec}(k), \sigma} \mathrm{Spec}(k), \quad 
\widetilde{X}' := \widetilde{X} \times_{\mathrm{Spec}\, W_2(k), \sigma} \mathrm{Spec}\, W_2(k).
\]
More generally, if \( \square \) is an object over \( k \), we write
\[
\square' := \square \times_{\mathrm{Spec}(k), \sigma} \mathrm{Spec}(k),
\]
and if \( \square \) is an object over \( W_2(k) \), we set
\[
\square' := \square \times_{\mathrm{Spec}\, W_2(k), \sigma} \mathrm{Spec}\, W_2(k).
\]


 We use the subscript $*$ to denote a complex and the superscript $\bullet$ to denote a graded module. That is, $\mathcal{F}^*$ refers to a complex, while $\mathcal{F}^\bullet$ denotes a graded module of the form $\bigoplus_{i = -\infty}^{\infty} \mathcal{F}^i$.

\subsection{Truncation}

Let \( \mathcal{F}^* \) be a complex of \( \mathcal{O}_X \)-modules on \( X \), and let \( a < b \) be integers.  
We denote the truncation \( \tau_{<b} \mathcal{F}^* \) by
\[
\cdots \longrightarrow \mathcal{F}^{b-3} \longrightarrow \mathcal{F}^{b-2} \longrightarrow \ker\left( \mathcal{F}^{b-1} \to \mathcal{F}^b \right) \longrightarrow 0,
\]
and the two-sided truncation \( \tau_{[a,b]} \mathcal{F}^* \) by
\[
0 \longrightarrow \frac{\mathcal{F}^a}{\mathrm{im}(\mathcal{F}^{a-1} \to \mathcal{F}^a)} 
\longrightarrow \mathcal{F}^{a+1} \longrightarrow \cdots \longrightarrow \mathcal{F}^{b-1} 
\longrightarrow \ker(\mathcal{F}^b \to \mathcal{F}^{b+1}) \longrightarrow 0.
\]


\subsection{$\lambda$-connections}
Let $S$ be a Noetherian scheme, and let $Y$ be a smooth $S$-scheme of pure relative dimension $d$ equipped with a simple normal crossing divisor $D_Y$ relative to $S$. Let $\lambda \in \Gamma(S, \mathcal{O}_S)$, and let $M$ be an $\mathcal{O}_Y$-module.
A \emph{$\lambda$-connection} on $M$ with pole along $D_Y$ is an $\mathcal{O}_S$-linear morphism
\[
\nabla \colon M \to M \otimes_{\mathcal{O}_Y} \Omega^1_{Y/S}(\log D_Y)
\]
satisfying the $\lambda$-Leibniz rule:
\[
\nabla(fm) = f \nabla(m) + m \otimes \lambda \, df, \quad \text{for all } f \in \mathcal{O}_Y,\, m \in M.
\]
This $\lambda$-connection associates to a $\lambda$-complex
$$
\Omega^*(M,\nabla):=[M\xrightarrow{\nabla}M\otimes\Omega^1_{Y/S}(\log D_Y)\xrightarrow{\nabla}\cdots\xrightarrow{\nabla}M\otimes\Omega^d_{Y/S}(\log D_Y)].
$$
It is customary to refer to $1$-complexes as de Rham complexes and to $0$-complexes as Higgs complexes.


We simply refer to the pair $(M, \nabla)$ as a $\lambda$-connection on $(Y, D_Y)/S$, and denote the corresponding category by $\lambda\mathrm{MIC}((Y, D_Y)/S)$. In accordance with standard convention, we denote the category corresponding to $\lambda = 0$ by $\mathrm{HIG}((Y, D_Y)/S)$, whose objects are called Higgs sheaves.

In this paper, we are interested in three special cases of $\lambda$-connections: namely, $1$-connections (i.e., integrable connections), $0$-connections (i.e., Higgs fields), and $p$-connections in the case where $S = \mathrm{Spec}\, W(k)$.


\subsection{Nilpotence}


 We say a Higgs sheaf $(E,\theta)$ on $(X,D)/k$ is nilpotent of level $\leq\ell$ if
$$
\prod_{i=1}^{\ell+1}(\mathrm{id}\otimes\partial_i)\theta=0\in\mathcal Hom_{\mathcal O_X}(E),~\partial_1,\cdots,\partial_{\ell+1}\in T_{X/k}(-\log D).
$$
Denote by $\mathrm{HIG}_\ell((X,D)/k)$ the category of all such nilpotent Higgs sheaves. 

\subsection{Inverse Cartier transform}
Let $C^{-1}_{(\tilde X,\tilde D)/W_2(k)}$ be the inverse Cartier transform associated to $(\tilde X,\tilde D)/W_2(k)$, which is an equivalence of categories from $\mathrm{HIG}_\ell((X,D)/k)$ to $\mathrm{MIC}_\ell((X,D)/k)$. For more details, see \cite{OV, Schepler, LSYZ, LSZ}.

For simplicity, we abbreviate \( C^{-1}_{(\widetilde{X}, \widetilde{D})/W_2(k)} \) by \( C^{-1} \). There are two constructions of the inverse Cartier transform \( C^{-1} \): one due to \cite{OV, Schepler}, and another given in \cite{Fa, LSYZ, LSZ}. In this paper, we adopt the latter construction, which we briefly recall below.

Let $\{U_i\}$ be an open affine covering $X$, and let $\tilde F_i:\tilde U_i\to\tilde U'_i$ be a Frobenius lifting of the relative Frobenius $F: U_i\to U_i'$, such that $\tilde F_i^*(\tilde D'|_{\tilde U_i'})=p\tilde D|_{\tilde U_i}$. These data give rise to the following two families of morphisms:
$$
\zeta_i=\frac{d\tilde F^*_i}{p}:F^*\Omega^1_{U'_i/k}\to \Omega^1_{U_i/k},\quad h_{ij}=\frac{\tilde F^*_j-\tilde F^*_i}{p}:F^*\Omega^1_{U'_i\cap U'_j/k}\to\Omega^1_{U_i\cap U_j/k}.
$$

Given any $(E,\theta)\in\mathrm{HIG}_{p-1}((X',D')/k)$. Then $(H,\nabla)$ is obtained by gluing the local integrable connections
$$
(H_i:=(F^*E)|_{U_i},(\mathrm{id}\otimes\zeta_{\tilde F_i})F^*\theta)
$$
via the transition isomorphisms
$$
G_{ij}=\mathrm{exp}((\mathrm{id}\otimes h_{ij})F^*\theta):H_j|_{U_i\cap U_j}\to H_i|_{U_i\cap U_j}.$$



\subsection{Hodge pairs of complexes in positive characteristic}

Let $K^*_\mathrm{Hig}$ be a complex of $\mathcal O_{X'}$-modules on $X'$, and let
$K^*_\mathrm{dR}$ be a complex of sheaves of $k$-vector spaces on $X$. We say the pair  $(K^*_\mathrm{Hig},K^*_\mathrm{dR})$ is a \emph{Hodge pair of complexes on $(X,D)/k$}, if it belongs to one of the following situations.


\emph{Hodge pair of complexes of type $(\mathrm{\uppercase\expandafter{\romannumeral1}},\ell)$}.
Given any $0\leq\ell\leq p-1$. Let $(E,\theta)\in\mathrm{HIG}_\ell((X,D)/k)$ and let $(H,\nabla)$ be its inverse Cartier transform.  Set either 
$$
K^*_{\mathrm{Hig}}:=\Omega^*(E,\theta),\quad
K^*_{\mathrm{dR}}:=\Omega^*(H,\nabla) $$
or
$$
K^*_{\mathrm{Hig}}:=W_i\Omega^*(E,\theta),\quad
K^*_{\mathrm{dR}}:=W_i\Omega^*(H,\nabla).
$$
In the later case, we additionally suppose that $(E,\theta)$ and $(H,\nabla)$ have no pole along $D$, namely
    $$
    \nabla(H)\subset H\otimes\Omega^1_{X/k},~\theta(E)\subset E\otimes\Omega^1_{X/k}.
    $$
 The notation $W_i$ means the $i$-th weight filtration (see e.g. \cite[Definition 8.3.19]{CEGT}), which is given by
 $$
 \begin{array}{c}
 W_i\Omega^q(E,\theta):=\left\{\begin{matrix}E\otimes\mathrm{im}(\Omega_{X/k}^i(\log D)\otimes\Omega^{q-i}_{X/k}\to\Omega_{X/k}^q(\log D)),&i\leq q,\\
E\otimes\Omega_{X/k}^q(\log D),&i>q
 \end{matrix}\right.\\
 \end{array}
 $$
 and
 $$
 \begin{array}{c}
 W_i\Omega^q(H,\nabla):=\left\{\begin{matrix}H\otimes\mathrm{im}(\Omega_{X/k}^i(\log D)\otimes\Omega^{q-i}_{X/k}\to\Omega_{X/k}^q(\log D)),&i\leq q,\\
H\otimes\Omega_{X/k}^q(\log D),&i>q.
 \end{matrix}\right.\\
 \end{array} $$

\emph{Hodge pair of complexes of type $(\mathrm{\uppercase\expandafter{\romannumeral2}},\ell)$}. Given any $0\leq\ell\leq p-1$. Let $(E,\theta)\in\mathrm{HIG}_\ell((X,D)/k)$ and let $(H,\nabla)$ be its inverse Cartier transform. Set
$$
K^*_{\mathrm{Hig}}:=\Omega_\mathrm{int}^*(E,\theta),\quad
K^*_{\mathrm{dR}}:=\Omega_\mathrm{int}^*(H,\nabla),
$$ 
Here, the subscript “int” refers to the intersection subcomplexes (see \cite[Definition 8.3.31]{CEGT} and \cite[\S2]{SZ2}). For completeness, we briefly recall the definition below. Without loss of generality, we may assume that $X/k$ admits a system of coordinates 
Let $t_1,\cdots,t_n$ and that $D:=(t_1\cdots t_n=0)$. Then
 $$
 \Omega^\bullet_\mathrm{int}(H,\nabla):=\bigoplus_{I\subset\{1,\cdots,n\}}(\sum_{J\subset I}t_J\nabla_{I-J}H)\otimes d_I\log t.
 $$
 Here $t_J:=\prod_{j\in J}t_j$ for $J\neq\emptyset$ and $t_\emptyset:=1$,  $\nabla_I:=\prod_{i\in I}\nabla_i$ for $I\neq\emptyset$ and $\nabla_\emptyset:=id$ ($\nabla=\sum_{i=0}^n\nabla_i\otimes d\log t_i$), and $d_I\log t:=d\log t_{i_1}\cdots d\log t_{i_r}$ for $I=\{i_1,\cdots,i_r\},i_1<\cdots<i_r$ and $d_\emptyset\log t:=1$. Similarly, one can define the intersection de Rham subcomplex $\Omega^*_\mathrm{int}(E,\theta)\subset\Omega^*(E,\theta)$.

\emph{Hodge pair of complexes of type $(\mathrm{\uppercase\expandafter{\romannumeral3}},\ell)$}. Given any $0\leq\ell\leq p-1$. Let $(E,\theta)\in\mathrm{HIG}_\ell((X,D)/k)$ and let $(H,\nabla)$ be its inverse Cartier transform.  Let $\tilde g:(\tilde X,\tilde D)\to(\mathbb A_{W_2(k)}^1,\tilde \Delta)$ be a semistble family over $W_2(k)$, which means that $\tilde g^*\tilde\Delta=\tilde D$ and $\tilde g$ is smooth outside $\tilde D$. Suppose that $\tilde\Delta$ containing $0$. Let $g:(X,D)\to(\mathbb A^1_k,\Delta)$ be the mod $p$ reduction of $\tilde g$. Set
$$
K^*_{\mathrm{Hig}}:=\Omega_{g'^{-1}}^*(E,\theta),\quad
K^*_{\mathrm{dR}}:=\Omega_{g^{-1}}^*(H,\nabla),
$$
which are Kontsevich subcomplexes of $\Omega^*(E,\theta)$ and $\Omega^*(H,\nabla)$, respectively. See \cite[Definition 2.11]{KKP} and \cite[\S1.3]{ESY} for more details. To facilitate later arguments, we give a more explicit description of the Kontsevich subcomplexes introduced above.
 Suppose that $X/k$ admits a system of coordinates $t_1,\cdots,t_n$ such that $g=t_1\cdots t_n$ and $D=(g=0)$. Then
$$
\Omega^i_{g^{-1}}(H,\nabla):=H\otimes\Omega^i_{g^{-1}},
$$
where
$$\Omega^i_{g^{-1}}:=\bigoplus_{I\subset\{2,\cdots,n\},|I|=i-1}\mathcal O_Xd\log g\wedge d_I\log t\oplus\bigoplus_{I\subset\{1,\cdots,n\},|I|=i}g\mathcal O_Xd_I\log t\subset\Omega^1_{X/k}(\log D).
$$
Similarly, one can define $\Omega_{g'^{-1}}^i(E,\theta)$.

\emph{Hodge pair of complexes of type $(\mathrm{\uppercase\expandafter{\romannumeral4}},0)$}. The last situation is a special case combining the first and the preceding ones. Let $\tilde{g}: (\tilde{X},\tilde D) \to (\mathbb{A}^1_{W_2(k)},\tilde\Delta)$ be a semistable family over $W_2(k)$ such that $\tilde\Delta=\sum_{i=1}^m[\tilde \lambda_i],~\tilde\lambda_i\in W_2(k)$. Let $g,\Delta,\lambda_i$ be the mod reductions of $\tilde g,\tilde\Delta,\tilde\lambda_i$, respectively. 

Set $K^*_{\mathrm{dR}} \subset \Omega^*_{X/k}(\log D)$ to be the subcomplex defined as follows: on $X-D$, define $K^*_{\mathrm{dR}} := \Omega^*_{X/k}$; around a fiber $g^{-1}(\lambda_i)$, define $K^*_{\mathrm{dR}} := W_j\Omega^*_{X/k}(\log D)$ or $\Omega^*_{(g - \lambda_i)^{-1}}$. Define $K^*_{\mathrm{Hig}}$ similarly.



\subsection{Exponential twisting of the Hodge pairs}
Inspired by Sabbah's work on twisted de Rham complexes (see \cite{Sah1}), we consider the exponential twisting of the Hodge pairs of complexes introduced in the previous subsection. Let $f\in\Gamma(X,\mathcal O_X)$ and let $(K^*_\mathrm{Hig},K^*_\mathrm{dR})$ be a Hodge pair of complexes. Define the $f$-exponential twisting of $(K^*_\mathrm{Hig},K^*_\mathrm{dR})$ by
$$
(e^{f'}K^*_\mathrm{Hig}:=(K^\bullet_\mathrm{Hig},\theta-df'\wedge),e^fK^*_\mathrm{dR}:=(K^\bullet_\mathrm{dR},\nabla+df\wedge)).
$$

To facilitate later arguments, we need the following observations:
$$
\begin{array}{cc}
     e^{f'}W_i\Omega^*(E,\theta)=W_i\Omega^*(E,\theta-df'\wedge),
     & e^fW_i\Omega^*(H,\nabla)=W_i\Omega^*(H,\nabla+df\wedge),\\
e^{f'}\Omega^*_\mathrm{int}(E,\theta)=\Omega^*_\mathrm{int}(E,\theta-df'\wedge),&e^f\Omega^*_\mathrm{int}(H,\nabla)=\Omega^*_\mathrm{int}(H,\nabla+df\wedge), \\
e^{f'}\Omega^*_{g'^{-1}}(E,\theta)=\Omega^*_{g'^{-1}}(E,\theta-df'\wedge),&e^f\Omega^*_{g^{-1}}(H,\nabla)=\Omega^*_{g^{-1}}(H,\nabla+df\wedge). 

\end{array}
$$









\iffalse
Let $X$ be a smooth projective algebraic variety over $\mathbb C$ endowed with a simple normal crossing divisor $D$. Let $f:X\to\mathbb P^1$ be a dominant  morphism of complex varieties such that $f^*[\infty]\subset D$. Let $\Omega^*_f$ be the Kontsevich complex associated to $(X,D,f)$ (see \cite[(1.3.1)]{ESY} or Definition \ref{Kontsevich complex}). Set $U:=X-D$. In \cite{ESY}, Esnault-Saito-Yu confirmed the following
\begin{conjecture}[Kontsevich]
For any $k\geq0$, the dimension of 
the hypercohomology $\mathbb H^k(X,(\Omega^\bullet_f,ud+vdf))$ is independent of $u,v\in\mathbb C$ and is equal to $\mathbb H^k(U,(\Omega^\bullet,d+df))$. In particular, the spectral sequence 
\begin{eqnarray}\label{E_1 Kontsevich}
E_1^{p,q}=H^q(X,\Omega^p_f)\Rightarrow\mathbb H^{p+q}(X,\Omega^*_f)
\end{eqnarray}
degenerates at $E_1$.
\end{conjecture}
\eqref{E_1 Kontsevich} was proved by two methods. The first one was suggested by Kontsevich and written down by Esnault in \cite[Appendix D]{ESY}. The second one was given by Saito in \cite[Appendix D]{ESY}. 
As a consequence of the conjecture, we have
\begin{eqnarray}\label{equality Kontsevich}
\dim_\mathbb C\mathbb H^k(U,(\Omega_U^\bullet,d+df))=\dim_\mathbb C\mathbb H^k(U,(\Omega_U^\bullet,-df)),~\forall k.
\end{eqnarray}
Using the nonabelian Hodge theory in positive characteristic, Ogus-Vologodsky gave another proof of the equality, see \cite[\S 4.7]{OV} for more details. Among other things, Li-Zhang established a rigid analogue of the equality in \cite{LZ}.

\cite[Appendix E]{ESY} showed that the Kontsevich complex $\Omega^*_f$ occurred in \eqref{E_1 Kontsevich}  is closely related to the nearby cycle complex $\psi_g\mathbb C_{X^\mathrm{an}}$ for $g=f^{-1}$. On the other hand, Sabbah \cite{Sah1} showed that the twisted de Rham complex $(\Omega^\bullet_U,d+df)$ occurred in \eqref{equality Kontsevich} is closely related to the vanishing cycle complexes $\phi_f\mathbb C_{X^\mathrm{an}}$.

This paper aims to investigate some positive characteristic analogues of the above discussions: (i) we extend \eqref{E_1 Kontsevich} and \eqref{equality Kontsevich} to twisted coefficients in positive characteristic; (ii) we define the vanishing/nearby cycle complexes of de Rham complexes in positive characteristic. 
\subsection{Twisted de Rham-Higgs comparison}
Let $k$ be an algebraically closed field of characteristic $p>0$, $X$ a smooth variety over $k$ endowed with a simple normal crossing $D$ and $f\in\Gamma(X,\mathcal O_X)$. Let $F$ be the Frobenius endomorphism of $X$.
Assume $(X,D,f)$ is $W_2(k)$-liftable and fix such one, say $(\tilde X,\tilde D,\tilde f)$. Using the lifting, one gets an inverse Cartier transform $C^{-1}$ from $\mathrm{HIG}_\ell((X,D)/k)$ to $\mathrm{MIC}_\ell((X,D)/k)$. Here we require $\ell<p$.

Let $(E,\theta)\in\mathrm{HIG}_\ell((X,D)/k)$ and set $(H,\nabla):=C^{-1}(E,\theta)$. Let
$\Omega^*(E,\theta)$ (resp. $\Omega^*(H.\nabla)$) be the Higgs complex (resp. de Rham complex) associated to $(E,\theta)$ (resp. $(H,\nabla)$).
Ogus-Vologodsky \cite[Corollary 2.27]{OV} and Schepler \cite[corollary 5.7]{Schepler} established the following de Rham-Higgs comparison:
\begin{eqnarray}\label{OV comparison}
F_*\tau_{<p-\ell}\Omega^*(H,\nabla)\cong\tau_{<p-\ell}\Omega^*(E,\theta)~\mathrm{in}~D(X).
\end{eqnarray}
Inspired by the complex Hodge theory, Sheng and the author \cite[\S2.1]{SZ2} introduced the intersection subcomplexes
$$
\Omega_\mathrm{int}^*(E,\theta)\subset\Omega^*(E,\theta),~\Omega_\mathrm{int}^*(H,\nabla)\subset\Omega^*(H,\nabla).
$$
Following the spirit of Deligne-Illusie \cite{DI}, \cite{SZ2} constructed another isomorphism between $F_*\tau_{p-\ell}\Omega^*(H,\nabla)$ and $\tau_{<p-\ell}\Omega^*(E,\theta)$ in $D(X)$ such that it induces
\begin{eqnarray}\label{SZ comparison}
F_*\tau_{<p-\ell}\Omega_\mathrm{int}^*(H,\nabla)\cong\tau_{<p-\ell}\Omega_\mathrm{int}^*(E,\theta)~\mathrm{in}~D(X).
\end{eqnarray}

The work of Sabbah \cite{Sah1} and Kontsevich's conjecture above enlighten us that \eqref{OV comparison} and \eqref{SZ comparison} may hold after twisting $\nabla$ and $\theta$ by $df$ and $-df$, respectively. This insight is completely confirmed for $\dim X=1$ in \S3.1. For $\dim X>1$, due to some technical problems, we can only confirm this insight for some special cases.  See \S3-\S4 for more details.
\begin{theorem}\label{Main theorem A intro}
Notation and assumption as above. For $\dim X=1$ or $\dim X>1$ and $f$ satisfies some technical condition, one has
\begin{eqnarray}\label{twisted log intro}
F_*\tau_{<p-\ell}\Omega^*(H,\nabla+df\wedge)\cong\tau_{<p-\ell}\Omega^*(E,\theta-df\wedge)~\mathrm{in}~D(X)
\end{eqnarray}
and
\begin{eqnarray}\label{twisted int intro}
F_*\tau_{<p-\ell}\Omega_\mathrm{int}^*(H,\nabla+df\wedge)\cong\tau_{<p-\ell}\Omega_\mathrm{int}^*(E,\theta-df\wedge)~\mathrm{in}~D(X).
\end{eqnarray}
\end{theorem}
When $D=\emptyset$ and $(E,\theta)=(\mathcal O_X,0)$, then \eqref{twisted log intro} was shown by \cite[Theorem 4.23 (2)]{OV} under some technical condition on $\Omega^*(\mathcal O_X,-df\wedge)$.
\subsection{Vanishing cycle de Rham/Higgs complexes in positive characteristic}
Notation and assumption as in \S1.1. Let $(\hat H,\hat\nabla):=(H,\nabla)_{/f^*p[0]}$ and $(\hat E,\hat\theta):=(E,\theta)_{/f^*[0]}$ which are formal completions of $(H,\nabla)$ and $(E,\theta)$ along $f^*p[0]$ and $f^*[0]$, respectively. Let $\mathfrak X:=X_{/f^*p[0]}=X_{/f^*[0]}$. Inspired by \cite[Remark 0.4]{Sah1}, we make the following
\begin{definition}
Define the vanishing cycle complex $\phi_f\Omega^*(H,\nabla)$ of $\Omega^*(H,\nabla)$ along $f$ by
$$
\phi_f\Omega^*(H,\nabla):=\Omega^*(\hat H,\hat\nabla-df\wedge).
$$
Similarly, we can define
$$
\begin{array}{c}
\phi_f\Omega_\mathrm{int}^*(H,\nabla):=\Omega_\mathrm{int}^*(\hat H,\hat\nabla-df\wedge),~\phi_f\Omega^*(E,\theta):=\Omega^*(\hat E,\hat\theta+df\wedge),\\
\phi_f\Omega_\mathrm{int}^*(E,\theta):=\Omega_\mathrm{int}^*(\hat E,\hat\theta+df\wedge).
\end{array}
$$
\end{definition}
As an immediate consequence of Theorem \ref{Main theorem A intro}, one gets the following de Rham-Higgs comparison for vanishing cycle complexes.
\begin{corollary}
Notation and assumption as in \S1.1. For $\dim X=1$ or $\dim X>1$ and $f$ satisfies some technical condition, one has
\begin{eqnarray}\label{twisted log intro}
F_*\tau_{<p-\ell}\phi_f\Omega^*(H,\nabla)\cong\tau_{<p-\ell}\phi_f\Omega^*(E,\theta)~\mathrm{in}~D(\mathfrak X)
\end{eqnarray}
and
\begin{eqnarray}\label{twisted int intro}
F_*\tau_{<p-\ell}\phi_f\Omega_\mathrm{int}^*(H,\nabla)\cong\tau_{<p-\ell}\phi_f\Omega_\mathrm{int}^*(E,\theta)~\mathrm{in}~D(\mathfrak X).
\end{eqnarray}
\end{corollary}
It is natural to study the monodromy operator $T$ of $\phi_f\Omega^*(H,\nabla)$. Unfortunately, we think it may hard to define the monodromy operator $T$ directly, but its logarithm $-\frac{\log T}{2\pi\sqrt{-1}} $ may possibly be defined. This insight is inspired by Sabbah-Saito \cite{SS}. Let us briefly recall their work.

Let $X$ temporarily be a complex algebraic variety and let $f\in\Gamma(X,\mathcal O_X)$. Let $\mathbb V$ be a $\mathbb C$-local system on $X^\mathrm{an}$. By Deligne's Riemann-Hilbert correspondence \cite{Del}, it corresponds to a locally free $\mathcal O_X$-module $H$ with an integrable connection $\nabla$ having regular singularity at infinity. Consider the formal microlocalization of $(H,\nabla)$ with respect to $f$, namely
$$
(H,\nabla)((u))^f=(H((u)),\nabla-\frac{df}{u}\wedge).
$$
Here $u$ is a variable and $H((u))$ is the sheaf of Laurent formal series with coefficients in $H$. One can check that the operator $u\partial_u+f/u$ acts on the de Rham complex $\Omega^*(H,\nabla)((u))^f$. On the other hand, for any $c\in\mathbb C$ and $k\in\mathbb Z$, consider the formal meromorphic connection
$$
\hat{\mathcal E}(H^{k-1},T_c):=(f^{-1}(c)^\mathrm{an},\phi_{f-c}\mathbb V)((u)),u\partial_u-\frac{\log T_c}{2\pi\sqrt{-1}}+cu^{-1}).
$$
Here $T_c$ is the monodromy operator of $\phi_{f-c}\mathbb V$ and the real part of the eigenvalues of $\frac{\log T_c}{2\pi\sqrt{-1}}$ are contained in $[0,1)$. The main result of \cite{SS} shows that there is an isomorphism of formal meromorphic connections
\begin{eqnarray}\label{SS}
(H^k(X,\Omega^*(H,\nabla)((u))^f),u\partial_u+\frac{f}{u})\cong\bigoplus_{c\in\mathbb C}\hat{\mathcal E}(H^{k-1},T_c).\end{eqnarray}

Return to the setting in positive characteristic. To avoid some technical problems, we adjust the setting as follows:
\begin{itemize}
\item $(E,\theta)$ is a Hodge bundle on $X/k$, namely 
$E=\bigoplus_{i=0}^{p-1}E_i,~\theta(E_i)\subset E_{i-1}$
and each $E_i$ is a locally free $\mathcal O_X$-module;
\item $X/k$ admits a system of coordinates $(t_1,\cdots,t_n)$ such that
$$
f=t_1^{a_1}+\cdots+t_n^{a_n},~p\nmid a_1,\cdots,p\nmid a_n
$$
and $Z(t_1,\cdots,t_n)=\{x\}$ for some $x\in X$ ($x$ is called a Brieskorn-Pham isolated singularity of $f$);
\item $(H,\nabla):=(F^*E,\nabla_\mathrm{can}+(\mathrm{id}\otimes\zeta)F^*\theta))$, where
$$
\zeta:F^*\Omega^1_{X/k}\to\Omega^1_{X/k},~dt_i\mapsto t_i^{p-1}dt_i,~i=1,\cdots,n.
$$
\end{itemize}
Under the setting above, we prove
\begin{theorem}
(i)  The amplitudes of $\Omega^*(H,\nabla)((u))^f$ and $\phi_f\Omega^*(H,\nabla)$ are $\{n\}$, i.e.,
$$
\mathcal H^i(\Omega^*(H,\nabla)((u))^f)=0,~\mathcal H^i(\phi_f\Omega^*(H,\nabla))=0
$$
for $i\neq n$. Moreover, the cohomology
$$
\mathcal H^n(\Omega^*(H,\nabla)((u))^f),~\mathcal H^n(\phi_f\Omega^*(H,\nabla))
$$
are supported on $\{x\}$.

(ii) There is an isomorphism of formal meromorphic connections
\begin{eqnarray}\label{Jordan decomposition intro}
(\Omega^*(H,\nabla)((u))^f,u\partial_u+\frac{f}{u})\cong(\phi_f\Omega^*(H,\nabla)((u)),u\partial_u+M).
\end{eqnarray}
Moreover, $M$ is induced by an endomorphism of $\phi_f\Omega^*(H,\nabla)$ and the Jordan decomposition $M=M_s+M_\mathrm{nil}$ of $M$ can be characterized as follows:
\begin{itemize}
\item The  eigenvalues of  $M_s$  together with multiplicities are given by
$$
\sum_{1\leq i\leq n}-\frac{b_i}{a_i},~1\leq b_i\leq a_i-1;
$$
\item $M_n$ involves the Higgs field $\theta$ and $M_\mathrm{nil}=0$ for $\theta=0$.
\end{itemize}
\end{theorem}
(i) may be regarded as a positive characteristic analogue of the fact about vanishing cycle complexes over $\mathbb C$: if $f^{-1}(0)$ contains only isolated singularities of $f$, then the amplitude of $\phi_f\mathbb V$ is $\{\dim X\}$. Take $(E,\theta)=(\mathcal O_X,0)$, then \eqref{Jordan decomposition intro} may be regarded as a positive characteristic analogue of the classical Brieskorn-Pham theorem (see e.g. \cite[Theorem 2.19]{Max}). In general, it may be regarded as a derived version of \eqref{SS} in positive characteristic. We believe that \eqref{Jordan decomposition intro} may hold in a more general situation.
\subsection{De Rham-Higgs comparison for Kontsevich complexes}
Let $k$ be an algebraically closed field of characteristic $p>0$, let $X$ be a smooth variety over $k$ endowed with a simple normal crossing divisor $D$, and let $f:X\to\mathbb P^1$ be a morphism of $k$-varieties such that $f^*[\infty]\subset D$. Assume $(X,D,f)$ is $W_2(k)$-liftable and fix such one, say $(\tilde X,\tilde D,\tilde f)$. In \cite{Kon1}, Kontsevich considered the following 
\begin{definition}
The Kontsevich de Rham complex $\Omega^*_f$ associated to $(X,D,f)$ is defined by 
$$
\Omega^*_f\subset\Omega^*_{X/k}(\log D),~\Omega^i_f:=\{\omega\in\Omega^i_{X/k}(\log D)|\beta\wedge df\in\Omega^{i+1}_{X/k}(\log D)\}.
$$
\end{definition}
Similarly, one can define the Kontsevich Higgs complex $(\Omega_f^\bullet,0)$. In \cite{Kon1} and \cite{Kon2}, Kontsevich sketched an approach to solving his conjecture.  The key of this approach is \eqref{E_1 Kontsevich}, which can be proved by using the method of Deligne-Illusie \cite{DI} (see \cite[Appendix D]{ESY} for more details). The essential input in positive characteristic is \cite[Proposition D.1.5]{ESY}: with the notation and assumption above,
there is an isomorphism 
\begin{eqnarray}\label{ESY decomposition}
\tau_{<p}F_*\Omega_f^*\cong\bigoplus_{i=0}^{p-1}\Omega_f^i[-i]~\mathrm{in}~D(X).
\end{eqnarray}
Using the method presented in \cite{SZ2}, we generalize this result to twisted coefficients.
\begin{definition}
Let $(H,\nabla)$ be a module with integrable connection on $(X,D)/k$. Define the Kontsevich de Rham complex $\Omega_f^*(H,\nabla)$ of $(H,\nabla)$ with respect to $f$ by 
$$
H\otimes\Omega^0_f\xrightarrow{\nabla}H\otimes\Omega^1_f\xrightarrow{\nabla}\cdots\xrightarrow{\nabla}H\otimes\Omega_f^n.
$$
Similarly, for a Higgs module $(H,\nabla)$ on $(X,D)/k$, one defines the Kontsevich Higgs complex of $(E,\theta)$ with respect to $f$ by
$$
E\otimes\Omega^0_f\xrightarrow{\theta}E\otimes\Omega^1_f\xrightarrow{\theta}\cdots\xrightarrow{\theta}E\otimes\Omega_f^n.
$$
\end{definition}
Our generalization is
\begin{theorem}\label{Kontsevich de Rham-Higgs intro}
Notation and assumption as above. Let $(E,\theta)$ be a nilpotent Higgs sheaf of level $\leq\ell$ on $(X,D)/k$ and let $(H,\nabla):=C^{-1}(E,\theta)$. Then there is an isomorphism
$$
F_*\tau_{<p-\ell}\Omega_f^*(H,\nabla)\cong\tau_{<p-\ell}\Omega_f^*(E,\theta)~\mathrm{in}~D(X).
$$
\end{theorem}
Next we assume that $X$ is proper over $k$. Let $(H,\nabla,\mathrm{Fil},\psi)$ be a Fontaine module over $(\tilde X,\tilde D)/W_2(k)$ (see \cite[Definition 4.16]{OV}) and let $(E,\theta):=\mathrm{Gr}_\mathrm{Fil}(H,\nabla)$. Suppose that $\mathrm{rank}(H)+n\leq p$. Clearly, the filtration $\mathrm{Fil}$ on $H$ induces a filtration on the Kontsevich de Rham complex $\Omega_f^*(H,\nabla)$ which is again denoted by $\mathrm{Fil}$. As immediate consequences of the theorem above, we have 
\begin{corollary}
 The Hodge to de Rham spectral sequence 
 $$
 E_1^{i,j}=\mathbb H^{i+j}(X,\mathrm{Gr}^i_\mathrm{Fil}\Omega^*_f(H,\nabla))\Rightarrow\mathbb H^{i+j}(X,\Omega^*_f(H,\nabla))
 $$
 degenerates at $E_1$. Suppose that $L$ is an ample line bundle on $X$, then
 $$
 \mathbb H^i(X,\Omega^*(E,\theta)\otimes L)=0,~i>n.
 $$
\end{corollary}
\subsection{Nearby cycle de Rham/Higgs complexes in positive characteristic}
Inspired by \cite[Appendix E]{ESY}, the Kontsevich de Rham complexes may be used to define the nearby cycle complexes in positive characteristic. Let us briefly recall loc.cit.. 
Let $X$ be a smooth algebraic variety over $\mathbb C$, let $D\subset X$ be a simple normal crossing divisor, and let $f:X\to\mathbb P^1$ be a projective morphism of complex varieties such that $P:=f^*[\infty]\subset D$. Set $U:=X-D$ and set $j:U\hookrightarrow X$ and $j':X-U\hookrightarrow X-P$ to be the natural inclusions. Set $g:=f^{-1}$.

Consider the short exact sequence of complexes
$$
0\to\Omega^*_f\to\Omega^*_X(\log D)\xrightarrow{\rho}\overline{\Omega}^*_X(\log D)\to 0.
$$
Clearly, it gives rise to a distinguished triangle 
$$
\overline{\Omega}^*_X(\log D)[n-1]\to\Omega^*_f[n]\to\Omega^*_X(\log D)[n]\xrightarrow{+1}.
$$
\cite[Theorem E.3]{ESY} shows that this distinguished triangle is the image of the short exact sequence of mixed Hodge modules
$$
0\to\psi_g
(j_*\mathbb Q_{h,U}[n])\to\Xi_g(j'_*\mathbb Q_{h,U}[n])\to j_*\mathbb Q_{h,U}[n]\to0
$$
under $\mathrm{DR}_X$. Here $\mathbb Q_{h,U}[n]$ is the pure Hodge module with underlying perverse sheaf $\mathbb Q_{U^\mathrm{an}}[n]$ and $\Xi_g$ is Beilinson's maximal extension functor. In particular, one has
$$
\mathrm{DR}_X(\psi_g(j_*\mathbb Q_{h,U}[n]))\cong\overline{\Omega}^*_X(\log D)[n-1].
$$
The isomorphism is our starting point to define the nearby cycle complexes in positive  characteristic.

Return to the setting in positive characteristic.
\begin{definition}
Let $(H,\nabla)$ be a bundle with integrable connection on $(X,D)/k$. We define the nearby cycle de Rham complex $\psi_g\Omega^*(H,\nabla)$ by
$$
H\otimes\overline{\Omega}^0_X(\log D)\xrightarrow{\nabla}H\otimes\overline{\Omega}^1_X(\log D)\xrightarrow{\nabla}\cdots\xrightarrow{\nabla}H\otimes\overline{\Omega}^n_X(\log D).
$$
Similarly, for a Higgs module $(H,\nabla)$ on $(X,D)/k$, one defines the nearby cycle Higgs complex $\psi_g\Omega^*(E,\theta)$ by
$$
E\otimes\overline{\Omega}^0_X(\log D)\xrightarrow{\theta}E\otimes\overline{\Omega}^0_X(\log D)\xrightarrow{\theta}\cdots\xrightarrow{\theta}E\otimes\overline{\Omega}^0_X(\log D).
$$


\end{definition}
As an immediate consequence of Theorem \ref{Kontsevich de Rham-Higgs intro}, one has
\begin{corollary}
Notation and assumption as in Theorem \ref{Kontsevich de Rham-Higgs intro}. There is a de Rham-Higgs comparison for nearby cycle complexes
$$
F_*\tau_{<p-\ell}\psi_g\Omega^*(H,\nabla)\cong\tau_{<p-\ell}\psi_g\Omega^*(E,\theta)~\mathrm{in}~D(g^*[0]).
$$
\end{corollary}
As we have done in \S1.2, we define the logarithmic monodromy operator of $\psi_g$ in positive characteristic. Consider the short exact sequence of complexes
$$
0\to\psi_g\Omega^*(H,\nabla)[-1]\xrightarrow{d\log g\wedge}\Omega^*(H,\nabla)\otimes\mathcal O_{g^*[0]}\xrightarrow{\rho}\psi_g\Omega^*(H,\nabla)\to0.
$$
Inspired by \cite[\S11.2.4]{CP}, one makes the following
\begin{definition}
Define the logarithmic monodromy operator $N_\mathrm{dR}$ of $\psi_g\Omega^*(H,\nabla)$ by the extension class of the short exact sequence above.
\end{definition}
\begin{theorem}
$N_\mathrm{dR}$ is nilpotent if $(H,\nabla)=C^{-1}(E,\theta)$ and $(E,\theta)$ is a nilpotent Higgs bundle on $(X,D)/k$ of level $\leq p-n$.
\end{theorem}
We believe that the arguments in \cite[\S3.1-3.2]{OV} may imply the theorem above.
\subsection{Further discussions}
It is natural to ask
\begin{question}
Assume $f:(X,D)\to(\mathbb A^1,[0])$ is a semistable morphism. Let $(H,\nabla,\mathrm{Fil},\kappa)$ be a Fontaine module on $(\tilde X,\tilde D)/W_2(k)$.
Could we define the distinguished triangle 
$$
i^*\Omega^*(H,\nabla)\xrightarrow{\mathrm{sp}}\psi_f\Omega^*(H,\nabla)\xrightarrow{\mathrm{can}}\phi_f\Omega^*(H,\nabla)\xrightarrow{+1}
$$
and $\mathrm{Var}:\phi_f\Omega^*(H,\nabla)\to\psi_f\Omega^*(H,\nabla)$ such that
$$
\mathrm{Var}\circ\mathrm{can}=N_\mathrm{dR},~\mathrm{can}\circ\mathrm{Var}=N_\mathrm{dR}?
$$
\end{question}
\fi


\section{De Rham-Higgs comparison ($df\neq0$)}

Let $k$ be a perfect field of characteristic $p$, let $X$ be a smooth variety of dimension $n$ over $k$ which endowed with a simple normal crossing divisor $D$, and let $f,g\in\Gamma(X,\mathcal O_X)$ such that $(g=0)$ is a subdivisor of $D$. We introduce the following assumptions:
\begin{itemize}
    \item[(S1)\ ] $X$ admits a system of coordinates $t_1,\cdots,t_n$ such that $D=(t_1\cdots t_r=0)$ and $g=t_1\cdots t_s$;
    \item[(S2)\ ] there exist $W_2(k)$-liftings $\tilde X,\tilde t_i,\tilde g$ of $X,t_i,g$, respectively and a Frobenius lifting $\tilde F:\tilde X\to \tilde X$ such that $\tilde F^*(\tilde t_i)=\tilde t_i^p$ and $\tilde g=\tilde t_1\cdots\tilde t_s$. 
\end{itemize}
Under the first assumption, we have a natural identification
$$
F_*\mathcal O_X=\mathcal O_X\{t_1^{i_1}\cdots t_n^{i_n}:0\leq i_1,\cdots,i_n<p\}.
$$
It follows that 
\begin{eqnarray}\label{decom f}
f=\sum_{0\leq i_1,\cdots,i_n<p}f_{i_1,\cdots,i_n},~f_{i_1,\cdots,i_n}=h_{i_1,\cdots,i_n}^pt_1^{i_1}\cdots t_n^{i_n},~h_{i_1,\cdots,i_n}\in\Gamma(X,\mathcal O_X).
\end{eqnarray}
Without loss of generality, we may assume that $f_{0,\cdots,0}=0$. Let $Z$ be the closed subscheme of $X$ defined by all $f_{i_1,\cdots,i_n}$. Let $(E,\theta)$ be a Higgs bundle of nilpotent level $\leq p-1$ on $(X,D)/k$ and set $(H,\nabla):=C^{-1}_{\tilde F}(E,\theta)=(F^*E,\nabla_{can}+(id\otimes\zeta_{\tilde F})F^*\theta)$. Set 
\begin{eqnarray}\label{MHM complex}
K^*_{dR}:=\star(H,\nabla),~K^*_{Hig}:=\star(E,\theta),~\star=\Omega^*,~\Omega^*_{int},~\Omega^*_{g^{-1}},~W_i\Omega^*.
\end{eqnarray}



The main purpose of this section is to verify that the slogan in the introduction is true after restricting to the formal neighborhood of $X$ along $Z$. More precisely, we have
\begin{theorem}\label{main thm s2}
With the notation and assumptions above. Set $\mathfrak X:=X_{/Z}$. Then there is a quasi-isomorphism of complexes
\begin{eqnarray}\label{Cartier quasi}
\varphi_{\tilde F}:(K^\bullet_{Hig},\theta-df\wedge)\otimes\mathcal O_{\mathfrak X}\to F_*(K^\bullet_{dR},\nabla+df\wedge)\otimes\mathcal O_{\mathfrak X}.
\end{eqnarray}
\end{theorem}
We call this kind of quasi-isomorphism \emph{the Cartier quasi-isomorphism}. For $(E,\theta)=(\mathcal O_X,0), f=0,\star=\Omega^*$, it recovers the classical Cartier quasi-isomorphism. For a general $(E,\theta)$ and $f=0,\star=\Omega^*$, Ogus-Vologodsky\cite{OV} and Schepler\cite{Schepler} established an isomorphism between $F_*\Omega^*(H,\nabla)$ and $\Omega^*(E,\theta)$ in the derived category $D(X)$. For a general $(E,\theta)$ and $f=0,\star=\Omega^*_{int}$, it was constructed by Sheng and the author\cite{SZ2}. For $(E,\theta)=(\mathcal O_X,0),D=\emptyset,\star=\Omega^*$ and some special $f$ ($df\neq0$), Ogus-Vologodsky\cite{OV} again and Arinkin-C\u{a}ld\u{a}raru-Hablicsek\cite{ACH} respectively constructed isomorphisms between $F_*\Omega^*(\mathcal O_X,d+df\wedge)$ and $\Omega^*(\mathcal O_X,-df\wedge)$ in the derived category $D(X)$.

 The theorem above has two immediate corollaries. If we take $f=0$, then $Z=X$. It follows that
\begin{corollary}
The Frobenius lifting $\tilde F$ above gives rise to a quasi-isomorphism of complexes
$$
\varphi_{\tilde F}:K_{Hig}^*\to F_*K^*_{dR}
$$
\end{corollary}



Keep the data $X/k,D,f,g,(E,\theta)$ above, but we do not require the assumptions (S1) and (S2). Instead, we only assume that there exist $W_2(k)$-liftings $\tilde X,\tilde D,\tilde g$ of $X,D,g$, respectively such that $(\tilde g=0)$ is a subdivisor of $\tilde D$. Set $(H,\nabla):=C^{-1}_{(\tilde X,\tilde D)/W_2(k)}(E,\theta)$ and construct $K^*_{dR},K^*_{Hig}$ as \eqref{MHM complex}.
 \begin{corollary}\label{corollary 2}
The supports of $\mathcal H^i(K^\bullet_{dR},\nabla+df\wedge)$ and $\mathcal H^i(K^\bullet_{Hig},\theta-df\wedge)$ coincide. In particular, if the support of $(K^\bullet_{Hig},\theta-df\wedge)$ is a finite set, then 
$$
F_*(K^\bullet_{dR},\nabla+df\wedge)\cong(K^\bullet_{Hig},\theta-df\wedge)~\mathrm{in}~D(X).
$$     
\end{corollary}
\begin{proof}
For any closed point $x\in X$, we can choose a system of coordinates $t_1,\cdots,t_n$ around $x$ such that $t_(x)=\cdots=t_n(x)=0$ and $g=t_1\cdots t_s$. Let $\tilde t_i$ be a lifting of $t_i$ and let $\tilde F$ be a Frobenius lifting around $x$ such that $\tilde F(\tilde t_i)=\tilde t_i^p$.
Write $f$ as \eqref{decom f} around $x$ and assume $f_{0,\cdots,0}=0$.  Let $Z$ be the locally closed subscheme of $X$ defined by all $f_{i_1,\cdots,i_n}$ and let $\mathfrak X=X_{/Z}$. According to the theorem above, coupled with the flatness of $\mathcal O_{\mathfrak X}$ over $\mathcal O_X$, we have
$$
\mathcal H^i(K^\bullet_{Hig},\theta-df\wedge)\otimes\mathcal O_{\mathfrak X}\cong\mathcal H^i(F_*(K^\bullet_{dR},\nabla+df\wedge))\otimes\mathcal O_{\mathfrak X}.
$$
Tensoring it with $\otimes_{\mathcal O_{\mathfrak X}}\hat{\mathcal O}_{X,x}$, one gets
$$
\mathcal H^i(K^\bullet_{Hig},\theta-df\wedge)_x\otimes\hat{\mathcal O}_{X,x}\cong\mathcal H^i(F_*(K^\bullet_{dR},\nabla+df\wedge))_x\otimes\hat{\mathcal O}_{X,x}.
$$
Since $\hat{\mathcal O}_{X,x}$ is fully faithful flat over $\mathcal O_{X,x}$, the isomorphism above implies that
$$
\mathcal H^i(K^\bullet_{Hig},\theta-df\wedge)_x\neq0\Longleftrightarrow\mathcal H^i(K^\bullet_{dR},\nabla+df\wedge)_x\neq0.
$$
From which, the first statement of this corollary follows. 

For the second statement, it suffices to assume that the support of $(K^\bullet_{Hig},\theta-df\wedge)$ is $\{x\}$. Keep the notation and assumption in the above paragraph. Consider the  diagram
$$
\xymatrix{
(K^\bullet_{Hig},\theta-df\wedge)\ar@{.>}[r]\ar[d]&F_*(K^\bullet_{dR},\nabla+df\wedge)\ar[d]\\
(\mathcal K^\bullet_{Hig},\theta-df\wedge)\otimes\hat{\mathcal O}_{X,x}\ar[r]^-{\varphi_{\tilde F}}& F_*(K^\bullet_{dR},\nabla+df\wedge)\otimes\hat{\mathcal O}_{X,x},
}
$$
where the vertical arrows are the natural morphisms of complexes. Combing the first statement with the assumption on the support of $(K^\bullet_{Hig},\theta-df\wedge)$, we conclude that they are quasi-isomorphisms. Consequently, there is a unique dotted arrow in $D(X)$ which makes the diagram above commutative. Obviously, it is an isomorphism. This completes the proof of the second statement.
\end{proof}

In the remaining part of this section, we dedicate to proving Theorem \ref{main thm s2}. Roughly speaking, the approach to prove Theorem \ref{main thm s2} contains two steps: the first step is to deform $\theta-df\wedge$ to an appropriate Higgs field $\Theta$ such that their Higgs complex are isomorphic and the inverse Cartier transform of $\Theta$ is isomorphic to $\nabla-df\wedge$; the second step is to show that $F_*\star(C_{\tilde F}^{-1}(\hat E,\Theta))$ is quasi-isomorphic to $\star(\hat E,\Theta)$, where $\star$ is defined as \eqref{MHM complex}. 

\subsection{Proof of Theorem \ref{main thm s2} (step 1)}
We need introduce some notions on formal schemes. Set 
$$(\mathfrak X,\mathfrak D):=(X,D)_{/Z},~\Omega^1_{\mathfrak X/k}(\log\mathfrak D):=\Omega^1_{X/k}(\log D)\otimes\mathcal O_{\mathfrak X},~\hat E:=E_{/Z},~\hat .$$
It is obvious that the notions of Higgs field, integrable connection and its $p$-curvature on $(X,D)/k$ still work on $(\mathfrak X,\mathfrak D)/k$.
Clearly, the Higgs field $\theta$ on $E$ induces a Higgs field on $\hat E$, namely
$$
\hat\theta:\hat E\to\hat E\otimes\Omega^1_{\mathfrak X/k}(\log\mathfrak D).
$$
Let $\mathcal I$ be the ideal subsheaf of $\mathcal O_X$ corresponding to $Z$. One may check that the logarithmic integrable connection $\nabla$ on $H=F^*E$ induces a morphism
$$
\nabla_i:H/(F^*\mathcal I)^iH\to H/(F^*\mathcal I)^iH\otimes\Omega^1_{X/k}(\log D).
$$
Note that the inverse limit of $H/(F^*\mathcal I)^iH$ is $\hat H$ and the inverse limit of $\nabla_i$ gives rise to a logarithmic connection on $\hat H$, namely
$$
\hat\nabla:\hat H\to\hat H\otimes\Omega^1_{\mathfrak X/k}(\log\mathfrak D).
$$
For convenient, the Frobenius of $\mathfrak X$ is denoted by $F$ again. Let
$$
\hat{\zeta}_{\tilde F}:F^*\Omega^1_{\mathfrak X/k}(\log\mathfrak D)\to\Omega^1_{\mathfrak X/k}(\log\mathfrak D)
$$
be the $\mathcal O_{\mathfrak X}$-linear morphism induced by 
$$\zeta_{\tilde F}:F^*\Omega^1_{X/k}(\log D)\to\Omega^1_{X/k}(\log D).$$

\begin{definition}
Let $\Theta$ be any Higgs field on $\hat E$ (not necessarily $\hat\theta$). Define the inverse Cartier transform of $(\hat E,\Theta)$ with respect to $\tilde F$ by 
$$
C^{-1}_{\tilde F}(\hat E,\Theta):=(F^*\hat E,\nabla_{can}+(id\otimes\hat{\zeta}_{\tilde F})F^*\Theta).
$$
\end{definition}
It is obvious that $C^{-1}_{\tilde F}(\hat E,\hat\theta)=C^{-1}_{\tilde F}(E,\theta)_{/Z}=(\hat H,\hat\nabla)$. By \cite[Theorem 2.8, 3]{OV}, we know that the $p$-curvature of $\nabla$ is $-F^*\theta$. By careful computation, one may check that the $p$-curvature of $\nabla+df\wedge$ is $F^*(-\theta+df\wedge)$. Unfortunately, this does not imply that $C_{\tilde F}^{-1}(E,\theta-df\wedge)$ is isomorphic to $(H,\nabla+df\wedge)$.
\begin{example}
Take $(E,\theta)=(\mathcal O_X,0)$ and $f=t_1$. One may check that the $p$-curvature of 
$$C^{-1}_{\tilde F}(\mathcal O_X,-dt_1\wedge)=(\mathcal O_X,d-t_1^{p-1}dt_1)$$
is $(\mathcal O_X,F^*dt_1-F^*(t_1^{p-1}dt_1))$. Consequently, the extra term $-F^*(t_1^{p-1}dt_1)$ prevents the existence of an isomorphism between $C^{-1}_{\tilde F}(\mathcal O_X,-dt_1\wedge)$ and $(\mathcal O_X,d+df\wedge)$.
\end{example}

Nevertheless, after restricting $\theta-df\wedge$ to the formal neighborhood of $Z$ in $X$, the resulting Higgs field $\hat\theta-df\wedge$ can be deformed to an appropriate Higgs field $\Theta$ such that its inverse cartier transform is $\hat\nabla+df\wedge$.
\begin{lemma}
Set
$$
\Theta:=\hat\theta-\sum_{0\leq i_1,\cdots,i_n<p}\sum_{j=0}^{\infty}f_{i_1\cdots,i_n}^{p^j-1}df_{i_1,\cdots,i_n}\wedge.
$$
Then its inverse Cartier transform is isomorphic to $\hat\nabla+df\wedge$.
\end{lemma}
\begin{proof}
It suffices to construct an isomorphism between the de Rham bundles $(\hat H, \hat\nabla+df\wedge)$ and $C^{-1}_{\tilde F}(\hat E,\Theta)$. Inspired by the proof of \cite[Theorem 4.28]{OV}, we consider the Artin-Hasse exponential of $g_{i_1,\cdots,i_n}$
$$
\mathrm{AH}(f_{i_1,\cdots,i_n})=\mathrm{exp}(\sum_{i=0}^\infty\frac{ f^{p^i}_{i_1,\cdots,i_n}}{p^i}).
$$
It is a power series of $g_{i_1,\cdots,i_n}$ with $p$-adically integral coefficients. Set 
$$
G:=\prod_{0\leq i_1,\cdots,i_n<p}\mathrm{AH}(g_{i_1,\cdots,i_n}).
$$
Then we claim that multiplication by $G$ gives rise to the desired isomorphism, i.e.,
$$
\xymatrix{
\hat H\ar[rrrr]^-{\nabla_{can}+(id\otimes\hat{\zeta}_{\tilde F})F^*\hat\theta+df\wedge}\ar[d]^-{G}&&&&\hat H\otimes\Omega^1_{\mathfrak X/k}(\log\mathfrak D)\ar[d]^-{G}\\
\hat H\ar[rrrr]^-{\nabla_{can}+(id\otimes\hat{\zeta}_{\tilde F})F^*\Theta}&&&&\hat H\otimes\Omega^1_{\mathfrak X/k}(\log\mathfrak D).
}
$$
The truth of the claim is equivalent to the equality between
\begin{eqnarray}\label{R}
(G\otimes id)(\nabla_{can}+(id\otimes\hat{\zeta}_{\tilde F})F^*\hat\theta+df\wedge)(F^*\hat e)
\end{eqnarray}
and 
\begin{eqnarray}\label{L}
(\nabla_{can}+(id\otimes\hat{\zeta}_{\tilde F})F^*\Theta)(GF^*\hat e)
\end{eqnarray}
for any $\hat e\in\hat E$. By an easy computation, \eqref{R} becomes
$$
(id\otimes\hat{\zeta}_{\tilde F})F^*\hat\theta(GF^*\hat e)+F^*\hat e\otimes Gdf
$$
and \eqref{L} becomes
$$
F^*\hat e\otimes dG+(id\otimes\hat{\zeta}_{\tilde F})F^*\hat\theta(GF^*\hat e)-F^*\hat e\otimes G\sum_{0\leq i_1,\cdots,i_n<p}\sum_{j=1}^{\infty}g_{i_1\cdots,i_n}^{p^j-1}dg_{i_1,\cdots,i_n}.
$$
Consequently, \eqref{R}=\eqref{L} is simplified to 
$$
F^*\hat e\otimes dG=F^*\hat e\otimes G\sum_{0\leq i_1,\cdots,i_n<p}\sum_{j=0}^{\infty}g_{i_1\cdots,i_n}^{p^j-1}dg_{i_1,\cdots,i_n}.
$$
Clearly, the last equality follows from 
$$
 dG=G\sum_{0\leq i_1,\cdots,i_n<p}\sum_{j=0}^{\infty}g_{i_1\cdots,i_n}^{p^j-1}dg_{i_1,\cdots,i_n}.
$$
This completes the proof.
\end{proof}

Next, we show that
\begin{lemma}
The complexes $\star(\hat E,\hat\theta-df\wedge)$ and $\star(\hat E,\Theta)$ are isomorphic.   
\end{lemma}
\begin{proof}
The proof contains three steps: the first step is to prove this lemma for $\star=\Omega^*$; the second step is for $\Omega^*_{int},W_i\Omega^*$; the last step is for $\star=\Omega^*_{g^{-1}}$.

\emph{Step 1}. Recall that the assumption (S1) gives rise to a system of coordinates $t_1,\cdots,t_n$ on $X/k$ such that $D=(t_1\cdots t_r=0)$ and $g=t_1\cdots t_s$. Write
\begin{eqnarray}\label{Theta cor}
\hat\theta=\sum_{i=1}^r\hat\theta_i\otimes d\log t_i+\sum_{i=r+1}^n\hat\theta_i\otimes dt_i,~\Theta=\sum_{i=1}^r\Theta_i\otimes d\log t_i+\sum_{i=r+1}^n\Theta_i\otimes dt_i
\end{eqnarray}
and the K\"ahler differential 
$$d=\sum_{i=1}^rt_i\partial_i\otimes d\log t_i+\sum_{i=r+1}^n\partial_i\otimes dt_i.$$ 
Clearly, we have natural isomorphisms
$$
\Omega^*(\hat E,\hat\theta-df\wedge)\cong\mathrm{Kos}(\hat E;\hat\theta_1-t_1\partial_1f,\cdots,\hat\theta_r-t_r\partial_rf,\hat\theta_{r+1}-\partial_{r+1}f,\cdots,\hat\theta_n-\partial_nf)
$$
and
$$
\Omega^*(\hat E,\Theta)\cong\mathrm{Kos}(\hat E;\Theta_1,\cdots,\Theta_n).
$$
To show $\Omega^*(\hat E,\hat\theta-df\wedge)$ is isomorphic to $\Omega^*(\hat E,\Theta)$, it suffices to check that 
$$
\Theta_i=\left\{
\begin{matrix}
  \vartheta_i(\hat\theta_i-t_i\partial_if),&1\leq i\leq r;\\
\vartheta_i(\hat\theta_i-\partial_if),&r+1\leq r\leq n,
\end{matrix}
\right.
$$
where each $\vartheta_i$ is an $\mathcal O_{\mathfrak X}$-linear automorphism of $\hat E$. In fact, we may construct an isomorphism of complexes 
$$
\psi:\Omega^*(\hat E,\hat\theta-df\wedge)\to\Omega^*(\hat E,\Theta) 
$$
as follows: set 
\begin{eqnarray}\label{omega}
\omega_1:=d\log t_1,\cdots,\omega_r:=d\log t_r,~\omega_{r+1}:=dt_{r+1},\cdots,\omega_n:=dt_n
\end{eqnarray}
and put
$$
\psi(\hat e):=\hat e,~\psi(\hat e\otimes\omega_{i_1}\wedge\cdots\wedge\omega_{i_q}):=\vartheta_{i_1}\cdots\vartheta_{i_q}(\hat e)\otimes\omega_{i_1}\wedge\cdots\wedge\omega_{i_q},~\hat e\in\hat E.
$$

Next, let us construct $\vartheta_i$. Using the facts
$$
t_q\partial_qf_{i_1,\cdots,i_n}=i_qf_{i_1,\cdots,i_n};~i_q=i_q^p\mod p;~x^p+y^p=(x+y)^p,~x,y\in\mathcal O_{\mathfrak X},
$$
we get
\begin{eqnarray*}
\sum_{0\leq i_1,\cdots,i_n<p}\sum_{j=0}^{\infty}f_{i_1\cdots,i_n}^{p^j-1}df_{i_1,\cdots,i_n}=\sum_{0\leq i_1,\cdots,i_n<p}\sum_{j=0}^{\infty}\sum_{q=1}^nf_{i_1\cdots,i_n}^{p^j-1}f_{i_1,\cdots,i_n}t_q\partial_qf_{i_1,\cdots,i_n}d\log t_q\\
=\sum_{0\leq i_1,\cdots,i_n<p}\sum_{j=0}^{\infty}\sum_{q=1}^n(i_qf_{i_1\cdots,i_n})^{p^j}d\log t_q=\sum_{q=1}^n\sum_{j=0}^{\infty}(t_q\partial_qf)^{p^j}d\log t_q.~~~~~~~~~~
\end{eqnarray*}
It follows that
$$
\Theta_i=\hat\theta_i-\sum_{q=1}^n\sum_{j=0}^{\infty}(t_q\partial_qf)^{p^j},~1\leq i\leq r;~t_i\Theta_i=t_i\hat\theta_i-\sum_{q=1}^n\sum_{j=0}^{\infty}.(t_q\partial_qf)^{p^j},~r+1\leq i\leq n.
$$
Consequently, for $1\leq i\leq r$, we may set
\begin{eqnarray*}
\vartheta_i:=\frac{\hat\theta_i-\sum_{j=0}^\infty(t_i\partial_i f)^{p^j}}{\hat\theta_i-t_i\partial_i f}=1-\frac{\sum_{j=1}^\infty(t_i\partial_i f)^{p^j}}{\hat\theta_i-t_i\partial_i f}~~~~~~\\
=1+\sum_{j=1}^\infty\sum_{q=0}^{p-1}(t_i\partial_i f)^{p^j-1}(\frac{\hat\theta_i}{t_i\partial_if})^q=1+\sum_{j=1}^\infty\sum_{q=0}^{p-1}(t_i\partial_i f)^{p^j-q-1}\hat\theta_i^q.
\end{eqnarray*}
Here the third equality follows from the Taylor expansion
$\frac{1}{1-x}=\sum_{i=0}^\infty x^i$
and the nilpotent condition $\hat\theta_i^p=0$. Since each term $(t_i\partial_i f)^{p^j-s-1}\hat\theta_i^s$ is a topological nilpotent endomorphism of $\hat E$, we conclude that $\vartheta_i$ is an automorphism of $\hat E$. Similarly, for $r+1\leq i\leq n$, we may set
$$
\vartheta_i:=1+\sum_{j=1}^\infty\sum_{q=0}^{p-1}(t_i\partial_i f)^{p^j-q-1}(t_i\hat\theta_i)^q.
$$
Clearly, it is an automorphism of $\hat E$ again.This completes the construction of $\vartheta_i$.

\emph{Step 2}. Using the following easily checked facts
$$
\psi(\Omega_{int}^j(\hat E,\hat\theta-df\wedge))=\Omega_{int}^j(\hat E,\Theta),~\psi(W_i\Omega^j(\hat E,\hat\theta-df\wedge))=W_i\Omega^j(\hat E,\Theta),
$$
we deduce that the isomorphism $\psi$ restricts to isomorphisms of subcomplexes
$$
\psi_{int}:\Omega_{int}^*(\hat E,\hat\theta-df\wedge)\to\Omega_{int}^*(\hat E,\Theta),~W_i\psi:W_i\Omega^*(\hat E,\hat\theta-df\wedge)\to W_i\Omega^*(\hat E,\Theta).
$$ 

\emph{Step 3}. Unfortunately, the isomorphism $\psi$ constructed in \emph{step 1} may not restricts to an isomorphism of complexes 
$$\psi_{g^{-1}}:\Omega^*_{g^{-1}}(\hat E,\hat\theta-df\wedge)\to\Omega^*_{g^{-1}}(\hat E,\Theta).$$
Nevertheless, we may construct the desired isomorphism $\psi_{g^{-1}}$ by using a similar idea as $\psi_{int}$ or $W_i\psi$. Write
$$
\hat\theta=\hat\theta'_1\otimes d\log g+\sum_{i=2}^r\hat\theta'_2\otimes d\log t_i+\sum_{i=r+1}^n\hat \theta'_i\otimes dt_i.
$$
and
$$
d=t_1\partial_1\otimes d\log g+\sum_{i=2}^s(t_i\partial_i-t_1\partial_1)\otimes d\log t_i+\sum_{i=s+1}^rt_i\partial_i\otimes d\log t_i+\sum_{i=r+1}^n\partial_i\otimes dt_i.
$$
it is easy to see that
$$
\hat\theta_1'=\hat\theta_1,~\hat\theta_2'=\hat\theta_2-\hat\theta_1,\cdots,\hat\theta_s'=\hat\theta_s-\hat\theta_1,~\hat\theta_{s+1}'=\hat\theta_{s+1},\cdots,\hat\theta_n'=\hat\theta_n
$$
and 
$$
df=t_1\partial_1f\otimes d\log g+\sum_{i=2}^s(t_i\partial_i-t_1\partial_1)f\otimes d\log t_i+\sum_{i=s+1}^rt_i\partial_if\otimes d\log t_i+\sum_{i=r+1}^n\partial_if\otimes dt_i.
$$
For our purpose, we construct another isomorphism of complexes 
$$
\psi':\Omega^*(\hat E,\hat\theta-df\wedge)\to\Omega^*(\hat E,\Theta) 
$$
as follows: set
$$
\omega'_1:=d\log g,~\omega'_2:=d\log t_2,\cdots,\omega'_r:=d\log t_r,~\omega'_{r+1}:=dt_{r+1},\cdots,\omega'_n:=dt_n
$$
and put
$$
\psi'(\hat e)=\hat e,~\psi'(\hat e\otimes\omega_{i_1}\wedge\cdots\omega_{i_q})=\vartheta'_{i_1}\cdots\vartheta'_{i_q}(\hat e)\otimes\omega'_{i_1}\wedge\cdots\wedge\omega'_{i_q}.
$$
Here 
$$
\vartheta_i':=\left\{
\begin{matrix}
1+\sum_{j=1}^\infty\sum_{q=0}^{p-1}((t_i\partial_i-t_1\partial_1) f)^{p^j-q-1}\hat\theta_i^q,&2\leq i\leq s,\\
\vartheta_i,&\mathrm{otherwise}.
\end{matrix}
\right.
$$
By the very construction of $\Omega^*_{g^{-1}}$, it is easy to see that $\psi'$ restricts to an isomorphism of subcomplexes 
$$
\psi'_{g^{-1}}:\Omega^*_{g^{-1}}(\hat E,\hat\theta-df\wedge)\to\Omega^*_{g^{-1}}(\hat E,\Theta).
$$
This completes the proof. 
\end{proof}


\subsection{Proof of Theorem \ref{main thm s2} (step 2)}

The key to our construction of a quasi-isomorphism between $F_*\star(C_{\tilde F}^{-1}(\hat E,\Theta))$ and $\star(\hat E,\Theta)$ is that $F_*\star(C_{\tilde F}^{-1}(\hat E,\Theta))$ can be decomposed into a direct sum of several acyclic complexes and $\star(\hat E,\Theta)$ for $\star=\Omega^*,\Omega^*_{int},\Omega^*_{g^{-1}}$. This idea is brought from log geometry, see e.g. the proof of \cite[Theorem (4.12)]{Kato}. Next, we construct the desired quasi-isomorphism case by case. 

\emph{Case of $\star=\Omega^*$}.  By definition, we a natural identification
\begin{eqnarray}\label{Frobenius push}
 F_*\Omega^q C^{-1}_{\tilde F}(\hat E,\Theta)\cong \hat E\otimes F_*\Omega^q_{\mathfrak X/k}(\log\mathfrak D). 
\end{eqnarray}
Meanwhile, one may check that
\begin{eqnarray}\label{decom F_*omega}
F_*\Omega^q_{\mathfrak X/k}(\log\mathfrak D)\cong\mathcal O_{\mathfrak X}\{t^{\boldsymbol{\alpha}}\omega_I:\boldsymbol{\alpha}\in\{0,\cdots,p-1\}^n,~I\subset\{1,\cdots,n\},~|I|=q\}.
\end{eqnarray}
Here $\omega_i$ is defined as \eqref{omega}, 
$$\omega_\emptyset:=1,~\omega_I:=\omega_{i_1}\wedge\cdots\wedge\omega_{i_q},~I=\{i_1,\cdots,i_q\},~i_1<\cdots<i_q$$
and $t^{\boldsymbol{\alpha}}=t_1^{\alpha_1}\cdots t_n^{\alpha_n}$ for $\boldsymbol{\alpha}=(\alpha_1,\cdots,\alpha_n)$.
\begin{definition}\label{Def w}
Let $B$ be the set consisting of all $t^{\boldsymbol{\alpha}}\omega_I$ with $\boldsymbol{\alpha}\in\{0,\cdots,p-1\}^n$ and $I\subset\{1,\cdots,n\}$. Define a function
$$
V:B\to \mathbb F_{p}^n,~t^{\boldsymbol{\alpha}}\omega_I\mapsto \boldsymbol{\alpha}+\boldsymbol{\alpha}_I \mod p.$$
Here $\boldsymbol{\alpha}_I=(\alpha_{1,I},\cdots,\alpha_{n,I})$ is defined as 
$$
\alpha_{i,I}:=\left\{
\begin{matrix}
1,&i\in I\cap \{r+1,\cdots,n\},\\
0,&\mathrm{otherwise}
\end{matrix}
\right.
$$
and $\boldsymbol{\beta}\mod p$ for $\boldsymbol{\beta}=(\beta_1,\cdots,\beta_n)\in\mathbb N^n$ is defined as 
$(\beta_1\mod p,\cdots,\beta_n\mod p)$.
\end{definition}
Write $\Theta$ as \eqref{Theta cor} and make the identification 
$$
F_*\Omega^q C^{-1}_{\tilde F}(\hat E,\Theta)\cong\hat E\{t^{\boldsymbol{\alpha}}\omega_I:\boldsymbol{\alpha}\in\{0,\cdots,p-1\}^n,~I\subset\{1,\cdots,n\},~|I|=q\}.
$$
Since
$$
\hat\nabla_\Theta:=\nabla_{can}+(id\otimes\zeta_{\tilde F})F^*\Theta=\sum_{i=1}^r(t_i\partial_i+F^*\Theta_i)\otimes d\log t_i+\sum_{i=r+1}^n(\partial_i+t_i^{p-1}F^*\Theta_i)\otimes dt_i,
$$
coupled with the identification \eqref{Frobenius push}, we conclude that $\hat\nabla_\Theta(\hat e\otimes t^{\boldsymbol{\alpha}}\omega_I)$ equals to
\begin{eqnarray*}
\sum_{i\leq r}(\alpha_iid+\Theta_i)\hat e\otimes d\log t_i\wedge t^{\boldsymbol{\alpha}}\omega_I+\sum_{i>r,\alpha_i>0}(\alpha_iid+t_i\Theta_i)\hat e\otimes d\log t_i\wedge t^{\boldsymbol{\alpha}}\omega_I\\
+\sum_{i>r,\alpha_i=0}\Theta_i\hat e\otimes t_i^{p-1}dt_i\wedge t^{\boldsymbol{\alpha}}\omega_I.
\end{eqnarray*}
Set 
$$
\gamma_i:=\left\{
\begin{matrix}
d\log t_i\wedge t^{\boldsymbol{\alpha}}\omega_I,&i\leq r~\mathrm{or}~i>r~\mathrm{and}~\alpha_i>0;\\
 t_i^{p-1}dt_i\wedge t^{\boldsymbol{\alpha}}\omega_I,&i>r~\mathrm{and}~\alpha_i=0.

\end{matrix}
\right.
$$ 
It is easy to see that if $\gamma_i\neq0$ (i.e., $i\notin I$), then either $\gamma_i$ or $-\gamma_i$ belongs to $B$. By the definition of $w$ above, we get the key observation
$$
V(t^{\boldsymbol{\alpha}}\omega_I)=V(\gamma_i)~\mathrm{or}~V(-\gamma_i),~i\leq i\leq n.
$$
It easily follows that
\begin{lemma}
For any $v\in\mathbb F_p^n$, the graded $\mathcal O_{\mathfrak X}$-submodule $\hat EV^{-1}(v)\subset\hat E\otimes F_*\Omega^\bullet_{\mathfrak X/k}(\log\mathfrak D)$ is closed under the differential $\hat\nabla_\Theta$. In other words, there is a subcomplex $K_v^*\subset F_*\Omega^*C_{\tilde F}^{-1}(\hat E,\Theta)$ such that $K_v^\bullet=\hat EV^{-1}(v)$.  Moreover, we have a decomposition of complexes
$$
F_*\Omega^*C_{\tilde F}^{-1}(\hat E,\Theta)\cong\bigoplus_{v\in\mathbb F_p^n}K_v^*.
$$
\end{lemma}
Note that $w(t^{\boldsymbol{\alpha}}\omega_I)=0$ if and only if 
$$t^{\boldsymbol{\alpha}}\omega_I\in\wedge^\bullet\{d\log t_1,\cdots,d\log t_r,t_{r+1}^{p-1}dt_{r+1},\cdots,t_n^{p-1}dt_n\}.$$
Hence for such $t^{\boldsymbol{\alpha}}\omega_I$, we have
$$
\hat\nabla_\Theta(\hat e\otimes t^{\boldsymbol{\alpha}}\omega_I)=\sum_{i=1}^r\Theta\hat e\otimes d\log t_i\wedge t^{\boldsymbol{\alpha}}\omega_I+\sum_{i=r+1}^n\Theta_i\hat e\otimes t_i^{p-1}dt_i\wedge t^{\boldsymbol{\alpha}}\omega_I.
$$
Consequently, we have an isomorphism of complexes
\begin{eqnarray}\label{varphi star=omega}
\begin{array}{c}
\varphi_{\tilde F}:\Omega^*(\hat E,\Theta)\to K_0^*,\\
\hat e\mapsto \hat e,\\
\hat e\otimes\omega_{i_1}\wedge\cdots\wedge\omega_{i_q}\mapsto \hat e\otimes\zeta_{\tilde F}(\omega_{i_1})\wedge\cdots\zeta_{\tilde F}(\omega_{i_q}). 
\end{array}
\end{eqnarray}

It remains to show that 
\begin{lemma}
$K_v^*$ is acyclic for $v\neq0$. 
\end{lemma}
\begin{proof}
To prove this lemma, we need the fact that $K_v^*$ is a Koszul complex, i.e., 
\begin{eqnarray}\label{Kos v}
K_v^*\cong\mathrm{Kos}(\hat E;v_1id+t_1^{\epsilon_1}\Theta_1,\cdots,v_nid+t_n^{\epsilon_n}\Theta_n),
\end{eqnarray}
where
\begin{eqnarray}\label{epsilon_i}
\epsilon_i:=\left\{
\begin{matrix}
 0,&i\leq r~\mathrm{or}~i>r~\mathrm{and}~v_i=0,\\
 1,&~i>r~\mathrm{and}~v_i\neq 0
\end{matrix}
\right.
\end{eqnarray}
and $v=(v_1,\cdots,v_n)$. Since each $\Theta_i$ is topologically nilpotent and at least one of $v_i$ is non-zero for $v\neq0$, we obtain that at least one of $v_iid+t_i^{\epsilon_i}\Theta_i$ is an automorphism of $\hat E$ for $v\neq0$. It follows that for $v\neq0$,  the right hand side of \eqref{Kos v} is acyclic, and hence its left hand side is also acyclic.
\end{proof}
The discussion above shows that the composite of 
\begin{eqnarray}\label{varphi tilde F}
\Omega^*(\hat E,\Theta)\stackrel{\varphi_{\tilde F}}{\to}K_0^*\subset\bigoplus_{v\in\mathbb F_p^n}K_v^*\cong F_*\Omega^*C_{\tilde F}^{-1}(\hat E,\Theta)
\end{eqnarray}
is a quasi-isomorphism.

\emph{Case of $\star=\Omega^*_{g^{-1}}$}. Recall that $g=t_1\cdots t_s$, $D=(t_1\cdots t_r=0)$ and $s\leq r$. By definition, we have a natural identification
$$
F_*\Omega^q_{g^{-1}}C^{-1}_{\tilde F}(\hat E,\Theta)\cong\hat E\otimes F_*\Omega^q_{g^{-1},\mathfrak X/k}(\log\mathfrak D).
$$
Since $\Omega^\bullet_{g^{-1},\mathfrak X/k}(\log\mathfrak D)$ is generated by
$$
d\log g\wedge^\bullet\{\omega_2,\cdots,\omega_n\},~g\wedge^\bullet\{\omega_2,\cdots,\omega_n\}
$$
over $\mathcal O_{\mathfrak X}$, there is an identification
$$
F_*\Omega^\bullet_{g^{-1}}C_{\tilde F}^{-1}(\hat E,\Theta)\cong\hat E\{t^{\boldsymbol{\alpha}}d\log g\wedge\omega_I,gt^{\boldsymbol{\alpha}}\omega_J:\boldsymbol{\alpha}\in\{0,\cdots,p-1\}^n,~I,J\subset\{2,\cdots,n\}\}.
$$
\begin{definition}
Let $B_g$ be the set consisting of all $t^{\boldsymbol{\alpha}}d\log g\wedge\omega_I$ and $gt^{\boldsymbol{\alpha}}\omega_J$. Define a function 
$$
V_g:B_g\to\mathbb F_p^n,~
t^{\boldsymbol{\alpha}}d\log g\wedge\omega_I\mapsto \boldsymbol{\alpha}+\boldsymbol{\alpha}_I,~t^{\boldsymbol{\alpha}}g\omega_J\mapsto\boldsymbol{\alpha}+\boldsymbol{\alpha}_I+\boldsymbol{\alpha}_g,
$$
where $\boldsymbol{\alpha}_I$ is defined as Definition \ref{Def w} and $\boldsymbol{\alpha}_g=(1,\cdots,1,0,\cdots,0)$ with the first $s$ components are $1$.
\end{definition}
By a similar argument as the case of $\star=\Omega^*$, we have
\begin{lemma}
There is a decomposition of complexes
$$
F_*\Omega_{g^{-1}}^*C_{\tilde F}^{-1}(\hat E,\Theta)\cong\bigoplus_{v\in\mathbb F_p^n}K_{g,v}^*,
$$
where $K_{g,v}^\bullet=\hat EV^{-1}_g(v)$ and 
$$K_{g,v}^*\cong\mathrm{Kos}(\hat E;\vartheta_1,\cdots,\vartheta_n),~\vartheta_i:=\left\{
\begin{matrix}
 v_1id+t_1^{\tau_1}\cdots t_s^{\tau_s}\Theta_1,&i=1,\\
 (v_i-v_1)id+t_i^{\epsilon_i}\Theta_i,&1<i\leq s,\\
 v_iid+t_i^{\epsilon_i}\Theta_i,&i>s.
\end{matrix}
\right.
$$
Here $\epsilon_i$ is defined as \eqref{epsilon_i} and
$$
\tau_i:=\left\{
\begin{matrix}
1,&v_i=0,\\
0,&v_i\neq0.
\end{matrix}
\right.
$$
Moreover, \eqref{varphi star=omega} induces an isomorphism of complexes
$$
\varphi_{g,\tilde F}:\Omega^*_{g^{-1}}(\hat E,\Theta)\to K^*_{g,0}
$$
and $K^*_{g,v}$ is acyclic for $v\neq0$.
\end{lemma}
Clearly, the lemma above shows that the composite of
$$\Omega_{g^{-1}}^*(\hat E,\Theta)\stackrel{\varphi_{g,\tilde F}}{\to}K_{g,0}^*\subset\bigoplus_{v\in\mathbb F_p^n}K_{g,v}^*\cong F_*\Omega_{g^{-1}}^*C_{\tilde F}^{-1}(\hat E,\Theta)$$ 
is a quasi-isomorphism.

\emph{Case of $\star=\Omega^*_{int}$}. For any $\boldsymbol{w}\in\{0,1\}^r$, set
$$
E_w:=\sum_{\boldsymbol{w}^1+\boldsymbol{w}^2=\boldsymbol{w},w_1,w_2\in\{0,1\}^r}t^{\boldsymbol{w}_1}\Theta^{\boldsymbol{w}_2}\hat E.
$$
Here
$$
t^{\boldsymbol{w}}:=t_1^{w_1}\cdots t_r^{w_r},~\Theta^{\boldsymbol{w}}:=\Theta_1^{w_1}\cdots \Theta_r^{w_r},~\boldsymbol{w}=(w_1,\cdots,w_r)\in\{0,1\}^r.
$$
Define 
$$W:B\to \{0,1\}^r,~t^{\boldsymbol{\alpha}}\omega_I\mapsto(\delta_0^{\alpha_1},\cdots,\delta_0^{\alpha_r}),~\boldsymbol{\alpha}=(\alpha_1,\cdots,\alpha_n)\in\{0,\cdots,p-1\}^n,$$
where $\delta_i^j:=1$ for $i=j$ and $0$ otherwise. 
By \cite[Lemma 3.14]{SZ2}, there is a decomposition
$$
F_*\Omega^\bullet_{int}C_{\tilde F}^{-1}(\hat E,\Theta)\cong\bigoplus_{\beta\in B}\hat E_{W(\beta)}.
$$
Oh the other hand, we have $B=\cup_{v\in\mathbb F_p^n}V^{-1}(v)$. It follows that
\begin{lemma}
There is decomposition of complexes
$$
F_*\Omega^*_{int}C_{\tilde F}^{-1}(\hat E,\Theta)\cong\bigoplus_{v\in\mathbb F_p^n}K^*_{int,v}
$$
with $K^\bullet_{int,v}=\bigoplus_{\beta\in V^{-1}(v)}\hat E_{W(\beta)}$. Moreover, $\varphi_{\tilde F}$ induces an isomorphism of subcomplexes
$$
\varphi_{int,\tilde F}:\Omega^*_{int}(\hat E,\Theta)\to K^*_{int,0}
$$
and $K_{v,int}^*$ is acyclic for $v\neq0$.
\end{lemma}
Clearly, the lemma above shows that the composite of
$$\Omega_{int}^*(\hat E,\Theta)\stackrel{\varphi_{int,\tilde F}}{\to}K_{int,0}^*\subset\bigoplus_{v\in\mathbb F_p^n}K_{int,v}^*\cong F_*\Omega_{int}^*C_{\tilde F}^{-1}(\hat E,\Theta)$$ 
is a quasi-isomorphism.

\emph{Case of $\star=W_q\Omega^*$}. Recall that in this situation,  we have 
$$(E,\theta)\in\mathrm{HIG}(X/k)\subset\mathrm{HIG}((X,D)/k).$$
Write 
$$\theta=\sum_{i=1}^n\theta'_i\otimes dt_i=\sum_{i=1}^rt_i\theta'_i\otimes d\log t_i+\sum_{i=r+1}^n\theta'_i\otimes dt_i.$$
For any $\emptyset\neq I\subset\{1,\cdots,r\}$, set $D_I$ to be the scheme-theoretic intersection of $D_i,i\in I$. Let $\theta_{D_I}$ be the Higgs field on $E_{D_I}:=E|_{D_I}$ induced by $\theta$, namely
$$
\theta_{D_I}:=\sum_{i\in I^c}\theta'_{i,D_I}\otimes dt_{i,D_I},~\theta'_{i,D_I}:=\theta'_i|_{D_I},~t_{i,D_I}:=t_i|_{D_I},~I^c:=\{1,\cdots,n\}-I.
$$
Set $\mathfrak D_I:=D_{I/Z\cap D_I}$ and $f_{D_I}:=f|_{D_I}$. Let $\tilde F_{D_I}:\tilde D_I\to\tilde D_I$ be the Frobenius lifting of $F:D_I\to D_I$ induced by $\tilde F$. In Step 1, replacing the data 
$$
t_i,i\in\{1,\cdots,n\},~f,~(E,\theta),~Z,~\tilde F
$$
by
$$
t_{i,D_I},i\in I^c,~(E_{D_I},\theta_{D_I}),~Z\cap D_I,~\tilde F_{D_I},
$$
then we get a Higgs filed $\Theta_{D_I}$ on $\hat E_{D_I}$.

By abuse of notation, denote by $\varphi_{\tilde F}$ again the composite of morphisms in \eqref{varphi tilde F}. For any $q\geq0$, it is easy to see that $\varphi_{\tilde F}$ induces a morphism of subcomplexes
$$
W_q\varphi_{\tilde F}:W_q\Omega^\bullet_{\log}(\hat E,\Theta)\to W_q\Omega^\bullet_{\log}C^{-1}_{\tilde F}(\hat E,\Theta).
$$
Moreover, we have the following diagram
$$
\xymatrix{
W_{q-1}\Omega^*_{\log}(\hat E,\Theta)\ar[d]^-{W_{q-1}\varphi_{\tilde F}}\ar[r]&W_q\Omega^*_{\log}(\hat E,\Theta)\ar[d]^-{W_q\varphi_{\tilde F}}\ar[r]^-{\mathrm{Res}^H_q}&\bigoplus_{|I|=q}\Omega^*_{\mathfrak D_I/k}(\hat E_{D_I},\Theta_{D_I})\ar[d]^-{\bigoplus_{|I|=q}\varphi_{\tilde F_{D_I}}}[-q]\\
W_{q-1}\Omega^*_{\log}C^{-1}_{\tilde F}(\hat E,\Theta)\ar[r]&W_q\Omega^*_{\log}C^{-1}_{\tilde F}(\hat E,\Theta)\ar[r]^-{\mathrm{Res}^D_q}&\bigoplus_{|I|=q}\Omega^*_{\mathfrak D_I/k}C^{-1}_{\tilde F_{D_I}}(\hat E_{D_I},\Theta_{D_I})[-q].
}
$$
Here $\mathrm{Res}^D_q$ is defined similar to $\mathrm{Res}^H_q$ and the later is the $q$-th surjective residue map defined by 
$$
\hat e\otimes(\sum_{|I|=q}\omega_I\wedge\beta_I+\beta_{q-1})\mapsto((-1)^q\beta_{I,D_I})_{|I|=q},
$$
where $\beta_I\in\Omega^\bullet_{\mathfrak X/k}(\hat E,\Theta)$ and $\beta_{q-1}\in W_{q-1}\Omega^\bullet_{(\mathfrak X,\mathfrak D)/k}(\hat E,\Theta)$. Since we have proved that $\varphi_{\tilde F_{D_I}}$ is a quasi-isomorphism, coupled with the $5$-lemma, we conclude that $W_q\varphi_{\tilde F}$ is a quasi-isomorphism. 

\iffalse
Let $\psi$ be the $p$-curvature of $\nabla$. Endow each term of the  $F$-Higgs complex 
$$
\Omega^*(H,-\psi)=[H\stackrel{-\psi}{\longrightarrow}H\otimes  F^*\Omega^1_{X/k}(\mathrm{log}~D)\stackrel{-\psi\wedge}{\longrightarrow}\cdots\stackrel{-\psi\wedge}{\longrightarrow}H\otimes  F^*\Omega^n_{X/k}(\mathrm{log}~D)]$$
with the tensor product connection, using the connection on $H$ and the Cartier descent connection $\mathrm{id}\otimes d$ on $F^*\Omega^i_{X/k}(\mathrm{log}~D)$. According to the following basic property of $p$-curvatures
$$
\psi(D_1)\nabla(D_2)=\nabla(D_2)\psi(D_1),~D_1,D_2\in T_{X/k}(-\log D),
$$
it is easy to see that there is a natural morphism of Lie algebras
$$
T_{X/k}(-\log D)\to\mathcal End_{C(\mathcal O_X)}(\Omega^*(H,-\psi)).
$$
In particular, for any $i$ one obtains a connection on $\mathcal H^i(\Omega^*(H,-\psi))$ with logarithmic pole along $D$, which is denoted by $\nabla$ again.

By the very construction of $(H,\nabla)$, Ogus-Vologodsky\cite[Theorem 2.13]{OV} shows that 
$$
(H,\psi)\cong F^*(E,-\theta).
$$
It follows that
$$
\Omega^*(H,-\psi)\cong F^*\Omega^*(E,\theta).
$$
On the other hand, by a general result of Ogus\cite[Theorem 1.2.1]{Ogus04}, we have 
\begin{eqnarray}\label{Ogus iso}
\mathcal H^i(\Omega^*(H,\nabla))\cong\mathcal H^i(\Omega^*(H,-\psi))^\nabla,~\forall i.
\end{eqnarray}
Combing these results, we demonstrate that there is a connection $\nabla$ on $F^*\mathcal H^i(\Omega^*(E,\theta))$ induced by the one on $\mathcal H^i(\Omega^*(H,-\psi))$ (it does not necessarily coincide with the Cartier descent connection)





\begin{remark}\label{explicit construction}
Assume that there is a global system of coordinates $(t_1,\cdots,t_n)$ on $X/k$ such that $D=(t_1\cdots t_n=0)$. Set
\begin{eqnarray}\label{de Rham F-Higgs}
\begin{array}{c}
\eta:\Omega^*(H,\nabla)\to\Omega^*(H,-\psi),\\
\eta^i: H\otimes\Omega_{X/k}^i(\mathrm{log}~D)\to H\otimes F^*\Omega_{X'/k}^i(\mathrm{log}~D'),\\
h\otimes\omega_I\mapsto\prod_{i\in I}(\mathrm{id}-\nabla_i^{p-1})(h)\otimes F^*\omega'_I
\end{array}
\end{eqnarray}
and
\begin{eqnarray}\label{F-Higgs de Rham}
\begin{array}{c}
\varphi: \Omega^*(H,-\psi)\to\Omega^*(H,\nabla),\\
\varphi^i: H\otimes F^*\Omega_{X'/k}^i(\mathrm{log}~D')\to H\otimes\Omega_{X/k}^i(\mathrm{log}~D),\\
h\otimes F^*\omega'_I\mapsto\prod_{i\in\{1,\cdots,n\}-I}(\mathrm{id}-\nabla_i^{p-1})(h)\otimes\omega_I.
\end{array}
\end{eqnarray}
Here $\nabla_i:=\nabla(t_i\partial_{t_i})$, $\omega_\emptyset:=1$ and $\omega_I:=d\mathrm{log}~t_{i_1}\wedge\cdots\wedge d\mathrm{log}~t_{i_s}$ for $I=\{i_1,\cdots,i_s\}$ with $1\leq i_1<\cdots<i_s\leq n$. Let 
$$\mathcal R:=k[\nabla_1,\cdots,\nabla_n]\subset D_{X_{\log}/k}.$$ 
Endow $\Omega^*(H,\nabla)$ an $\mathcal R$-action as follows:
$$
R(h\otimes\omega_I)=(Rh)\otimes\omega_I,~R\in\mathcal R.
$$
One can check that $\psi,\varphi$ are $\mathcal R$-equivariant and $\mathcal R$ acts on $\mathcal H^i(\Omega^*(H,\nabla))$ as zero.

In particular, we have morphisms of $\mathcal R$-horizontal cohomology sheaves
\begin{eqnarray}\label{coordinate C}
\mathcal H^i(\eta):\mathcal H^i(\Omega^*(H,\nabla))\to\mathcal H^i(\Omega^*(H,-\psi))^\mathcal{R}
\end{eqnarray}
and
\begin{eqnarray}\label{coordinate C^{-1}}
\mathcal H^i(\varphi):\mathcal H^i(\Omega^*(H,-\psi))^\mathcal{R}\to\mathcal H^i(\Omega^*(H,\nabla)).
\end{eqnarray}
Clearly, $\mathcal H^i(\Omega^*(H,-\psi))^\mathcal{R}=\mathcal H^i(\Omega^*(H,-\psi))^\nabla$. Ogus proved that these morphisms coincide with $C_{X/k}$ and $C^{-1}_{X/k}$, respectively.
\end{remark}
Suppose that $(H,\nabla)=C^{-1}_{\mathcal X/\mathcal S}(E,\theta)$ for a nilpotent Higgs module $(E,\theta)$ on $(X',D')/k$ of level $\leq p-n-1$. OV provided another construction of \eqref{Ogus iso}, we briefly recall it below. 

{\itshape Step 1.} By the construction of $C^{-1}_{\mathcal X/\mathcal S}$, $(H,-\psi)$ can be obtained by gluing the family $\{(F^*E_\alpha,F^*\theta_\alpha)\}$ by $\{G_{\alpha\beta}\}$. 

{\itshape Step 2.} Endow each term of $\Omega^*(F^*E_\alpha,F^*\theta_\alpha)$ with the tensor product connection, using the connection $\nabla_\alpha$ on $F^*E_\alpha$ and the Cartier descent connection $\mathrm{id}\otimes d$ on $F^*\Omega^i_{U_\alpha'/k}(\mathrm{log}~D'|_{U'_\alpha})$.  It follows that the aforementioned complex can be regarded as complex of modules with integrable connection on $(X,D)/k$. One can further check that $G_{\alpha\beta}$ is a morphism of such complexes. 

{\itshape Step 3.} One can check that the integrable connection on
$$
\mathcal H^i(\Omega^*(F^*E_\alpha,F^*\theta_\alpha))\cong F^*\mathcal H^i(\Omega^*(E,\theta))|_{U'_\alpha}
$$ 
is the Cartier descend connection $\mathrm{id}\otimes d$. This fact follows from $\nabla_\alpha=\nabla_\mathrm{can}+(\mathrm{id}\otimes\zeta_\alpha)F^*\theta$ and $\theta(\partial')$ acts on $\mathcal H^i(\Omega^*(E,\theta))$ as zero for any $\partial'\in T_{X'_{\log}/k}$.
Therefore, its horizontal subsheaf
can be identified with $\mathcal H^i(\Omega^*(E,\theta))|_{U_\alpha'}$. Moreover, the transition morphism $G_{\alpha\beta}$ on $\mathcal H^i(\Omega^*(E,\theta))|_{U_{\alpha\beta}'}$ is the identity map.  Consequently, one has
$$
F_*\mathcal H^i(\Omega^*(H,-\psi))^\nabla\cong\mathcal H^i(\Omega^*(E,\theta)).
$$

{\itshape Step 4.} The de Rham-Higgs comparison of OV shows that 
$$
F_*\mathcal H^i(\Omega^*(H,\nabla))\cong\mathcal H^i(\Omega^*(E,\theta)).
$$
It follows that
$$
\mathcal H^i(\Omega^*(H,\nabla))\cong\mathcal H^i(\Omega^*(H,-\psi))^\nabla,
$$
which coincides with \eqref{Ogus iso}.

Next we extend the last three steps to an intersection version.  To this end, we need to introduce a new notion. Clearly, there is an inclusion $F^*\Omega_\mathrm{int}^*(E,\theta)\subset\Omega^*(F^*E,F^*\theta)$, from which we make the following
\begin{definition}
For any $F$-Higgs module $(E,\psi)$ on $(X,D)/k$, we define its intersection complex $\Omega_\mathrm{int}^*(E,\psi)$ by local systems of coordinates $(t_1,\cdots,t_n)$ on $(U,D|_U)/k$ such that $D|_U=(t_1\cdots t_n=0)$:
$$
\Omega^i_\mathrm{int}(E,\nabla):=\bigoplus_{I\subset\{1,\cdots,n\},|I|=i}(\sum_{J\subset I}t_J^p\psi_{I-J}E)\otimes F^*\omega'_I,
$$
where $t_J:=\prod_{j\in J}t_j$ and $\psi_{I-J}:=\prod_{i\in I-J}(\nabla(t_i\partial_{t_i})^p-\nabla(t_i\partial_{t_i}))$. 
\end{definition}
\begin{remark}
 (i) Let $(E,\theta)$ be a Higgs module on $(X',D')/k$. Then by definition one has
$$
F^*\Omega_\mathrm{int}^*(E,\theta)=\Omega_\mathrm{int}^*(F^*E,F^*\theta).
$$

(ii) If $X$ is a curve over $k$, then
$$
\Omega_\mathrm{int}^*(E,\psi)=[E\stackrel{\psi}{\longrightarrow}E\otimes F^*\Omega^1_{X'/k}+\mathrm{Im}(\psi)].
$$
\end{remark}
We claim that {\itshape Steps 2-4} are compatible with the definition. 

{\itshape Intersection version of step 2.} It means that 
$$
\Omega_\mathrm{int}^*(F^*E_\alpha,F^*\theta_\alpha)\subset\Omega^*(F^*E_\alpha,F^*\theta_\alpha)
$$
is a subcomplex of modules with integrable connection on $(X,D)/k$ and 
$$
G_{\alpha\beta}:\Omega_\mathrm{int}^*(F^*E_\beta,F^*\theta_\beta)|_{U_{\alpha\beta}}\to\Omega_\mathrm{int}^*(F^*E_\alpha,F^*\theta_\alpha)|_{U_{\alpha\beta}}
$$
is an isomorphism.
To check these facts, it suffices to observe that
$$
[\nabla_\alpha(\partial), F^*\theta(\partial')]=0,~[\nabla_\alpha(\partial),\lambda^p]=0,
$$
where $\partial\in T_{U_{\alpha,\log}/k},\partial'\in T_{U_{\alpha,\log}'/k},\lambda\in\mathcal O_{U_\alpha}$.

{\itshape Intersection version of step 3.} It means that the integrable connection on
$$
\mathcal H^i(\Omega_\mathrm{int}^*(F^*E_\alpha,F^*\theta_\alpha))\cong F^*\mathcal H^i(\Omega_\mathrm{int}^*(E,\theta))|_{U'_\alpha}
$$ 
is the Cartier descend connection $\mathrm{id}\otimes d$. This fact follows from
\begin{lemma}
$\theta(\partial')$ acts on $\mathcal H^i(\Omega_\mathrm{int}^*(E,\theta))$ as zero for any $\partial'\in T_{X'_{\log}/k}$.
\end{lemma}
\begin{proof}
The problem is local, we may assume that there is a global system of coordinates $(t_1,\cdots,t_n)$ on $(X,D)/k$ such that $D=(t_1\cdots t_n=0)$. Without loss of generality, we may assume that $\partial'=t_1'\partial_{t_1'}$. Recall that there is a homotopy of
$$
\theta(\partial'):\Omega^*(E,\theta)\to\Omega^*(E,\theta)
$$
defined as
$$
\begin{array}{c}
s_i:\Omega^i(E,\theta)\to\Omega^{i-1}(E,\theta),~0<i\leq n,\\
\sum_{|I|=i}e_I\otimes\omega'_I\mapsto\sum_{|J|=i-1,1\notin J}e_{\{1\}\cup J}\otimes\omega'_J,
\end{array}
$$
where $\omega_I$ is given as in Theorem \ref{Cartier Ogus}.
Clearly, $s_i(\Omega_\mathrm{int}^i(E,\theta))\subset\Omega_\mathrm{int}^{i-1}(E,\theta)$, it follows that
$$
\theta(\partial'):\Omega_\mathrm{int}^*(E,\theta)\to\Omega_\mathrm{int}^*(E,\theta)
$$
is homotopic to zero. This completes the proof.
\end{proof}

Therefore, the horizontal subsheaf of $\mathcal H^i(\Omega_\mathrm{int}^*(F^*E_\alpha,F^*\theta_\alpha))$ can be identified with $\mathcal H^i(\Omega_\mathrm{int}^*(E,\theta))|_{U_\alpha'}$.
 
{\itshape Intersection version of step 4.} It means that there are isomorphisms
$$
F_*\mathcal H^i(\Omega_\mathrm{int}^*(H,\nabla))\cong\mathcal H^i(\Omega_\mathrm{int}^*(E,\theta)),
$$
which are cohomological consequences of the de Rham-Higgs comparison of \cite{SZ2}.

By combing the discussions above, one has
\begin{theorem}\label{SZ Cartier}
If $(H,\nabla)=C^{-1}_{\mathcal X/\mathcal S}(E,\theta)$ for $(E,\theta)$ is a nilpotent Higgs module on $(X',D')/k$ of level $\leq p-n-1$, then there are natural isomorphisms
\begin{eqnarray}\label{Ogus int iso}
\mathcal H^i(\Omega^*_\mathrm{int}(H,\nabla))\cong\mathcal H^i(\Omega_\mathrm{int}^*(H,-\psi))^\nabla.
\end{eqnarray}
\end{theorem}
The isomorphisms can be generalized to a larger class of modules with connection on $(X,D)/k$, but not for all. The following is a counterexample.
\begin{example}
Let $X=\mathrm{Spec}(k[t]), D=(t=0).$ Set
$$
H:=\mathcal O_X/t^p\mathcal O_X,~\nabla:=d+\mathrm{id}\otimes d\log t.
$$
By the definition of intersection complexes, one has 
$$
\Omega^1_\mathrm{int}(H,\nabla)=\Omega^1(H,\nabla)\stackrel{\mathrm{id}\otimes t\partial_t}{\cong}H.$$
Note that 
$$F_*H\cong(\mathcal O_{X'}/t'\mathcal O_{X'})\{1,t,\cdots,t^{p-1}\},~(\mathrm{id}\otimes t\partial_t)\nabla=t\partial_t+\mathrm{id}.$$
It follows that
$$
F_*\mathcal H^1(\Omega^*_\mathrm{int}(H,\nabla))\cong F_*\mathrm{coker}(t\partial_t+\mathrm{id})\cong(\mathcal O_{X'}/t'\mathcal O_{X'})t^{p-1}\cong\mathcal O_{X'}/t'\mathcal O_{X'}.
$$
On the other hand, one can check that
$$
\psi=0,~\Omega^1_\mathrm{int}(H,-\psi)=0,~\mathcal H^1(\Omega^*_\mathrm{int}(H,-\psi))^\nabla=0.
$$
Consequently, $(H,\nabla)$ violates \eqref{Ogus int iso} for $i=1$.
\end{example}
\begin{definition}\label{good}
Let $(H,\nabla)$ be a module with integrable connection on $(X,D)/k$. We say it is {\bfseries good} if for each $x\in D$, there is an open subset $U\subset X$ containing $x$ such that $\nabla|_U=\Delta+\vartheta$. Here
$$
\Delta:H|_{U}\to H|_{U}\otimes\Omega^1_{U/k},~\vartheta:H|_{U}\to H|_{U}\otimes\Omega^1_{U/k}(\mathrm{log}~D|_U)
$$
satisfying $[\Delta(\partial),\vartheta(\partial)]=0$ and $\mathrm{id}-\vartheta(\partial)^{p-1}$ is an isomorphism for any $\partial\in T_{U_{\log}/k}$.
\end{definition}
Clearly, if $(H,\nabla)$ is given as in Theorem \ref{SZ Cartier}, then it is good.
\iffalse
\begin{lemma}
\eqref{Ogus curve iso} coincides with \eqref{Cartier int}.
\end{lemma}
\begin{proof}
The problem is local, we may assume that $X/k$ admits a global coordinate $t$ such that $D=(t=0)$ and there is a global logarithmic lifting $\tilde F$ of $F$. Since the case of $i=0$ is obvious, it suffices to take $i=1$. Clearly, one has 
$$(H,\nabla)\cong(F^*E,\nabla_\mathrm{can}+(\mathrm{id}\otimes\zeta_{\tilde F})F^*\theta).$$
\begin{eqnarray}\label{de Rham F-Higgs int}
\alpha_\mathrm{int}:\Omega_\mathrm{int}^*(H,\nabla)\to\Omega_\mathrm{int}^*(H,-\psi)
\end{eqnarray}
and \eqref{Cartier int} is the morphism of $\mathcal H^i(\alpha_\mathrm{int})$ onto its image.
\end{proof}

 Here
$
\Delta=\sum_{i=1}^n\Delta_i\otimes dt_i,~\vartheta=\sum_{i=1}^n\vartheta_i\otimes d\mathrm{log}~t_i.
$
\fi
\begin{theorem}\label{Ogus comp int}
If $(H,\nabla)$ is good, then \eqref{Ogus int iso} holds locally.
\end{theorem}
\begin{proof}
Suppose that there is a global system of coordinates $(t_1,\cdots,t_n)$ on $X/k$ such that $D=(t_1\cdots t_n=0)$. Keep the notations in Remark \ref{explicit construction}. In priori, we have inclusions
$$
\varphi\eta(\Omega^*_\mathrm{int}(H,\nabla))\subset\Omega^*_\mathrm{int}(H,\nabla),~\eta\varphi(\Omega^*_\mathrm{int}(H,-\psi))\subset\Omega^*_\mathrm{int}(H,-\psi).
$$
On the other hand, one can check that the homotopy of 
$$
\mathrm{id}-\varphi\eta:\Omega^*_\mathrm{int}(H,\nabla)\to\Omega^*_\mathrm{int}(H,\nabla)
$$
given by Ogus restricts to a homotopy of
$$
\mathrm{id}-\varphi\eta:\Omega^*_\mathrm{int}(H,\nabla)\to\Omega^*_\mathrm{int}(H,\nabla).
$$
It follows that $\mathcal H^i(\varphi\eta)$ acts on $\mathcal H^i(\Omega_\mathrm{int}^*(H,\nabla))$ as identity. Clearly, $\mathcal H^i(\eta\varphi)$ acts on $\mathcal H^i(\Omega_\mathrm{int}^*(H,-\psi))^\nabla$ as identity. 

We claim that there are inclusions 
\begin{eqnarray}\label{inclusion harder}
\eta(\Omega^*_\mathrm{int}(H,\nabla))\subset\Omega_\mathrm{int}^*(H,-\psi).
\end{eqnarray}
and
\begin{eqnarray}\label{inclusion easer}
\varphi((\Omega_\mathrm{int}^*(H,-\psi)))\subset\Omega^*_\mathrm{int}(H,\nabla).
\end{eqnarray}
Note that the subcomplexes
$$
\Omega^*_\mathrm{int}(H,\nabla)\subset\Omega^*(H,\nabla),~\Omega^*_\mathrm{int}(H,-\psi)\subset\Omega^*(H,-\psi)
$$
are $\mathcal R$-closed and $\mathcal R$ acts on $\mathcal H^i(\Omega_\mathrm{int}^*(H,\nabla))$ as zero. In particular, \eqref{inclusion harder} and \eqref{inclusion easer} imply that there are morphisms of $\mathcal R$-horizontal cohomology sheaves 
$$
\mathcal H^i(\eta):\mathcal H^i(\Omega_\mathrm{int}^*(H,\nabla))\to\mathcal H^i(\Omega_\mathrm{int}^*(H,-\psi))^\nabla
$$
and
$$
\mathcal H^i(\varphi):\mathcal H^i(\Omega_\mathrm{int}^*(H,-\psi))^\nabla\to\mathcal H^i(\Omega_\mathrm{int}^*(H,\nabla)).
$$
By combing them with the assertions obtained in previous paragraph, the theorem follows. It remains to show  \eqref{inclusion harder} and \eqref{inclusion easer}. Let $\nabla=\Delta+\vartheta$ be a decomposition as in  Definition \ref{good}. Set 
$$
\Delta_i:=\Delta(t_i\partial_{t_i}),~\vartheta_i:=\vartheta(t_i\partial_{t_i}),~\psi_i:=\psi(F^*(t'_i\partial_{t'_i})).
$$

{\itshape The proof of \eqref{inclusion harder}.}  It is equivalent to show that
$$
\prod_{i\in I}(\mathrm{id}-\nabla_i^{p-1})(t_{I-J}\nabla_JH)\subset\sum_{K\subset I}t_{I-K}^p\psi_KH.
$$
One can check that this inclusion is derived from
\begin{eqnarray}\label{intersection eta}
(\mathrm{id}-\nabla_i^{p-1})t_iH\subset t_i^pH+\psi_iH.
\end{eqnarray}
Since 
\begin{eqnarray}\label{decomposition i}
\nabla_i=t_i\Delta_i+\vartheta_i,
\end{eqnarray}
 coupled with Deligne's formula, we get
$$
\psi_i=\nabla_i^p-\nabla_i=t_i^p\Delta_i^p+\vartheta_i^p-\vartheta_i.
$$
According to the second assumption on $\nabla$ and $H$ is coherent, we conclude that 
$$(\vartheta_i^{p-1}-\mathrm{id})H=H$$ 
around $D$. It follows that 
$$
t_i^pH+\psi_iH=t_i^p+\vartheta_iH.
$$
By combing the equality and \eqref{decomposition i} again, \eqref{intersection eta} can be reduced to
\begin{eqnarray}\label{t_i H}
(\mathrm{id}-(t_i\Delta_i)^{p-1})(t_iH)\subset t_i^pH+\psi_iH.
\end{eqnarray}
 
Consider an operator 
$$P_i=\sum_{j=0}^{p-1}\frac{(-t_i)^j}{j!}\Delta_i^j\in\mathrm{End}_k(H).
$$
One can check the following properties
$$t_i\Delta_iP_i\equiv0,~\mathrm{id}\equiv\sum_{j=0}^{p-1}\frac{t_i^j}{j!}P_i\Delta_i^j~\mod~t_i^p\mathrm{End}_k(H).$$
Using these properties, we have
$$
\begin{array}{rcl}
(\mathrm{id}-(t_i\Delta_i)^{p-1})t_i&\equiv&\sum_{j=0}^{p-1}(\mathrm{id}-(t_i\Delta_i)^{p-1})\frac{t_i^{j+1}}{j!}P_i\Delta_i^j\\
&\equiv&\sum_{j=0}^{p-2}(\frac{t_i^{j+1}}{j!}-\frac{(\mathrm{ad}_{t_i\Delta_i})^{p-1}(t_i^{j+1})}{j!})P_i\Delta_i^j\\
&\equiv&\sum_{j=0}^{p-2}(1-(j+1)^{p-1})\frac{t_i^{j+1}}{j!}P_i\Delta_i^j\\
&\equiv&0~\mod~t_i^p\mathrm{End}_k(H),
\end{array}
$$
from which \eqref{t_i H} follows. This completes the proof of \eqref{inclusion harder}.

{\itshape The proof of \eqref{inclusion easer}.} It follows immediately from the definitions of intersection complexes. 
\end{proof}
Ogus proved that the two universally cohomological $\partial$-functors from $\mathrm{MIC}(X_{\log}/k)$ to $F$-$\mathrm{HIG}(X_{\log}/k)$
$$\{\mathcal H_{DR}^i\circ\Omega^*: i\in\mathbb N\},~\{\mathcal H^{i\nabla}_{-\psi}\circ\Omega^*:i\in\mathbb N\}$$  
are isomorphic. It guarantees that the isomorphisms \eqref{Ogus int iso} induced by \eqref{de Rham F-Higgs} is independent of choices of coordinate systems. Unfortunately, the following example shows that $\Omega^*_\mathrm{int}$ is not exact (similar phenomenon in complex Hodge theory is well-known, namely the intermediate extension is not exact.) It leads to $\{\mathcal H_{DR}^i\circ\Omega_\mathrm{int}^*: i\in\mathbb N\}$ is not even a well-defined cohomological $\partial$-functors. We believe that the theorem above is globally true, but it can not be shown by the aforementioned trick.
\begin{example} 
Let $X=\mathrm{Spec}(k[t]), ~D=(t=0)$. Set 
$$
(H',\nabla')=(H'',\nabla''):=(\mathcal O_X,d),~(H,\nabla)\in\mathrm{MIC}(X_{\log}/k).
$$
Here $H=\mathcal O_X\bigoplus\mathcal O_X$ and 
$$
\nabla(x,y)=(dx+yd\log t,dy),~x,y\in\mathcal O_X.
$$
Consider the following exact sequence
$$
0\to(H',\nabla')\to(H,\nabla)\to(H'',\nabla'')\to0,
$$
where the first (resp. second) morphism is given by $x\mapsto(x,0)$ (resp. $(x,y)\mapsto y$). By direct computation, one has
$$
\Omega^1_\mathrm{int}(H',\nabla')=\Omega^1_\mathrm{int}(H'',\nabla'')=\Omega^1_{X/k},~\Omega^1_\mathrm{int}(H,\nabla)=\Omega^1_{X_{\log}/k}\bigoplus\Omega^1_{X/k}.
$$
Consequently, the sequence 
$$
0\to\Omega^1_\mathrm{int}(H',\nabla')\to\Omega^1_\mathrm{int}(H,\nabla)\to\Omega^1_\mathrm{int}(H'',\nabla'')\to0
$$
is not exact.
\end{example}
Nevertheless, for $\mathcal H^0,\mathcal H^1$, one has
\begin{corollary}
If $(H,\nabla)$ is good, then \eqref{Ogus int iso} holds globally for $i=0,1$.
\end{corollary}
\begin{proof}
It suffices to observe that 
$$\mathcal H^1(\Omega_\mathrm{int}^*(H,\nabla))\subset\mathcal H^1(\Omega_\mathrm{int}^*(H,\nabla)),~\mathcal H^1(\Omega_\mathrm{int}^*(H,-\psi))\subset\mathcal H^1(\Omega_\mathrm{int}^*(H,-\psi)).$$
\end{proof}
\fi

\section{De Rham-Higgs comparison ($df=0$)}
Let $k$ be a perfect field of characteristic $p$, let $X$ be a smooth variety of dimension $n$ over $k$, and let $D\subset X$ a simple normal crossing divisor on $X$. Assume that $(X,D)/k$ admits a $W_2(k)$-lifting, say $(\tilde X,\tilde D)/W_2(k)$. Let $(E,\theta)$ be a nilpotent Higgs sheaf on $(X,D)/k$ of level $\leq\ell$ ($\ell<p$) and set $(H,\nabla):=C^{-1}_{(\tilde X,\tilde D)/W_2(k)}(E,\theta)$. Then we have the following \emph{de Rham-Higgs comparison}, which vastly generalizes the decomposition theorem of Deligne-Illusie\cite[Théorème 2.1]{DI}.
\begin{theorem}[{{Ogus-Vologodsky\cite[Corollary 2.27]{OV}, Schepler\cite[Corollary 5.7]{Schepler}}}]\label{OV}
In the derived category $D(X)$, we have
$$
\tau_{<p-\ell}F_*\Omega^*(H,\nabla)\cong\tau_{<p-\ell}\Omega^*(E.\theta).
$$   
\end{theorem} 
As shown by Petrov\cite[Corollary 10.7]{P1}, the truncation $\tau_{<p-\ell}$ is necessary. However, to derive nice geometric consequences, such as the $E_1$-degeneration and vanishing theorem, this truncation is not expected to appear. Obviously, it can be automatically dropped when $\dim X<p-\ell$. Recently, a beautiful result of Petrov in \cite[Corollary 4.6]{P2} demonstrated that $\tau_{<p}$ can be dropped, if $X$ is quasi-$F$-split, $D=\emptyset$, and $(E,\theta)=(\mathcal O_X,0)$. In this section, we show that this truncation can be dropped by additionally requiring that the logarithmic cotangent bundle $\Omega^1_{X/k}(\log D)$ splits into a direct sum of subbundles of rank $< p-\ell$. 

\subsection{Higher homotopies}
Theorem \ref{OV} was re-proven by Sheng and the author\cite{SZ2} using a different approach, inspired by the $\mathrm{\check{C}}$ech construction of Deligne-Illusie\cite{DI}. More precisely, we constructed an isomorphism between $\tau_{<p-\ell}\Omega^*(E,\theta)$ and $\tau_{<p-\ell}F_*\Omega^*(H,\nabla)$ in the derived category $D(X)$ as follows:
$$
\tau_{<p-\ell}\Omega^*(E,\theta)\xrightarrow{\varphi_\mathrm{SZ}}\check{\mathcal C}(\mathcal U,\tau_{<p-\ell}F_*\Omega^*(H,\nabla))\xleftarrow{\rho}\tau_{<p-\ell}F_*\Omega^*(H,\nabla).
$$
Here $\mathcal U=\{U_i\}_{i\in I}$ is an open affine covering of $X$, $\rho$ is the natural augmentation morphism and $\varphi$ is a quasi-isomorphism of complexes. The construction of $\varphi$ is intricate, we outline it below. For any $i\in I$, we take a Frobenius lifting $\tilde F_{U_i}:\tilde U_i\to\tilde U_i$ of $F_{U_i}:U_i\to U_i$ such that $\tilde F^*_{\tilde U}(\tilde D|_{\tilde U})=p\tilde D|_{\tilde U}$. As considered in \eqref{varphi star=omega}, we construct a local Cartier quasi-isomorphism
$$
\varphi_{\tilde F_{U_i}}:\Omega^*(E,\theta)|_{U_i}\to F_*\Omega^*(H,\nabla)|_{U_i},~e\otimes\omega_1\wedge\cdots\wedge\omega_s\mapsto \iota_{\tilde F_{U_i}}(F^*e)\otimes\zeta_{\tilde F_{U_i}}(\omega_1)\wedge\cdots\wedge\zeta_{\tilde F_{U_i}}(\omega_s).
$$
 If we only regard $\varphi$ as a morphism of graded $\mathcal O_X$-modules, then
$$
\varphi_\mathrm{SZ}=\bigoplus_{r=0}^\infty\varphi_\mathrm{SZ}^r,~\varphi_\mathrm{SZ}^r:\tau_{<p-\ell}\Omega^\bullet(E,\theta)\to\check{\mathcal C}^r(\mathcal U,\tau_{<p-\ell}F_*\Omega^\bullet(H,\nabla)[-r]).
$$
In particular, $\varphi^0$ is induced by the family $\{\tau_{<p-\ell}\varphi_{\tilde F_{U_i}}\}_{i\in I}$. We say \emph{the family $\{\varphi_\mathrm{SZ}^r\}_{r>0}$ is a family of higher homotopies of $\varphi^0_{\mathrm{SZ}}$.}

Next, we further assume that the logarithmic cotangent bundle $\Omega^1_{X/k}(\log D)$ splits into a direct sum of subbundles of rank $<p-\ell$:
\begin{eqnarray}\label{splitting cotangent bundle}
\Omega=\bigoplus_{i=1}^\beta\Omega_i,~\mathrm{rank}(\Omega_i)<p-\ell,~\Omega:=\Omega^1_{X/k}(\log D).
\end{eqnarray}
Using this splitting, we construct a quasi-isomorphism of complexes
\begin{eqnarray}\label{Cech quasi-iso without truncation}
\varphi_{\underline \Omega}:\Omega^*(E,\theta)\to\check{\mathcal C}(\mathcal U,F_*\Omega^*(H,\nabla)),~\underline\Omega:=(\Omega_1,\cdots,\Omega_\beta).
\end{eqnarray}
In particular, it gives rise to a desired isomorphism between $\Omega^*(E,\theta)$ and $F_*\Omega^*(H,\nabla)$ (without truncation) in the derived category $D(X)$. Similar to $\varphi_\mathrm{SZ}$, $\varphi_Z^0$ is induced  by the family $\{\varphi_{\tilde F_{U_i}}\}_{i\in I}$, but its higher homotopies are different from $\varphi_\mathrm{SZ}$ in general.

We now give an informal but suggestive explanation for the existence of $\varphi_{\underline\Omega}$. To simplify the discussion, we may assume $(E, \theta) = (\mathcal{O}_X, 0)$ ($\ell=0$) and $D = \emptyset$. According to the splitting~\eqref{splitting cotangent bundle}, we may heuristically regard it as inducing a decomposition $X = \prod_{i=1}^\beta X_i$, such that each $\Omega_i$ is the pullback of the cotangent bundle of $X_i$. By the result of Deligne–Illusie~\cite{DI}, each de Rham complex $F_*\Omega^\bullet_{X_i/k}$ is decomposable, which in turn suggests that $F_*\Omega^\bullet_{X/k}$ is also decomposable.




\iffalse
In the construction of $\{\varphi_\mathrm{SZ}^r\}_{r>0}$, we need the canonical section of the projection $\pi:\Omega^1_{X/k}(\log D)^{\otimes q}\to\Omega^q_{X/k}(\log D)$, i.e.,
$$
\Omega^q_{X/k}(\log D)\to\Omega^1_{X/k}(\log D)^{\otimes q},~\omega_1\wedge\cdots\wedge\omega_q\mapsto\frac{1}{q!}\sum_{\sigma\in S_q}\mathrm{sgn}(\sigma)\omega_{\sigma(1)}\otimes\cdots\otimes\omega_{\sigma(q)}.
$$
The factor $1/q!$ partially answers why the truncation $\tau_{<p-\ell}$ is necessary. However, in the construction of $\{\varphi_Z^r\}_{r>0}$, we use a non-canonical section of the projection $\pi$. 
More precisely, for any splitting $\Omega^1_{X/k}(\log D)=\bigoplus_{i=1}^n \Omega_i$ with $\mathrm{rank}(\Omega_i)<p-\ell$, the decomposition
$$
\Omega^q_{X/k}(\log D)\cong\bigoplus_{q_1+\cdots+q_\beta=q,0\leq q_i\leq\mathrm{rank}(\Omega_i)}\wedge^{q_1}\Omega_1\otimes\cdots\otimes\wedge^{q_\beta}\Omega_\beta,~\wedge^0\Omega_i:=\mathcal O_X
$$
combined with the natural section
$$
\wedge^{q_i}L_i\to\otimes^{q_i}L_i,~\omega_1\wedge\cdots\wedge \omega_{q_i}\mapsto\frac{1}{q_i!}\sum_\sigma\mathrm{sgn}(\sigma)\omega_{\sigma(1)}\otimes\cdots\otimes \omega_{\sigma(q_i)},~\omega_1,\cdots,\omega_{q_i}\in \Omega_i,
$$
gives rise to a non-canonical section of $\pi$.
It is important to note that in this non-canonical section, the factor $1/q!$ is replaced by $1/q_i!$, enabling the possibility of a de Rham–Higgs comparison without truncation.\fi

To construct $\varphi_{\underline\Omega}$, we recall some notations from \cite[\S3.1]{SZ2}. Let $\mathcal L$ be the sheaf of local Frobenius liftings $\tilde F$ to $\tilde X$ such that $\tilde F^*(\tilde D)=p\tilde D$. 
Let $\Delta_\ast(\mathcal L)$ be the simplicial complex attached to $\mathcal L$, which is described as follows. For $r \geq 0$, $\Delta_r(\mathcal L)$ is the sheaf associated to the presheaf of abelian groups, which assigns to an open subset $U \subset X$ the free abelian group generated by elements of $\Gamma(U, \mathcal L^{r+1})$, where $\mathcal L^{r+1}$ is the direct product of $(r + 1)$ copies of $\mathcal L$. The boundary operator 
\[
\partial : \Delta_{r+1}(\mathcal L) \to \Delta_r(\mathcal L), \quad r \geq 0
\]
is given as usual:
\[
(\tilde F_0, \cdots, \tilde F_{r+1}) \mapsto \sum_{q=0}^{r+1} (-1)^q (\tilde F_0, \cdots, \widehat{\tilde F_q}, \cdots, \tilde F_{r+1}).
\]
Set $\Delta_r(\mathcal L) = 0$ for $r < 0$. For \( \mathcal F^\ast \), \(\mathcal  G^\ast \), two complexes of sheaves of \(\mathcal{O}_X\)-modules,  
we define \(\mathcal Hom^\ast_{\mathcal{O}_X}(\mathcal F^\ast, \mathcal G^\ast)\) to be the associated Hom complex.  
Explicitly, for \( r \in \mathbb{Z} \),  
\[
\mathcal Hom^r_{\mathcal{O}_X}(\mathcal F^\ast, \mathcal G^\ast) := \bigoplus_{i \in \mathbb{Z}} 
\mathcal Hom_{\mathcal{O}_X}(\mathcal F^i, \mathcal G^{i+r}),
\]
and
\[
d_{\mathcal Hom} : \mathcal Hom^r_{\mathcal{O}_X}(\mathcal F^\ast, \mathcal G^\ast) 
\to \mathcal Hom^{r+1}_{\mathcal{O}_X}(\mathcal F^\ast, \mathcal G^\ast)
\]
sends 
\[
\{f^i \in \mathcal Hom_{\mathcal{O}_Y}(\mathcal F^i, \mathcal G^{i+r})\}_{i \in \mathbb{Z}}
\]
to
\[
\{d_{\mathcal G^\ast} \circ (-1)^{i+r}f^i + (-1)^{i+r+1} f^{i+1} \circ d_{\mathcal F^\ast} \in 
\mathcal Hom_{\mathcal{O}_X}(\mathcal F^i,\mathcal G^{i+r+1})\}_{i \in \mathbb{Z}}.
\]


\begin{definition}
 An $\mathcal L$-indexed $\infty$-homotopy from $\mathcal F^*$ to $\mathcal G^*$ is a morphism of complexes of sheaves of abelian groups
\[
\mathrm{Ho} : \Delta_\ast(\mathcal L) \to \mathcal Hom^\ast_{\mathcal{O}_X}(\mathcal F^*, \mathcal G^*)
\]
such that $\mathrm{Ho}^0(\tilde F)$ is a quasi-isomorphism for any $\tilde F\in\mathcal L$.
\end{definition}

By \cite[Proposition 3.2]{SZ2}, the existence of $\varphi_{\underline\Omega}$ is equivalent to the existence of an $\mathcal L$-indexed homotopy 
\begin{eqnarray}\label{L-indexed homotopy theta to nabla}
\mathrm{Ho}_{\underline\Omega}: \Delta_\ast(\mathcal L) \to \mathcal Hom^\ast_{\mathcal{O}_X}(\Omega^*(E,\theta), F_*\Omega^*(H,\nabla)),~\tilde F\in\mathcal L\mapsto\varphi_{\tilde F}.
\end{eqnarray}


\begin{theorem}\label{de-Hig without truncation 2}
With the notation and assumptions introduced above Theorem \ref{OV}, suppose further that $\Omega^1_{X/k}(\log D)$ splits into a direct sum of subbundles of rank $<p-\ell$ as in \eqref{splitting cotangent bundle}. Then there is a desired $\mathcal L$-indexed $\infty$-homotopy \eqref{L-indexed homotopy theta to nabla}. In particular, in the derived category $D(X)$ we have
$$
 F_*\Omega^*(H,\nabla)\cong\Omega^*(E,\theta).
$$
 \end{theorem}



 
The construction of $\mathrm{Ho}_{\underline\Omega}$ requires some preliminaries. Suppose that $\Omega$ is free over $\mathcal O_X$ with basis $\omega_1,\cdots,\omega_n$, and that there is a nondecreasing surjective function 
\begin{eqnarray}\label{D}
D:\{1,\cdots,n\}\to\{1,\cdots,\beta\}
\end{eqnarray}
such that $\{\omega_j:j\in D^{-1}(i)\}$ forms a basis for $\Omega_i$ over $U$.  Write $\theta=\sum_{i=1}^n\theta_i\otimes\omega_i$. For any $r>0,s\geq0$ and any Frobenius liftings $\tilde F_0,\cdots,\tilde F_r\in\Gamma(X,\mathcal L)$, we will construct an $\mathcal O_X$-linear morphism of complexes
\begin{eqnarray}\label{varphi r s}
\varphi(r,s)_{\tilde F_0,\cdots,\tilde F_r}:\Omega^{r+s}(E,\theta)\to F_*\Omega^s(H,\nabla).
\end{eqnarray}
Its construction requires the following notation table.
\begin{itemize}
\item Let $T_{\underline\Omega}(r,s)$ be the set of all triples $(\underline i,\underline S,\underline{\overline j})$ satisfying the following conditions:
\begin{itemize}
    \item $\underline S:=(S_0,\cdots,S_r)$, where $S_0,\cdots,S_r\subset\{1,\cdots,n\}$ and 
    $$\#S_0+\cdots+\#S_r=\#(S_0\cup\cdots\cup S_r)=s;$$
    \item $\underline{\overline j}:=(j^1_1,\cdots,j^1_n;\cdots;j^r_1,\cdots,j^r_n)$, where $0\leq j^q_l<p$;
    \item $\lambda_1\geq\cdots\geq\lambda_r$, where $\lambda_q := \max D( S_q \cup \{ l : j^q_l \neq 0 \} \big)$;
    \item $\underline i:=(i_1,\cdots,i_r)\in D^{-1}(\lambda_1)\times\cdots\times D^{-1}(\lambda_r)$ and
    $$
    \#(\{i_1,\cdots,i_r\}\cup S_0\cup\cdots\cup S_r)=r+s;
    $$
    \item $j^q_{i_q}>0$ for any $1\leq q\leq r$.
\end{itemize}
    \item Let  \((\underline i,\underline S, \underline{\overline{j}})\in T_{\underline\Omega}(r,s) \).  For any \( 1 \leq i \leq \beta \), set  
\[
a_i := \max  \{q:\lambda_q = i \}  ,~ 
b_i := \min \{ q : \lambda_q = i \}.
\]
For any $1\leq q\leq r$, set 
$$s'_q := \# ( S_q \cap D^{-1}(\lambda_q)),$$
where we stipulate that $\mathrm{max}(\emptyset):=0$.
Now, set  
\begin{equation}\label{C}
C(\underline i,\underline S, \underline{\overline{j}}) := \prod_{i=1}^\beta c_i,
\end{equation}  
where  
\[
c_i := 
\begin{cases}
\prod_{b_i \leq q \leq a_i} \Big[ (j^q + s'_q) + \cdots + (j^{a_i} + s'_{a_i}) \Big]^{-1}, & \text{if } b_i > 0~\mathrm{and}~\sum_{q=b_i}^{a_i}(j^q+s_q')<p, \\
1, & \text{otherwise}.
\end{cases}
\]  
Here,  \( j^q := \sum_{l \in D^{-1}(D(\lambda_q))} j^q_l \) (we stipulate that $\sum_\emptyset:=0$). 

\item $\theta^{\underline{\overline j}}:=\theta_1^{j_1^1+\cdots+j^r_1}\cdots\theta_n^{j^1_n+\cdots+j^r_n}$. 
    \item $\zeta_{i,j}:=\zeta_{\tilde F_i}(F^*\omega_j)$ and $h_{i,j}:=h_{\tilde F_{i-1}\tilde F_i}(F^*\omega_j)$.
    \item $h^{[\underline{\overline j}]}:=h^{[j^1_1]}_{1,1}\cdots h^{[j^1_n]}_{1,n}\cdots h^{[j^r_1]}_{r,1}\cdots h^{[j^r_n]}_{r,n}$, where $a^{[j]}:=a^j/j!$.
    \item For \( S = \{i_1, \dots, i_m\} \subset \{1, \dots, n\} \) with \( i_1 < \cdots < i_m \), we define  
\[
\zeta_{0,S} := \zeta_{0,i_1} \wedge \cdots \wedge \zeta_{0,i_m}, \quad dh_{q,S} := dh_{q,i_1} \wedge \cdots \wedge dh_{q,i_m}.
\]  
By convention, we set \( \zeta_{0,\emptyset} := 1 \) and \( dh_{q,\emptyset} := 1 \). For subsets \( S_0, \dots, S_r \subset \{1, \dots, n\} \), define  
\[
\boldsymbol{\omega}_{\underline{S}} := \zeta_{0,S_0} \wedge dh_{1,S_1} \wedge \cdots \wedge dh_{r,S_r}, \quad \underline{S} = (S_0, \dots, S_r).
\]
\end{itemize}
\begin{construction}\label{construction higher homotopy}
Note that any section $\boldsymbol{e}\in\Omega^{r+s}(E,\theta)$ can be uniquely written as 
\begin{eqnarray}\label{section e}
\sum_{1\leq i_1<\cdots<i_{r+s}\leq n}e_{i_1,\cdots,i_{r+s}}\otimes(\omega_{i_1}\wedge\cdots\wedge\omega_{i_{r+s}}),~e_{i_1,\cdots,i_{r+s}}\in E.
\end{eqnarray}
Define
\begin{eqnarray}\label{key construction}
\varphi(r,s)_{\tilde F_0,\cdots,\tilde F_r}(\boldsymbol{e}):=\sum_{(\underline i,\underline S,\underline{\overline j})\in T_{\underline\Omega}(r,s)}C(\underline i,\underline S,\underline{\overline j})j^1_{i_1}\cdots j^r_{i_r}\iota_{\tilde F_0}(F^*(\theta^{\underline{\overline j}}\theta_{i_1}^{-1}\cdots\theta_{i_r}^{-1}e_{\underline S,\underline i}))(\boldsymbol{\omega}_{\underline S}h^{[\underline{\overline j}]}),
\end{eqnarray}
where $e_{\underline S,\underline i}:=e_{i^0_1,\cdots,i^0_{s_0},\cdots,i^r_1,\cdots,i^r_{s_r},i_1,\cdots,i_r}\footnote{By convention, we set \( e_{\sigma(1), \cdots, \sigma(r+s)} := \mathrm{sgn}(\sigma) e_{i_1, \cdots, i_{r+s}} \), where \( \sigma \) is a permutation of \( \{1, \dots, n\} \). If \( i_q = i_l \) for some \( q \neq l \), we set \( e_{i_1, \dots, i_{r+s}} := 0 \).
},
$ $\underline S=(S_0,\cdots,S_r), S_q=\{i^q_1,\cdots,i^q_{s_q}\}$ with $i^q_1<\cdots<i^q_{s_q}$ (if $S_q\neq\emptyset$).


\iffalse Be aware that \eqref{key construction} applies only when $s_0,\cdots,s_r>0$. If any $s_q=0$, it should be appropriately modified. For example, if $r=s=2$ and $s_0=s_1=0$, then \eqref{key construction} becomes 
$$
\sum Cj^1_{i_1}j^2_{i_2}\iota_{\tilde F_0}(F^*(\theta^{\underline{\overline j}}\theta_{i_1}^{-1}\theta_{i_2}^{-1}e_{i^2_1,i^2_2,i_1,i_2}))dh_{2;i^2_1,i^2_2}h^{[\underline{\overline j}]}.
$$
Here, the summation is taken over $i^2_1,i^2_2,i_1,i_2,\underline{\overline j}$, subject to the following conditions:
\begin{itemize}
    \item $i_q=\mathrm{max}\{j^q_l:1\leq l\leq n\}$;
    \item The indices $i^2_1,i^2_2,i_1,i_2$ are pairwise distinct and
    $$
     i^2_1<i^2_2,~D(i^2_2)\leq D(i_2),~D(i_1)\geq D(i_2).
    $$
\end{itemize}\fi
\end{construction}
\begin{lemma}\label{indep of basis}
The construction of $\varphi(r,s)_{\tilde F_0,\cdots,\tilde F_r}$ is independent of the choice of basis $\omega_1,\cdots,\omega_n$. 
\end{lemma}

Clearly, the splitting \eqref{splitting cotangent bundle} gives rise to a natural isomorphism
\begin{eqnarray}\label{splitting omega r+s}
\Omega^{r+s}\cong\bigoplus_{(i_1,\cdots,i_\beta)\in I_{r+s}}\bigwedge^{i_1}\Omega_1\otimes\cdots\otimes\bigwedge^{i_\beta}\Omega_\beta,
\end{eqnarray}
where $I_{r+s}$ consiting of $(i_1,\cdots,i_\beta)$ satisfying $i_1+\cdots+i_\beta=r+s$ and $0\leq i_l\leq\mathrm{rank}(\Omega_l)$. Since the rank of any $\Omega_q$ is strictly smaller than $p$, we have natural inclusions
\begin{eqnarray}\label{inclusion omega}
\bigoplus_{(i_1,\cdots,i_\beta)\in I_{r+s}}\bigwedge^{i_1}\Omega_1\otimes\cdots\otimes\bigwedge^{i_\beta}\Omega_\beta\subset\bigoplus_{(i_1,\cdots,i_\beta)\in I_{r+s}}\Omega_1^{\otimes i_1}\otimes\cdots\otimes\Omega_\beta^{\otimes i_\beta}\subset\Omega^{\otimes(r+s)}.
\end{eqnarray}
Denote by $\tilde\Omega^{\otimes(r+s)}$ the middle term of the above inclusions. Using the right inclusion, we obtain a projection $\pi:\tilde\Omega^{\otimes(r+s)}\to\Omega^{r+s}$. Combing \eqref{splitting omega r+s} and the left inclusion, we obtain a section of $\pi$, denoted by $\mathrm{sec}_{\underline\Omega}$, which is independent of the choice of the basis $\omega_1,\cdots,\omega_n$.

To prove this lemma, it suffices to lift $\varphi(r,s)_{\tilde F_0,\cdots,\tilde F_r}$ to $E\otimes\tilde\Omega^{\otimes(r+s)}$, namely a morphism
$$
\tilde\varphi(r,s)_{\tilde F_0,\cdots,\tilde F_r}:E\otimes\tilde\Omega^{\otimes(r+s)}\to F_*\Omega^s(H,\nabla)
$$
whose composite with $\mathrm{sec}_{\underline\Omega}$ is $\varphi(r,s)_{\tilde F_0,\cdots,\tilde F_r}$, and
which is independent of the choice of the basis $\omega_1,\cdots,\omega_n$ again. The following notation table is necessary for the construction of the lifting.
\begin{itemize}

    \item Let $\Theta_i$ be the composite of
    $$
    E\stackrel{\theta}{\to}E\otimes(\Omega_1\oplus\cdots\oplus\Omega_\beta)\xrightarrow{\mathrm{proj}_i} E\otimes\Omega_i.
    $$
    Let $H_i^q$ be the composite of
    $$
    F^*\Omega_i\hookrightarrow F^*\Omega\xrightarrow{h_q=h_{\tilde F_{q-1}\tilde F_q}}\mathcal O_X.
    $$
    Set $\sigma^q_i:=(id\otimes H^q_i)F^*\Theta_l\in\mathcal End_{\mathcal O_X}(F^*E)$.
    
    \item Let $U(r,s)$ be the set consisting of triples $(\overline{\underline j},\overline{\underline s},\underline s)$, subject to the following conditions:
\begin{itemize}
    \item $\overline{\underline j}=(j^1_1,\cdots,j^1_\beta;\cdots;j^r_1,\cdots,j^r_\beta)$, where $0\leq j^q_i<p$;
    \item $\underline s=(s_1,\cdots,s_\beta)$, where $s_1+\cdots+s_\beta=r$ and $s_i\geq0$;
 \item $\overline{\underline s}=(s^0_1,\cdots,s^0_\beta;\cdots;s^r_1,\cdots,s^r_\beta)$, where $s^q_i\geq 0$;
 \item $s^q_i=0$ for $1\leq i\leq \beta$ and $q>s_i+\cdots+s_\beta$, $\sum_{q=1}^rj^q_i+\sum_{q=0}^rs_i^q<p$ for any $1\leq i\leq \beta$, and 
     $$(\sum_{q=0}^rs^q_1+s_1,\cdots,\sum_{q=0}^rs^q_\beta+s_\beta)\in I_{r+s},~\sum_{q,i}s^q_i+\sum_l s_i=r+s.$$ 
       \end{itemize}    
    
 \item
    For any $s>0$ and any $0\leq q\leq r$, define 
 $$
 \wedge^sdh_q:\Omega^{\otimes s}\to F_*\Omega^s,~\xi_1\otimes\cdots\otimes\xi_s\mapsto dh_{\tilde F_{q-1}\tilde F_q}(F^*\xi_1)\wedge\cdots\wedge dh_{\tilde F_{q-1}\tilde F_q}(F^*),
 $$
 where $h_{\tilde F_{-1}\tilde F_0}:=\zeta_0$. Moreover, we stipulate that $\wedge^0dh_q$ is the Frobenius morphism $F^*:\mathcal O_X\to F_*\mathcal O_X$.
   
\item Given any $\underline s$ as in the construction of $U(r,s)$. Define 
$$
\wedge^rh:\Omega^{\otimes r}\to F_*\mathcal O_X,~\xi_1\otimes\cdots\otimes\xi_r\mapsto h_{\tilde F_0\tilde F_1}(F^*\xi_1)\cdots h_{\tilde F_{r-1}\tilde F_r}(F^*\xi_r)
$$
Moreover, we stipulate that $\wedge^0h$ is the Frobenius morphism $F^*:\mathcal O_X\to F_*\mathcal O_X$.
    
\item Given any $\overline{\underline j}$ as in the construction of $U(r,s)$. Set 
$$
\sigma^{\underline{\overline j}}:=\prod_{1\leq q\leq r,1\leq i\leq\beta}(\sigma^q_i)^{[j^q_i]}\in\mathcal End_{\mathcal O_X}(F^*E).
$$
\item Given any $(\underline{\overline s},\underline s)$ as in the construction of $U(r,s)$. Rewrite the increasing sequence $1,\cdots,r+s$ by 
$$
\begin{array}{c}
i^0_{1,1},\cdots,i^0_{1,s^0_1},\cdots,i^r_{1,1},\cdots,i^r_{1,s^r_1},i_{1,1},\cdots,i_{1,s_1},\\
\cdots\\
i^0_{\beta,1},\cdots,i^0_{\beta,s^0_\beta},\cdots,i^r_{\beta,1},\cdots,i^r_{\beta,s^r_\beta},i_{\beta,1},\cdots,i_{\beta,s_\beta}.
\end{array}
$$
Here, if some $s^q_i=0$, then we stipulate that the sequence $i^q_{i,1},\cdots,i^q_{i,s^q_i}$ is empty.
Denote by $\epsilon(\underline{\overline s},\underline s)$ the sign of the rearrangement of $1,\cdots,r+s$ given by
$$
\begin{array}{c}
  i^0_{1,1},\cdots,i^0_{1,s^0_1},\cdots,i^0_{\beta,1},\cdots,i^0_{\beta,i^0_\beta},\\
  \cdots\\
  i^r_{1,1},\cdots,i^r_{1,s^0_1},\cdots,i^r_{\beta,1},\cdots,i^r_{\beta,i^0_\beta},\\
  i_{1,1},\cdots,i_{1,s_1},\cdots,i_{\beta,1},\cdots,i_{\beta,s_\beta}.
\end{array}
$$
\item 
 Given any $(\overline{\underline j},\underline{\overline s},\underline s)\in U(r,s)$.  For any \( 1 \leq i \leq \beta \), set  
\[
a_i :=s_i+\cdots+s_\beta,~ 
b_i := s_{i+1}+\cdots+s_\beta.
\]
Here, we stipulate that $b_\beta=0$.
Set  
$$C(\underline{\overline{j}},\underline{\overline s},\underline s) := \epsilon(\underline{\overline s},\underline s)\frac{\prod_{1\leq i\leq\beta}(\sum_{q=0}^rs^q_i+s_i)!}{\prod_{0\leq q\leq r,1\leq i\leq\beta}s^q_i!}\prod_{i=1}^\beta c_i,$$
where  
\[
c_i := 
\begin{cases}
\prod_{b_i<q \leq a_i} \Big[ (j^q_i+ s^q_i+1) + \cdots + (j^{a_i}_i+ s^{a_i}_i+1) \Big]^{-1}, & \text{if } s_i > 0, \\
1, & \text{otherwise}.
\end{cases}
\]
    \end{itemize}

\begin{construction}\label{tilde varphi}
  Given any $(\overline{\underline j},\underline{\overline s},\underline s)\in U(r,s)$ and set
$l_i:=\sum_{q=0}^rs^q_i+s_i,~1\leq i\leq\beta$. Define
$$
\tilde\varphi_{(\overline{\underline j},\overline{\underline s},\underline s)}:E\otimes\Omega_1^{\otimes l_1}\otimes\cdots\otimes\Omega_\beta^{\otimes l_\beta}
\to F_*\Omega^s(H,\nabla)
$$
by the setting the section
$$
e\otimes(\xi^0_1\otimes\cdots\otimes\xi^r_1\otimes\xi_1)\otimes\cdots\otimes(\xi^0_\beta\otimes\cdots\otimes\xi^r_\beta\otimes\xi_\beta),~e\in E,~\xi^q_i\in\Omega_i^{\otimes s^q_i},~\xi_i\in\Omega_i^{\otimes s_i}
$$
to
$$
 C(\overline{\underline j},\overline{\underline s},\underline s)(\wedge^sh)(\xi_\beta\otimes\cdots\otimes\xi_1)\iota_{\tilde F_0}(\sigma^{\underline{\overline j}}(F^*e))\otimes\Xi,
$$
where
$$
\Xi:=[(\wedge^{s^0_1}dh_0)(\xi^0_1)\wedge\cdots\wedge(\wedge^{s^0_\beta}dh_0)(\xi^0_\beta)]\wedge\cdots\wedge[(\wedge^{s^r_1}dh_r)(\xi^r_1)\wedge\cdots\wedge(\wedge^{s^r_\beta}dh_r)(\xi^r_\beta)].
$$
\end{construction}
\emph{The proof of Lemma \ref{indep of basis}}. Set
$$
(\tilde\varphi(r,s)_{\tilde F_0,\cdots,\tilde F_r}:=\bigoplus_{(\overline{\underline j},\underline{\overline s},\underline s)\in U(r,s)}\tilde\varphi_{(\overline{\underline j},\overline{\underline s},\underline s)}):E\otimes\tilde\Omega^{\otimes(r+s)}\to F_*\Omega^s(H,\nabla)
$$
By construction, we know that it is independent of the choice of the basis $\omega_1,\cdots,\omega_n$. On the other hand, a straightforward computation shows that 
$$
\varphi(r,s)_{\tilde F_0,\cdots,\tilde F_r}=\tilde\varphi(r,s)_{\tilde F_0,\cdots,\tilde F_r}\circ\mathrm{sec}_{\underline\Omega},
$$
from which the lemma follows. \qed

\begin{remark}
\upshape

When $(E,\theta)=(\mathcal O_X,0)$, the above construction had been essentially constructed in the proof of \cite[Theorem 3.9]{AS} for $s=1,2$ and $\mathrm{rank}(\Omega_i)=1$. 

 Let $ K^*$ be an abstract Koszul complex on a ringed topos $(X,\mathcal O)$ (see \cite[Definirion 2.1]{AS}) and set $T:=\mathcal H^1(K^*)^{\vee}$. Suppose that the gerbe $\tau_{\leq 1}K^*$ is trivial and $p\mathcal O=0$. Let $(E,\theta)$ be a nilpotent $T$-Higgs bundle of level $\leq\ell$ ($\ell<p$), namely $(\mathrm{Sym}^{\ell+1}T)E=0$. Using the nonabelian Hodge theory of Ogus-Vologodsky \cite{OV} (see also Lan-Sheng-Zuo \cite{LSZ}), we can twist $K^*$ by $(E,\theta)$, denote by $K^*_\theta$ the resulting complex. If there is a splitting $T=\bigoplus_{i=1}^\beta T_i$ with $\mathrm{rank}(T_i)<p-\ell$, we may use \eqref{key construction} to establish an isomorphism between the Higgs complex $\Omega^*(E,\theta)$ and $K^*_\theta$. In particular, it follows that Theorem \ref{de-Hig without truncation 2} can be generalized to logarithmic geometry.
\end{remark}


\iffalse Note that for any $0<q\leq n$, the splitting $\Omega^1_{X/k}(\log D)=\bigoplus_{i=0}^\beta E_i$ induces a natural isomorphism
$$
\Omega^q_{X/k}(\log D)\cong\bigoplus_{q_1+\cdots+q_\beta=q,0\leq q_i\leq\mathrm{rank}(E_i)}\wedge^{q_1}E_1\otimes\cdots\otimes\wedge^{q_\beta}E_\beta.
$$
Here, we stipulate that $\wedge^0E_i=\mathcal O_X$. On the other hand, for any vector bundle $E$ on $X$ and any $0\leq q\leq\mathrm{min}\{p-1,\mathrm{rank}(E_i)\}$, the projection $\otimes^qE\to\wedge^qE$ has a natural section
$$
\wedge^qE\to\otimes^qE,~e_1\wedge\cdots\wedge e_q\mapsto\frac{1}{q!}\sum_\sigma\mathrm{sgn}(\sigma)e_{\sigma(1)}\otimes\cdots\otimes e_{\sigma(q)},~e_1,\cdots,e_q\in E.
$$
It follows that if $\mathrm{rank}(E_i)<p$ holds for any $i$, then there are natural inclusions
$$
\Omega^q_{X/k}(\log D)\subset\bigoplus_{q_1+\cdots+q_\beta=q,0\leq q_i\leq\mathrm{rank}(E_i)}(\otimes^{q_1}E_1)\otimes\cdots\otimes(\otimes^{q_\beta}E_\beta)\subset\Omega^1_{X/k}(\log D)^{\otimes q}.
$$
Using the non-canonical section of $\pi$ induced by the splitting $\Omega^1_{X/k}(\log D)=\bigoplus_{i=1}^\beta \Omega_i$, coupled with Construction \ref{construction higher homotopy},
one can check that there is an $\mathcal O_U$-linear morphism 
$$
\tilde\varphi(r,s)_{\tilde F_0,\cdots,\tilde F_r}:E\otimes\Omega^1_{X/k}(\log D)^{\otimes (r+s)}|_U\to F_*\Omega^s(H,\nabla)|_U
$$
such that its composite with the inclusion 
$$
\Omega^{r+s}(E,\theta)|_U\hookrightarrow E\otimes\Omega^1_{X/k}(\log D)^{\otimes(r+s)}|_U
$$
is precisely $\varphi(r,s)_{\tilde F_0,\cdots,\tilde F_r}$.\fi



\emph{The proof of Theorem \ref{de-Hig without truncation 2}}. For any $r>0$, we construct 
$$
\mathrm{Ho}^{-r}:\Delta_r(\mathcal L)\to \mathcal Hom^{-r}_{\mathcal{O}_X}(\Omega^*(E,\theta), F_*\Omega^*(H,\nabla))
$$
as follows. For any $0\leq s\leq n-r$, any open subset $U\subset X$ and any local liftings $\tilde F_0,\cdots,\tilde F_r\in\Gamma(U,\mathcal L)$, set
$$
\mathrm{Ho}^{-r}(\tilde F_0,\cdots,\tilde F_r):=\{\varphi(r,s)_{\tilde F_0,\cdots,\tilde F_r}\}_{0\leq s\leq n-r}.
$$
To verify the family $\mathrm{Ho}:=\{\mathrm{Ho}^{-r}\}_{r\geq0}$ constructed above forms an $\mathcal L$-indexed $\infty$-homotopy, it suffices to check that
\begin{eqnarray}\label{infty homotopy relation}
\nabla\varphi(r,s-1)_{\tilde F_0,\cdots,\tilde F_r}+\sum_{q=0}^r(-1)^{q+s}\varphi(r-1,s)_{\tilde F_0,\cdots,\widehat{\tilde F_q},\cdots,\tilde F_r}=\varphi(r,s)_{\tilde F_0,\cdots,\tilde F_r}\theta.
\end{eqnarray}
Without loss of generality, we may assume that $\Omega^1_{X/k}(\log D)$ admits a basis $\omega_1,\cdots,\omega_n$ over $U$ as described above \eqref{varphi r s}, and Write $\theta|_U=\sum_{i=1}^n\theta_i\otimes\omega_i$ again. Take a section $\boldsymbol{e}$ of $\Omega^{r+s-1}(E,\theta)|_U$ and write it as in \eqref{section e}, namely
$$
\boldsymbol{e}=\sum_{1\leq i_1<\cdots<i_{r+s-1}\leq n}e_{i_1,\cdots,i_{r+s-1}}\otimes(\omega_{i_1}\wedge\cdots\wedge\omega_{i_{r+s-1}}),~e_{i_1,\cdots,i_{r+s-1}}\in E|_U.
$$

\emph{The computation of $\nabla\varphi(r,s-1)_{\tilde F_0,\cdots,\tilde F_r}(\boldsymbol{e})$}. According to the identification of integrable connections
$$
\iota_{\tilde F_0}:((F^*E)|_U,\nabla_\mathrm{can}+\sum_{i=1}^nF^*\theta_i\otimes\zeta_{0,i})\cong (H,\nabla)|_U,
$$
we see that $\nabla\varphi(r,s-1)_{\tilde F_0,\cdots,\tilde F_r}(\boldsymbol{e})$ equals
\begin{eqnarray}\label{L nabla}
 \sum_{\underline S,\underline{\overline j},q,i}
f_{\underline S}^{\underline{\overline j}}[(\textcolor{red}{j^q_i}\textcolor{red}{dh_{q,i}}\wedge\boldsymbol{\omega}_{\underline S})\frac{h^{[\underline{\overline j}]}}{\textcolor{red}{h_{q,i}}}]+\sum_{\underline S,\underline{\overline j},i} \textcolor{red}{(F^*\theta_i)}(f_{\underline S}^{\underline{\overline j}})[(\textcolor{red}{\zeta_{0,i}}\wedge\boldsymbol{\omega}_{\underline S})h^{[\underline{\overline j}]}].
\end{eqnarray}
Here, $h^{[\underline{\overline j}]}/h_{q,i}$ is well-defined  for $j^q_i>0$ and we stipulate that $h^{[\underline{\overline j}]}/h_{q,i}:=1$ for $j^q_i=0$. The summation $\sum_{\underline S,\underline{\overline j},q,i}$ is taken over 
$$
\underline S:=(S_0,\cdots,S_r),~\underline{\overline j}:=(j^1_1,\cdots,j^1_n;\cdots;j^r_1,\cdots,j^r_n),~q,i,
$$ 
subject to the following conditions:
\begin{itemize}
    \item $\#S_0+\cdots+\#S_r=s-1$;
    \item $0\leq j^q_l<p$;
    \item $0\leq q\leq r$ and $1\leq i\leq n$.
\end{itemize}
The another summation $\sum_{\underline S,\underline{\overline j},i}$ is defined similar to $\sum_{\underline S,\underline{\overline j},q,i}$.

By grouping and simplifying like terms in \eqref{L nabla}, we have
$$\nabla\varphi(r,s-1)_{\tilde F_0,\cdots,\tilde F_r}(\boldsymbol{e})=\sum_{(\underline S,\underline{\overline j})\in T(r,s)}g_{\underline S}^{\underline{\overline j}}(\boldsymbol{\omega}_{\underline S}h^{[\underline{\overline j}]}),$$
where
$$
g_{\underline S}^{\underline{\overline j}}:=\sum_{q,l}\textcolor{red}{\epsilon}f_{(S_0,\cdots,S_q-\{i^q_l\},\cdots,S_r)}^{\underline{\overline j}(q,l)}C\underline{\overline j}(q,l)^1_{i_1}\cdots \underline{\overline j}(q,l)^r_{i_r}\iota_{\tilde F_0}(F^*(\theta^{\underline{\overline j}}\theta_{i_1}^{-1}\cdots\theta_{i_r}^{-1}\textcolor{red}{\theta_{i^q_l}}e_{i^0_1,\cdots,\textcolor{red}{\widehat{i^q_l}},\cdots,i^r_{s_r},\underline i})).
$$
Here, $\underline{\overline j}(q,l)^{q'}_{l'}:=j^q_l$ for $(q',l')\neq(q,l)$ and $j^q_l+1$ for $(q',l')=(q,l)$.
If $S_q\neq\emptyset$, then we write $S_q=\{i_1^q,\cdots,i^q_{s_q}\}$ with $i^q_1<\cdots<i^q_{s_q}$. Set
$$C:=C(S_0,\cdots,S_q-\{i^q_l\},\cdots,S_r,\underline{\overline j}(q,l)),~\epsilon:=\frac{\mathrm{sgn}(i^0_1,\cdots,i^r_{s_r})}{\mathrm{sgn}(i_l^q,i^0_1,\cdots,\widehat{i^q_l},\cdots,i^r_{s_r})}.$$
Clearly, one can check that $C=C(S_0,\cdots,S_r,\underline{\overline j})$. The summation $\sum_{q,l}$ is taken over $q,l$, subject to the following conditions:
\begin{itemize}
    \item $1\leq q\leq r$ such that $S_q\neq\emptyset$;
    \item $1\leq l\leq s_q$.
\end{itemize}

By the very construction of $\underline{\overline j}(q,l)$, $g_{\underline{S}}^{\underline{\overline j}}$ is the sum of two parts. The first part equals
\begin{eqnarray}\label{f'}
f_{\underline S}'^{\underline{\overline j}}:=\sum_{\underline i,q,l}\epsilon Cj^1_{i_1}\cdots j^r_{i_r}\iota_{\tilde F_0}(F^*(\theta^{\underline{\overline j}}\theta_{i_1}^{-1}\cdots\theta_{i_r}^{-1}\textcolor{red}{\theta_{i^q_l}}e_{i^0_1,\cdots,\textcolor{red}{\widehat{i^q_l}},\cdots,i^r_{s_r},\underline i})).
\end{eqnarray}
Here, the summation $\sum_{\underline i,q,l}$ is taken over $\underline i=(i_1,\cdots,i_r),q,l$, subject to the following conditions:
\begin{itemize}
\item $1\leq i_1<\cdots<i_r\leq n$;
    \item $i_q\in D^{-1}(D(\lambda_q)),~\lambda_q=\mathrm{max}(\{l:j^q_l\neq0\}\cup S_q),~D(\lambda_1)\geq\cdots\geq D(\lambda_r)$.;
    \item $0\leq q\leq r$ such that $S_q\neq\emptyset$, and $1\leq l\leq s_q$.
\end{itemize}
The second part equals
\begin{eqnarray}\label{f''}
f_{\underline S}''^{\underline{\overline j}}:=\sum_{q=1}^r\sum_{i_1,\cdots,{\textcolor{red}{\widehat{i_q}}},\cdots,i_q}\textcolor{red}{(-1)^{s+q}} C \textcolor{red}{s'_q}j_{i_1}^1\cdots\textcolor{red}{\widehat{j^q_{i_q}}}\cdots j^r_{i_r}\iota_{\tilde F_0}(F^*(\theta^{\underline{\overline j}}\theta_{i_1}^{-1}\cdots\textcolor{red}{\widehat{\theta_{i_q}^{-1}}}\cdots\theta_{i_r}^{-1}e_{\underline S,i_1,\cdots,\textcolor{red}{\widehat{i_q}},\cdots,i_r})),
\end{eqnarray}
where 
$$
s_q'=\#D^{-1}(D(\lambda_q))\cap S_q,~\lambda_q=\mathrm{max}( S_q\cup\{l:j^q_l\neq0\}).
$$
Note that
the summation $\sum_{\underline S,\underline{\overline j}}$ has been described above. The summation $\sum_{i_1,\cdots,\widehat{i_q},\cdots,i_q}$ is taken over $i_1,\cdots,i_{q-1},i_{q+1},\cdots,i_r$, subject to the following conditions:
\begin{itemize}
\item $1\leq i_1<\cdots<i_{q-1}<i_{q+1}<\cdots<i_r\leq n$;
    \item For any $u\in\{1,\cdots,r\}-\{q\}$, we have $i_u\in D^{-1}(D(\lambda_q))$;
    \item $D(\lambda_1)\geq\cdots\geq D(\lambda_r)$.
\end{itemize}

\emph{The computation of $\sum_{q=0}^r(-1)^{q+s}\varphi(r-1,s)_{\tilde F_0,\cdots,\textcolor{red}{\widehat{\tilde F_q}},\cdots,\tilde F_r}(\boldsymbol{e})$}. If $r=1$, there is nothing to do. Thus, we may assume $r>1$. By \eqref{key construction}, we have
\begin{eqnarray}\label{r-1,s,q}
\varphi(r-1,s)_{\tilde F_0,\cdots,\textcolor{red}{\widehat{\tilde F_q}},\cdots,\tilde F_r}(\boldsymbol{e})=\sum_{(\underline S,\underline{\overline j})\in T(r-1,s)}f_{\underline S}^{\underline{\overline j},\textcolor{red}{\hat q}}(\boldsymbol{\omega}_{\underline S}^{\textcolor{red}{\hat{q}}}h^{[\underline{\overline j}],\textcolor{red}{\hat{q}}}).
\end{eqnarray}
 To explain the right hand side of the equality above, we introduce an additional notation table below.
\begin{itemize}
    \item For $q=0$, set $$\boldsymbol{\omega}_{\underline S}^{\hat 0}:=\zeta_{1,S_0}\wedge dh_{2,S_1}\wedge\cdots\wedge dh_{r,S_{r-1}},~\underline S=(S_0,\cdots,S_{r-1}).$$ 
    For $q=r$, set 
    $$\boldsymbol{\omega}_{\underline S}^{\hat r}:=\zeta_{0,S_0}\wedge dh_{1,S_1}\wedge\cdots\wedge dh_{r-1,S_{r-1}}.$$
     For $0<q<r$, set 
    $$\boldsymbol{\omega}_{\underline S}^{\hat q}:=\zeta_{0,S_0}\wedge dh_{1,S_1}\wedge\cdots\wedge dh_{q-1,S_{q-1}}\wedge dh_{\tilde F_{q-1}\tilde F_{q+1},S_q}\wedge dh_{q+2,S_{q+1}}\wedge\cdots dh_{r,S_{r-1}}.$$
    Here, for \( S = \{i_1, \dots, i_m\} \subset \{1, \dots, n\} \) with \( i_1 < \cdots < i_m \), we define  
\[
dh_{\tilde F_{q-1}\tilde F_{q+1},S} :=dh_{\tilde F_{q-1}\tilde F_{q+1},i_1}\wedge \cdots \wedge dh_{\tilde F_{q-1}\tilde F_{q+1},i_m},~h_{\tilde F_{q-1}\tilde F_{q+1},i}:=h_{\tilde F_{q-1}\tilde F_{q+1}}(F^*\omega_i).
\]  
By convention, we set \( dh_{q+1/2,\emptyset} := 1 \). 
\item $f_{\underline S}^{\underline{\overline j},\hat q}:=\left\{
\begin{matrix}
 \sum_{\underline i}C(S_0,\cdots,S_{r-1},\underline{\overline j})j^1_{i_1}\cdots j^{r-1}_{i_{r-1}}\iota_{\tilde F_1}(F^*(\theta^{\underline{\overline j}}\theta_{i_1}^{-1}\cdots\theta_{i_{r-1}}^{-1}e_{\underline S,\underline i})),&q=0,\\
\sum_{\underline i}C(S_0,\cdots,S_{r-1},\underline{\overline j})j^1_{i_1}\cdots j^{r-1}_{i_{r-1}}\iota_{\tilde F_0}(F^*(\theta^{\underline{\overline j}}\theta_{i_1}^{-1}\cdots\theta_{i_{r-1}}^{-1}e_{\underline S,\underline i})),&0<q\leq r.
\end{matrix}
\right.$
\item $h^{[\underline{\overline j}],\hat q}:=a_qb_qc_q$, where
$$
\begin{array}{c}
a_q:=\left\{
\begin{matrix}
1,&q=0,1,\\
(h_{1,1}^{[j^1_1]}\cdots h_{1,n}^{[j^1_n]})\cdots (h_{q-1,1}^{[j^{q-1}_1]}\cdots h_{q-1,n}^{[j^{q-1}_n]}),&1<q\leq r,
\end{matrix}
\right.\\
b_q:=\left\{
\begin{matrix}
1,&q=0,r,\\
h_{\tilde F_{q-1}\tilde F_{q+1},1}^{[j^q_1]}\cdots h_{\tilde F_{q-1}\tilde F_{q+1},n}^{[j^q_n]},&0<q<r,
\end{matrix}
\right.\\
c_q:=\left\{
\begin{matrix}
1,&q=r-1,r,\\
 (h_{q+2,1}^{[j^{q+1}_1]}\cdots h_{q+2,n}^{[j^{q+1}_n]})\cdots(h_{r,1}^{[j_1^{r-1}]}\cdots h_{r,n}^{[j_n^{r-1}]}),&0\leq q\leq r-2.
\end{matrix}
\right.
\end{array}
$$
\end{itemize}

Using the basic relations
$$
\begin{array}{c}
\zeta_{1,i}=\zeta_{0,i}+dh_{1,i},~\iota_{\tilde F_1}(F^*e)=\iota_{\tilde F_0}(\exp (\sum_{i=1}^nh_{1,i}F^*\theta_i)(F^*e))~(e\in E|_U),\\
h_{\tilde F_{q-1}\tilde F_{q+1},i}=h_{q,i}+h_{q+1,i},
\end{array}
$$
it is easy to see that for any $0\leq q<r$, we have
$$
\varphi(r-1,s)_{\tilde F_0,\cdots,\widehat{\tilde F_q},\cdots,\tilde F_r}(\boldsymbol{e})=\sum_{(\underline S,\underline{\overline j})\in T(r,s)}f_{\underline S\textcolor{red}{(\cup q)}}^{\underline{\overline j}\textcolor{red}{(\cup q)},\hat q}(\boldsymbol{\omega}_{\underline S}h^{[\underline{\overline j}]}),
$$
where
\begin{eqnarray}\label{cup j}
\underline{\overline j}(\cup q):=\left\{
\begin{matrix}
(j^2_1,\cdots,j^2_n;\cdots;j^r_1,\cdots,j^r_n),&q=0,\\
(\cdots;j^{q-1}_1,\cdots,j^{q-1}_n;j^q_1+j^{q+1}_n,\cdots,j^q_n+j^{q+1}_n;j^{q+2}_1,\cdots,j^{q+2}_n;\cdots),&0\leq q<r,
\end{matrix}
\right.
\end{eqnarray}
and
\begin{eqnarray}\label{cup s}
\underline S(\cup q):=(S_0,\cdots,S_{q-1},S_q\cup S_{q+1},S_{q+2},\cdots,S_r).
\end{eqnarray}
Consequently, we have
$$
\sum_{q=0}^r(-1)^{q+s}\varphi(r-1,s)_{\tilde F_0,\cdots,\widehat{\tilde F_q},\cdots,\tilde F_r}(\boldsymbol{e})=\sum_{(\underline S,\underline{\overline j})\in T(r,s)}f_{\underline S}'''^{\underline{\overline j}}(\boldsymbol{\omega}_{\underline S}h^{[\underline{\overline j}]}),
$$
where
$$
f_{\underline S}'''^{\underline{\overline j}}:=\left\{
\begin{matrix}
  \sum_{q=0}^{r-1}(-1)^{q+s}f_{\underline S(\cup q)}^{\underline{\overline j}(\cup q),\hat q},&\mathrm{if}~\lambda_r>0,\\
 \sum_{q=0}^{r-1}(-1)^{q+s}f_{\underline S(\cup q)}^{\underline{\overline j}(\cup q),\hat q}+(-1)^{r+s}f_{(S_0,\cdots,S_{r-1})}^{(j^1_1,\cdots,j^1_n;\cdots;j^{r-1}_1,\cdots,j^{r-1}_n),\hat r},&\mathrm{otherwise}.
\end{matrix}
\right.
$$

One can check that if $\lambda_r=0$, then
 $$
 f_{\underline S(\cup 0)}^{\underline{\overline j}(\cup 0),\hat 0}=\cdots=f_{\underline S(\cup r-2)}^{\underline{\overline j}(\cup r-2),\widehat{r-2}}=0,\cdots,f_{\underline S(\cup r-1)}^{\underline{\overline j}(\cup r-1),\widehat{r-1}}=f_{(S_0,\cdots,S_{r-1})}^{(j^1_1,\cdots,j^1_n;\cdots;j^{r-1}_1,\cdots,j^{r-1}_n),\hat r}.
 $$
 On the other hand, if $\lambda_r>0$ and $D(\lambda_q)<D(\lambda_{q+1})$ for some $0\leq q<r$, then either all
$$
f_{\underline S(\cup 0)}^{\underline{\overline j}(\cup 0),\hat 0},\cdots,f_{\underline S(\cup r-1)}^{\underline{\overline j}(\cup r-1),\widehat{r-1}}
$$
are $0$ or two adjacent terms coincide while the rest are zero. It follows that $f_{\underline S}'''^{\underline{\overline j}}=0$ if $\lambda_r=0$ or  $D(\lambda_q)<D(\lambda_{q+1})$ for some $0\leq q<r$. Next, we assume that
$$
D(\lambda_1)\geq\cdots\geq D(\lambda_r)>0.
$$
Using the following observations:
\begin{itemize}
    \item If $D(\lambda_q)=D(\lambda_{q+1})$, then
    $$
     C(\cdots,S_{q-1}\cup S_q,\cdots,\underline{\overline j})-C(\cdots,S_q\cup S_{q+1},\cdots,\underline{\overline j})=(j^q+s_q')C(S_0,\cdots,S_r,\underline{\overline j});
    $$
    \item If  $D(\lambda_q)>D(\lambda_{q+1})$, then
    $$
    C(\cdots,S_{q-1}\cup S_q,\cdots,\underline{\overline j})=(j^q+s_q')C(S_0,\cdots,S_r,\underline{\overline j}),
    $$
\end{itemize}
we deduce that $f_{\underline S}'''^{\underline{\overline j}}$ equals
\begin{eqnarray}\label{f'''}
\sum_{q=1}^r\sum_{i_1,\cdots,{\textcolor{red}{\widehat{i_q}}},\cdots,i_q}\textcolor{red}{(-1)^{s+q-1}} C \textcolor{red}{(j^q+s'_q)}j_{i_1}^1\cdots\textcolor{red}{\widehat{j^q_{i_q}}}\cdots j^r_{i_r}\iota_{\tilde F_0}(F^*(\theta^{\underline{\overline j}}\theta_{i_1}^{-1}\cdots\textcolor{red}{\widehat{\theta_{i_q}^{-1}}}\cdots\theta_{i_r}^{-1}e_{\underline S,i_1,\cdots,\textcolor{red}{\widehat{i_q}},\cdots,i_r})).
\end{eqnarray}
Since the summation above is similar to \eqref{f''}, we omit further explanation.

Combing \eqref{f'}, \eqref{f''}, and \eqref{f''}, we obtain that the image of the left hand side of \eqref{infty homotopy relation} under $\boldsymbol{e}$ is 
\begin{eqnarray}\label{LHS infty homotopy relation}
\sum_{(\underline S,\underline{\overline j})\in T(r,s)}(f_{\underline S}'^{\underline{\overline j}}+f_{\underline S}''^{\underline{\overline j}}+f_{\underline S}'''^{\underline{\overline j}})(\omega_{\boldsymbol{\underline S}}h^{[\underline{\overline j}]}).
\end{eqnarray}
By direct computation, we have
$$
f_{\underline S}''^{\underline{\overline j}}+f_{\underline S}'''^{\underline{\overline j}}=\sum_{q=1}^r\sum_{\underline i}\textcolor{red}{(-1)^{s+q-1}} C(S_0,\cdots,S_r,\underline{\overline j}) j_{i_1}^1\cdots j^r_{i_r}\iota_{\tilde F_0}(F^*(\theta^{\underline{\overline j}}\theta_{i_1}^{-1}\cdots\theta_{i_r}^{-1}\textcolor{red}{\theta_{i_q}}e_{\underline S,i_1,\cdots,\textcolor{red}{\widehat{i_q}},\cdots,i_r})).
$$
Consequently, we demonstrate that \eqref{LHS infty homotopy relation} equals the image of the right hand side of \eqref{infty homotopy relation} under $\boldsymbol{e}$. This completes the proof. \hfill\qed


\subsection{Comparison between different higher homotopies}

Following the previous subsection, we continue to use the same notation and assumptions. It is natural to ask 
\begin{question}
Is the $\infty$-homotopy $\mathrm{Ho}_{\underline\Omega}$ described in the proof of Theorem \ref{de-Hig without truncation 2} independent of the choice of the splitting $\Omega=\bigoplus_{i=1}^\beta\Omega_i$, up to homotopy? More precisely, if $\Omega=\bigoplus_{i=1}^{\beta'} \Omega'_i$, where $\mathrm{rank}(\Omega'_i)<p-\ell$, is another splitting, is $\mathrm{Ho}_{\underline\Omega}$ homotopic to $\mathrm{Ho}_{\underline\Omega'}$ ($\Omega':=(\Omega'_1,\cdots,\Omega_{\beta'}')$)?
\end{question}
This question can be resolved in the following case: there exists $1\leq o<\beta$ such that
\begin{eqnarray}\label{case i}
\Omega_i':=\left\{
\begin{matrix}
 \Omega_i,&1\leq i<o,\\
 \Omega_o\bigoplus\Omega_{o+1},&i=o,\\
 \Omega_{i+1},&o<i\leq\beta-1,
\end{matrix}
\right.
\end{eqnarray}
and $\mathrm{rank}(\Omega_o)+\mathrm{rank}(\Omega_{o+1})<p-\ell$.
\iffalse
\item[(ii)] All $\Omega_i$ are line bundles and there exists $1\leq o<\beta$ such that
\begin{eqnarray}\label{case ii}
\Omega_i':=\left\{
\begin{matrix}
 \Omega_i,&i\neq o,o+1,\\
 \Omega_{o+1},&i=o,\\
 \Omega_o,&i=o+1.
\end{matrix}
\right.
\end{eqnarray}
\fi

\begin{proposition}\label{Ho^2}
With notation and assumptions as above, $\mathrm{Ho}_{\underline\Omega}$ is homotopic to $\mathrm{Ho}_{\underline\Omega'}$.
\end{proposition}
To prove this lemma, for any $r>0$ and 
$s\geq0$, we utilize the following notation table, which is built upon the notable table presented in the previous subsection.
\begin{itemize}
    \item $\underline{\overline j}:=(j^1_1,\cdots,j^1_n;\cdots;j^r_1,\cdots,j^r_n)$, where $0<\sum_{l=0}^nj^q_l<p$ for any $1\leq q\leq r$.
    \item $\underline{S}:=(S_0,\cdots,S_r)$, where $S_0,\cdots,S_r\subset\{1,\cdots,n\}$ and 
    $$\#S_0+\cdots+\#S_r==\#(S_0\cup\cdots\cup S_r)=s.$$
    \item Given any pair $(\underline S,\underline{\overline j})$.
    \begin{itemize}
    
        \item For $1\leq q\leq r$, set 
        $$\lambda_q:=\max D(S_q\cup\{l: j^q_l\neq0\}).$$
        
        \item  For $1\leq i\leq \beta$ and $i\neq o,o+1$, set
$$
   a_i:=\max\{l:\lambda_l=i\},~b_i:=\min\{l:\lambda_l=i\}.
   $$ 
        \item Set
   $$
   a_\star:=\max\{l:\lambda_l\in \{o,o+1\}\},~b_\star:=\min\{l:\lambda_l\in \{o,o+1\}\}.
   $$
   
   \item  For any $1\leq q\leq r$, set
   $$
   j^q:=\left\{
  \begin{matrix}
      \sum_{i\in D^{-1}(\lambda_q)}j^q_i,&q\notin [b_\star,a_\star],\\
      \sum_{i\in D^{-1}(\{o,o+1\})}j^q_i,&q\in[b_\star,a_\star],
  \end{matrix}
  \right.   
  $$
  and
  $$
  s_q':=\left\{
  \begin{matrix}
\#S_q\cap D^{-1}(\lambda_q)),&q\notin [b_\star,a_\star],\\
\#S_q\cap D^{-1}(\{o,o+1\}),&q\in [b_\star,a_\star].
\end{matrix}
  \right.
  $$
\item   For any $b_\star\leq q\leq a_\star$, set
  $$
  \Delta_q:=\sum_{i\in D^{-1}(o+1)}j^q_i+\#S_q\cap D^{-1}(o+1).
  $$
    \end{itemize} 
    \item Let $H_{\underline\Omega,o}(r,s)$ be the set of $5$-tuples $(q,i,\underline{S},\underline i,\underline{\overline j})$, subject to the following conditions:
    \begin{itemize}

\item $1\leq q\leq r$;

\item $\underline S,\overline{\underline j}$ are defined as above;
    
 \item $\lambda_q=o+1$. Moreover, for any $1\leq l\leq r$, either $\lambda_l\geq\lambda_{l+1}$ or $\lambda_l=o, \lambda_{l+1}=o+1$;
    
    \item $i\in D^{-1}(o)\cap\{w:j^q_w\neq0\}$, $\underline i:=(i_1,\cdots,i_r)$ satisfying
 $$
  i_l\in\left\{
  \begin{matrix}
      D^{-1}(\lambda_l)\cap\{w:j^l_w\neq0\},&l\notin [b_\star,a_\star],\\
      D^{-1}(\{o,o+1\})\cap\{w:j^l_w\neq0\}, &l\in[b_\star,a_\star]-\{q\},\\
      D^{-1}(o+1)\cap\{w:j^q_w\neq0\},&l=q,
  \end{matrix}
  \right.
  $$   
  and $\#(\{i_1,\cdots,i_r\}\cup S_0\cup\cdots\cup S_r)=r+s$.
    \end{itemize}
  
  \item Given any $(q,i,\underline S,\underline i,\underline{\overline j})\in H_{\underline\Omega,o}(r,s)$. Set
  $$
  a_\star':=\max\{l:D(i_l)=o+1\},~b'_\star := \min \{l : \lambda_l=\cdots=\lambda_{a_\star'} = o+1,l\leq a'_\star\}.
  $$
  Define 
  $$
  C_o(q,i,\underline S,\underline i,\underline{\overline j}):=\left\{
\begin{matrix}
c_1\cdots c_{o-1}c_\star c_{o+2}\cdots c_\beta,&\mathrm{if}~b_\star'\leq q\leq a_\star',\\
0,&\mathrm{otherwise}.
\end{matrix}
  \right.
  $$
Here, for any $w=1,\cdots,o-1,o+2,\cdots,\beta$, we set
   \[
c_w:= 
\begin{cases}
\prod_{b_w \leq u \leq a_w}[(j^u+s'_u)+ \cdots +(j^{a_w}+s'_{a_w})]^{-1}, & \text{if } b_w > 0,~\sum_{u=b_w}^{a_w}(j^u+s'_u)<p, \\
1, & \text{otherwise}.
\end{cases}
\]  
Set
$$
c_\star:=\left\{
\begin{matrix}
 f_\star/g_\star,&\mathrm{if}~b'_\star\leq q\leq a'_\star,
~\sum_{u=b_\star}^{a_\star}(j^u+s'_u)<p,~\mathrm{and}~\lambda_u=o~\mathrm{for}~a_\star'<u\leq a_\star,\\
0,&\mathrm{otherwise},
\end{matrix}
\right.
$$
where
$$
\begin{array}{c}
f_\star:=f_{b_\star'}+\cdots+f_q,~f_l:=f_l'f_l'',\\
f_l':=\left\{
\begin{matrix}
\prod_{l<u\leq a_\star'}[(j^u+s_u')+\cdots+(j^{a_\star}+s'_{a_\star})],&b_\star'\leq l<a_\star',\\
1,&l=a_\star',
\end{matrix}
\right.\\
f_l'':=\left\{
\begin{matrix}
\prod_{b_\star'\leq u<l}(\Delta_u+\cdots+\Delta_{a'_\star}),&b_\star'<l\leq a_\star',\\
1,&l=b_\star',
\end{matrix}
\right.
\end{array}
$$
and
$$
g_\star:=\prod_{b_\star\leq u\leq a_\star}[(j^u+s_u')+\cdots+(j^{b_\star}+s'_{b_\star})]\prod_{b_\star'\leq u\leq a_\star'}(\Delta_u+\cdots+\Delta_{a'_\star}).
$$
\end{itemize}

Given any $r>0,s\geq0$ such that $r+s+1\leq n$, and any Frobenius liftings $\tilde F_0,\cdots,\tilde F_r\in\Gamma(X,\mathcal L)$. As in \eqref{D} and its corresponding neighborhood, let $\omega_1,\cdots,\omega_n$ be a basis for $\Omega^1_{X/k}(\log D)$ over $U$ associated to a nondecreasing function $D$. Let $\boldsymbol{e}\in\Omega^{r+s+1}_{X/k}(\log D)$ be any section which can be uniquely written as 
\begin{eqnarray}\label{bold e r+s+1}
\sum_{1\leq i_1<\cdots<i_{r+s+1}\leq n}e_{i_1,\cdots,i_{r+s+1}}\otimes(\omega_{i_1}\wedge\cdots\wedge\omega_{i_{r+s+1}}),~e_{i_1,\cdots,i_{r+s+1}}\in E.
\end{eqnarray}

\begin{construction}\label{case 1}
Suppose that $\Omega_i'$ is defined as in \eqref{case i}. Define $\psi(r,s)_{\tilde F_0,\cdots,\tilde F_r}(\boldsymbol{e})$ to be
\begin{eqnarray}\label{homotopy case 1}
\sum_{(q,i,\underline S,\underline i,\underline{\overline j})\in H_{\underline\Omega,o}(r,s)} C_o(q,i,\underline S,\underline i,\underline{\overline j})j^q_ij^1_{i_1}\cdots j^r_{i_r}\iota_{\tilde F_0}(F^*(\theta^{\underline{\overline j}}\theta_i^{-1}\theta_{i_1}^{-1}\cdots\theta_{i_r}^{-1}e_{i,\underline S,\underline i}))\boldsymbol{\omega}_{\underline S}h^{[\underline{\overline j}]}.
\end{eqnarray}
\iffalse
where the summation $\sum_{q,i,\underline i,\underline S,\underline{\overline j}}$ is taken over $q,i,\underline i=(i_1,\cdots,i_r),\underline S=(S_0,\cdots,S_r)$, $\underline{\overline j}=(j^1_1,\cdots,j^1_n;\cdots;j^r_1,\cdots,j_n^r)$, subject to the following conditions:
\begin{itemize}
    \item $(\underline S,\underline{\overline j})\in T_o(r,s)$, $\underline i\in Q_o(\underline S,\underline{\overline j})$;
    \item $1\leq q\leq r$ and $i\in D^{-1}(o)$.\fi
    \iffalse
    \item $1\leq i_1,\cdots,i_r\leq n$. For any $1\leq j<r$, either $D(i_j)\geq D(i_{j+1})$ or $D(i_j)=o$ and $D(i_{j+1})=o+1$;
    \item $1\leq q\leq r$, $D(i_q)=o+1$ and $i\in D^{-1}(o)$;
    \item For any $1\leq j\leq r$, if $D(i_j)\neq o$, then $D(\mathrm{max}(S_j))\leq D(i_j)$;
    \item If $D(i_j)=o$ for some $1\leq j\leq r$, then $D(\mathrm{max}(S_j))\leq o+1$.
    \fi

\iffalse
(ii) Suppose that $\Omega_i'$ is defined as in \eqref{case ii}. Define 
\begin{eqnarray}\label{homotopy case 2}
\psi(r,s)_{\tilde F_0,\cdots,\tilde F_r}(\boldsymbol{e}):=\sum_{(\underline S,\underline{\overline j})\in L_o(r,s)}\theta^{\underline{\overline j}}\theta_o^{-1}\theta_{i_1}^{-1}\cdots\theta_{i_r}^{-1}\iota_{\tilde F_0}(F^*e_{o,\underline S,\underline i})\boldsymbol{\omega}_{\underline S}h^{\underline{\overline j}},
\end{eqnarray}
where $\underline i:=(i_1,\cdots,i_r)$ and $i_l:=\lambda_q'(=\max\{u:j^l_u\neq0\})$.\fi
\iffalse
Moreover, we set $h_{\underline i}:=h_{1,i_1}\cdots h_{r,i_r}$ and
$$
 C(\underline S,\underline i,q):=c_1\cdots c_{o-1}c_\star c_{o+2}\cdots c_\beta.
$$
Here, for any $j=1,\cdots,o-1,o+2,\cdots,\beta$, we define
   \[
c_j := 
\begin{cases}
\prod_{b_j \leq\ell \leq a_j}(s'_\ell+ \cdots + s'_{a_j}) ^{-1}, & \text{if } b_j > 0~\mathrm{and}~\sum_{\ell=b_j}^{a_j}(1+s'_\ell)<p, \\
1, & \text{otherwise}.
\end{cases}
\]  
Here, $a_j,b_j$ and $s'_\ell$ are defined as follows:
\begin{itemize}
    \item $a_j := \max \{ \ell : D(i_\ell) = j \}  ,~ 
b_j := \min \{ \ell : D(i_\ell) = j\}$;
\item For any $b_j\leq\ell\leq a_j$, set 
$$
s'_\ell := \# ( S_\ell \cap D^{-1}(j))+1.
$$
\end{itemize}



Meanwhile, for $j=\star$, we set
$$
\begin{array}{c}
a_\star := \max \{ \ell : D(i_\ell) = o~\mathrm{or}~o+1\},~b_\star := \min \{ \ell : D(i_\ell) = o~\mathrm{or}~o+1\},\\
b'_\star := \max \{ \ell : D(i_\ell) = o~\mathrm{and}~\mathrm{max}\{D(i_{\ell+1}),D(\mathrm{max}(S_\ell))\}=o+1\},\\
a_\star':=\max\{\ell:\ell>b'_\star~\mathrm{and}~D(i_\ell)=o+1\}.
\end{array}
$$
Define
$$
c_\star:=\left\{
\begin{matrix}
 \frac{f_\star}{\prod_{b_\star\leq\ell\leq a_\star}(j^\ell+\cdots+j^{b_\star})\prod_{b_\star'<\ell\leq a_\star'}(s_\ell'+\cdots+s'_{a'_\star})},&\mathrm{if}~b'_\star< \min\{a'_\star, q\},
~\mathrm{and}~\sum_{\ell=b_\star}^{a_\star}j^\ell<p,\\
0,&\mathrm{otherwise}.
\end{matrix}
\right.
$$
Here, $j^\ell,s'_\ell$ and $f_\star$ are defined as follows: 
\begin{itemize}
\item For any $b_\star\leq\ell\leq a_\star$, set 
    $$
    j^\ell:=\#(S_\ell\cap(D^{-1}(o)\cup D^{-1}(o+1)))+1+\delta^q_\ell,~\delta^q_\ell:=1~\mathrm{for}~\ell=q~\mathrm{and}~0~\mathrm{otherwise};
    $$
\item For any $b_\star'<\ell\leq a_\star'$, set
$$
s_\ell':=\#(S_\ell\cap D^{-1}(o+1))+1;
$$
\item If $b_\star'<q$, then set
$$
f_\star:=f_{b_\star'+1}+\cdots+f_q,
$$
where $f_\ell:=f_\ell'f_\ell''$ with
$$
f_\ell':=\left\{
\begin{matrix}
\prod_{\ell<\ell'\leq a_\star'}(j^{\ell'}+\cdots+j^{a_\star}),&b_\star'<\ell<a_\star',\\
1,&\ell=a_\star'
\end{matrix}
\right.
$$
and
$$
f_\ell'':=\left\{
\begin{matrix}
\prod_{b_\star'<\ell'<\ell}(s_{\ell'}'+\cdots+s_{a_\star}'),&b_\star'+1<\ell\leq a_\star',\\
1,&\ell=b_\star'+1.
\end{matrix}
\right.
$$
\end{itemize}
\fi
\end{construction}

\begin{lemma}
The construction above is independent of the choice of basis $\omega_1,\cdots,\omega_n$.
\end{lemma}
\begin{proof}
Its proof is similar to that of Lemma \ref{indep of basis}, and we omit the details.
\end{proof}

\emph{Proof of Proposition \ref{Ho^2}}. To construct a homotopy from $\mathrm{Ho}_{\underline\Omega}$ to $\mathrm{Ho}_{\underline\Omega'}$, we need to define a family of morphisms
$$
\psi_r:\Delta_r(\mathcal L)\to\mathcal Hom^{-r-1}_{\mathcal O_X}(\Omega^*(E,\theta),F_*\Omega^*(H,\nabla)),~r\geq0.
$$
They can be given as follows:
\begin{itemize}
    \item $\psi_0:=0$;
    \item For any $r>0,s\geq0$, any open subset $U\subset X$, and any $\tilde F_0,\cdots,\tilde F_r\in\Gamma(U,\mathcal L)$, set
    $$
    \psi_r(\tilde F_0,\cdots,\tilde F_r)_s:=\psi(r,s)_{\tilde F_0,\cdots,\tilde F_r}:\Omega^{r+s+1}(E,\theta)\to F_*\Omega^s(H,\nabla).
    $$
\end{itemize}
We claim that the family $\psi:=\{\psi_r\}_{r\geq0}$ constructed above is the desired homotopy. Equivalently, for any $r>0,s\geq-1$, the two expressions
\begin{eqnarray}\label{psi L}
  \nabla\psi(r,s)_{\tilde F_0,\cdots,\tilde F_r}+\sum_{q=0}^r(-1)^{s+q+1}\psi(r-1,s+1)_{\tilde F_0,\cdots,\widehat{\tilde F_q},\cdots,\tilde F_r}+\psi(r,s+1)_{\tilde F_0,\cdots,\tilde F_r}\theta\end{eqnarray}
and
\begin{eqnarray}\label{psi R}
    \mathrm{Ho}_{\underline\Omega}^{-r}(\tilde F_0,\cdots,\tilde F_r)_{s+1}-\mathrm{Ho}_{\underline\Omega'}^{-r}(\tilde F_0,\cdots,\tilde F_r)_{s+1}
    \end{eqnarray}
coincide.

Given any section \(\boldsymbol{e} \in \Omega^{r+s+1}(E, \theta)\), written as in \eqref{bold e r+s+1}, to verify that \eqref{psi L} and \eqref{psi R} coincide, it suffices to show that their images of \(\boldsymbol{e}\) are identical. Moreover, without loss of generality, we may further assume that $\underline\Omega=(\Omega_1,\Omega_2)$ and $o=1$.
\iffalse
For the sake of clarity in the verification process, we introduce the following notation and mappings.
\begin{itemize}
    \item Let $V_{r,s}$ be the $\mathbb F_p$-vector space of formal sums
$\sum_{\eta\in H_{\underline\Omega,1}(r,s)}c_\eta\eta$, where $c_\eta\in\mathbb F_p$.

\item Let $U_{\underline\Omega,1}(r,s)$ be the set of $6$-tuples $(l,q,i,\underline i,\underline S,\overline{\underline j})$ satisfying the following conditions:
\begin{itemize}

\item $1\leq l,q\leq r$ and $l\neq q$;

    \item $\underline S:=(S_0,\cdots,S_r)$, where $S_0,\cdots,S_r\subset\{1,\cdots,n\}$ and 
    $$\#S_0+\cdots+\#S_r=\#(S_0\cup\cdots\cup S_r)=s.$$
    
    \item $\underline{\overline j}:=(j^1_1,\cdots,j^1_n;\cdots;j^r_1,\cdots,j^r_n)$, where $0\leq j^q_l<p$;
    
 \item $i\in D^{-1}(1)$ and $j^q_i>0$;
 
    \item $\underline i:=(i_1,\cdots,\widehat{i_l},\cdots,i_r)$ satisfying $i_q\in D^{-1}(2)$, $j^w_{i_w}>0$ for any $w\neq l$, and 
    $$
   \#(\{i\}\cup\{i_1,\cdots,\widehat{i_l},\cdots,i_r\}\cup S_0\cup\cdots\cup S_r)=r+s.
    $$
\end{itemize}

\item Let $H'_{\underline\Omega,1}$ be the disjoint union of $U_{\underline\Omega,1}$ and $T_\Omega(r,s)$, where $T_\Omega(r,s)$ was defined in the notation table above Construction \ref{construction higher homotopy}, and the subscript $\Omega$ stands for the trivial splitting of $\Omega$ (i.e., $\beta=1$).

 \item Let $V'_{r,s}$ be the $\mathbb F_p$-vector space of formal sums
$\sum_{\eta\in H'_{\underline\Omega,1}(r,s)}c_\eta\eta$, where $c_\eta\in\mathbb F_p$.
\item 
\end{itemize}
\fi


\emph{Computation of the image of $\boldsymbol{e}$ under \eqref{psi L}}. It contains four steps. The first step is the expansion of $\nabla\psi(r,s)_{\tilde F_0,\cdots,\tilde F_r}(\boldsymbol{e})$, which equals the sum $M_{11}+M_{12}$, where
$$
M_{11}:=\sum\textcolor{red}{\sum_u} C_1j^q_ij^1_{i_1}\cdots j^r_{i_r}\iota_{\tilde F_0}(F^*(\theta^{\underline{\overline j}}\theta_i^{-1}\theta_{i_1}^{-1}\cdots\theta_{i_r}^{-1}\textcolor{red}{\theta_u}e_{i,\underline S,\underline i}))(\textcolor{red}{\zeta_{0,u}\wedge}\boldsymbol{\omega}_{\underline S})h^{[\underline{\overline j}]},$$
$$
M_{12}:=\sum\textcolor{red}{\sum_{w,u}}C_1j^q_ij^1_{i_1}\cdots j^r_{i_r}\iota_{\tilde F_0}(F^*(\theta^{\underline{\overline j}}\textcolor{red}{\theta^{-1}_u}\theta_i^{-1}\theta_{i_1}^{-1}\cdots\theta_{i_r}^{-1}\textcolor{red}{\theta_u}e_{i,\underline S,\underline i}))(\textcolor{red}{dh_{w,u}\wedge}\boldsymbol{\omega}_{\underline S})(\textcolor{red}{j^w_uh^{-1}_{w,u}}h^{[\underline{\overline j}]}).
$$
\iffalse
$$
M_{13}:=\sum\textcolor{red}{\sum_w}C_1j^q_ij^1_{i_1}\cdots j^r_{i_r}(\theta^{\underline{\overline j}}\textcolor{red}{\theta^{-1}_{i_w}})\theta_i^{-1}\theta_{\underline i}^{-1}\textcolor{red}{\theta_{i_w}}\iota_{\tilde F_0}(F^*e_{i,\underline S,\underline i})(\textcolor{red}{dh_{w,i_w}\wedge}\boldsymbol{\omega}_{\underline S})(\textcolor{red}{j^w_{i_w}h^{-1}_{w,i_w}}h^{[\underline{\overline j}]}),
$$
and 
$$
M_{14}:=\sum\textcolor{red}{\sum_{u=i,i_q}}C_1j^q_ij^1_{i_1}\cdots j^r_{i_r}(\theta^{\underline{\overline j}}\textcolor{red}{\theta^{-1}_u})\theta_i^{-1}\theta_{\underline i}^{-1}\textcolor{red}{\theta_u}\iota_{\tilde F_0}(F^*e_{i,\underline S,\underline i})(\textcolor{red}{dh_{q,u}\wedge}\boldsymbol{\omega}_{\underline S})(\textcolor{red}{j^q_uh^{-1}_{q,u}}h^{[\underline{\overline j}]}).
$$\fi
Here, the first summation $\sum$ in either $M_{11}$ or $M_{12}$ is taken over $(q,i,\underline i,\underline S,\underline{\overline j})\in H_{\underline\Omega,1}$, and $C_1:=C_1(q,i,\underline i,\underline S,\underline{\overline j})$. 
\iffalse
subject to the following conditions:
\begin{itemize}
    \item $1\leq w\leq r$ and $1\leq u\leq n$;
    \item $u\neq\left\{
\begin{matrix}
i_w,&w\neq q,\\
i~\mathrm{or}~i_q,& w=q.
\end{matrix}
\right.$
\end{itemize}
\fi
The red summation in $M_{13}$ is taken over $w$ satisfying $1\leq w\leq n$ and $w\neq q$.


The second step is the expansion of $\sum_{q=0}^r(-1)^{s+q+1}\psi(r-1,s+1)_{\tilde F_0,\cdots,\widehat{\tilde F_q},\cdots,\tilde F_r}(\boldsymbol{e})$.
Using the transition morphism $G_{\tilde F_0\tilde F_1}$, $\psi(r-1,s+1)_{\textcolor{red}{\widehat{\tilde F_0}},\tilde F_1,\cdots,\tilde F_r}(\boldsymbol{e})$ equals the sum
$$
 M_{20}:=\sum C_1(q,i,\underline i,\underline S\textcolor{red}{(\cup0)},\underline{\overline j}\textcolor{red}{(\cup0)})j^{q+1}_ij^2_{i_2}\cdots j^r_{i_r}\iota_{\tilde F_0}(F^*(\theta^{\underline{\overline j}}\theta_i^{-1}\theta^{-1}_{i_2}\cdots\theta^{-1}_{i_r}e_{i,\underline S,\underline i})\boldsymbol{\omega}_{\underline S}h^{\underline{\overline j}},$$
where $\underline S(\cup 0),\underline{\overline j}(\cup 0)$ are defined as in \eqref{cup j} and \eqref{cup s}, respectively. Here, the summation is taken over $$q,i,\underline i:=(i_2,\cdots,i_r),~\underline S:=(S_0,\cdots,S_r),~\underline{\overline j}:=(j^1_1,\cdots,j^1_n;\cdots;j^r_1,\cdots,j^r_n),$$ 
subject to the following conditions:
\begin{itemize}
    \item $\#S_0+\cdots+\#S_r=s+1$;
    \item $\#(\{i\}\cup\{i_2,\cdots,i_r\}\cup S_0\cup\cdots\cup S_r)=r+s+1$;
    \item $(q,i,\underline i,\underline S(\cup0),\overline{\underline j}(\cup0))\in H_{\underline\Omega,1}(r-1,s+1)$.
\end{itemize}

For $0<u<r$, using the binomial expansion $(h_{u,l}+h_{u+1,l})^{[m]}=\sum_{m_1+m_2=m}h_{u,l}^{[m_1]}h_{u+1,l}^{[m_2]}$, we obtain that $\psi(r-1,s+1)_{\tilde F_0,\cdots,\textcolor{red}{\widehat{\tilde F_u}},\cdots,\tilde F_r}(\boldsymbol{e})$ equals the sum $M_{2u}'+M_{2u}''+M_{2u}'''$, where
$$
M'_{2u}:=\sum C_1'j^q_ij^1_{i_1}\cdots\textcolor{red}{\widehat{j^u_{i_u}}}\cdots j^r_{i_r}\iota_{\tilde F_0}(F^*(\theta^{\underline{\overline j}}\theta_i^{-1}\theta_{i_1}^{-1}\cdots\textcolor{red}{\widehat{\theta_{i_u}^{-1}}}\cdots\theta_{i_r}^{-1}e_{i,\underline S,\underline i}))\boldsymbol{\omega}_{\underline S} h^{\underline{\overline j}},$$
$$
M_{2u}'':=\sum C_1''j^q_ij^1_{i_1}\cdots\textcolor{red}{\widehat{j^{u+1}_{i_{u+1}}}}\cdots j^r_{i_r}\iota_{\tilde F_0}(F^*(\theta^{\underline{\overline j}}\theta_i^{-1}\theta_{i_1}^{-1}\cdots\textcolor{red}{\widehat{\theta_{i_{u+1}}^{-1}}}\cdots\theta_{i_r}^{-1}e_{i,\underline S,\underline i}))\boldsymbol{\omega}_{\underline S} h^{\underline{\overline j}},$$
$$
M_{2u}''':=\sum C_1'''j^1_{i_1}\cdots j^r_{i_r}\iota_{\tilde F_0}(F^*(\theta^{\underline{\overline j}}\theta_{\underline i}e_{\underline S,\underline i}))\boldsymbol{\omega}_{\underline S} h^{\underline{\overline j}}.
$$
In $M_{2u}'$ or $M_{2u}''$, we write $\underline i$ as $(i_1,\cdots,\textcolor{red}{\widehat{i_u}},\cdots,i_r)$ or $(i_1,\cdots,\textcolor{red}{\widehat{i_{u+1}}},\cdots,i_r)$, respectively. Moreover, the first two summations are taken over $q,i,\underline i,\underline S,\underline{\overline j}$,
subject to the following conditions:
\begin{itemize}
    \item $\#S_0+\cdots+\#S_r=s+1$;
    \item $\#(\{i\}\cup\{\underline i\}\cup S_0\cup\cdots\cup S_r)=r+s+1$, where $\{\underline i\}$ is the set consisting of entries of $\underline i$;
    \item $(q,i,\underline i,\underline S(\cup u),\overline{\underline j}(\cup u))\in H_{\underline\Omega,1}(r-1,s+1)$.
\end{itemize}
The last summation is taken over
$$\underline i:=(i_1,\cdots,i_r),~\underline S:=(S_0,\cdots,S_r),~\underline{\overline j}:=(j^1_1,\cdots,j^1_n;\cdots;j^r_1,\cdots,j^r_n),$$
subject to the following conditions:
\begin{itemize}
\item  $\{D(i_u),D(i_{u+1})\}=\{1,2\}$;
    \item $\#S_0+\cdots+\#S_r=s+1$;
    \item $\#(\{i_1,\cdots,i_r\}\cup S_0\cup\cdots\cup S_r)=r+s+1$;
    \item Let $i\in\{i_u,i_{u+1}\}\cap D^{-1}(1)$, then
    $$(\textcolor{red}{u},\textcolor{red}{i},(i_1,\cdots,\textcolor{red}{\widehat{i}},\cdots,i_r),\underline S(\cup u),\overline{\underline j}(\cup u))\in H_{\underline\Omega,1}(r-1,s+1).$$ 
\end{itemize}
 The three coefficients $C_1',C_1'',C_1'''$ are defined as follows:
 $$
 C_1'=C_1'':=C_1(q,i,\underline i,\underline S(\cup u),\overline{\underline j}(\cup u))
 $$
and
$$
C_1''':=\left\{
\begin{matrix}
(-1)^{s+u}C_1(u,i_u,(i_1,\cdots,\widehat{i_u},\cdots,i_r),\underline S(\cup u),\overline{\underline j}(\cup u)),&i_u\in D^{-1}(1),\\
(-1)^{s+u+1}C_1(u,i_{u+1},(i_1,\cdots,\widehat{i_{u+1}},\cdots,i_r),\underline S(\cup u),\overline{\underline j}(\cup u)),&i_{u+1}\in D^{-1}(1).
\end{matrix}
\right.
$$

Set $M_{2r}:=\psi(r-1,s+1)_{\tilde F_0,\cdots,\tilde F_{r-1},\textcolor{red}{\widehat{\tilde F_r}}}(\boldsymbol{e})$. By definition, we have
$$
\sum_{(q,i,\underline S,\underline i,\underline{\overline j})\in H_{\underline\Omega,1}(r-1,s+1)} C_1(q,i,\underline S,\underline i,\underline{\overline j})j^q_ij^1_{i_1}\cdots j^r_{i_r}\iota_{\tilde F_0}(F^*(\theta^{\underline{\overline j}}\theta_i^{-1}\theta_{i_1}^{-1}\cdots\theta_{i_r}^{-1}e_{i,\underline S,\underline i}))\boldsymbol{\omega}_{\underline S}h^{\underline{\overline j}}.
$$

The third step is the expansion of $\psi(r,s+1)\theta(\boldsymbol{e})$. Define
$$
(\theta\boldsymbol{e})_{i_1,\cdots,i_{r+s+2}}:=\sum_{q=1}^{r+s+2}(-1)^{q-1}\theta_qe_{i_1,\cdots,\textcolor{red}{\widehat{i_q}},\cdots,i_{r+s+2}}.
$$
It is easy to see that
$$
\theta\boldsymbol{e}=\sum_{1\leq i_1,\cdots,i_{r+s+2}\leq n}(\theta\boldsymbol{e})_{i_1,\cdots,i_{r+s+2}}\otimes\omega_{i_1}\wedge\cdots\wedge\omega_{i_{r+s+2}},
$$
and $$(\theta\boldsymbol{e})_{i_{\sigma(1)},\cdots,i_{\sigma(r+s+2)}}=\mathrm{sgn}(\sigma)(\theta\boldsymbol{e})_{i_1,\cdots,i_{r+s+2}}$$ holds for any permutation $\sigma$ of $\{1,\cdots,r+s+2\}$. Recall that $\psi(r,s+1)\theta(\boldsymbol{e})$ equals 
\begin{eqnarray}\label{psi theta r s+1}
\sum_{(q,i,\underline S,\underline i,\underline{\overline j})\in H_{\underline\Omega,1}(r,s+1)} C_1(q,i,\underline S,\underline i,\underline{\overline j})j^q_ij^1_{i_1}\cdots j^r_{i_r}\iota_{\tilde F_0}(F^*(\theta^{\underline{\overline j}}\theta_i^{-1}\theta_{i_1}^{-1}\cdots\theta_{i_r}^{-1}(\theta\boldsymbol{e})_{i,\underline S,\underline i}))\boldsymbol{\omega}_{\underline S}h^{[\underline{\overline j}]}.
\end{eqnarray}
Dividing $(\theta\boldsymbol{e})_{i,\underline S,\underline i}$ into the following three parts:
\begin{itemize}
    \item $\theta_1e_{\underline S,\underline i}+(-1)^{s+q+1}\theta_qe_{i,\underline S,i_1,\cdots,\textcolor{red}{\widehat{i_q}},\cdots,i_r}$;
    \item $\sum_{0\leq w\leq r,u\in S_w}(-1)^{\#S_0+\cdots+\#S_{w-1}+\#\{l\in S_w:l<i_w\}+1}\theta_ue_{i,S_0,\cdots,\textcolor{red}{S_w-\{u\}},\cdots,S_r,\underline i}$;
    \item $\sum_{w\neq q}(-1)^{s+w+1}e_{i,\underline S,i_1,\cdots,\textcolor{red}{\widehat{i_w}},\cdots,i_r}$.
\end{itemize}
Let \( M_{31}, M_{32}, M_{33} \) denote the three parts of \eqref{psi theta r s+1} obtained by replacing \( (\theta\boldsymbol{e})_{i,\underline{S},\underline{i}} \) with the three corresponding parts described above.

The first three steps show that the image of $\boldsymbol{e}$ under \eqref{psi L} is
\begin{eqnarray}\label{huge summation}
\begin{array}{c}
 (M_{11}+M_{12})\\
 +[(-1)^{s+1}M_{20}+\sum_{0<u<r}(-1)^{s+1+u}(M_{2u}'+M_{2u}''+M_{2u}''')+(-1)^{s+1+r}M_{2r}]\\
 +(M_{31}+M_{32}+M_{33}).
 \end{array}
\end{eqnarray}
The final step is to restructure the summation above in a finer form. To this end, we introduce two sets. The first is $T'_{\underline\Omega}(r,s+1)$, consisting of $6$-tuples $(l,q,i,\underline i,\underline S,\overline{\underline j})$ that meet the following conditions:
\begin{itemize}

\item $1\leq l,q\leq r$ and $l\neq q$;

    \item $\underline S:=(S_0,\cdots,S_r)$, where $S_0,\cdots,S_r\subset\{1,\cdots,n\}$ and 
    $$\#S_0+\cdots+\#S_r=\#(S_0\cup\cdots\cup S_r)=s+1;$$
    
    \item $\underline{\overline j}:=(j^1_1,\cdots,j^1_n;\cdots;j^r_1,\cdots,j^r_n)$, where $0\leq j^q_l<p$;
    
 \item $i\in D^{-1}(1)$ and $j^q_i>0$;
 
    \item $\underline i:=(i_1,\cdots,\widehat{i_l},\cdots,i_r)$ satisfying $i_q\in D^{-1}(2)$, $j^w_{i_w}>0$ for any $w\neq l$, and 
    $$
   \#(\{i\}\cup\{i_1,\cdots,\widehat{i_l},\cdots,i_r\}\cup S_0\cup\cdots\cup S_r)=r+s+1.
    $$
\end{itemize}
The second is $T''_{\underline\Omega}(r,s+1)$, consisting of $6$-tuples $(u,q,i,\underline i,\underline S,\overline{\underline j})$ that meet the following conditions:
\begin{itemize}
    \item $(q,i,\underline S,\underline i,\overline{\underline j})\in H_{\underline\Omega,1}(r,s+1)$;
    \item $u\in S_0\cup\cdots\cup S_r$.
\end{itemize}
Clearly, \eqref{huge summation} can be reformulated as the sum of the following three parts:
$$
\sum_{(\underline S,\underline i,\underline{\overline j})\in T_{\Omega}(r,s+1)}c^L_{\underline i,\underline S,\underline{\overline j}}j^1_{i_1}\cdots j^r_{i_r}\iota_{\tilde F_0}(F^*(\theta^{\underline{\overline j}}\theta_{i_1}^{-1}\cdots\theta_{i_r}^{-1}e_{\underline S,\underline i}))(\boldsymbol{\omega}_{\underline S}h^{[\underline{\overline j}]}),
$$
$$
\sum_{(l,q,i,\underline S,\underline i,\underline{\overline j})\in T'_{\underline\Omega}(r,s+1)} c^L_{l,q,i,\underline S,\underline i,\underline{\overline j}}j^q_ij^1_{i_1}\cdots\textcolor{red}{\widehat{j^l_{i_l}}}\cdots j^r_{i_r}\iota_{\tilde F_0}(F^*(\theta^{\underline{\overline j}}\theta_i^{-1}\theta_{i_1}^{-1}\cdots\theta_{i_r}^{-1}e_{i,\underline S,(i_1,\cdots,\textcolor{red}{\widehat{i_l}},\cdots,i_r)}))\boldsymbol{\omega}_{\underline S}h^{[\underline{\overline j}]},
$$
$$
\sum_{(l,q,i,\underline S,\underline i,\underline{\overline j})\in T''_{\underline\Omega}(r,s+1)} c^L_{u,q,i,\underline S,\underline i,\underline{\overline j}}j^q_ij^1_{i_1}\cdots j^r_{i_r}\iota_{\tilde F_0}(F^*(\theta^{\underline{\overline j}}\theta_i^{-1}\theta_{i_1}^{-1}\cdots\theta_{i_r}^{-1}e_{i,\underline S\textcolor{red}{-\{u\}},\underline i}\boldsymbol{\omega}_{\underline S}h^{[\underline{\overline j}]}.
$$
In the first summation, \(T_\Omega(r,s+1)\) is the set of all triples $(\underline S,\underline i,\underline{\overline j})$ that meet the following conditions:
\begin{itemize}
    \item $\underline S:=(S_0,\cdots,S_r)$, where $S_0,\cdots,S_r\subset\{1,\cdots,n\}$ and 
    $$\#S_0+\cdots+\#S_r=\#(S_0\cup\cdots\cup S_r)=s+1;$$
    \item $\underline{\overline j}:=(j^1_1,\cdots,j^1_n;\cdots;j^r_1,\cdots,j^r_n)$, where $0\leq j^q_l<p$;
    \item $\underline i:=(i_1,\cdots,i_r)$ such that
    $$
    \#(\{i_1,\cdots,i_r\}\cup S_0\cup\cdots\cup S_r)=r+s+1.
    $$
\end{itemize}
In the last summation, $\underline S-\{u\}$ means $(S_0,\cdots,S_w-\{u\},\cdots,S_r)$ if $u\in S_w$.

The coefficients $c^L_{\underline S,\underline i,\underline{\overline j}},~c^L_{l,q,i,\underline S,\underline i,\underline{\overline j}}$ and $c^L_{u,q,i,\underline S,\underline i,\underline{\overline j}}$ can be initially written as
$$
\begin{array}{c}
c^L_{\underline S,\underline i,\underline{\overline j}}=\sum_{v=1}^r(c^{L,v,1}_{\underline S,\underline i,\underline{\overline j}}+c^{L,v,2}_{\underline S,\underline i,\underline{\overline j}}+c^{L,v,3}_{\underline S,\underline i,\underline{\overline j}}),\\
c^L_{l,q,i,\underline S,\underline i,\underline{\overline j}}=c^{L,1}_{l,q,i,\underline S,\underline i,\underline{\overline j}}+c^{L,2}_{l,q,i,\underline S,\underline i,\underline{\overline j}}+c^{L,3}_{l,q,i,\underline S,\underline i,\underline{\overline j}},\\
c^L_{u,q,i,\underline S,\underline i,\underline{\overline j}}=c^{L,1}_{u,q,i,\underline S,\underline i,\underline{\overline j}}+c^{L,3}_{u,q,i,\underline S,\underline i,\underline{\overline j}},
\end{array}
$$
where the superscripts $1$, $2$, and $3$ denote the contributions arising from the first, second, and third parts of \eqref{huge summation}, respectively. The explicit descriptions for the right-hand sides of the above equalities are given below. 
\begin{itemize}
    \item  If $i_v\in D^{-1}(2)$, then 
    $$c^{L,v,1}_{\underline S,\underline i,\underline{\overline j}}:=\sum \underline{\overline j}(w,u)^v_uC_1(v,u,\underline S-\{u\},\underline i,\underline{\overline j}(w,u)),$$ 
    where the summation is taken over $u$, subject to the following conditions:
    \begin{itemize}
        \item $u\in D^{-1}(1)\cap S_w$ for some $0\leq w\leq r$;
        \item $(v,u,\underline S-\{u\},\underline i,\underline{\overline j}(w,u))\in H_{\underline\Omega,1}(r,s)$.   
    \end{itemize}
    
If $i_v\in D^{-1}(1)$, then 
    $$c^{L,v,1}_{\underline S,\underline i,\underline{\overline j}}:=\sum \underline{\overline j}(w,u)^v_uC_1(v,i_v,\underline S-\{u\},\underline i(v, u),\underline{\overline j}(w,u)),$$ 
    where $\underline{i}(v,u)$ (resp. $\underline{\overline j}(w,u)$) denotes the result of substituting the $v$-th (resp. $(w,u)$-th) entry of $\underline{i}$ (resp. $\underline{\overline j}$) by $u$ (resp. $j^w_u+1$).
 The summation is taken over $u$, subject to the following conditions:
    \begin{itemize}
        \item $u\in D^{-1}(2)\cap S_w$ for some $0\leq w\leq r$;
        \item $(v,i_v,\underline S-\{u\},\underline i(v,u),\underline{\overline j}(w,u))\in H_{\underline\Omega,1}(r,s)$.
\end{itemize}
\item Set $c^{L,0,2}_{\underline S,\underline i,\underline{\overline j}}:=0$ for $v=0$ or $D(i_v)=D(i_{v-1})$, $0<v\leq r$. Otherwise, set
$$
c^{L,v,2}_{\underline S,\underline i,\underline{\overline j}}:=\left\{
\begin{matrix}
-C_1(v-1,i_{v-1},\underline S,\underline i(\widehat{v-1}),\overline{\underline j}(\cup(v-1))),&\mathrm{if}~i_{v-1}\in D^{-1}(1), i_v\in D^{-1}(2),\\

   C_1(v-1,i_v,\underline S,\underline i(\hat v),\overline{\underline j}(\cup(v-1))),&\mathrm{if}~i_{v-1}\in D^{-1}(2), i_v\in D^{-1}(1), 
\end{matrix}
\right.
$$
where $\underline i(\widehat{v}):=(i_1,\cdots,\widehat{i_v},\cdots,i_r)$ and
$$
\overline{\underline j}(\cup(v-1)):=(j^1_1,\cdots,j^1_n;\cdots;j^{v-1}_1+j^v_1,\cdots,j^{v-1}_n
+j^v_n;\cdots;j^r_1,\cdots,j^r_n).
$$
\item For $i_v\in D^{-1}(2)$, we set
$$c^{L,v,3}_{\underline S,\underline i,\underline{\overline j}}:=\sum j^v_uC_1(v,u,\underline S,\underline i,\overline{\underline j}).$$
Here, the summation is taken over $u$, subject to the following conditions:
\begin{itemize}
    \item $u\in D^{-1}(1)-S_0\cup\cdots\cup S_r$;
    \item $(v,u,\underline S,\underline i,\overline{\underline j})\in H_{\underline\Omega,1}(r,s+1)$.
\end{itemize}
For $i_v\in D^{-1}(1)$, we set
$$c^{L,v,3}_{\underline S,\underline i,\underline{\overline j}}:=-\sum j^v_uC_1(v,i_v,\underline S,\underline i(v,u),\overline{\underline j}).$$
Here, the summation is taken over $u$, subject to the following conditions:
\begin{itemize}
    \item $u\in D^{-1}(1)-S_0\cup\cdots\cup S_r$;
    \item $(v,u,\underline S,\underline i,\overline{\underline j})\in H_{\underline\Omega,1}(r,s+1)$.
\end{itemize}

\item Set 
$$c^{L,1}_{l,q,i,\underline S,\underline i,\underline{\overline j}}:=(-1)^{s+l-1}\sum \overline{\underline j}(w,u)^l_uC_1(q,i,\underline S-\{u\},(l,u,\underline i),\overline{\underline j}(w,u)),$$
where $(l,u,\underline i)$ is obtained by inserting $u$ into $\underline i$ as the $l$-th entry. The summation is taken over $u$, subject to the following conditions:
\begin{itemize}
    \item $u\in S_w$ for some $0\leq w\leq r$;
    \item $(q,i,\underline S-\{u\},(l,u,\underline i),\overline{\underline j}(w,u))\in H_{\underline\Omega,1}(r,s)$.
\end{itemize}

\item For $0\leq l<q$, set $c^{L,2}_{l,q,i,\underline S,\underline i,\underline{\overline j}}$ to be the sum 
$$
(-1)^{s+l}C_1(q-1,i,\underline S(\cup(l-1)),\underline i,\overline{\underline j}(\cup(l-1)))+(-1)^{s+l+1}C_1(q-1,i,\underline S(\cup l),\underline i,\overline{\underline j}(\cup l)).
$$
For $q<l<r$, set $c^{L,2}_{l,q,i,\underline S,\underline i,\underline{\overline j}}$ to be the sum 
$$
(-1)^{s+l}C_1(q,i,\underline S(\cup(l-1)),\underline i,\overline{\underline j}(\cup(l-1)))+(-1)^{s+l+1}C_1(q,i,\underline S(\cup l),\underline i,\overline{\underline j}(\cup l)).
$$
For $l=r$, set $c^{L,2}_{r,q,i,\underline S,\underline i,\underline{\overline j}}$ to be the sum 
$$
(-1)^{s+r}C_1(q,i,\underline S(\cup(r-1)),\underline i,\overline{\underline j}(\cup(r-1)))+(-1)^{s+r+1}\epsilon C_1(q,i,\underline S(\cup l),\underline i,\overline{\underline j}(\cup l)).
$$
Here, we stipulate that 
$$
\epsilon:=\left\{
\begin{matrix}
1,&\mathrm{if}~S_r=\emptyset, j^r_1=\cdots=j^r_n=0;\\
0,&\mathrm{otherwise}.
\end{matrix}
\right.
$$

\item Set 
$$c^{L,3}_{r,q,i,\underline S,\underline i,\underline{\overline j}}=\sum (-1)^{s+l-1}j^l_uC_1(q,i,\underline S,(l,u,\underline i),\overline{\underline j}),$$
where the summation is taken over $u\in\{1,\cdots,n\}-S_0\cup\cdots\cup S_r$.

\item Let $u\in S_w$ for some $0\leq w\leq r$, and then define
$$c^{L,1}_{u,q,i,\underline S,\underline i,\underline{\overline j}}:=(-1)^{\#S_0+\cdots+\#S_{w-1}+\#\{u'\in S_w:u'<u\}}C_1(q,i,\underline S-\{u\},\underline i,\overline{\underline j}).
$$

\item Let $u\in S_w$ for some $0\leq w\leq r$, and then define 
$$c^{L,3}_{u,q,i,\underline S,\underline i,\underline{\overline j}}:=(-1)^{\#S_0+\cdots+\#S_{w-1}+\#\{u'\in S_w:u'<u\}+1}C_1(q,i,\underline S-\{u\},\underline i,\overline{\underline j}).
$$
\end{itemize}

\emph{Computation of the image of $\boldsymbol{e}$ under \eqref{psi R}}.
\iffalse
the sum of
$$
M_{31}:=\sum_{q,i,\underline i,\underline S,\underline{\overline j}} C_1(q,\underline i,\underline S,\underline{\overline j})j^q_ij^1_{i_1}\cdots j^r_{i_r}\theta^{\underline{\overline j}}\theta_{\underline i}^{-1}\iota_{\tilde F_0}(F^*e_{\textcolor{red}{\hat i},\underline S,\underline i})\boldsymbol{\omega}_{\underline S}h^{\underline{\overline j}},
$$
$$
M_{32}:=\sum_{q,i,\underline i,\underline S,\underline{\overline j}}\textcolor{red}{\sum_{w,u}}\textcolor{red}{\epsilon}C_1(q,\underline i,\underline S,\underline{\overline j})j^q_ij^1_{i_1}\cdots j^r_{i_r}\theta^{\underline{\overline j}}\theta_i^{-1}\theta_{\underline i}^{-1}\textcolor{red}{\theta_u}\iota_{\tilde F_0}(F^*e_{i,S_0,\cdots,\textcolor{red}{S_w-\{u\}},\cdots,S_r,\underline i})\boldsymbol{\omega}_{\underline S}h^{\underline{\overline j}},
$$
$$
M_{33}:=\sum_{q,i,\underline i,\underline S,\underline{\overline j}}\textcolor{red}{\sum_w}\textcolor{red}{\epsilon}C_1(q,\underline i,\underline S,\underline{\overline j})j^q_ij^1_{i_1}\cdots j^r_{i_r}\theta^{\underline{\overline j}}\theta_{i_w}^{-1}\theta_{\underline i}^{-1}\textcolor{red}{\theta_{i_w}}\iota_{\tilde F_0}(F^*e_{i,S_0,\cdots,\textcolor{red}{S_w-\{i_w\}},\cdots,S_r,\underline i})\boldsymbol{\omega}_{\underline S}h^{\underline{\overline j}},
$$
$$
M_{34}:=\sum_{q,i,\underline i,\underline S,\underline{\overline j}}\textcolor{red}{\sum_{u=i,i_q}}\textcolor{red}{\epsilon}C_1(q,\underline i,\underline S,\underline{\overline j})j^q_ij^1_{i_1}\cdots j^r_{i_r}\theta^{\underline{\overline j}}\theta_i^{-1}\theta_{\underline i}^{-1}\textcolor{red}{\theta_u}\iota_{\tilde F_0}(F^*e_{i,S_0,\cdots,\textcolor{red}{S_q-\{u\}},\cdots,S_r,\underline i})\boldsymbol{\omega}_{\underline S}h^{\underline{\overline j}},
$$
$$
M_{35}:=\sum_{q,i,\underline i,\underline S,\underline{\overline j}}\textcolor{red}{\sum_{w\neq q}}\textcolor{red}{(-1)^{s+w+1}}C_1(q,\underline i,\underline S,\underline{\overline j})j^q_ij^1_{i_1}\cdots j^r_{i_r}\theta^{\underline{\overline j}}\theta_{\underline i}^{-1}\textcolor{red}{\theta_{i_w}}\iota_{\tilde F_0}(F^*e_{i,\underline S,i_1,\cdots,\textcolor{red}{\widehat{i_w}},\cdots,i_r})\boldsymbol{\omega}_{\underline S}h^{\underline{\overline j}},
$$
and
$$
M_{36}:=\sum_{q,i,\underline i,\underline S,\underline{\overline j}}\textcolor{red}{(-1)^{s+q+1}}C_1(q,\underline i,\underline S,\underline{\overline j})j^q_ij^1_{i_1}\cdots j^r_{i_r}\theta^{\underline{\overline j}}\theta_{\underline i}^{-1}\textcolor{red}{\theta_{i_q}}\iota_{\tilde F_0}(F^*e_{i,\underline S,i_1,\cdots,\textcolor{red}{\widehat{i_q}},\cdots,i_r})\boldsymbol{\omega}_{\underline S}h^{\underline{\overline j}}.
$$
Here, we set
$$
\epsilon:=\left\{
\begin{matrix}
(-1)^{\#S_0+\cdots+\#S_{r-1}+\#\{l\in S_w:l<u\}+1}~\mathrm{in}~M_{32},\\
(-1)^{\#S_0+\cdots+\#S_{w-1}+\#\{l\in S_w:l<i_w\}+1}~\mathrm{in}~M_{33},\\
(-1)^{\#S_0+\cdots+\#S_{q-1}+\#\{l\in S_q:l<u\}+1}~\mathrm{in}~M_{34}.
\end{matrix}
\right.
$$
The summation $\sum_{q,i,\underline i,\underline S,\underline{\overline j}}$ is clear. The red summation in $M_{32}$ is taken over $w,u$, subject to the following conditions:
\begin{itemize}
    \item $0\leq w\leq r$, $1\leq u\leq n$, and $u\in S_w$;
    \item If $1\leq w\leq r$, we require that
    $$
u\neq\left\{
\begin{matrix}
i_w,&w\neq q,\\
i~\mathrm{or}~i_q,& w=q.
\end{matrix}
\right.
$$
\end{itemize}
The red summation in $M_{33}$ is taken over $w$ satisfying $1\leq w\leq r$ and $w\neq q$.

By direct computation, we have the following equalities: 
\begin{itemize}
    \item $M_{11}+M_{12}+M_{32}=0$;
    \item  $M_{13}+(M_{20}'+\sum_{0<u<r}(M_{2u}+M'_{2u})+M_{2r})+(M_{33}+M_{35})=0$.
\end{itemize}
It follows that the image of \eqref{psi L} under $\boldsymbol{e}$ equals
$$
M_{14}+\sum_{0<u<r}M_{2u}'+(M_{31}+M_{34}+M_{36})=\sum_{\underline i,\underline S,\underline{\overline j}}C_L(\underline i,\underline S,\underline{\overline j})j^1_{i_1}\cdots j^r_{i_r}\theta^{\underline{\overline j}}\theta_{\underline i}^{-1}\iota_{\tilde F_0}(F^*e_{\underline S,\underline i})\boldsymbol{\omega}_{\underline S}h^{\underline{\overline j}},
$$
where $C_L(\underline i,\underline S,\underline{\overline j})$ is the sum of 
$$
-\sum_{q\in Q_1}\sum_{i\in D^{-1}(2)}(j^q_i+\delta_{S_q}^i)C_1(q,\underline i(q,i),\underline S,\underline{\overline j}),~
\sum_{q\in Q_2}\sum_{i\in D^{-1}(1)}(j^q_i+\delta_{S_q}^i)C_1(q,\underline i,\underline S,\underline{\overline j}),
$$
$$
-\sum_{q\in Q_3}C_1(q,\underline i(\cup q),\underline S(\cup\underline q),\underline{\overline j}(\cup q)),~\sum_{q\in Q_4}C_1(q-1,\underline i(\cup q),\underline S(\cup q-1),\underline{\overline j}(\cup q-1)).
$$
Here, we set 
\begin{itemize}
    \item $Q_1:=\{q:i_q=1\},Q_2=\{q:i_q=2\}$,
    \item $Q_3:=\{q:D(i_q)=1,D(i_{q+1})=2\}, Q_4:=\{q:D(i_q)=2,D(i_{q+1})=1\}$,
    \item $\delta_{S_q}^i:=1$ if $i\in S_q$ and $0$ if $i\notin S_q$,
    \item $\underline i(q,i)$ is obtained by replacing the $q$-th entry $i_q$ of $\underline i$ by $i$.
\end{itemize}\fi
It is straightforward to verify that this image is equal to
$$
\sum_{(\underline S,\underline i,\underline{\overline j})\in T_\Omega(r,s+1)}c^R_{\underline S,\underline i,\underline{\overline j}}j^1_{i_1}\cdots j^r_{i_r}\iota_{\tilde F_0}(F^* 
(\theta^{\underline{\overline j}}\theta_{i_1}^{-1}\cdots\theta_{i_r}^{-1}
e_{\underline S,\underline i}))\boldsymbol{\omega}_{\underline S}h^{\underline{\overline j}}.
$$
Here, the coefficient $c^R_{\underline S,\underline i,\underline{\overline j}}$ are given as follows.
\begin{itemize}
    \item If there exists some $1\leq a\leq r$ such that 
    $D(i_u)=2$ for $u\leq a$ and $1$ for $u>a$.
    then 
    $$
    c^R_{\underline S,\underline i,\underline{\overline j}}:=\frac{1}{\prod_{1\leq u\leq a
    }\sum_{l=u}^a\Delta_l\prod_{a<u\leq r}\sum_{l=u}^a(j^l+s_l)}-\frac{1}{\prod_{1\leq u\leq r}\sum_{l=u}^r(j^l+s_l)}.
    $$
    Here, we stipulate that $\prod_{a<u\leq a}\sum_{l=u}^a(j^l+s_l):=1$.
    \item If $D(i_1)=\cdots=D(i_r)=1$ and $j^q_i=0$ for any $1\leq q\leq r$ and any $i\in D^{-1}(2)$, then 
    $$
    c^R_{\underline S,\underline i,    \underline{\overline j}}:=0;
    $$
    \item In the remaining situation, we set
    $$
    c^R_{\underline S,\underline i,\underline{\overline j}}:=-\frac{1}{\prod_{1\leq u\leq r}\sum_{l=u}^r(j^l+s_l)}.
    $$
\end{itemize}

\emph{Comparison between \eqref{psi L} and \eqref{psi R}}.
To verify that the images of \eqref{psi L} and \eqref{psi R} under \( \boldsymbol{e} \) coincide, it suffices to check that
\[
c^L_{\underline{S},\underline{i},\underline{\overline{j}}}
= c^R_{\underline{S},\underline{i},\underline{\overline{j}}}, \quad
c^L_{l,q,i,\underline{S},\underline{i},\underline{\overline{j}}} = 0, \quad
c^L_{u,q,i,\underline{S},\underline{i},\underline{\overline{j}}} = 0
\]
hold for all \( (\underline{S}, \underline{i}, \underline{\overline{j}}) \in T_{\Omega}(r, s+1) \), all \( (l, q, i, \underline{S}, \underline{i}, \underline{\overline{j}}) \in T'_{\underline{\Omega}}(r, s+1) \), and all \( (u, q, i, \underline{S}, \underline{i}, \underline{\overline{j}}) \in T''_{\underline{\Omega}}(r, s+1) \), respectively. The last equality is obvious. 

For the first equality, we check a typical situation, namely  \( (\underline{S}, \underline{i}, \underline{\overline{j}}) \in T_{\Omega}(r, s+1) \) meets the following condition: there exist $1< a,b,c<r$ such that 
\begin{itemize}
\item $c+2\leq b\leq a$,
    \item $D(i_{c-1})=1$, $D(i_c)=\cdots=D(i_{b-2})=2$, $D(i_{b-1})=1$, $D(i_b)=\cdots=D(i_a)=2$, $D(i_{a+1})=\cdots=D(i_r)=1$, and
    \item $j^q_i=0$ for $a<q\leq r,i\in D^{-1}(2)$.
\end{itemize}
In this situation, $c^L_{\underline{S},\underline{i},\underline{\overline{j}}}$ is equal to the sum of terms
$$
c^{L,a+1,2}_{\underline{S},\underline{i},\underline{\overline{j}}}=z\frac{\sum_{b\leq u\leq a}f_u}{\prod_{1\leq u\leq r}\sum_{u\leq l\leq r}(x_l+y_l)\prod_{c\leq u\leq a}\sum_{u\leq l\leq a}y_l},
$$
$$
c^{L,q,1}_{\underline{S},\underline{i},\underline{\overline{j}}}+c^{L,q,3}_{\underline{S},\underline{i},\underline{\overline{j}}}=x_q\frac{\sum_{b\leq u\leq q}f_u}{\prod_{1\leq u\leq r}\sum_{u\leq l\leq r}(x_l+y_l)\prod_{c\leq u\leq a}\sum_{u\leq l\leq a}y_l},~b\leq q\leq a,$$
$$
c^{L,b,2}_{\underline{S},\underline{i},\underline{\overline{j}}}+
c^{L,b-1,2}_{\underline{S},\underline{i},\underline{\overline{j}}}=\frac{y_{b-1}\sum_{c\leq u\leq b-1}f_u-(\sum_{b\leq u\leq a}y_u)f_{b-1}}{\prod_{1\leq u\leq r}\sum_{u\leq l\leq r}(x_l+y_l)\prod_{c\leq u\leq a}\sum_{u\leq l\leq a}y_l},$$
$$
c^{L,b-1,1}_{\underline{S},\underline{i},\underline{\overline{j}}}+c^{L,b-1,3}_{\underline{S},\underline{i},\underline{\overline{j}}}=-y_{b-1}\frac{\sum_{c\leq u\leq b-1}f_u}{\prod_{1\leq u\leq r}\sum_{u\leq l\leq r}(x_l+y_l)\prod_{c\leq u\leq a}\sum_{u\leq l\leq a}y_l}.$$
Here, we set 
$$
f_u:=\prod_{c\leq q<u}\sum_{l=q}^ay_l\prod_{u<q\leq a}\sum_{l=q}^r(x_l+y_l),~\prod_{c\leq q<c}\sum_{l=q}^ay_l:=1,~\prod_{a<q\leq a}\sum_{l=q}^r(x_l+y_l):=1$$
$$
x_u:=\sum_{i\in D^{-1}(1)}j^u_i+\#(S_u\cap D^{-1}(1)),~y_u:=\sum_{i\in D^{-1}(2)}j^u_i+\#(S_u\cap D^{-1}(2)),~1\leq u\leq r,$$
and 
$$z=\sum_{a<u\leq r}(x_l+y_l).$$
Using the identity
$$
(\prod_{u=c}^{b-1}\sum_{u\leq v\leq a}y_v)(\prod_{u=b}^a\sum_{u\leq l\leq r}(x_l+y_l)-\prod_{u=b}^a\sum_{u\leq l\leq a}y_l)=z\sum_{b\leq u\leq a}f_u+\sum_{b\leq u\leq a}y_u\sum_{b\leq l\leq u}f_l,
$$
the first equality $c^L_{\underline{S},\underline{i},\underline{\overline{j}}}
= c^R_{\underline{S},\underline{i},\underline{\overline{j}}}$ follows.

For the second equality, we check a typical situation, namely $(l, q, i, \underline{S}, \underline{i}, \underline{\overline{j}}) \in T'_{\underline{\Omega}}(r, s+1)$ meets the following condition: there exist $1<a,b<r$ such that
\begin{itemize}
    \item $b+2<l+2\leq q<a$, 
    \item $D(i_{b-1})=1$, $D(i_b)=\cdots=D(i_a)=2$, $D(i_{a+1})=\cdots=D(i_r)=1$, and
    \item $j^w_v=0$ for $a<w\leq r$ and $v\in D^{-1}(2)$.
\end{itemize}
In this situation, $c^L_{l, q, i, \underline{S}, \underline{i}, \underline{\overline{j}}}$ is equal to the sum of terms
$$
c^{L,l+1,2}_{l, q, i, \underline{S}, \underline{i}, \underline{\overline{j}}}=(-1)^{s+l+1}\frac{\sum_{l+1\leq u\leq r}(x_u+y_u)\sum_{l+2\leq u\leq q}f_u+\sum_{l+1\leq u\leq a}y_u\sum_{b\leq u\leq l}f_u}{\prod_{1\leq u\leq r}\sum_{u\leq v\leq r}(x_v+y_v)\prod_{b\leq u\leq a}\sum_{u\leq v\leq a}y_v},$$
$$
c^{L,l,2}_{l, q, i, \underline{S}, \underline{i}, \underline{\overline{j}}}=(-1)^{s+l}\frac{\sum_{l\leq u\leq r}(x_u+y_u)\sum_{l+1\leq u\leq q}f_u+\sum_{l\leq u\leq a}y_u\sum_{b\leq u\leq l-1}f_u}{\prod_{1\leq u\leq r}\sum_{u\leq v\leq r}(x_v+y_v)\prod_{b\leq u\leq a}\sum_{u\leq v\leq a}y_v},$$
$$
c^{L,l,1}_{l, q, i, \underline{S}, \underline{i}, \underline{\overline{j}}}+c^{L,l,3}_{l, q, i, \underline{S}, \underline{i}, \underline{\overline{j}}}=(-1)^{s+l+1}\frac{(x_l+y_l)\sum_{l+1\leq u\leq q}f_u+y_l\sum_{b\leq u\leq l}f_u}{\prod_{1\leq u\leq r}\sum_{u\leq v\leq r}(x_v+y_v)\prod_{b\leq u\leq a}\sum_{u\leq v\leq a}y_v}.$$
Here, we set $x_u,y_u$ as in the verification of the first equality and reset 
$$
f_u=\prod_{b\leq v<u}\sum_{w=v}^ay_w\prod_{u<v\leq a}\sum_{w=v}^r(x_w+y_w),~\prod_{b\leq v<b}\sum_{w=v}^ay_w:=1,~\prod_{a<v\leq a}\sum_{w=v}^r(x_w+y_w):=1.$$
By direct computation, we have $c^L_{l, q, i, \underline{S}, \underline{i}, \underline{\overline{j}}}=0$, the second equality follows.

To check the last equality, it suffices to observe that 
$$
c^{L,1}_{u, q, i, \underline{S}, \underline{i}, \underline{\overline{j}}}=-c^{L,3}_{u, q, i, \underline{S}, \underline{i}, \underline{\overline{j}}}.
$$
This completes the proof of the proposition. \qed

\subsection{Compatibility}
Set $\Omega := \Omega^1_{X/k}(\log D)$. In \cite{SZ2}, Sheng and the author proved that the morphism $\mathrm{Ho}_\Omega$ is compatible with the intersection subcomplexes, namely $\Omega_{\mathrm{int}}(E, \theta)$ and $\Omega_{\mathrm{int}}(H, \nabla)$. In the previous two subsections, we introduced several constructions that depend on splittings of \( \Omega \). It is therefore natural to expect that these constructions are compatible with the intersection subcomplexes again, and even possibly with the complexes introduced in \S2.6.
 However, the incompatibility between the divisor $D$ and a general ordered splitting $\underline{\Omega}$ may cause the $\infty$-homotopy $\mathrm{Ho}_{\underline{\Omega}}$ to fail to be compatible with the intersection subcomplexes. It is therefore necessary to impose appropriate compatibility conditions on the pair $(D, \underline{\Omega})$.

\begin{definition}
Let $\underline{\Omega} = (\Omega_1, \ldots, \Omega_\beta)$ be an ordered splitting of $\Omega$. We say that $\underline{\Omega}$ is compatible with the divisor $D = \sum_{i \in I} D_i$ (where each $D_i$ is an irreducible component of $D$) if the index sets $I_1, \ldots, I_\beta$ are pairwise disjoint, where
\[
I_j := \left\{ i \in I \mid \mathrm{Res}_{D_i}|_{\Omega_j} \neq 0 \right\}, \quad 1 \leq j \leq \beta.
\]
Suppose that $g:(X,D)\to(\mathbb A_k^1,\Delta)$ is a semistable family. We say that $\underline{\Omega}$ is compatible with $g$ if there exists some $1 \leq i \leq \beta$ such that
$dg \in \Gamma(X, \Omega_i)$.

\end{definition}
Let $(K^*_\mathrm{Hig},K^*_\mathrm{dR})$ be a Hodge pair of complexes of type $(\star,\ell)$ on $(X,D)/k$, where $\star\in\{\mathrm{\uppercase\expandafter{\romannumeral1}}, \mathrm{\uppercase\expandafter{\romannumeral2}} 
 ,\mathrm{\uppercase\expandafter{\romannumeral3}}\}
$ and $0\leq \ell<p$ or $\star=\mathrm{\uppercase\expandafter{\romannumeral4}}$ and $\ell=0$ (see \S2.6).
Let $\underline{\Omega}$ be an ordered splitting of $\Omega$, and let $\underline{\Omega}'$ be as defined in \eqref{case i}. Suppose that in types $(\mathrm{\uppercase\expandafter{\romannumeral1}},\ell)$ and $(\mathrm{\uppercase\expandafter{\romannumeral2}},\ell)$, the splitting $\underline{\Omega}$ is compatible with the divisor $D$, while in types $(\mathrm{\uppercase\expandafter{\romannumeral3}},\ell)$ and $(\mathrm{\uppercase\expandafter{\romannumeral4}},0)$, it is compatible with $g:(X,D)\to(\mathbb A^1_k,\Delta)$, which is the semistable family appearing in the construction of the pair $(K^*_{\mathrm{Hig}},K^*_{\mathrm{dR}})$. 

\begin{definition}
Assume that $(K^*_{\mathrm{Hig}},K^*_{\mathrm{dR}})$ has type $(\mathrm{\uppercase\expandafter{\romannumeral4}},0)$.Let $\Delta=\sum_{i=1}^m[\lambda_i]$ with each $\lambda_i\in k$, and let
$D_i:=g^*[\lambda_i]$.
Define the formal $g$-transform along $D$ of  $(K^*_{\mathrm{Hig}},K^*_{\mathrm{dR}})$ to be the pair $(K^*_{\mathrm{Hig},f},K^*_{\mathrm{dR},f})$, where
$$
\begin{array}{c}
K^*_{\mathrm{Hig},f}:=\bigoplus_{i=1}^m((K^\bullet_{\mathrm{Hig}})_{/D_i'},-\sum_{j=0}^\infty(f'-\lambda'_i)^{p^j-1}df'\wedge),\\
K^*_{\mathrm{dR},f}:=\bigoplus_{i=1}^m((K^\bullet_{\mathrm{dR}})_{/pD_i},d-\sum_{j=1}^\infty(f'-\lambda'_i)^{p^j-1}df'\wedge).
\end{array}
$$

\end{definition}







\begin{theorem}\label{compatibility MHM}
Notation and assumptions as above. 
 Then the $\mathcal L$-indexed $\infty$-homotopy 
 $\mathrm{Ho}_{\underline\Omega}$ induces an $\mathcal L_\star$-indexed $\infty$-homotopy 
 \begin{eqnarray}\label{compatibility 1}
 \mathrm{Ho}_{\underline\Omega,\star}:\Delta_*(\mathcal L_\star)\to\left\{\begin{matrix}
 \mathcal Hom^*(\tau_{<q}K^*_\mathrm{Hig},\tau_{<q}F_*K^*_\mathrm{dR}),&\star=\mathrm{\uppercase\expandafter{\romannumeral1}}, \mathrm{\uppercase\expandafter{\romannumeral2}}, 
 \mathrm{\uppercase\expandafter{\romannumeral3}},\\
 \mathcal Hom^*(\tau_{<q}K^*_{\mathrm{Hig},f},\tau_{<q}F_*K^*_{\mathrm{dR},f}),&\star=\mathrm{\uppercase\expandafter{\romannumeral4}},
 \end{matrix}
 \right.
 \end{eqnarray}
where $q = \infty$ if $\mathrm{rank}(\underline\Omega) < p - \ell$, and $q = p - \ell$ otherwise. Moreover, $\mathrm{Ho}_{\underline\Omega,\underline\Omega'}$ induces a graded morphism  
\begin{eqnarray}\label{compatibility 2}
\mathrm{Ho}_{\underline\Omega,\underline\Omega',\star}:\Delta_\bullet(\mathcal L_\star)\to\left\{\begin{matrix}
 \mathcal Hom^{\bullet-1}(\tau_{<q}K^*_\mathrm{Hig},\tau_{<q}F_*K^*_\mathrm{dR}),&\star=\mathrm{\uppercase\expandafter{\romannumeral1}}, \mathrm{\uppercase\expandafter{\romannumeral2}}, 
 \mathrm{\uppercase\expandafter{\romannumeral3}},\\
 \mathcal Hom^{\bullet-1}(\tau_{<q}K^*_{\mathrm{Hig},f},\tau_{<q}F_*K^*_{\mathrm{dR},f}),&\star=\mathrm{\uppercase\expandafter{\romannumeral4}},
 \end{matrix}
 \right.
    \end{eqnarray}
    where $q = \infty$ if both $\mathrm{rank}(\underline{\Omega})$ and $\mathrm{rank}(\underline{\Omega}')$ are less than $p - \ell$, and $q = p - \ell$ otherwise. 

In particular, if the pair $(K^*_\mathrm{Hig},K^*_\mathrm{dR})$ has 
type $(\mathrm{\uppercase\expandafter{\romannumeral1}},\ell), (\mathrm{\uppercase\expandafter{\romannumeral2}},\ell),(\mathrm{\uppercase\expandafter{\romannumeral3}},\ell),$
or
$(\mathrm{\uppercase\expandafter{\romannumeral4}},0)$, then
$\underline\Omega$ and $\underline\Omega'$ induce the same de Rham-Higgs comparison
\begin{eqnarray}\label{de Rham-Higgs 4.3}
\tau_{<q}F_*K^*_\mathrm{dR}\cong\tau_{<q}K^*_\mathrm{Hig}~\mathrm{in}~D(X'),
\end{eqnarray}
where $q$ is as defined below \eqref{compatibility 1} (note that in the last type, $\ell=0$).
\end{theorem}

\iffalse
\emph{Case 1}. For $\star=\Omega^*_\mathrm{int},W_i\Omega^*$, the compatibility means the following.
\begin{itemize}
    \item 
$\mathrm{Ho}_{\underline \Omega}$ restricts to an $\mathcal L_{\uppercase\expandafter{\romannumeral1}}$-indexed $\infty$-homotopy 
    \begin{eqnarray}\label{compatibility 1}
    \Delta_*(\mathcal L)\to\mathcal Hom^*(\tau_{<q}\star(E,\theta),\tau_{<q}F_*\star(H,\nabla)),
    \end{eqnarray}
\item  $\mathrm{Ho}_{\underline\Omega,\underline\Omega'}$ restricts to a graded morphism  
   \begin{eqnarray}\label{compatibility 2}
\Delta_\bullet(\mathcal L)\to\mathcal Hom^{\bullet-1}(\tau_{<q}\star(E,\theta),\tau_{<q}F_*\star(H,\nabla)),
    \end{eqnarray}
    where $q = \infty$ if both $\mathrm{rank}(\underline{\Omega})$ and $\mathrm{rank}(\underline{\Omega}')$ are less than $p - \ell$, and $q = p - \ell$ otherwise.
\end{itemize}

\emph{Case 2}.
Let $f : (X, D) \to (\mathbb{A}^1_k, E)$ be a semistable family over $k$ along some effective divisor $E$ on $\mathbb{A}^1_k$ containing $0$. Suppose that there exists a semistable family $\tilde{f} : (\tilde{X}, \tilde{D}) \to (\mathbb{A}^1_{W_2(k)}, \tilde{E})$ on $W_2(k)$ along an effective divisor $\tilde{E}$ on $\mathbb{A}^1_{W_2(k)}$, such that the reductions modulo $p$ of $\tilde{f}, \tilde{E}$ are $f, E$, respectively. Suppose further that $df$ is contained in $\Omega_i$ for some $1 \leq i \leq \beta$. Set $\mathcal{L}_f \subset \mathcal{L}$ to be the subsheaf consisting of sections $\tilde{F} \in \mathcal{L}$ satisfying $\tilde{F}(\tilde{f}) = \tilde{f}^p$, and define
$\mathcal{L}_{f/D} := \mathcal{L}_f|_{\mathfrak{X}}$ with $\mathfrak{X} := X_{/D}$.
 The remaining part of compatibility means the following.
\begin{itemize}
    \item $\mathrm{Ho}_{\underline \Omega}$ restricts to an $\mathcal L_f$-indexed $\infty$-homotopy
    \begin{eqnarray*}
    \Delta_*(\mathcal L_f)\to\mathcal Hom^*(\tau_{<q}\Omega^*_{f^{-1}}(E,\theta),\tau_{<q}F_*\Omega^*_{f^{-1}}(H,\nabla)),
    \end{eqnarray*}  
    where \( q \) is as defined immediately after \eqref{compatibility 1}.
   \item $\mathrm{Ho}_{\underline\Omega,\underline\Omega'}$ restricts to a graded morphism
   \begin{eqnarray*}
\Delta_\bullet(\mathcal L_f)\to\mathcal Hom^{\bullet-1}(\tau_{<q}\Omega^*_{f^{-1}}(E,\theta),\tau_{<q}F_*\Omega^*_{f^{-1}}(H,\nabla)),
    \end{eqnarray*}
    where \( q \) is as defined immediately after \eqref{compatibility 2}.
  \item For $(E, \theta) = (\mathcal{O}_X, 0)$, let $K^*_{\mathrm{dR}} \subset \Omega^*(\mathcal{O}_X, d)$ denote the subcomplex defined just before Corollary~\ref{intro tau<p structure sheaf}. Then
  $\mathrm{Ho}_{\underline \Omega}$ restricts to an $\mathcal L_{f/D}$-indexed $\infty$-homotopy
    \begin{eqnarray*}
    \Delta_*(\mathcal L_{f/D})\to\mathcal Hom^*(\tau_{<q}(K^\bullet_\mathrm{dR}\otimes\mathcal O_{\mathfrak X},-\sum_{i=0}^\infty f^{p^i-1}df\wedge),\tau_{<q}F_*(K^\bullet_\mathrm{dR}\otimes\mathcal O_{\mathfrak X},d+df\wedge)),
    \end{eqnarray*}  
    $q = \infty$ if $\mathrm{rank}(\underline\Omega) < p $, and $q = p$ otherwise.
    \item Notation as above. Then $\mathrm{Ho}_{\underline \Omega,\underline\Omega'}$ restricts to a graded morphism
    \begin{eqnarray*}
    \Delta_\bullet(\mathcal L_{f/D})\to\mathcal Hom^{\bullet-1}(\tau_{<q}(K^\bullet_\mathrm{dR}
    \otimes\mathcal O_{\mathfrak X},-\sum_{i=0}^\infty f^{p^i-1}df\wedge),\tau_{<q}F_*(K^\bullet_\mathrm{dR}
    \otimes\mathcal O_{\mathfrak X},d+df\wedge)),
    \end{eqnarray*}   
    where $q = \infty$ if both $\mathrm{rank}(\underline{\Omega})$ and $\mathrm{rank}(\underline{\Omega}')$ are less than $p$, and $q = p$ otherwise. 
    \end{itemize}
\fi
\begin{proof}
Clearly, if the Hodge pair $(K^*_{\mathrm{Hig}}, K^*_{\mathrm{dR}})$ is of type $(\mathrm{\uppercase\expandafter{\romannumeral1}}, \ell)$, $(\mathrm{\uppercase\expandafter{\romannumeral2}}, \ell)$, or $(\mathrm{\uppercase\expandafter{\romannumeral3}}, \ell)$, then \eqref{de Rham-Higgs 4.3} is a direct consequence of \eqref{compatibility 1} and \eqref{compatibility 2}.
If it is of type $(\mathrm{\uppercase\expandafter{\romannumeral4}}, 0)$, then there are isomorphisms of complexes
\begin{eqnarray}\label{f-exponential}
\begin{array}{ll}
K^*_{\mathrm{Hig}} \cong K^*_{\mathrm{Hig}} \otimes \mathcal{O}_{\mathfrak{X}}, & 
F_* K^*_{\mathrm{dR}} \cong F_* K^*_{\mathrm{dR}} \otimes \mathcal{O}_{\mathfrak{X}}, \\
K^*_{\mathrm{Hig}} \otimes \mathcal{O}_{\mathfrak{X}} \cong K^*_{\mathrm{Hig}, f}, & 
F_* K^*_{\mathrm{dR}} \otimes \mathcal{O}_{\mathfrak{X}} \cong F_* K^*_{\mathrm{dR}, f}.
\end{array}
\end{eqnarray}
The first isomorphism follows from the fact that $K^*_{\mathrm{Hig}}$ is acyclic outside $D$. Combining it with Corollary~\ref{corollary 2}, we get the second isomorphism. 
The third isomorphism is given by multiplication by $\left( \sum_{i=0}^\infty f^{p^i - 1} \right)^j$ on $K^j_{\mathrm{Hig}} \otimes \mathcal{O}_{\mathfrak{X}}$, and the fourth isomorphism is given by multiplication by the Artin–Hasse exponential of $f$, namely,
\[
\mathrm{AH}(f) = \exp\left( \sum_{i=0}^\infty f^{p^i}/p^i \right),
\]
on $K^j_{\mathrm{dR}} \otimes \mathcal{O}_{\mathfrak{X}}$.
Combining these isomorphisms with \eqref{compatibility 1} and \eqref{compatibility 2} again, we conclude that \eqref{de Rham-Higgs 4.3} also holds in the case of type $(\mathrm{\uppercase\expandafter{\romannumeral4}}, 0)$.

By Theorem~\ref{main thm s2} and its proof, if \eqref{compatibility 1} is well-defined, then it is an $\infty$-homotopy. Hence, it remains to verify that \eqref{compatibility 1} and \eqref{compatibility 2} are well-defined. Furthermore, it suffices to check that \eqref{compatibility 1} is well-defined in the case of $\star = \mathrm{\uppercase\expandafter{\romannumeral4}}$, since the remaining cases can be treated similarly.

We now verify this specific case. Without loss of generality, we may assume that \( X \) admits a system of coordinates \( t_1, \cdots, t_n \) such that \( f = t_1 \cdots t_n \) and \( D = (f = 0) \), and that
\[
K^*_{\mathrm{Hig},f} = \left( \Omega^\bullet_{\mathrm{int}}(\hat E, \hat\theta), \hat\theta - \sum_{i=0}^\infty f^{p^i - 1} \, df \wedge \right), \quad
K^*_{\mathrm{dR},f} = \left( \Omega^\bullet_{\mathrm{int}}(\hat H, \hat\nabla), \hat\nabla - \sum_{i=1}^\infty f^{p^i - 1} \, df \wedge \right),
\]
where \( (\hat E, \hat\theta) := (E, \theta)_{/D} \) and \( (\hat H, \hat\nabla) := (H, \nabla)_{/D} \).
Given any \( r > 0 \), \( s \geq 0 \), and any \( \tilde F_0, \cdots, \tilde F_r \in \Gamma(\mathfrak U, \mathcal L_{\mathrm{\uppercase\expandafter{\romannumeral4}}}) \) for any open subset \( \mathfrak U \subset \mathfrak X \), it suffices to check that
\[
\varphi(r,s)_{\tilde F_0, \tilde F_1, \cdots, \tilde F_r;\underline\Omega}(K^{r+s}_{\mathrm{Hig},f}) \subset F_* K^s_{\mathrm{dR},f}.
\]
Without loss of generality, this inclusion can be checked by showing that
\begin{eqnarray}\label{compatibility inclusion}
\varphi(r,s)_{\tilde F_0, \tilde F_1, \cdots, \tilde F_r;\underline\Omega} \left( \hat\theta_1 \cdots \hat\theta_{r+s} \hat E \otimes d\log t_1 \wedge \cdots \wedge d\log t_{r+s} \right) \subset F_* K^s_{\mathrm{dR},f}.
\end{eqnarray}
where \( \theta = \sum_{i=1}^n \theta_i \otimes d\log t_i \) and \( \hat\theta_i := (\theta_i)_{/D} \). Since \( \underline\Omega \) is compatible with \( D \), we have indices \( 1 \leq i_1, \cdots, i_n \leq \beta \) such that
\[
d\log t_j = \omega^j_1 + \cdots + \omega^j_\beta, \quad \omega^j_q \in \Omega_q \cap \Omega^1_{X/k} \text{ for } q \neq i_j.
\]
Combining this decomposition of \( d\log t_j \) with the base-free description presented in the proof of Lemma \ref{indep of basis}, the left-hand side of \eqref{compatibility inclusion} is contained in
\[
\begin{array}{c}
F_* \iota_{\tilde F_0} \left( \hat\theta_1 \cdots \hat\theta_{r+s} \hat E \right) \otimes \left[ \bigoplus_{i=0}^{r+s} \wedge^i \left\{ \zeta_{\tilde F_0}(\omega_{i_1}), \cdots, \zeta_{\tilde F_0}(\omega_{i_{r+s}}) \right\} \otimes \Omega^{r+s-i}_{X/k} \right] \\
\subset F_* \iota_{\tilde F_0} \left( \hat\theta_1 \cdots \hat\theta_{r+s} \hat E \right) \otimes \left[ \bigoplus_{i=0}^{r+s} \wedge^i \left\{ \zeta_{\tilde F_0}(d\log t_1), \cdots, \zeta_{\tilde F_0}(d\log t_{r+s}) \right\} \otimes \Omega^{r+s-i}_{X/k} \right] \\
\subset F_* K^s_{\mathrm{dR},f}.
\end{array}
\]
This completes the proof.
\end{proof}
\begin{remark}
    We hope that \eqref{de Rham-Higgs 4.3} holds even when \( (K^*_{\mathrm{Hig}}, K^*_{\mathrm{dR}}) \) has type \( (\mathrm{\uppercase\expandafter{\romannumeral4}}, \ell) \) with \( \ell \neq 0 \) (in other words, \( \theta \neq 0 \)). To this end, the key point is to find an isomorphism between \( \Omega^*(\hat E, \hat\theta - df \wedge) \) and \( \Omega^*(\hat E, \hat\theta - f^{p^i - 1} df \wedge) \). Note that the desired isomorphism always exists locally due to the existence of a nice system of coordinates, as shown in the proof of Theprem \ref{main thm s2}. To glue together these local isomorphisms, we need to construct an \( \mathcal{L}_\mathrm{\uppercase\expandafter{\romannumeral4}} \)-indexed \( \infty \)-homotopy between the aforementioned complexes that is compatible with the subcomplexes, inspired by the theory of mixed Hodge modules.
\end{remark}




\subsection{Two-term truncations}

In general, it is not reasonable to expect that the logarithmic cotangent bundle $\Omega^1_{X/k}(\log D)$ splits into a direct sum of line bundles. However, this expectation can be realized locally. Once we have such an ordered splitting on some Zariski open subset $U\subset X$, say $\underline\Omega=(\Omega_1,\cdots,\Omega_n)$ with $\mathrm{rank}(\Omega_i)=1$, then \S3.1 showed that for any nilpotent Higgs sheaf $(E,\theta)$ on $(X,D)/k$ of level $\leq p-2$, there is an $\mathcal L|_U$-indexed $\infty$-homotopy
\begin{eqnarray}\label{homotopy|_U}
\mathrm{Ho}_{\underline{\Omega}}:\Delta_*(\mathcal L|_U)\to\mathcal Hom_{\mathcal O_U}^*(\Omega^*(E|_U,\theta|_U),F_*\Omega^*(H|_U,\nabla|_U)).
\end{eqnarray}
Inspired by the refinements of the classical Deligne-Illusie decomposition theorem (see Achinger-Suh \cite{AS}, Drinfeld, Bhatt-Lurie \cite{BL}, Li-Mondal \cite{LM}, and Ogus \cite{O24}), we observe that \eqref{homotopy|_U} is independent of the choice of \( \underline\Omega \) after applying the two-term truncation \( \tau_{[a, a+1]} \) and strengthening the nilpotency condition on \( (E, \theta) \) to \( \leq p-3 \). This leads to the following.

\iffalse
More precisely, if there is another splitting $\Omega^1_{X/k}(\log D)|_U=\bigoplus_{i=1}^{\beta'}\Omega'_i$ with $\mathrm{rank}(\Omega'_i)=1$, then $\mathrm{Ho}_{\underline{\Omega},a}=\mathrm{Ho}_{\underline{\Omega'},a}$. Here, 
$$
\mathrm{Ho}_{\underline{\Omega},a}(\mathrm{resp.}~\mathrm{Ho}_{\underline{\Omega}',a}):\Delta_*(\mathcal L|_U)\to\mathcal Hom_{\mathcal O_U}^*(\tau_{[a,a+1]}\Omega^*(E|_U,\theta|_U),\tau_{[a,a+1]}F_*\Omega^*(H|_U,\nabla|_U))
$$
is induced by $\mathrm{Ho}_{\underline\Omega}$ (resp. $\mathrm{Ho}_{\underline\Omega'}$). In particular, we have an $\mathcal L$-indexed $\infty$-homotopy
$$\mathrm{Ho}_a:\Delta_*(\mathcal L)\to\mathcal Hom_{\mathcal O_X}^*(\tau_{[a,a+1]}\Omega^*(E,\theta),\tau_{[a,a+1]}F_*\Omega^*(H,\nabla)).$$
Furthermore, for $\star=\Omega^*_{\mathrm{int}},W_i\Omega^*$, $\mathrm{Ho}_a$ restricts to an $\mathcal L_\star$-indexed $\infty$-homotopy
$$
\mathrm{Ho}_{a,\star}:\Delta_*(\mathcal L_\star)\to\mathcal Hom^*(\tau_{[a,a+1]}\star(E,\theta),\tau_{[a,a+1]}F_*\star(H,\nabla));
$$
for $\star=\Omega^*_{f^{-1}}$, $\mathrm{Ho}_a$ restricts to an $\mathcal L_f$-indexed $\infty$-homotopy
$$
\mathrm{Ho}_{a,\star}:\Delta_*(\mathcal L_f)\to\mathcal Hom^*(\tau_{[a,a+1]}\star(E,\theta),\tau_{[a,a+1]}F_*\star(H,\nabla));
$$
for $\star=W_i\Omega^*,\Omega^*_{f^{-1}}$ and $(E,\theta)=(\mathcal O_X,-df\wedge)$, $\mathrm{Ho}_a$ induces an $\mathcal L_{df\wedge}$-indexed $\infty$-homotopy
$$
\mathrm{Ho}_{a,\star}:\Delta_*(\mathcal L_{df\wedge})\to\mathcal Hom^*(\tau_{[a,a+1]}\star(\mathcal O_{\mathfrak X},-df\wedge),\tau_{[a,a+1]}F_*\star(\mathcal O_{\mathfrak X},d+df\wedge)).
$$It follows that\fi

\begin{theorem}\label{two sided truncation}
Suppose that $p\geq 3$ and given any $a\geq0$. Let $(K^*_\mathrm{Hig},K^*_\mathrm{dR})$ be a Hodge pair of complexes of type $(\star,p-3)$ $(\star\in\{\mathrm{\uppercase\expandafter{\romannumeral1}}, \mathrm{\uppercase\expandafter{\romannumeral2}} 
 ,\mathrm{\uppercase\expandafter{\romannumeral3}}\})$ or $(\mathrm{\uppercase\expandafter{\romannumeral4}},0)$. Then
$$
\tau_{[a,a+1]}F_*K^*_\mathrm{dR}\cong \tau_{[a,a+1]}K^*_\mathrm{Hig}~\mathrm{in}~D(X').
$$
\end{theorem}
\begin{proof}
We verify the last type; the first three types can be checked in a similar manner. Without loss of generality, we may assume that \( X \) admits a system of coordinates \( t_1, \cdots, t_n \) such that \( f = t_1 \cdots t_n \) and \( D = (f = 0) \), and that
\[
K^*_{\mathrm{Hig}} = (\Omega^\bullet_{f^{-1}}, -df \wedge), \quad K^*_{\mathrm{dR}} = (\Omega^\bullet_{f^{-1}}, d + df \wedge).
\]
Set $\underline\Omega:=(\Omega_1,\cdots,\Omega_n)$, where $\Omega_1:=\mathcal O_Xd\log f$ and $\Omega_i=\mathcal O_Xd\log t_i$ for $2\leq i\leq n$. Obviously, $\underline\Omega$ is compatible with $f$. Taking into account Theorem \ref{compatibility MHM}, we have an \( \mathcal{L}_{\mathrm{\uppercase\expandafter{\romannumeral4}}} \)-indexed \( \infty \)-homotopy
$$
\mathrm{Ho}_{\underline\Omega,\mathrm{\uppercase\expandafter{\romannumeral4}}}:\Delta_*(\mathcal L_\mathrm{\uppercase\expandafter{\romannumeral4}})\to\mathcal Hom^*(K^*_{\mathrm{Hig},f},F_*K^*_{\mathrm{dR},f}),
$$
which induces another \( \mathcal{L}_{\mathrm{\uppercase\expandafter{\romannumeral4}}} \)-indexed \( \infty \)-homotopy
$$
\mathrm{Ho}_{\underline\Omega,\mathrm{\uppercase\expandafter{\romannumeral4}},a}:\Delta_*(\mathcal L_\mathrm{\uppercase\expandafter{\romannumeral4}})\to\mathcal Hom^*(\tau_{[a,a+1]}K^*_{\mathrm{Hig},f},\tau_{[a,a+1]}F_*K^*_{\mathrm{dR},f}).
$$
To prove this theorem, it suffices to show that this $\infty$-homotopy is independent of the choice of $\underline\Omega$.

Let $m_1,\cdots,m_n$ be another system of coordinates on $X/k$ such that $f=m_1\cdots m_n$. Set $\underline\Omega':=(\Omega_1',\cdots,\Omega_n')$, where $\Omega_1':=\mathcal O_Xd\log f$ and $\Omega_i':=\mathcal O_Xd\log m_i$ for $2\leq i\leq n$.  Now, the aforementioned independence reduces to the equality 
$$
\mathrm{Ho}_{\underline\Omega,\mathrm{\uppercase\expandafter{\romannumeral4}},a}=\mathrm{Ho}_{\underline\Omega',\mathrm{\uppercase\expandafter{\romannumeral4}},a}.
$$
One can check that if there exists a homotopy between \( \mathrm{Ho}_{\underline\Omega, \mathrm{\uppercase\expandafter{\romannumeral4}}} \) and \( \mathrm{Ho}_{\underline\Omega', \mathrm{\uppercase\expandafter{\romannumeral4}}} \), then the above equality follows. According to Theorem \ref{compatibility MHM} again, to construct such a homotopy, it suffices to find a chain of \( f \)-compatible ordered splittings of \( \Omega^1_{X/k}(\log D) \) between \( \underline\Omega \) and \( \underline\Omega' \), say
\begin{eqnarray}\label{chain ordered splitting}
\underline\Omega_0 := \underline\Omega, \ \underline\Omega_1, \ \cdots, \ \underline\Omega_c := \underline\Omega',
\end{eqnarray}
such that \( \mathrm{rank}(\underline\Omega) \leq 2 \) for \( 1 \leq i \leq c \), and for \( 1 \leq i < c \), either \( \underline\Omega_i \) is constructed from \( \underline\Omega_{i+1} \) or vice versa, via \eqref{case i}.

We show the existence of the chain in \eqref{chain ordered splitting}. Set
\[
(d\log f, d\log m_2, \cdots, d\log m_n) = (d\log f, d\log t_2, \cdots, d\log t_n) A, \quad A \in \mathrm{GL}_n(\Gamma(X, \mathcal{O}_X)).
\]
By basic linear algebra and shrinking \( X \) if necessary, we may assume that \( A \) can be decomposed as a product of elementary matrices, namely
\begin{eqnarray}\label{decomposition A}
A = A_1 \cdots A_s.
\end{eqnarray}
Moreover, we can further assume that each \( A_i \) is one of the following:
\begin{itemize}
    \item The elementary matrix obtained by adding \( \lambda \) times the first column of the identity matrix \( I_n \) to its second column;
    \item The elementary matrix obtained by interchanging the \( j \)-th and \( (j+1) \)-th columns of \( I_n \) (\( 2 \leq j < n \));
    \item The elementary matrix obtained by adding \( \lambda \) times the \( j \)-th column of the identity matrix \( I_n \) to its \( (j+1) \)-th column.
\end{itemize}
Consequently, the decomposition of \( A \) yields a chain of bases for \( \Omega^1_{X/k}(\log D) \) over \( \mathcal{O}_X \), namely
\[
(\omega_1^i, \cdots, \omega_n^i), \quad 0 \leq i \leq s,
\]
subject to the following conditions:
\begin{itemize}
    \item \( (\omega_1^0, \cdots, \omega_n^0) = (d\log f, d\log t_2, \cdots, d\log t_n) \);
    \item \( (\omega_1^s, \cdots, \omega_n^s) = (d\log f, d\log m_2, \cdots, d\log m_n) \);
    \item \( (\omega_1^i, \cdots, \omega_n^i) = (\omega_1^{i-1}, \cdots, \omega_n^{i-1}) A_i \), where \( 1 \leq i \leq s \).
\end{itemize}
Clearly, the basis \( (\omega_1^i, \cdots, \omega_n^i) \) gives rise to an $f$-compatible ordered splitting of \( \Omega^1_{X/k}(\log D) \), denoted by \( \underline\Omega_{2i} \). For \( 1 \leq i \leq s \), denote by \( \underline\Omega_{2i-1} \) the ordered splitting constructed from \( \underline\Omega_{2i} \) via \eqref{case i} by setting
\[
o := \left\{
\begin{matrix}
1, & A_i \text{ satisfies the first case}; \\
j, & \text{ otherwise}.
\end{matrix}
\right.
\]
Finally, it can be verified that \( \underline\Omega_0, \cdots, \underline\Omega_{2s} \) forms the desired chain of \( f \)-compatible ordered splittings. This completes the proof.
\end{proof}


\section{Further consequences}

\iffalse
\subsection{Gerbes of splittings}

In this subsection, we do not assume that $(X,D)/k$ is $W_2(k)$-liftable. In \cite[\S3]{DI}, Deligne and Illusie considered the following gerbes:
\begin{itemize}
    \item (See also \cite[\S3.2]{AS}) The gerbe of splittings of $\mathcal K_0^*:=\tau_{\leq 1}F_*\Omega^*_{X/k}(\log D)$, denoted by $\mathrm{sc}(\mathcal K_0^*)$. It is defined as the stackification of the prestack \(\text{sc}'(\mathcal K_0^*)\), whose objects are local splittings
\( s: \mathcal H^1(\mathcal K_0^*) \to \mathcal K_0^1 \)
of the projection \( \mathcal K_0^1  \to \mathcal H^1(\mathcal K_0^*) \), where morphisms \( s \to s' \) are maps
\( h: \mathcal H^1(\mathcal K_0^*) \to \mathcal K_0^0 \)
such that \( dh = s' - s \). The automorphisms of an object \( s \) are then identified with sections of \( \mathcal Hom(\mathcal H^1(\mathcal K_0^*), \mathcal H^0(\mathcal K_0^*)) \), and this makes \( \text{sc}(\mathcal K_0^*) \) into a gerbe under \(\mathcal Hom(\mathcal H^1(\mathcal K_0^*), \mathcal H^0(\mathcal K_0^*))=T_{X/k}(-\log D)\);
\item The gerbe of lifts of \( (X,D)\) over \( W_2(k) \), denoted by \( \text{Rel}((X,D), W_2(k)) \), is the gerbe over $X$, whose objects over an open subset \( U \) are pairs \( (Y,D_Y) \), where $Y$ is a smooth \( W_2(k) \)-sheme, $D_Y\subset Y$ is a simple normal crossing divisor relative to $W_2(k)$, and its reduction modulo \( p \) is \( (U,D|_U) \). A morphism \( (Y',D_Y') \to (Y'',D_Y'') \) is a $W_2(k)$-morphism $u:Y'\to Y''$, with $u^*D_Y''=D_Y'$ and reduction modulo \( p \) equal to the identity. The sheaf of automorphisms of any lift of \( U \) is the abelian sheaf \( T_{X/k}(-\log D)\).
\end{itemize}
In \cite[Proposition 3.3]{DI}, the authors proved that the class of the gerbe $\mathrm{sc}(\mathcal K_0^*)$ is equal to $-\delta^1_{\mathcal K^*_0}$, where $\delta^1_{\mathcal K^*_0}$ fits into the distinguished triangle 
$$
\mathcal O_X\to\mathcal K_0^*\to\Omega^1_{X/k}(\log D)[-1]\stackrel{\delta^1_{\mathcal K^*_0}}{\to}\mathcal H^0(\mathcal K^*_0)[1].$$
Moreover, as shown in \cite[Th\'eorème 3.5]{DI}, there is a natural equivalence between the above gerbes, which associates to a lift $(Y,D_Y)$ the induced splitting of $\mathcal K_0^*$.

For any $a\geq0$, set $\mathcal K^*_a:=\tau_{[a,a+1]}\Omega^*_{X/k}(\log D)$. In \cite[\S 3]{AS}, Achinger and Suh considered the gerbe of splittings of $\mathcal K^*_a:=\tau_{[a,a+1]}\Omega^*_{X/k}(\log D)$ under $\mathcal Hom_{\mathcal O_X}(\Omega^{a+1}_{X/k}(\log D),\Omega^a_{X/k}(\log D))$.  They constructed a functor $\wedge^a$ from $\mathrm{sc}(\mathcal K^*_0)$ to $\mathrm{sc}(\mathcal K^*_a)$, and established a relation between the classes of $\mathrm{sc}(\mathcal K^*_0)$ and $\mathrm{sc}(\mathcal K^*_a)$.

Next, we take a global section $f\in\Gamma(X,\mathcal O_X)$ such that $(f=0)\subset D$. We demonstrate that all the discussion above can be extended to $\Omega^*_{f^{-1}}$. More precisely, we consider the following  gerbes:
\begin{itemize}

    \item The gerbe of splittings of $\mathcal G_a^*:=\tau_{[a,a+1]}\Omega^*_{f^{-1}}$ under $\mathcal Hom_{\mathcal O_X}(\Omega^{a+1}_{f^{-1}},\Omega^a_{f^{-1}})$ for $a\geq 0$;
    
    \item The stack of lifts of $(X,D,f)$ over $X$, denoted by $\mathrm{Rel}((X,D,f),W_2(k))$, is the stackification of the prestack $\mathrm{Rel}'((X,D,f),W_2(k))$.  Objects in the prestack over an open subset \( U \) are triples \( (Y,D_Y, f_Y) \), where $Y$ is a smooth \( W_2(k) \)-sheme, $D_Y\subset Y$ is a simple normal crossing divisor relative to $W_2(k)$, $f_Y\in\Gamma(X,\mathcal O_Y)$ satisfying $(f_Y=0)\subset D_Y$, and its reduction modulo \( p \) is \( (U,D|_U,f|_U) \). A morphism \( (Y',D_{Y'},f_{Y'}) \to (Y'',D_{Y''},f_{Y''}) \) is a $W_2(k)$-morphism $u:Y'\to Y''$ satisfying $u^*D_{Y''}=D_{Y'}$, $u^*f_{Y''}\in f_{Y'}(1+pf_{Y'}\mathcal O_{Y'})$, and reduction modulo \( p \) equal to the identity. Since the sheaf of automorphisms of any lift of \( U \) is the abelian sheaf \( \mathcal Hom_{\mathcal O_X}(\Omega^1_f,f\mathcal O_X)\), $\mathrm{Rel}((X,D,f),W_2(k))$ forms a gerbe under it.
\end{itemize}
These gerbes are related by the following
\begin{theorem}\label{functors between stacks}
There are functors of stacks
$$
\mathrm{Rel}((X,D,f),W_2(k))\stackrel{\Phi}{\to}\mathrm{sc}(\mathcal G^*_0)\stackrel{\wedge^a}{\to}\mathrm{sc}(\mathcal G^*_a)
$$
and 
$$
\mathrm{Rel}((X,D,f),W_2(k))\xrightarrow{\mathrm{Forget}}\mathrm{Rel}((X,D),W_2(k)).
$$
Here, 
\begin{itemize}
    \item the equivalence of gerbes $\Phi$ is given by Theorem \ref{two sided truncation},
    \item  $\wedge^a$ sends a section $s:\Omega^1_{f^{-1}}\to Z^1F_*\Omega_{f^{-1}}^*$ to
$$
\wedge^as:\Omega^a_{f^{-1}}\to Z^aF_*\Omega^*_{f^{-1}},~\omega_1\wedge\cdots\wedge\omega_a\mapsto s(\omega_1)\wedge\cdots\wedge s(\omega_a),
$$
where $\omega_1\in\Omega^1_{f^{-1}},\omega_2,\cdots,\omega_a\in\Omega^1_{X/k}(\log D)$, and
\item $\mathrm{Forget}$ is the forget functor.
\end{itemize}
\end{theorem} 

As an immediate consequence of the above theorem, we have
\begin{corollary}
The splitting of $\mathcal G^*_0$ implies the splitting of $\mathcal G^*_a$ and $\mathcal K^*_a$.
\end{corollary}

For the classes of the above gerbes, we have 

\begin{theorem}
$\mathrm{(i)}$ There are equalities 
 $$\mathrm{cl}~\mathcal G^*_0=\mathrm{cl}~\mathrm{Rel}((X,D,f),W_2(k))=-\delta^1_{\mathcal G^*_0},$$
where $\delta^1_{\mathcal G^*_0}$ fits into the distinguished triangle 
$$
f\mathcal O_X\to\mathcal G^*_0\to\Omega^1_{f^{-1}}[-1]\xrightarrow{\delta^1_{\mathcal G^*_0}}f\mathcal O_X[1].
$$

$\mathrm{(ii)}$ $\mathrm{cl}~\mathrm{Rel}((X,D),W_2(k))$ is the image of $\mathrm{cl}~\mathrm{Rel}((X,D,f),W_2(k))$ under the map
$$
H^2(X,Hom_{\mathcal O_X}(\Omega^1_f,f\mathcal O_X))\to H^2(X,Hom_{\mathcal O_X}(\Omega^1_{X/k}(\log D),\mathcal O_X)),
$$
which is induced by the inclusion $\mathcal Hom_{\mathcal O_X}(\Omega^1_f,f\mathcal O_X)\subset\mathcal Hom_{\mathcal O_X}(\Omega^1_{X/k}(\log D),\mathcal O_X)$.

$\mathrm{(iii)}$ Taking an open affine covering $\{U_i\}$ of $X$. Then
$$
\eta_a:=\mathrm{cl}~\mathrm{sc}(\mathcal G^*_a)\in H^2(X,\mathcal Hom_{\mathcal O_X}(\Omega^1_f,f\mathcal O_X))\cong\mathrm{Hom}_{K(\mathcal O_X)}(\Omega^{a+1}_{f^{-1}},\mathcal {\check C}^*(\{U_i\},\Omega^a_{f^{-1}})[2]),
$$
where $K(\mathcal O_X)$ is the homotopy category of $\mathcal O_X$-modules. Moreover,  
$$
\eta_a(\omega_1\wedge\cdots\wedge\omega_{a+1})=\sum_{i=0}^{a+1}(-1)^{i-1}\eta_0(\omega_i)\omega_1\wedge\cdots\wedge\textcolor{red}{\widehat{\omega_i}}\wedge\cdots\wedge\omega_{a+1},
$$
where $\omega_1\in\Omega^1_{f^{-1}},\omega_2,\cdots,\omega_{a+1}\in\Omega^1_{X/k}(\log D)$.
\end{theorem}
\fi
\subsection{Kunneth formula}
Given $(\tilde X_i,\tilde D_i)/W_2(k)$ for $i=1,2$, where $\tilde X_i$ is a smooth $W_2(k)$-scheme, $\tilde D_i\subset\tilde X_i$ a simple normal crossing divisor relative to $W_2(k)$. Let $(X_i,D_i)/k$ be the mod $p$ reduction of $(\tilde X_i,\tilde D_i)/W_2(k)$ and suppose that $\dim X_1+\dim X_2<p$. Set 
$$\tilde X_0:=\tilde X_1\times_{W_2(k)}\tilde X_2,~\tilde D_0:=\tilde D_1\times_{W_2(k)}\tilde X_2+\tilde X_1\times_{W_2(k)}\tilde D_2,$$
and $(X_0,D_0)$ to be the mod $p$ reduction of $(\tilde X_0,\tilde D_0)$. Clearly, there are natural isomorphisms
\begin{eqnarray}\label{kunneth DR}
\Omega^*_{X_0/k}(\log D_0)\cong\Omega^*_{X_1/k}(\log D_1)\boxtimes\Omega^*_{X_2/k}(\log D_2)
\end{eqnarray}
and
\begin{eqnarray}\label{kunneth Higgs}
\bigoplus_{i=0}^{\infty}\Omega^i_{X_0/k}(\log D_0)[-i]\cong\bigoplus_{i=0}^{\infty}\Omega^i_{X_1/k}(\log D_1)[-i]\boxtimes\bigoplus_{i=0}^{\infty}\Omega^i_{X_2/k}(\log D_2)[-i].
\end{eqnarray}
On the other hand, \cite[Th\'eorème 2.1]{DI} shows that for $i=0,1,2$, there are decompositions
$$
\Phi_{(\tilde X_i,\tilde D_i)}:F_*\Omega^*_{X_i/k}(\log D_i)\cong\bigoplus_{i=0}^{\infty}\Omega^i_{X_i/k}(\log D_i)[-i]~\mathrm{in}~D(X_i).
$$

Inspired by the Kunneth formula in any good cohomology theory, we establish the Kunneth formula for the Deligne-Illlusie decomposition. 
\begin{theorem}
Notation and assumptions as above. Under the isomorphisms \eqref{kunneth DR} and \eqref{kunneth Higgs}, we have
$$
\Phi_{(\tilde X_0,\tilde D_0)}=\Phi_{(\tilde X_1,\tilde D_1)}\boxtimes\Phi_{(\tilde X_2,\tilde D_2)}.
$$
\end{theorem}
\begin{proof}
Let $\mathcal U_i:=\{(U_{i,j}\}_{j\in J_i}$ be an open affine covering of $X_i$ and let $\{\tilde F_{i,j}:(\tilde U_{i,j},\tilde D_i|_{\tilde U_{i,j}})\to(\tilde U_{i,j},\tilde D_i|_{\tilde U_{i,j}})\}_{j\in J_i}$ be a family of local logarithmic Frobenius liftings on $(\tilde X_i,\tilde D_i)$ for $i=1,2$. Set 
$$\mathcal U_0:=\{U_{1,j_1}\times U_{2,j_2}\}_{(j_1,j_2)\in J_1\times J_2},~\{\tilde F_{j_1,j_2}:=\tilde F_{1,j_1}\times_{W_2(k)}\tilde F_{2,j_2}\}_{(j_1,j_2)\in J_1\times J_2}
$$ to be the induced open affine covering of $X$ and  the induced family of local logarithmic Frobenius liftings on $(\tilde X_0,\tilde D_0)$, respectively.
For $i=1,2,3$, we set
$$
\mathcal H^*_i:=\bigoplus_{j=0}^\infty\Omega^j_{X_i/k}(\log D_i)[-j],~\mathcal K_i^*:=\check{\mathcal C}(\mathcal U_i,F_*\Omega^*_{X_i/k}(\log D_i)),~\mathcal R^*_i:=F_*\Omega^*_{X_i/k}(\log D_i).
$$
Let $\pi_i:X_0\to X_i$ be the natural projection for $i=1,2$. 

Consider the commutative diagram in $D(X_0)$:
$$
\xymatrix{
\mathcal H^*_0 \ar[r]\ar[rrrd]_-{\varphi}
  & \pi_1^*\mathcal H^*_1 \otimes \pi_2^*\mathcal H^*_2 
    \ar[rr]^-{\pi_1^*\varphi_1 \otimes \pi_2^*\varphi_2} 
  && \pi_1^*\mathcal K^*_1 \otimes \pi_2^*\mathcal K^*_2 
    \ar[d]^-{-\cup-} 
  && \pi_1^*\mathcal R_1^* \otimes \pi_2^*\mathcal R_2^* 
    \ar[ll]_-{\pi_1^*\rho_1 \otimes \pi_2^*\rho_2} 
  & \mathcal R^*_0 
    \ar[l]
    \ar@/_3pc/[llllll]_-{\Phi_{(\tilde X_1,\tilde D_1)} \boxtimes \Phi_{(\tilde X_2,\tilde D_2)}} 
    \ar[llld]^-{\rho_0} \\
&&& \mathcal K^*_0&&& 
},
$$
where $\varphi_i$ denotes the quasi-isomorphism $\varphi_{(\tilde X_i,\tilde D_i)}$ attached to the family $\{\tilde F_j\}_{j \in J_i}$, $\rho_i$ denotes the natural augmentation, and the cup product $-\cup-$ is defined as 
$$
\begin{array}{c}
\pi^*_1\check{\mathcal C}^{r_1}(\mathcal U_1,F_*\Omega^{s_1}_{X_1/k}(\log D_1))\otimes\pi_2^*\check{\mathcal C}^{r_2}(\mathcal U_2,F_*\Omega^{s_2}_{X_2/k}(\log D_2))\to\check{\mathcal C}^{r_1+r_2}(\mathcal U_0,F_*\Omega^{s_1+s_2}_{X_0/k}(\log D_0)),\\
\pi_1^*\{\beta_{j^0_1,\cdots,j_1^{r_1}}\in F_*\Omega^{s_1}_{X_1/k}(\log D_1)\}_{j^0_1,\cdots,j_1^{r_1}\in J_1}\bigotimes\pi_2^*\{\beta_{j_2^0,\cdots,j_2^{r_2}}\in F_*\Omega^{s_2}_{X_2/k}(\log D_2)\}_{j_2^0,\cdots,j_2^{r_2}\in J_2}\\
\tikz[baseline={(0,-0.5ex)}]{
  \draw[|->](0,0) -- (0,-0.5);
}\\
\{(-1)^{r_1s_2}\pi_1^*\beta_{j^0_1,\cdots,j_1^{r_1}}\bigwedge\pi_2^*\beta_{j_2^{r_1},\cdots,j_2^{r_1+r_2}}\in F_*\Omega^{s_1+s_2}_{X_0/k}(\log D_0)\}_{(j^0_1,j_2^0),\cdots,(j_1^{r_1+r_2},j_2^{r_1+r_2})\in J_1\times J_2}.
\end{array}
$$
It follows that $\Phi_{(\tilde X_1,\tilde D_1)}\boxtimes\Phi_{(\tilde X_2,\tilde D_2)}$ can be alternatively represented by the diagram
$$
\mathcal H^*_0\stackrel{\varphi}{\rightarrow}\mathcal K^*_0\stackrel{\rho_0}{\leftarrow}\mathcal R^*_0.
$$

We claim that $\varphi$ is nothing but $\varphi_{\underline\Omega}$ attached to the family $\{\tilde F_{(j_1,j_2)}\}_{(j_1,j_2)\in J_1\times J_2}$, where $\underline{\Omega}:=(\Omega_1,\Omega_2)$ with
$$
\Omega_1:=\pi^*_2\Omega^1_{X_2/k}(\log D_2)\subset\Omega^1_{X_0/k}(\log D_0),~\Omega_2:=\pi^*_1\Omega^1_{X_1/k}(\log D_1)\subset\Omega^1_{X_0/k}(\log D_0).
$$
Combining this fact with Proposition \ref{Ho^2}, we conclude that $\varphi$ is homotopic to $\varphi_{(\tilde X_0,\tilde D_0)}$, the quasi-isomorphism attached to the family $\{\tilde F_{(j_1,j_2)}\}_{(j_1,j_2)\in J_1\times J_2}$ again. Consequently, the desired equality follows. 

It remains to verify the claim. Let $r,s>0$, $(j_1^0,j_2^0),\cdots,(j^r_1,j^r_2)\in J_1\times J_2$, and let
$$\omega_1^1,\cdots,\omega_1^{m_1}\in\Omega^1_{X_1/k}(\log D_1),~\omega_2^1,\cdots,\omega_2^{m_2}\in\Omega^1_{X_2/k}(\log D_2),~m_1+m_2=r+s.
$$
Using the definition of the cup product $-\cup-$ described above, together with the following easily checked facts:
$$
\begin{array}{c}
\pi_1^*h_{\tilde F_{j_1^i},\tilde F_{j_1^{i+1}}}(\omega_1^l)=h_{\tilde F_{(j_1^i,j_2^i)},\tilde F_{(j_1^{i+1},j_2^{i+1})}}(\pi^*_1\omega_1^l),~\pi_2^*h_{\tilde F_{j_2^i},\tilde F_{j_2^{i+1}}}(\omega_2^l)=h_{\tilde F_{(j_1^i,j_2^i)},\tilde F_{(j_1^{i+1},j_2^{i+1})}}(\pi^*_2\omega_2^l),\\
\pi_2^*\zeta_{\tilde F_{j_2^i}}(\omega_2^l)=\zeta_{\tilde F_{(j^0_1,j_2^0)}}(\pi_2^*\omega_2^l)+h_{\tilde F_{(j_1^0,j_2^0)},\tilde F_{(j_1^1,j_2^1)}}(\pi^*_2\omega_2^l)+\cdots+h_{\tilde F_{(j_1^{i-1},j_2^{i-1})},\tilde F_{(j_1^i,j_2^i)}}(\pi^*_2\omega_2^l),
\end{array}
$$
we have
$$
\begin{array}{c}
\varphi(r,s)_{(j_1^0,j_2^0),\cdots,(j_1^r,j_2^r)}(\pi_1^*\omega_1^1\wedge\cdots\wedge\pi_1^*\omega_1^{m_1}
\wedge\pi_2^*\omega_2^1\wedge\cdots\wedge\pi_2^*\omega_2^{m_2})\\
=\bigoplus\pi_1^*\varphi_1(r_1,s_1)_{j^0_1,\cdots,j^{r_1}_1}(\omega_1^1\wedge\cdots\wedge\omega_1^{m_1})\bigcup\pi_2^*\varphi_2(r_2,s_2)_{j^{r_1}_2,\cdots,j^{r_1+r_2}_2}(\omega_1^1\wedge\cdots\wedge\omega_1^{m_2})\\
=\varphi_{\underline\Omega}(r,s)_{\tilde F_{(j^0_1,j^0_2)},\cdots,\tilde F_{(j^r_1,j^r_2)}}(\pi_1^*\omega_1^1\wedge\cdots\wedge\pi_1^*\omega_1^{m_1}
\wedge\pi_2^*\omega_2^1\wedge\cdots\wedge\pi_2^*\omega_2^{m_2}).
\end{array}
$$
Here, the direct sum in the second equality is taken over $r_1,s_1,r_2,s_2$, subject to the following conditions:
$$
r_1+r_2=r,~s_1+s_2=s,~r_1+s_1=m_1,~r_2+s_2=m_2.
$$
This completes the proof.
\end{proof}
\iffalse
More general, let $(E_i,\theta_i)$ be a nilpotent Higgs bundle on $(X_i,D_i)/k$ of level $\leq p-\dim X_i-1$ for $i=1,2$. Set $(H_i,\nabla_i):=C^{-1}_{(\tilde X_i,\tilde D_i)}(E_i,\theta_i)$ for $i=1,2$ and
$$(E_0,\theta_0):=(E_1,\theta_1)\boxtimes(E_2,\theta_2),~(H_0,\nabla_0)=(H_1,\nabla_1)\boxtimes(H_2,\nabla_2).$$
Similar to \eqref{kunneth DR} and \eqref{kunneth Higgs}, for $\star=\Omega^*,\Omega^*_\mathrm{int}$, we have natural isomorphisms
\begin{eqnarray}\label{kunneth theta}
\star(H_0,\nabla_0)\cong\star(H_1,\nabla_1)\boxtimes\star(H_2,\nabla_2),~\star(E_0,\theta_0)\cong\star(E_1,\theta_1)\boxtimes\star(E_2,\theta_2).
\end{eqnarray}
By combing \cite[Corollary 2.27]{OV}, \cite[Corollary 5.7]{Schepler} and \cite[Corollary 3.6]{SZ2}, for $i=0,1,2$ and $\star$ as above, there are isomorphisms
$$
\Phi_{(\tilde X_i,\tilde D_i),\star}^{(E_i,\theta_i)}:F_*\star(H_i,\nabla_i)\cong\star(E_i,\theta_i)~\mathrm{in}~D(X_i).
$$
\begin{theorem}
Notation and assumption as above. Under the isomorphisms \eqref{kunneth theta}, we have
$$
\Phi_{(\tilde X_0,\tilde D_0),\star}^{(E_0,\theta_0)}=\Phi_{(\tilde X_1,\tilde D_1),\star}^{(E_1,\theta_1)}\boxtimes\Phi_{(\tilde X_2,\tilde D_2),\star}^{(E_2,\theta_2)}.
$$
\end{theorem}
Next, we assume that any $X_i$ is proper over $k$. Let $(H_i,\nabla_i,\mathrm{Fil}_i,\psi_i)$ be a Fontaine module on $(\tilde X_i,\tilde D_i)/W_2(k)$ with $\mathrm{rank}(H_i)\leq p-\dim X_i$ for $i=1,2$, and let 
$$
(H_0,\nabla_0,\mathrm{Fil}_0,\psi_0):=(H_1,\nabla_1,\mathrm{Fil}_1,\psi_1)\boxtimes(H_2,\nabla_2,\mathrm{Fil}_2,\psi_2).
$$
As an immediate consequence of the above theorem, we have the Kunneth formula for Fontaine-Laffallie modules coming from geometry.
\begin{corollary}
Notation and assumption as above. Then for $\star=\Omega^*,\Omega_\mathrm{int}$, we have an isomorphism of Fontaine-Laffallie modules
$$
(V_0^m,\mathrm{Fil}_0,\psi_{\tilde X_0})\cong\bigoplus_{i+j=m}(V^m_1,\mathrm{Fil}_1,\psi_{\tilde X_1})\bigotimes(V^m_2,\mathrm{Fil}_2,\psi_{\tilde X_2}),
$$
where $V^m_i:=\mathbb H^m(X_i,\star(H_i,\nabla_i))$.
\end{corollary}
\fi




\subsection{Inverse Cartier transform of $p$-connections}
 Let $Y$ be a smooth scheme over $W(k)$ of relative dimension $n$, 
 and let $\mathcal D_{Y/W(k)}$ be the sheaf of rings of relative differential operators on $Y/W(k)$. In \cite{Fa}, Faltings introduced the abelian category $\mathrm{MF}^\nabla_{[0,p-2]}(Y/W(k))$ of Fontaine-Faltings modules on $Y/W(k)$ of Hodge-Tate weight $\leq p-2$. A key ingredient in his construction is a procedure that relates $p$-connections to integrable connections. In what follows, we present a slight generalization of this procedure, which can be viewed as an inverse Cartier transform of 
$p$-connections.

\begin{definition}
Let $p\mathcal D_{Y/W(k)}\subset\mathcal D_{Y/W(k)}$ be the subring locally defined by
$$
p\mathcal D_{Y/W(k)}:=\bigoplus_{I\in\mathbb N^n}\mathcal O_Y(p\partial)^I,\quad (p\partial)^I:=\prod_{j=1}^n(p\partial_j)^{i_j},\quad I=(i_1,\cdots,i_n).
$$
Here, $\partial_1,\cdots,\partial_n$ are dual basis of $dt_1,\cdots,dt_n$ for a system of local coordinates $t_1,\cdots,t_n$ on $Y/W(k)$.

Let \( 0 \leq \ell \leq p - 2\). Define the subring \( p\mathcal{D}_{Y/W(k)}^\ell \subset \mathcal{D}_{Y/W(k)} \) locally by
\[
p\mathcal{D}_{Y/W(k)}^\ell := 
\left( \bigoplus_{|I| \leq \ell} \mathcal{O}_Y \cdot (p\partial)^I \right)
\oplus 
\left( \bigoplus_{|I| > \ell} \frac{1}{p^{|I| - \ell}} \mathcal{O}_Y \cdot (p\partial)^I \right).
\] 
\end{definition}
Clearly, giving a \( W(k) \)-linear \( p \)-connection on an \( \mathcal{O}_Y \)-module \( E \) is equivalent to endowing \( E \) with a  \( p\mathcal{D}_{Y/W(k)} \)-module structure.
\begin{definition}
Let \( E \) be a \( p \)-torsion \( \mathcal{O}_Y \)-module; that is, \( p^N E = 0 \) for some \( N > 0 \). We say that a \( W(k) \)-linear \( p \)-connection on \( E \) is \( p \)-nilpotent of level \( \leq \ell \) if the corresponding \( p\mathcal{D}_{Y/W(k)} \)-action on \( E \) extends to a \( p\mathcal{D}^\ell_{Y/W(k)} \)-action. Denote by \( p\mathrm{MIC}_{\leq \ell}(Y/W(k)) \) the category of \( p \)-torsion \( \mathcal{O}_Y \)-modules equipped with a \( p \)-connection that is \( p \)-nilpotent of level \( \leq \ell \).
\end{definition}
Let $\{U_\alpha\}$ be an open affine covering of $Y$, and let $F_\alpha:\hat U_\alpha\to \hat U_\alpha$ be a Frobenius lifting on the $p$-adic completion of $U_\alpha$. Let $(E,\theta)\in p\mathrm{MIC}_{\leq p-1}(Y/W(k))$. We glue the family \( \{ F_\alpha^* E|_{U_\alpha} \} \) using Taylor's formula. More precisely, for any $e\in E|_{U_i}$, we identify $F^*_\alpha e$ with a section of $F^*_\beta E|_{U_\beta}$ described as follows:
\begin{eqnarray}\label{Taylor's formula}
\sum_{I\in\mathbb N^n}h_{\beta\alpha}(t)^IF_\beta^*(\frac{\theta^I}{I!}(m)),
\end{eqnarray}
where
$$
h_{\beta\alpha}(t)^I:=\prod_{j=1}^nh_{\beta\alpha}(t_j)^{i_j},\quad h_{\beta\alpha}:=\frac{F_\alpha^*-F_\beta^*}{p},\quad \theta^I:=\prod_{j=1}^n\theta_j^{i_j},\quad I!:=\prod_{j=1}^ni_j!,\quad I=(i_1,\cdots,i_n).
$$
Here, we assume that $U_\alpha \cap U_\beta$ admits a system of coordinates $t_1, \dots, t_n$ over $W(k)$, and write the $p$-connection as $\theta = \sum_{j=1}^n \theta_j \otimes dt_j$. Let $H$ denote the $\mathcal{O}_Y$-module obtained by gluing the aforementioned family. We construct an integrable connection $\nabla$ on $H$ as follows: for any $e\in E|_{U_\alpha}$, set 
$$
\nabla(F^*_\alpha e):=(\mathrm{id}\otimes\frac{dF^*_\alpha}{p})(F^*_\alpha(\theta e)),
$$
where
$$
\frac{dF^*_\alpha}{p}:F^*_\alpha\Omega^1_{Y/W(k)}\to\Omega^1_{Y/W(k)},\quad dx\mapsto \frac{dF^*_\alpha(x)}{p},\quad x\in\mathcal O_Y.
$$

\begin{definition}
The pair $(H, \nabla)$ is called the inverse Cartier transform of $(E, \theta)$, and is denoted by $(H, \nabla) = C^{-1}(E, \theta)$. In other words, there is an inverse Cartier transform between categories 
$$
C^{-1}:p\mathrm{MIC}_{\leq p-1}(Y/W(k))\to\mathrm{MIC}(Y/W(k)).
$$
\end{definition}
This inverse Cartier transform is closely related to the work of Shiho \cite{Shi}, Xu \cite{X}, Ogus \cite{O24}, and Wang \cite{W}.

\begin{lemma}
Let $(H,F^\bullet H,\nabla,\varphi)\in\mathrm{MF}^\nabla_{[0,p-2]}(Y/W(k))$ and let $(E,\theta)$ be the corresponding torsion  $p$-connection. Then $(E,\theta)\in p\mathrm{MIC}_{\leq p-1}(Y/W(k))$ and $(H,\nabla)=C^{-1}(E,\theta)$.
\end{lemma}  
Motivated by \cite[Corollary 2.27]{OV}, one naturally expects a comparison between $\Omega^*(H,\nabla)$ and  $\Omega^*(E,\theta)$. In this paper, we establish such a comparison in the case where 
$$Y = Y_1 \times_{W(k)} \cdots \times_{W(k)} Y_n,$$ 
with each $Y_i$ a smooth $W(k)$-scheme of relative dimension one.

\begin{theorem}
Notation and assumption as above. Let $(E,\theta)\in p\mathrm{MIC}_{\leq p-2}(Y/W(k))$ and let $(H,\nabla)=C^{-1}(E,\theta)$. Then 
\begin{eqnarray}\label{de Rham-p comparison}
\Omega^*(H,\nabla)\cong\Omega^*(E,\theta)~\mathrm{in}~D(\mathrm{Ab}(Y)),
\end{eqnarray}
where $\mathrm{Ab}(Y)$ is the category of sheaves of abelian groups on $Y$.
\end{theorem}
\begin{proof}
Let $\mathcal U_i:=\{U_i^q\}_{q\in Q_i}$ be an open affine affine covering of $Y_i$, and let $F_i^q$ be a Frobenius lifting on the $p$-adic completion of $U_i^q$. These coverings give rise to an open affine covering of $Y$, namely
$$
\mathcal U:=\{U_{(q_1,\cdots,q_n)}\}_{(q_1,\cdots,q_n)\in Q_1\times\cdots\times Q_n},\quad U_{(q_1,\cdots,q_n)}:=U_1^{q_1}\times_{W(k)}\cdots\times_{W(k)}U_n^{q_n}.
$$
Let $F_{(q_1,\cdots,q_n)}:=F_1^{q_1}\times_{W(k)}\cdots\times_{W(k)}F_n^{q_n}$, which is a Frobenius lifting on $p$-adic completion of $U^{(q_1,\cdots,q_n)}$.
As usual, we construct the desired isomorphism \eqref{de Rham-p comparison} as follows:
$$
\xymatrix{\Omega^*(H,\nabla)\ar[r]^-{\rho}&\check{\mathcal C}(\mathcal U,\Omega^*(H,\nabla))&\Omega^*(E,\theta)\ar[l]_-{\varphi_{\underline{\Omega}}}},
$$
where $\rho$ is the natural augmentation and 
$$\underline\Omega:=(\pi_1^*\Omega^1_{Y_1/W(k)},\cdots,\pi_n^*\Omega^1_{Y_n/W(k)}),\quad \pi_i:Y\to Y_i.$$
Clearly, the construction of the quasi-isomorphism $ \varphi_{\underline{\Omega}}$ reduces to define
\begin{eqnarray}\label{varphi p-connection}
   \varphi_{\underline{\Omega}}(r,s)_{ F_{\underline q_1},\cdots, F_{\underline q_r}}:\Omega^{r+s}(E,\theta)|_{U^{\underline q_0}\cap\cdots\cap U^{\underline q_r}}\to\Omega^s(H,\nabla)|_{U^{\underline q_0}\cap\cdots\cap U^{\underline q_r}}
\end{eqnarray}
for any $r,s\geq0$ and any $\underline q_0,\cdots,\underline q_r\in Q_1\times\cdots\times Q_n$. 

This morphism was essentially given in Construction \ref{construction higher homotopy}. For the reader's convenience, we recall it below with a slight modification. Without loss of generality, we may assume  that each \( Y_i/W(k) \) admits a coordinate \( t_i \).
 Write $\theta=\sum_{i=1}^n\theta_i\otimes dt_i$. Given any section
$$
\boldsymbol{e}=\sum_{1\leq i_1<\cdots<i_{r+s}\leq n}e_{i_1,\cdots,i_{r+s}}\otimes dt_{i_1}\wedge\cdots\wedge dt_{i_{r+s}},\quad e_{i_1,\cdots,i_{r+s}}\in E|_{U^{\underline q_0}\cap\cdots\cap U^{\underline q_r}}.
$$
Then \eqref{varphi p-connection} is defined by
\begin{eqnarray}\label{varphi p r s}
\boldsymbol{e}\mapsto \sum_{(\underline i,\underline S,\underline{\overline j})\in T_{\underline\Omega}(r,s)}\iota_{F^{\underline q_0}}(F^{\underline q_0*}(\theta^{\underline{\overline j}}\theta_{i_1}^{-1}\cdots\theta_{i_r}^{-1}e_{\underline S,\underline i}))(\boldsymbol{\omega}_{\underline S}h^{[\underline{\overline j}]}),
\end{eqnarray}
where
$$
\zeta_{F_{\underline q_0}}(dt_j):=\frac{dF^*_{\underline q_0}(t_j)}{p},\quad h_{i,j}:=\frac{F^*_{\underline q_i}(t_j)-F^*_{\underline q_{i-1}}(t_j)}{p}.
$$

We need to show that \eqref{varphi p r s} is independent of the choice of coordinates on \( Y_i/W(k) \). Without loss of generality, we may assume that \( n = 1 \) and \( (E, \theta) = (\mathcal{O}_X, 0) \). It then suffices to check that \( \varphi_{\underline\Omega}(1, 0) \) is independent of the choice of coordinates on \( Y/W(k) \).
Since the verification can be carried out formally, we may reduce the independence to a statement over \( \mathbb{C} \): for any holomorphic function \( f \in \mathcal{O}_{\mathbb{C}^{\mathrm{an}}} \), any holomorphic coordinate functions \( t \) and \( w \) on \( \mathbb{C}^{\mathrm{an}} \), and any points \( P_1, P_2 \in \mathbb{C}^{\mathrm{an}} \), we have the identity
\[
\sum_{j=1}^\infty (\partial_w^{j-1} f)(P_1) \cdot (w(P_2) - w(P_1))^{[j]} 
= \sum_{j=1}^\infty \left( \partial_t^j \left( f \frac{\partial w}{\partial t} \right) \right)(P_1) \cdot (t(P_2) - t(P_1))^{[j]}.
\]
By the basic theory of one-variable complex analysis, there exists a holomorphic function \( F \) locally such that \( dF = f\,dw \). Combining this with Taylor's formula, the desired identity follows. 

It remains to show that $\varphi_{\underline{\Omega}}$ is a quasi-isomorphism. To this end, we need to check that
$$
\varphi_{\underline\Omega}(0,*)_{F_{\underline q}}:\Omega^*(E,\theta)|_{U_{\underline q}}\to\Omega^*(H,\nabla)|_{U_{\underline q}}
$$
is a quasi-isomorphism. Without loss of generality, we may assume that each 
$Y_i/W(k)$ admits a coordinate $t_i$ such that
$F^*_{\underline q}(t_i)=t_i^p$ and
$$
\Omega^1_{U_{\underline q}/W(k)}\cong\bigoplus_{i=1}^n\mathcal O_{U_{\underline q}}d\log t_i.
$$ 
Under these assumptions, we have
$$
F_{\underline q*}\Omega^*(H,\nabla)|_{\hat U_{\underline q}}\cong\bigoplus_{\beta_1,\cdots,\beta_n\in\{0,\cdots,p-1\}}\mathrm{Kos}(E|_{U_{\underline q}};\beta_1+\theta_1,\cdots,\beta_n+\theta_n),
$$
where $\theta=\sum_{i=1}^n\theta_i\otimes d\log t_i$.
Clearly, the Koszul complex $\mathrm{Kos}(E|_{U_{\underline q}};\beta_1+\theta_1,\cdots,\beta_n+\theta_n)$ is acyclic for $(\beta_1,\cdots,\beta_n)\neq0$, and $\varphi_{\underline\Omega}(0,*)_{F_{\underline q}}$ induces an isomorphism from $\Omega^*(E,\theta)|_{ U_{\underline q}}$ onto its image $\mathrm{Kos}(E|_{U_{\underline q}};\theta_1,\cdots,\theta_n)$. Consequently, we conclude that $\varphi_{\underline\Omega}(0,*)_{F_{\underline q}}$ is a quasi-isomorphism. This completes the proof.
\end{proof}



We proceed to extend the above theorem to the logarithmic setting. Let $D = \sum_{i=1}^n D_i$ be a divisor on $Y$, where each $D_i$ is a relative simple normal crossing divisor on $Y_i$ over $W(k)$. One can similarly define the category $p\mathrm{MIC}_{\leq p-2}((Y,D)/W(k))$ and the inverse Cartier transform between categories
$$
C^{-1}:p\mathrm{MIC}_{\leq p-2}((Y,D)/W(k))\to\mathrm{MIC}((Y,D)/W(k)).
$$

\begin{theorem}
Let $(E,\theta)\in p\mathrm{MIC}_{\leq p-2}((Y,D)/W(k))$ and let $(H,\nabla)=C^{-1}(E,\theta)$. Then 
\begin{eqnarray}\label{log de Rham-p}
\Omega^*(H,\nabla)\cong\Omega^*(E,\theta)~\mathrm{in}~D(\mathrm{Ab}(Y)),
\end{eqnarray}
which restricts to an isomorphism of intersection subcomplexes
\begin{eqnarray}\label{int de Rham-p}
\Omega_\mathrm{int}^*(H,\nabla)\cong\Omega_\mathrm{int}^*(E,\theta)~\mathrm{in}~D(\mathrm{Ab}(Y)).
\end{eqnarray}
\end{theorem}
\begin{proof}
Its proof is similar to that of the former theorem.
\end{proof}

\begin{corollary}
  Let $(H,F^\bullet H,\nabla,\varphi)\in\mathrm{MF}^\nabla_{[0,p-2]}((Y,D)/W(k))$. Then the spectral sequence
  $$
  E^{i,j}_1=\mathbb H^{i+j}(Y,\mathrm{Gr}_F^i\Omega^*_\mathrm{int}(H,\nabla))\Rightarrow\mathbb H^{i+j}(Y,\Omega^*_\mathrm{int}(H,\nabla))
  $$
  degenerates at the $E_1$-page.
\end{corollary}
\begin{proof}
Let $(E,\theta)$ be the $p$-connection corresponding to $(H,F^\bullet H,\nabla,\varphi)$.
Following a similar argument as in \cite[Theorem 4.12]{SZ2}, we obtain the following \emph{intersection adaptedness}: 
$$
\Omega_\mathrm{int}^*(E,\theta)\cong\mathrm{colim}(\cdots\rightarrow F^i\Omega_\mathrm{int}^*(H,\nabla)\xleftarrow{p}F^i\Omega_\mathrm{int}^*(H,\nabla)\rightarrow F^{i-1}\Omega_\mathrm{int}^*(H,\nabla)\xleftarrow{p}\cdots).
$$
Combing this fact with \eqref{int de Rham-p}, we get a $W(k)$-semilinear isomorphism from 
$$
\mathrm{colim}(\cdots\rightarrow R\pi_*F^i\Omega_\mathrm{int}^*(H,\nabla)\xleftarrow{p}R\pi_*F^i\Omega_\mathrm{int}^*(H,\nabla)\rightarrow R\pi_*F^{i-1}\Omega_\mathrm{int}^*(H,\nabla)\xleftarrow{p}\cdots).
$$
to $R\pi_*\Omega_\mathrm{int}^*(H,\nabla)$, where $\pi:Y\to\mathrm{Spec}(W(k))$ is the structural morphism. In other words, 
the triple
$$(R\pi_*\Omega_\mathrm{int}^*(H,\nabla),R\pi_*F^i\Omega_\mathrm{int}^*(H,\nabla),\varphi_{\underline{\Omega}})$$
is a Fontaine complex over $W(k)$ in the sense of \cite[Definition 5.3.6]{Ogus}. By \cite[Corollary 5.3.7]{Ogus}, for any $i$ and any $m$, the map
$$
H^m(R\pi_*F^i\Omega_\mathrm{int}^*(H,\nabla))\to H^m(R\pi_*\Omega_\mathrm{int}^*(H,\nabla))
$$
is injective. From which this theorem follows. 
\end{proof}




\subsection{$E_1$-degeneration and vanishing theorem}

\begin{theorem}
Keep the notation and assumption as in \S3.3 and additionally assume that $f$ is proper.
\begin{itemize}
\item[$\mathrm{(i)}$\ ] If the logarithmic cotangent bundle $\Omega^1_{X/k}(\log D)$ splits into a direct sum of subbundles of rank $<p$, then for $\star=W_i\Omega^*,\Omega^*_{f^{-1}}$ we have
$$
\dim_k\mathbb H^j(X,\star(\mathcal O_X,d+df\wedge))=\dim_k\mathbb H^j(X,\star(\mathcal O_X,-df\wedge))<\infty,~\forall j.
$$

\item[$\mathrm{(ii)}$\ ] Let $(H,\nabla,\mathrm{Fil},\psi)$ be a Fontaine module on $(\tilde X,\tilde D)/W_2(k)$ and set $(E,\theta):=\mathrm{Gr}_\mathrm{Fil}(H,\nabla)$. Suppose that $f$ has only isolated singularities, then
$$
\dim_k\mathbb H^j(X,\Omega^*(H,\nabla+df\wedge))=\dim_k\mathbb H^j(X,\Omega^*(E,\theta-df\wedge))<\infty,~\forall j.
$$
\end{itemize}
\end{theorem}
In the next two theorems, let $\tilde f:(\tilde X,\tilde D)\to(\mathbb P^1_{W_2(k)},\tilde E)$ be a projective semistable family over $W_2(k)$ along an effective divisor $\tilde E$ containing $\infty$ on $\mathbb P^1_{W_2(k)}$, and let $f:(X,D)\to(\mathbb P^1_k,E)$ be its mod $p$ reduction. Let $(H,\nabla,\mathrm{Fil},\psi)$ be a Fontaine module on $(\tilde X,\tilde D)/W_2(k)$ and set $(E,\theta):=\mathrm{Gr}_\mathrm{Fil}(H,\nabla)$.
\begin{theorem}
    If the logarithmic cotangent bundle $\Omega^1_{X/k}(\log D)$ splits into a direct sum of subbundles of rank $\leq p-\mathrm{rank}(H)$, then for $\star=W_i\Omega^*,\Omega^*_\mathrm{int},\Omega^*_f$ we have
$$
\dim_k\mathbb H^j(X,\star(H,\nabla))=\dim_k\mathbb H^j(X,\star(E,\theta)).
$$
In other words, the Hodge to de Rham spectral sequence 
$$
E_1^{r,s}:=\mathbb H^{r+s}(X,\mathrm{Gr}_\mathrm{Fil}^r\star(H,\nabla))\Rightarrow\mathbb H^{r+s}(X,\star(H,\nabla))
$$
degenerates at $E_1$.
\end{theorem} 
\begin{theorem}
If the logarithmic cotangent bundle $\Omega^1_{X/k}(\log D)$ splits into a direct sum of subbundles of rank $\leq p-\mathrm{rank}(H)$, then for $\star=W_i\Omega^*,\Omega^*_\mathrm{int},\Omega^*_f$  and any ample line bundle $L$ on $X$, we have
$$
\mathbb H^j(X,\star(E,\theta)\otimes L)=0,~j>\dim X.
$$
\end{theorem}

Finally, we turn to complex Hodge theory. Let $X$ be a smooth quasi-projective algebraic variety over $\mathbb C$ endowed with a simple normal crossing divisor $D$. Let $(H,\nabla,\mathrm{Fil})$ be a Gauss-Manin system associated to a smooth projective family $g:Y\to X$, namely there is some $m\geq0$ such that
$$
H=R^mg_*\Omega^*_{Y/X},~\nabla=\nabla^{GM},~\mathrm{Fil}^p=\mathrm{Im}(Rg_*\Omega_{Y/X}^{\geq p}\to Rg_*\Omega_{Y/X}^*).
$$
Set $(E,\theta):=\mathrm{Gr}_\mathrm{Fil}(H,\nabla)$.



By combing spreading-out technique and the above theorems, we can give the following vanishing theorems an alternative mod $p$ proof, respectively.


\begin{theorem}
 Assume that $D=\emptyset$. Let $f:X\to\mathbb A^1_\mathbb{C}$ be a projective morphism such that it has only isolated singularities. Then
 $$
 \dim_\mathbb{C}\mathbb H^i(X-D,\Omega^*(H,\nabla+df\wedge))=\dim_\mathbb{C}\mathbb H^i(X-D,\Omega^*(E,\theta-df\wedge)),~\forall i.
 $$
\end{theorem}

  

\begin{theorem}
Suppose that there is a semistable family $f:(X,D)\to(\mathbb P^1_{\mathbb C},E)$ along a divisor $E$ containing $\infty$ on $\mathbb P^1_{\mathbb C}$.
We regard $(E,\theta)$ as a Higgs bundle on $(X,D)/\mathbb C$. Then for $\star=W_i\Omega^*,\Omega^*_\mathrm{int},\Omega^*_f$ and any ample line bundle $L$ on $X$, we have
 $$
 \mathbb H^i(X,\star(E,\theta)\otimes L)=0,~i>\dim X.
 $$  


\end{theorem}




\iffalse
\section{De Rham-Higgs comparison for vanishing cycles \uppercase\expandafter{\romannumeral1}}
\subsection{Completion along one single point}
Let $(E,\theta)$ be a coherent nilpotent Higgs module on $(X',D')/k$ of level $\leq p-1$. Let $f\in\Gamma(X,\mathcal O_X)$ be any global function on $X$. Set $(H,\nabla):=C^{-1}(E,\theta)$. Given any $x\in X$ and set $Z$ to be the reduced closed subscheme of $X$ with support $\{x\}$. Set
$$(\mathfrak X,\mathfrak D):=(X,D)_{/F^*Z'},~(\mathfrak X',\mathfrak D'):=(X',D')_{/Z'}.$$ Put $(\hat E,\hat \theta):=(E,\theta)_{/Z'}$ and $(\hat H,\hat\nabla):=(H,\nabla)_{/F^*Z'}$.
\iffalse
We reduce the theorem to Proposition \ref{formal curve} below. Write
$$
Z_{df}+D=\sum_{i=1}^mn_ix_i,~n_i\in\mathbb Z_{>0},~\mathrm{Supp}(Z_{df}+D)=\{x_1,\cdots,x_m\}.
$$
Set $\mathcal I:=\mathcal O_{X'}(-Z_{df}-D)$ and $\mathcal I_i:=\mathcal O_X(-n_ix_i)$. Let $\mathfrak X:=X_{/(Z_{df}+D)}$ and $\mathfrak X_i:=X_{/(n_ix_i)}$ be formal completions of $X$ along $Z_{df}+D$ and $n_ix_i$, respectively.
Set 
$$
\begin{array}{c}
\hat E:=\varprojlim_jE/\mathcal I'^jE,~\hat H:=\varprojlim_jH/F^*\mathcal I'^jH,\\
\hat E_i:=\varprojlim_jE/\mathcal I_i'^jE,~\hat H_i:=\varprojlim_jH/F^*\mathcal I_i'^jH.
\end{array}
$$
Clearly, $\theta$ induces a Higgs field $\hat\theta$ (resp. $\hat\theta_i$) on $\hat E$ (resp. $\hat E_i$) and $\nabla$ induces  an integrable connection $\hat\nabla$ (resp. $\hat\nabla_i$) on $\hat H$ (resp. $\hat H_i$).  One can check that
$$
(\hat E,\hat\theta)\cong\bigoplus_{i=1}^m(\hat E_i,\hat\theta_i),~(\hat H,\hat\nabla)\cong\bigoplus_{i=1}^m(\hat H_i,\hat\nabla_i)
$$
and
$$
\begin{array}{c}
\Omega^*(\hat E,\hat\theta-df'\wedge)\cong\bigoplus_{i=1}^m\Omega^*(\hat E_i,\hat\theta_i-df'\wedge),\\
\Omega^*(\hat H,\hat\nabla+df\wedge)\cong\bigoplus_{i=1}^m\Omega^*(\hat H_i,\hat\nabla_i+df\wedge),\\
\Omega_{\mathrm{int}}^*(\hat E,\hat\theta-df'\wedge)\cong\bigoplus_{i=1}^m\Omega_\mathrm{int}^*(\hat E_i,\hat\theta_i-df'\wedge),\\
\Omega_\mathrm{int}^*(\hat H,\hat\nabla+df\wedge)\cong\bigoplus_{i=1}^m\Omega_\mathrm{int}^*(\hat H_i,\hat\nabla_i+df\wedge).
\end{array}
$$
Since $\Omega^*( E,\theta-df'\wedge)$ is acyclic outside $Z_{df}+D$, coupled with Theorem \ref{Cartier Ogus}, we know that the natural morphisms
$$
\begin{array}{c}
\Omega^*(E,\theta-df'\wedge)\to\Omega^*(\hat E,\hat\theta-df'\wedge),~\Omega^*(H,\nabla+df\wedge)\to\Omega^*(\hat H,\hat\nabla+df\wedge),\\
\Omega_\mathrm{int}^*(E,\theta-df'\wedge)\to\Omega_\mathrm{int}^*(\hat E,\hat\theta-df'\wedge),~\Omega_\mathrm{int}^*(H,\nabla+df\wedge)\to\Omega_\mathrm{int}^*(\hat H,\hat\nabla+df\wedge)
\end{array}
$$
are quasi-isomorphisms. Therefore Theorem \ref{twisted de Rham-Higgs comparison curve} follows immediately from \fi
\begin{theorem}\label{formal de Rham-Higgs isolated}
With the notations above, we have
\begin{eqnarray}\label{formal log}
F_*\Omega^*(\hat H,\hat\nabla+df\wedge)\cong\Omega^*(\hat E,\hat\theta-df'\wedge)~\mathrm{in}~D(\mathfrak X')
\end{eqnarray}
and
\begin{eqnarray}\label{formal int}
F_*\Omega_\mathrm{int}^*(\hat H,\hat\nabla+df\wedge)\cong\Omega_\mathrm{int}^*(\hat E,\hat\theta-df'\wedge)~\mathrm{in}~D(\mathfrak X').\end{eqnarray}
\end{theorem}
\begin{proof}
 Without loss of generality, we may assume that $\mathrm{dim}(X)=2$ and that there is a system of coordinates $(t_1,t_2)$ around $x$ such that $D=(t_1=0)$ and $t_1(x)=t_2(x)=0$. By Theorems \ref{Cartier Ogus} and \ref{Ogus comp int}, if $\partial f/\partial t_2(x)=0$, then all complexes in \eqref{formal log} and \eqref{formal int} are acyclic. It follows that we may assume $\partial f/\partial t_2(x)\neq0$.
Under these assumptions, we prove \eqref{formal log} (resp. \eqref{formal int}) by three steps:
$$
\begin{array}{c}
\Omega^*(\hat E,\hat\theta-df'\wedge)\stackrel{\mathrm{step}~1}{\cong}A^*\stackrel{\mathrm{step}~2}{\cong}B^*\stackrel{\mathrm{step}~3}{\cong}F_*\Omega^*(\hat H,\hat\nabla+df\wedge)\\
(\mathrm{resp.}~\Omega_\mathrm{int}^*(\hat E,\hat\theta-df'\wedge)\stackrel{\mathrm{step}~1}{\cong}A_\mathrm{int}^*\stackrel{\mathrm{step}~2}{\cong}B_\mathrm{int}^*\stackrel{\mathrm{step}~3}{\cong}F_*\Omega_\mathrm{int}^*(\hat H,\hat\nabla+df\wedge)),
\end{array}
$$
where complexes $A^*,B^*$ (resp. $A^*_\mathrm{int},B^*_\mathrm{int}$) will be explicitly described below. 

{\itshape Step 1.} Since $k$ is an algebraically closed field, we may assume $f(x)=0$. Let $\hat{\mathcal O}_{X,x}$ be the completion with respect to its maximal ideal. It is well-know that $\hat{\mathcal O}_{X,x}\cong k[[t_1,t_2]]$ and $\mathfrak X=\mathrm{Spf}(k[[t_1,t_2]])$. Using the natural inclusion $\mathcal O_{X,x}\subset k[[t_1,t_2]]$, we can consider the Taylor expansion 
$$f=\sum_{\underline\alpha\in\mathbb N^2}f_{\underline\alpha},~f_{\underline\alpha}=a_{\underline\alpha}t_1^{\alpha_1}t_2^{\alpha_2},~\underline\alpha=(\alpha_1,\alpha_2),~a_{\underline\alpha}\in k.$$
Note that $f_{(0,0)}=0$ by assumption. Using the above expansion, we construct complexes of $\mathcal O_{\mathfrak X'}$-modules
$$
A^*:=\Omega^*(\hat E,\hat\theta-\sum_{\underline\alpha\in\mathbb N^2}\sum_{i\in\mathbb N} f_{\underline\alpha}'^{p^i-1}df'_{\underline\alpha}\wedge),~A_\mathrm{int}^*:=\Omega^*_\mathrm{int}(\hat E,\hat\theta-\sum_{\underline\alpha\in\mathbb N^2}\sum_{i\in\mathbb N} f_{\underline\alpha}'^{p^i-1}df'_{\underline\alpha}\wedge).
$$
We claim that $A^*$ and $A^*_\mathrm{int}$ are isomorphic to $\Omega^*(\hat E,\hat\theta-df'\wedge)$ and $\Omega_\mathrm{int}^*(\hat E,\hat\theta-df'\wedge)$, respectively. It is easy to see that $A^*,\Omega^*(\hat E,\hat\theta-df'\wedge)$
can be reformulated as Koszul complexes of $\mathcal O_{\mathfrak X'}$-modules
$$
K^*(\hat E;\hat\theta_1-\sum_{i\in\mathbb N}(t_1'\frac{\partial f'}{\partial t'_1})^{p^i},\hat\theta_2-t_2^{-1}\sum_{i\in\mathbb N}(t_2'\frac{\partial f'}{\partial t'_2})^{p^i})
,~K^*(\hat E;\hat\theta_1-t_1'\frac{\partial f'}{\partial t'_1},\hat\theta_2-\frac{\partial f'}{\partial t'_2}).$$
Let us do some formal computations:

Clearly, the quotients are automorphisms of $\hat E$ which give rise to an isomorphism of the Koszul complexes above and hence an isomorphism 
\begin{eqnarray}\label{step 1}
\Omega^*(\hat E,\hat\theta-df'\wedge)\cong A^*,~
\hat e\otimes\omega'_I\mapsto(\prod_{i\in I}q_i)(\hat e)\otimes\omega'_I.
\end{eqnarray}
Here $I\subset\{1,2\},~\prod_{i\in I}q_i
:=\mathrm{id}_{\hat E}$ and 
$$ \omega_\emptyset:=1,~\omega_{\{1\}}:=d\log t_1,~\omega_{\{2\}}:=dt_2,~\omega_{\{1,2\}}:=d\log t_1\wedge dt_1.$$
By the definition of intersection complexes, one easily sees that  \eqref{step 1} induces an isomorphism
$$
A_\mathrm{int}^*\cong\Omega_\mathrm{int}^*(\hat E,\hat\theta-df'\wedge).
$$
This completes {\itshape Step 1.}
\iffalse
Let $\partial'\in T_{X'/k}$ be the dual of $dt'\in\Omega^1_{X'/k}$.  Using the identification 
$$
t'\partial':\Omega^1_{X'/k}(\mathrm{log}~D')\cong\mathcal O_{X'},
$$
it follows that (by abuse of notation, we rewrite $[\mathrm{id}\otimes(t'\partial')]\hat\theta$ by $\hat\theta$)
$$
\begin{array}{rcl}

\\cong&[\hat E\xrightarrow{\hat\theta-t'\partial'f'}\hat E],\\
\Omega^*(\hat E,\hat\theta-\sum^\infty_{i=0}f'^{p^i-1}df'\wedge)&\cong&[\hat E\xrightarrow{\hat\theta-t'\partial'f'\sum^\infty_{i=0}f'^{p^i-1}}\hat E].
\end{array}
$$
Let us do some formal computations:
$$
\begin{array}{rcl}
\frac{\hat\theta-t'\partial'f'-t'\partial'f'\sum_{i=1}^\infty f'^{p^i-1}}{\hat\theta-t'\partial'f'}&=&1+\frac{\sum_{i=1}^\infty f'^{p^i-1}}{1-\frac{\hat\theta}{t'\partial'f'}}\\
&=&1+(\sum_{i=0}^\infty f'^{p^i-1})^p\sum_{i=0}^{p-2} f'^{p-1-i}(\frac{f'}{t'\partial'f'})^i\hat\theta^i.
\end{array}
$$
Since $\mathrm{ord}_{x'}(t'\partial'f')=\mathrm{ord}_{x'}(f')$, then  $f'/(t'\partial'f')\in\mathcal O^*_{X',x'}\subset\mathcal O^*_{\mathfrak X'}$. Denote by $\hat h$ the expression above,  it is obvious that $\hat h\in\mathrm{Aut}_{\mathcal O_{\mathfrak X}}(\hat E)$. One can check that
$$
\xymatrix{
\hat E\ar[r]^-{P}\ar[d]^-{\mathrm{id}}&\hat E\ar[d]^-{\hat h}\\
\hat E\ar[r]^-{Q}&\hat E}
$$
is an isomorphism of complexes, where
$$
P=\hat\theta-t'\partial'f',~Q=\hat\theta-t'\partial'f'\sum^\infty_{i=0}f'^{p^i-1}.$$
Moreover, it induces an isomorphism of subcomplexes
$$
\xymatrix{
\hat E\ar[r]^-{P}\ar[d]^-{\mathrm{id}}&t\hat E+P\hat E\ar[d]^-{\hat h}\\
\hat E\ar[r]^-{Q}&t\hat E+Q\hat E.}
$$
Consequently, we obtain two isomorphisms of complexes
\begin{eqnarray}\label{step 1.1}
\Omega^*(\hat E,\hat\theta-df'\wedge)&\cong&\Omega^*(\hat E,\hat\theta-\sum^\infty_{i=0}f'^{p^i-1}df'\wedge)\end{eqnarray}
and
\begin{eqnarray}\label{step 1.2}
\Omega_\mathrm{int}^*(\hat E,\hat\theta-df'\wedge)&\cong&\Omega_\mathrm{int}^*(\hat E,\hat\theta-\sum^\infty_{i=0}f'^{p^i-1}df'\wedge).
\end{eqnarray}
\fi

{\itshape Step 2.} Set
$$
\hat\vartheta:=\hat\theta-\sum_{\underline\alpha\in\mathbb N^2}\sum_{i\in\mathbb N} f_{\underline\alpha}'^{p^i-1}df'_{\underline\alpha}\wedge
$$
and write
$$
\hat\vartheta=\hat\vartheta_1\otimes d\log t'_1+\hat\vartheta_2\otimes d t_2'.
$$
Since $\partial f/\partial t_2(x)=0$ by the assumption located at the very beginning of the proof, we deduce that $\hat\vartheta_1,\hat\vartheta_2$ are topologically nilpotent endomorphisms of $\hat E$. One constructs a module with integrable connection on $(\mathfrak X,\mathfrak D)/k$ by
\begin{eqnarray}\label{inverse Cartier transform}
\begin{array}{c}
C_\zeta^{-1}(\hat E,\hat\vartheta)=(F^*\hat E,\nabla_\mathrm{can}+(\mathrm{id}\otimes\zeta)F^*\hat\vartheta)=\\
(F^*\hat E,\nabla_\mathrm{can}+(\mathrm{id}\otimes\zeta)F^*\theta-\sum_{\underline\alpha\in\mathbb N^2}\sum_{i=1}^\infty f_{\underline\alpha}^{p^i-1}df_{\underline\alpha}\wedge),
\end{array}
\end{eqnarray}
where $\nabla_\mathrm{can}$ is the Cartier descent connection and 
$$\zeta:F^*\Omega^1_{\mathfrak X'/k}(\log\mathfrak D')\to\Omega^1_{\mathfrak X/k}(\log\mathfrak D)$$ is given by $F^*d\log t'_1\mapsto d\log t_1$ and $F^*dt_2'\mapsto t_2^{p-1}dt_2$. Set
$$
B^*:=F_*\Omega^*(F^*\hat E,\nabla_\mathrm{can}+(\mathrm{id}\otimes\zeta)F^*\hat\vartheta),~B^*_\mathrm{int}:=F_*\Omega_\mathrm{int}^*(F^*\hat E,\nabla_\mathrm{can}+(\mathrm{id}\otimes\zeta)F^*\hat\vartheta).$$
We claim that the morphism of complexes of $\mathcal O_{\mathfrak X'}$-modules 
\begin{eqnarray}\label{varphi}
\varphi:A^*\to B^*,~\hat e\otimes\omega'_I\mapsto F^*\hat e\otimes\zeta(\omega'_I)
\end{eqnarray}
is a quasi-isomorphism. Moreover, it induces a quasi-isomorphism of subcomplexes 
\begin{eqnarray}\label{varphi int}
\varphi_\mathrm{int}:A_\mathrm{int}^*\to B_\mathrm{int}^*.
\end{eqnarray}
These statements can be proved by using the similar arguments presented in the proof of \cite[Proposition 3.8]{SZ2}. Here we give another proof following the spirit of logarithmic algebraic  geometry, see e.g. \cite{Kato}, \cite{O18} for more details.

We firstly prove \eqref{varphi} is a quasi-isomorphism. Let us construct a family of Koszul complexes of $\mathcal O_{\mathfrak X'}$-modules
$$
K^*_{\underline v}:=K^*(\hat E;v_1\mathrm{id}+\hat\vartheta_1,v_2\mathrm{id}+t_2\hat\vartheta_2),~\underline v=(v_1,v_2)\in\{0,1,\cdots,p-1\}^2.
$$
Clearly, there are inclusions $K^*_{\underline v}\subset B^*$ of complexes given by 
$$
\hat E_I\subset B^{|I|},~\hat e\mapsto t_1^{v_1}t_2^{u_2}(F^*\hat e)\eta_I,~I\subset\{1,2\}.
$$
Here $u_2:=v_2$ for $v_2-\delta_2^I\geq0$, $p-1$ for $v_2=0,2\in I$, and
$$\eta_\emptyset:=1,~\eta_{\{1\}}:=d\log t_1,~\eta_{\{2\}}:=d\log t_2,~\eta_{\{1,2\}}=d\log t_1\wedge d\log t_2.
$$
Under the inclusions, one can check that we have a decomposition
\begin{eqnarray}\label{decomposition}
B^*=\bigoplus_{\underline v\in\{0,1,\cdots,p-1\}^2}K^*_{\underline v}.
\end{eqnarray}
Note that $\varphi$ is injective whose image is $K^*_{(0,0)}$. Moreover, $K^*_{\underline v}$ is acyclic for $\underline v\neq(0,0)$. In fact, assume $v_1\neq0$, then we can construct a homotopy on $K^*_{\underline v}$ between $v_1\mathrm{id}_{K^*_{\underline v}}$ and $0$:
\begin{eqnarray}\label{homotopy}
\begin{array}{c}
h_{\underline v}:\bigoplus_I\hat E_I\to\bigoplus_I\hat E_I,\\
(\hat e_\emptyset,\hat e_{\{1\}},\hat e_{\{2\}},\hat e_{\{1,2\}})\mapsto (\hat e_{\{1\}},0,\hat e_{\{1,2\}},0).
\end{array}
\end{eqnarray}
It follows that \eqref{varphi} is a quasi-isomorphism.

Proceeding to prove \eqref{varphi int} is a quasi-isomorphism. Set
$$
\begin{array}{c}
\Phi:=\{t_1^{v_1}t_2^{v_2}|v_1,v_2\in\{0,1,\cdots,p-1\}\}\cdot\{1,d\log t_1.dt_2,d\log t_1\wedge dt_2\}\\
\Phi_0:=\{t_2^{v_2}|v_2\in\{0,1,\cdots,p-1\}\}\cdot\{d\log t_1,d\log t_1\wedge dt_2\},~\Phi_1:=\Phi-\Phi_0.
\end{array}
$$
According to the proof of \cite[Lemma 3.14]{SZ2}, one has
$$
B^\bullet_\mathrm{int}=(\bigoplus_{\phi\in\Phi_0}(\hat\vartheta\hat E+t_1\hat E)\otimes\phi)\bigoplus(\bigoplus_{\phi\in\Phi_1}\hat E\otimes\phi).
$$
It implies that $B^*_\mathrm{int}$ is compatible with the decomposition \eqref{decomposition}, i.e.,
$$
B_\mathrm{int}^*=\bigoplus_{\underline v\in\{0,1,\cdots,p-1\}^2}B^*_\mathrm{int}\cap K^*_{\underline v}.
$$
By the definition of intersection complexes, one can check that $\varphi_\mathrm{int}$ is injective whose image is $B^*_\mathrm{int}\cap K^*_{(0,0)}$. Moreover, $B^\bullet_\mathrm{int}\cap K^\bullet_{\underline v}$  is an invariant subspace of \eqref{homotopy}. It follows that \eqref{varphi int} is a quasi-isomorphism.
\iffalse
Without loss of generality, we may assume $g(x)=1$. By assumption we know that $(p,r)=1$. Consider the Taylor expansion
$$
g'^{1/r}=(1+(1-g'))^{1/r}=\sum_{i=0}^\infty\left(\begin{matrix}
1/r\\
i
\end{matrix}
\right)(1-g')^i\in\mathcal O^*_{\mathfrak X'},
$$
where
$$
\left(\begin{matrix}
1/r\\
i
\end{matrix}
\right)=\frac{\frac{1}{r}\cdots(\frac{1}{r}-i+1)}{i!}.
$$
One can check that this number is a $p$-adic integer. Set $T:=g^{1/r}t$. We have the following facts:
\begin{itemize}
\item[(i)\ ] $\mathcal O_{\mathfrak X}\cong k[[t]],\mathcal O_{\mathfrak X'}\cong k[[t']]$ and $F^*:\mathcal O_{\mathfrak X'}\to F_*\mathcal O_{\mathfrak X}$ sends $\lambda\in k,t'$ to $\lambda^p,t^p$, respectively;
\item[(ii)\ ] $F_*\mathcal O_{\mathfrak X}\cong\mathcal O_{\mathfrak X'}\{1,T,\cdots,T^{p-1}\}$;
\item[(iii)\ ] let $\mathfrak D'=(t'=0)$ and $\mathfrak D=(t=0)$, then
$$
\begin{array}{rl}
\Omega^1_{\mathfrak X'/k}\cong\mathcal O_{\mathfrak X'}\cdot dT',&\Omega^1_{\mathfrak X/k}\cong\mathcal O_{\mathfrak X}\cdot dT,\\
\Omega^1_{\mathfrak X'/k}(\mathrm{log}~\mathfrak D')\cong\mathcal O_{\mathfrak X'}\cdot d\mathrm{log}~T',&\Omega^1_{\mathfrak X/k}(\mathrm{log}~\mathfrak D)\cong\mathcal O_{\mathfrak X}\cdot d\mathrm{log}~T.
\end{array}
$$
\end{itemize}
Using the third isomorphism in (iii), we have an identification 
\begin{eqnarray}\label{identification T}
\Omega^1_{\mathfrak X'/k}(\mathrm{log}~\mathfrak D')\cong\mathcal O_{\mathfrak X'}.
\end{eqnarray}
Consequently, one has
$$
\Omega^*(\hat E,\hat\theta-\sum^\infty_{i=0}f'^{p^i-1}df'\wedge)\cong [\hat E\xrightarrow{\hat\theta-r\sum^\infty_{i=0}f'^{p^i}}\hat E].
$$
According to the facts above, one can check that the following morphism of complexes
\begin{eqnarray}\label{step 2.1}
\xymatrixcolsep{5pc}\xymatrix{
\hat E\ar[r]^-{Q_T}\ar[d]^-{\mathrm{id}\otimes1}&\hat E\ar[d]^-{\mathrm{id}\otimes1}\\
\hat E\otimes F_*\mathcal O_{\mathfrak X}\ar[r]^-{T\partial_T+Q_T\otimes\mathrm{id}}&\hat E\otimes F_*\mathcal O_{\mathfrak X}
}
\end{eqnarray}
is a quasi-isomorphism, where $Q_T=\hat\theta-r\sum^\infty_{i=0}f'^{p^i}$ and
$$T\partial_T(\hat e\otimes T^j)=\hat e\otimes jT^j,~\hat e\in\hat E,~0\leq j<p.$$
 In fact, the morphism is an inclusion of complexes and its quotient is acyclic.
 
By direct computation, one has
$$
T(\hat E\otimes F_*\mathcal O_{\mathfrak X})+\mathrm{Im}(T\partial_T+Q_T\otimes\mathrm{id})
=(T'\hat E+Q_T\hat E)\otimes1+\sum_{i=1}^{p-1}\hat E\otimes T^i.
$$
It follows that \eqref{step 2.1} induces a quasi-isomorphism of complexes
\begin{eqnarray}\label{step 2.2}
\xymatrixcolsep{5pc}\xymatrix{
\hat E\ar[r]^-{Q_T}\ar[d]^-{\mathrm{id}\otimes1}&T'\hat E+Q_T\hat E\ar[d]^-{\mathrm{id}\otimes1}\\
\hat E\otimes F_*\mathcal O_{\mathfrak X}\ar[r]^-{T\partial_T+Q_T\otimes\mathrm{id}}&T(\hat E\otimes F_*\mathcal O_{\mathfrak X})+\mathrm{Im}(T\partial_T+Q_T\otimes\mathrm{id}).
}
\end{eqnarray}
By virtue of (iii) again, one gets
$$
\begin{array}{rcl}
\Omega_\mathrm{int}^1(\hat E,\hat\theta-\sum^\infty_{i=0}f'^{p^i-1}df'\wedge)&=&\hat E\otimes\Omega^1_{\mathfrak X'/k}+(\hat\theta-\sum^\infty_{i=0}f'^{p^i-1}df'\wedge)(\hat E)\\
&\stackrel{\eqref{identification T}}{\cong}&T'\hat E+(\hat\theta-r\sum^\infty_{i=0}f'^{p^i})(\hat E)\\
&=&T'\hat E+Q_T\hat E.
\end{array}
$$
 Therefore the upper complex in \eqref{step 2.2} is isomorphic to the intersection complex $\Omega_\mathrm{int}^*(\hat E,\hat\theta-\sum^\infty_{i=0}f'^{p^i-1}df'\wedge)$.
\fi 

{\itshape Step 3.} Since $(H,\nabla)=C^{-1}(E,\theta)$, coupled with the explicit construction of $C^{-1}$, one has
$$
(\hat H,\hat\nabla+df\wedge)=(F^*\hat E,\nabla_\mathrm{can}+(\mathrm{id}\otimes\zeta)F^*\hat\theta+df\wedge).
$$
Consider the Artin-Hasse exponential of $f_{\underline{\alpha}}$, which is given formally by
$$
\mathrm{AH}(f_{\underline{\alpha}})=\mathrm{exp}(f_{\underline{\alpha}}+f_{\underline{\alpha}}^p/p+f_{\underline{\alpha}}^{p^2}/p^2+\cdots)
$$
and which in fact has $p$-adically integral coefficients. Set
$$
g:=\prod_{\underline\alpha\in\mathbb N^2}\mathrm{AH}(f_{\underline\alpha})\in\mathcal O^*_{\mathfrak X}.
$$
Clearly, multiplication by $g$ gives rise to an isomorphism of modules with integrable connection
$$
g:(\hat H,\hat\nabla+df\wedge)\to C_\zeta^{-1}(\hat E,\hat\vartheta).
$$
In particular, we have isomorphisms of complexes of $\mathcal O_{\mathfrak X'}$-modules
$$
B^*\cong F_*\Omega^*(\hat H,\hat\nabla+df\wedge),~B_\mathrm{int}^*\cong F_*\Omega_\mathrm{int}^*(\hat H,\hat\nabla+df\wedge).$$
This completes the proof of the theorem.
\iffalse
\begin{eqnarray}\label{step 3.1}
\xymatrixcolsep{5pc}\xymatrix{
E\otimes F_*\mathcal O_{\mathfrak X}\ar[d]^-{\mathrm{AH}(f)}\ar[r]^-{T\partial_T+\hat\theta\otimes\mathrm{id}+rf}&E\otimes F_*\mathcal O_{\mathfrak X}\ar[d]^-{\mathrm{AH}(f)}\\
E\otimes F_*\mathcal O_{\mathfrak X}\ar[r]^-{T\partial_T+Q_T\otimes\mathrm{id}}&E\otimes F_*\mathcal O_{\mathfrak X}
}
\end{eqnarray}
is an isomorphism. It follows that the induced morphism of subcomplexes
\begin{eqnarray}\label{step 3.2}
\xymatrixcolsep{5pc}\xymatrix{
E\otimes F_*\mathcal O_{\mathfrak X}\ar[d]^-{\mathrm{AH}(f)}\ar[r]^-{T\partial_T+\hat\theta\otimes\mathrm{id}+rf}&E\otimes F_*\mathcal O_{\mathfrak X}\ar[d]^-{\mathrm{AH}(f)}+\mathrm{Im}(T\partial_T+\hat\theta\otimes\mathrm{id}+rf)\\
E\otimes F_*\mathcal O_{\mathfrak X}\ar[r]^-{T\partial_T+Q_T\otimes\mathrm{id}}&E\otimes F_*\mathcal O_{\mathfrak X}+\mathrm{Im}(T\partial_T+Q_T\otimes\mathrm{id})
}
\end{eqnarray}
is an isomorphism. Note that the upper complex in \eqref{step 3.2} is isomorphic to  $F_*\Omega_\mathrm{int}^*(\hat H,\hat\nabla+df\wedge)$.

By combing \eqref{step 1.1}, \eqref{step 2.1} and \eqref{step 3.1}, we conclude that 
$$
\begin{array}{c}
A^*:=\Omega^*(\hat E,\hat\theta-\sum^\infty_{i=0}f'^{p^i-1}df'\wedge),\\
B^*:=[\hat E\otimes F_*\mathcal O_{\mathfrak X}\xrightarrow{T\partial_T+Q_T\otimes\mathrm{id}}\hat E\otimes F_*\mathcal O_{\mathfrak X}]
\end{array}
$$
are the desired complexes in \eqref{outline}. Similarly, by combing \eqref{step 1.2}, \eqref{step 2.2} and \eqref{step 3.2}, we obtain \eqref{formal int}. This completes the proof.\fi
\end{proof}

As immediate consequences of the theorem above, we have
\begin{corollary}\label{C1}
Assume that $\Omega^*(E,\theta+df'\wedge)$ (resp. $\Omega_\mathrm{int}^*(E,\theta+df'\wedge)$) is acyclic outside a finite subset of $X'$. Then we have 
\begin{eqnarray}\label{twisted de Rham-Higgs comparison isolated sing}
\begin{array}{c}
F_*\Omega^*(H,\nabla+df\wedge)\cong\Omega^*(E,\theta-df'\wedge)~\mathrm{in}~D(X')\\
(\mathrm{resp.}~F_*\Omega_\mathrm{int}^*(H,\nabla+df\wedge)\cong\Omega_\mathrm{int}^*(E,\theta-df'\wedge)~\mathrm{in}~D(X').)
\end{array}
\end{eqnarray}
\end{corollary}
\begin{corollary}
Assume that $X$ is a smooth curve over $k$. Then we have 
$$
F_*\Omega^*(H,\nabla+df\wedge)\cong\Omega^*(E,\theta-df'\wedge)~\mathrm{in}~D(X')
$$
and
$$
F_*\Omega_\mathrm{int}^*(H,\nabla+df\wedge)\cong\Omega_\mathrm{int}^*(E,\theta-df'\wedge)~\mathrm{in}~D(X').
$$
\end{corollary}
\begin{definition}[{\cite[Definition 2.3.2]{Ogus04}}]
Let $(E,\theta)$ be any coherent Higgs module on $(X',D')/k$. Define the critical set of $\Omega^*(E,\theta)$ (resp. $\Omega_\mathrm{int}^*(E,\theta)$) to be
$$
\begin{array}{c}
\mathrm{CS}(E,\theta):=\{x\in X|\Omega^*(E,\theta)_x~\mathrm{is}~\mathrm{not}~\mathrm{acyclic}\}\\
(\mathrm{resp.}~\mathrm{CS}_\mathrm{int}(E,\theta):=\{x\in X|\Omega_\mathrm{int}^*(E,\theta)_x~\mathrm{is}~\mathrm{not}~\mathrm{acyclic}\}).
\end{array}
$$
\end{definition}
It is natural to ask
\begin{question}\label{question}
Is \eqref{twisted de Rham-Higgs comparison isolated sing} holds if $\mathrm{CS}(E,\theta-df'\wedge)$ (resp. $\mathrm{CS}_\mathrm{int}(E,\theta-df'\wedge)$) has positive dimension?
\end{question}
 Let $Z$ be the reduced closed subscheme of $X$ with support $\mathrm{CS}(E,\theta-df'\wedge)$ (resp. $\mathrm{CS}_\mathrm{int}(E,\theta-df'\wedge)$). Set 
 $$
 \begin{array}{c}
 (\hat E,\hat\theta):=(E,\theta)_{/Z'},~(\hat H,\hat\nabla):=(H,\nabla)_{/F^*Z'},\\
 (\mathfrak X',\mathfrak D'):=(X',D')_{/Z'},~(\mathfrak X,\mathfrak D):=(X,D)_{/F^*Z'}.
 \end{array}
 $$
 We attack this problem by using a similar approach as presented in the proof of the theorem above. The key is to find a decomposition $f=\sum f_i$, $f_i|_Z=0$ such that we can run the following steps in $D(\mathfrak X')$.
 \begin{itemize}
\item[{\itshape Step 1}] (Twist Higgs Higgs field)
$$
\begin{array}{c}
\Omega^*(\hat E,\hat\theta-df'\wedge)\cong\Omega^*(\hat E,\hat\theta-\sum_i\sum_{j=0}^\infty f_i'^{p^j-1}df'_i\wedge),\\
(\mathrm{resp.}~\Omega_\mathrm{int}^*(\hat E,\hat\theta-df'\wedge)\cong\Omega_\mathrm{int}^*(\hat E,\hat\theta-\sum_i\sum_{j=0}^\infty f_i'^{p^j-1}df'_i\wedge));
\end{array}
$$
\item[{\itshape Step 2}] (De Rham-Higgs comparison)
$$
\begin{array}{c}
(\tau_<)F_*\Omega^*(\hat H,\hat\nabla-\sum_i\sum_{j=1}^\infty f_i^{p^j-1}df_i\wedge)\cong(\tau_<)\Omega^*(\hat E,\hat\theta-\sum_i\sum_{j=0}^\infty f_i'^{p^j-1}df'_i\wedge),\\
(\mathrm{resp.}~(\tau_<)\Omega_\mathrm{int}^*(\hat H,\hat\nabla-\sum_i\sum_{j=0}^\infty f_i^{p^j-1}df_i\wedge)\cong(\tau_<)\Omega_\mathrm{int}^*(\hat E,\hat\theta-\sum_i\sum_{j=0}^\infty f_i'^{p^j-1}df'_i\wedge)),
\end{array}
$$
wehere $(\tau_<)$ means a suitable truncation if necessary;
\item[{\itshape Step 3}] (Twist integrable connection) 
$$
(\hat H,\hat \nabla+df\wedge)\cong(\hat H,\hat\nabla-\sum_i\sum_{j=0}^\infty f_i^{p^j-1}df_i\wedge).
$$
\end{itemize}
The following and the next section are devoted to partially answer this question. 
 \begin{corollary}\label{C2}
 Suppose that there is a system of coordinates $(t_1,\cdots,t_n)$ around $\mathrm{CS}(E,\theta-df'\wedge)$ (resp. $\mathrm{CS}_\mathrm{int}(E,\theta-df'\wedge)$) such that $D=(t_1\cdots t_r=0)$, and that
there is decomposition $f=\sum_{i=1}^mf_i$, where
\begin{eqnarray}\label{coordinate term}
f_i=g_i^pt_1^{\alpha^i_1}\cdots t_n^{\alpha^i_n},~g_i\in\Gamma(X,\mathcal O_X),~\alpha^i_1,\cdots,\alpha^i_n\in\mathbb N
\end{eqnarray}
and $\{j|\alpha^i_j>0\}\cap\{j|\alpha^s_j>0\}=\emptyset$ for $i\neq s$. 
Then  the question above has an affirmative answer.
 \end{corollary}
\begin{proof}
  It suffices to prove that
$$
\begin{array}{c}
 F_*\Omega^*(\hat H,\hat\nabla+df\wedge)\cong\Omega^*(\hat E,\hat\theta-df'\wedge)~\mathrm{in}~D(\mathfrak X')\\
 (\mathrm{resp.}~F_*\Omega_\mathrm{int}^*(\hat H,\hat\nabla+df\wedge)\cong\Omega_\mathrm{int}^*(\hat E,\hat\theta-df'\wedge)~\mathrm{in}~D(\mathfrak X')).
 \end{array}
 $$
 The isomorphism can be verified by three steps as in the proof of Theorem \ref{formal de Rham-Higgs isolated}. Note that by \eqref{coordinate term}, the zero locus of $f_i$ contains $\mathrm{CS}(E,\theta-df'\wedge)$ (resp. $\mathrm{CS}_\mathrm{int}(E,\theta-df'\wedge)$). It allows us to make the following constructions:
 $$
 A^*:=\Omega^*(\hat E,\hat\theta-\sum_{i=1}^m\sum_{j=0}^\infty f_i'^{p^j-1}df_i')~(\mathrm{resp.}~A^*_\mathrm{int}:=\Omega_\mathrm{int}^*(\hat E,\hat\theta-\sum_{i=1}^m\sum_{j=0}^\infty f_i'^{p^j-1}df_i'))
 $$
 and 
 $$
 \begin{array}{c}
 B^*:=\Omega^*(F^*\hat E,\nabla_{\mathrm{can}}+(\mathrm{id}\otimes\zeta)F^*\hat\theta-\sum_{i=1}^m\sum_{j=1}^\infty f_i^{p^j-1}df_i)\\
 (\mathrm{resp.}~B_\mathrm{int}^*:=\Omega_\mathrm{int}^*(F^*\hat E,\nabla_{\mathrm{can}}+(\mathrm{id}\otimes\zeta)F^*\hat\theta-\sum_{i=1}^m\sum_{j=1}^\infty f_i^{p^j-1}df_i)).
 \end{array}
 $$
 Here  $\nabla_\mathrm{can}$ is the Cartier descent connection and 
$$\zeta:F^*\Omega^1_{\mathfrak X'/k}(\log\mathfrak D')\to\Omega^1_{\mathfrak X/k}(\log\mathfrak D)$$ is given by $F^*d\log t'_i\mapsto d\log t_i$ for $i\leq r$ and $F^*dt_i'\mapsto t_i^{p-1}dt_i$ for $i>r$. The remaining arguments are clear, we omit the details. This completes the proof.
\iffalse
  Let $(t_1,\cdots,t_n)$ be a system of local coordinates around $x$ such that $f=t_1\cdots t_n$  and $t_1(x)=\cdots=t_n(x)$. By Theorem \ref{Cartier Ogus}, it suffices to show that 
 \begin{eqnarray}\label{twisted de Rham-Higgs higher dimension}
 F_*\Omega^*(\hat H,\hat\nabla+df\wedge)\cong\Omega^*(\hat E,\hat\theta-df'\wedge)~\mathrm{in}~D(\mathfrak X').
 \end{eqnarray}
 Similar to the proof of Proposition \ref{formal curve}, the isomorphism can be proved by three steps. We sketch them below.
 
 {\itshape Step 1.}  Set $\omega_i:=dt_i$ and $\hat\theta_i:=\hat\theta(\partial_{t_i})$.  Clearly, one has
 $$
( \hat\theta-df'\wedge)(\partial_{t'_i}):=\hat\theta_i-f'/t_i,~(\hat\theta-\sum_{i=0}^\infty f'^{p^i-1}df'\wedge)(\partial_{t'_i})=\hat\theta_i-\sum_{i=0}^\infty f'^{p^i}/t_i'.
 $$
 Set
$$
\hat h_i:=\frac{\hat\theta_i-\sum_{i=0}^\infty f'^{p^i}/t_i'}{\hat\theta_i-f'/t_i'}=1+(\sum_{i=0}^\infty f'^{p^i-1})^p\sum_{j=0}^{p-2} f'^{p-1-j}t_i^j\hat\theta^j_i\in\mathrm{Aut}_{\mathcal O_{\mathfrak X}}(\hat E).
$$
These automorphisms give rise to an isomorphism of complexes 
 $$
 \begin{array}{c}
 \Omega^*(\hat E,\hat\theta-df'\wedge)\cong\Omega^*(\hat E,\hat\theta-\sum_{i=0}^\infty f'^{p^i-1}df'\wedge),\\
 \hat e\otimes\omega_I'\mapsto(\prod_{i\in I}h_i)(\hat e)\otimes\omega_I',\\
 \hat e\in\hat E,~I\subset\{1,\cdots,n\}.
 \end{array}
 $$
 
 {\itshape Step 2.} Let $s_i=t_i+1$ and let $\tilde F$ be a lifting of $F$ around $x$ such that $\tilde F(\tilde s_i')=\tilde s_i^p$. Consider a module with integrable connection 
 $$
C^{-1}_{\tilde F}(\hat E,\hat\theta-\sum_{i=0}^\infty f'^{p^i-1}df'\wedge)=(F^*\hat E,\nabla_\mathrm{can}
 +(\mathrm{id}\otimes\hat\zeta_{\tilde F})F^*\hat\theta-\sum_{i=1}^\infty f^{p^i-1}df\wedge).$$
 Since $\hat\theta-\sum_{i=0}^\infty f'^{p^i-1}df'\wedge$ is a topological nilpotent Higgs field on $\hat E$, there is a quasi-isomorphism of complexes
 $$
\varphi_{\tilde F}:\Omega^*(\hat E,\hat\theta-\sum_{i=0}^\infty f'^{p^i-1}df'\wedge)\to\Omega^*C^{-1}_{\tilde F}(\hat E,\hat\theta-\sum_{i=0}^\infty f'^{p^i-1}df'\wedge).$$

{\itshape Step 3.} The Artin-Hasse exponential $\mathrm{AH}(f)$ of $f$ gives rise to an isomorphism of modules with integrable connection
$$
(F^*\hat E,\nabla_\mathrm{can}
 +(\mathrm{id}\otimes\hat\zeta_{\tilde F})F^*\hat\theta+df\wedge)\cong C^{-1}_{\tilde F}(\hat E,\hat\theta-\sum_{i=0}^\infty f'^{p^i-1}df'\wedge).
 $$
 On the other hand, by construction we have
 $$
 (\hat H,\hat\nabla)\cong C^{-1}_{\tilde F}(\hat E,\hat\theta)=(F^*\hat E,\nabla_\mathrm{can}
 +(\mathrm{id}\otimes\hat\zeta_{\tilde F})F^*\hat\theta).$$
 
 Obviously, \eqref{twisted de Rham-Higgs higher dimension} follows immediately from the three steps above. \fi
 \end{proof}

\section{De Rham-Higgs comparison for vanishing cycles \uppercase\expandafter{\romannumeral2}}
\subsection{The formal inverse Cartier transform}
Let $(X,D)/k$ be as in previous section.  Let $Z$ be a closed subscheme of $X$ and then set
$$(\mathfrak X,\mathfrak D):=(X,D)_{/F^*Z},~(\mathfrak X',\mathfrak D'):=(X',D')_{/Z'}.$$
Let $\mathcal L_{\mathcal X/\mathcal S}$ be the sheaf of local liftings $\tilde F$ of $F$  such that $F^*\tilde D'=p\tilde D$. Assume that there is a subsheaf $\mathcal L\subset F_*\mathcal L_{\mathcal X/\mathcal S}$ such that $\mathcal L_{x'}\neq\emptyset$ holds for any $x\in Z'$.  Put 
$$I_{\mathcal L}:=\{(U',\tilde F)|\tilde F\in\Gamma(U',\mathcal L)\},~U'_\alpha:= U',~\tilde F_\alpha:=\tilde F~\mathrm{for}~\alpha=( U',\tilde F)\in I_\mathcal{L}.$$ 
Set $h_{\alpha\beta}:=h_{\tilde F_\alpha\tilde F_\beta},~\zeta_\alpha:=\zeta_{\tilde F_\alpha}.$
\iffalse
Suppose that there are two families 
$\{\zeta_\alpha:F^*\Omega^1_{U'_\alpha/S}\to\Omega^1_{U_\alpha/S}\}$ and $\{h_{\alpha\beta}:F^*\Omega^1_{U'_{\alpha\beta}/S}\to\mathcal O_{U_{\alpha\beta}}\}$
 such that
\begin{itemize}
\item $\zeta_\beta-\zeta_\alpha=dh_{\alpha\beta}$, and
\item $h_{\alpha\beta}+h_{\beta\gamma}=h_{\alpha\gamma}$.
\end{itemize}
\fi

This subsection contains two inputs. The simple one is that {\itshape $\mathcal L$ can be used to construct an inverse Cartier transform from $\mathrm{HIG}_{\leq p-1}((\mathfrak X',\mathfrak D')/k)$ to $\mathrm{MIC}_{\leq p-1}((\mathfrak X,\mathfrak D)/k)$.} The subtle one is that {\itshape in the construction, the nilpotent condition on Higgs field is not necessary which may be replaced by a weaker $(\mathcal L,h)$-nilpotent condition on Higgs field.}
\begin{definition}
For any non-negative integer $\ell$, we say a Higgs field $(E,\theta)$ on $(\mathfrak X',\mathfrak D')/k$ is $(\mathcal L,h)$-nilpotent of level $\leq\ell$ if
\begin{eqnarray}\label{nilpotent h}
\prod_{i=1}^{\ell+1}(\mathrm{id}\otimes h_{\alpha_i\beta_i})F^*\theta=0,~\forall\alpha_i,\beta_i\in I_{\mathcal L},~1\leq i\leq p.
\end{eqnarray}
Denote by $\mathrm{HIG}_{(\mathcal L,h)\leq\ell}((\mathfrak X',\mathfrak D')/k)$ the full subcategory of $\mathrm{HIG}((\mathfrak X',\mathfrak D')/k)$ consisting of objects described above.
\end{definition}
Clearly, $\mathrm{HIG}_{\leq\ell}((\mathfrak X',\mathfrak D')/k)$ is contained in $\mathrm{HIG}_{(\mathcal L,h)\leq\ell}((\mathfrak X',\mathfrak D')/k)$ and
$$
\mathrm{HIG}_{\leq\ell}((\mathfrak X',\mathfrak D')/k)=\mathrm{HIG}_{(\mathcal L_{\mathcal X/\mathcal S},h)\leq\ell}((\mathfrak X',\mathfrak D')/k).$$
Let $\omega\in\Gamma(X,\Omega^1_{X_{\log}/k})$ be a global non-zero one-form such that
$$
h_{\alpha\beta}(F^*\omega')=0,~\forall\alpha,\beta\in I_{\mathcal L}.
$$
We are interested in the following twisting functor
\begin{eqnarray}\label{twisting functor}
\begin{array}{c}
T_{\omega'}:\mathrm{HIG}_{\leq\ell}((\mathfrak X',\mathfrak D')/k)\to\mathrm{HIG}_{(\mathcal L,h)\leq \ell}((\mathfrak X',\mathfrak D')/k),\\
(E,\theta)\mapsto (E,\theta+\omega').
\end{array}
\end{eqnarray}
One can check that if $E$ is torsion-free over $\mathcal O_{X'}$, then $T_{\omega'}(E,\theta)$ does not belong to $\mathrm{HIG}_{\leq\ell}((\mathfrak X',\mathfrak D')/k)$.

The construction of the inverse Cartier transform on $\mathrm{HIG}_{\leq\ell}((\mathfrak X',\mathfrak D')/k)$ can be extended to $\mathrm{HIG}_{(\mathcal L,h)\leq \ell}((\mathfrak X',\mathfrak D')/k)$ verbatim. 
\begin{lemma}\label{gluing}
Given $(E,\theta)\in\mathrm{HIG}_{(\mathcal L,h)\leq p-1}((\mathfrak X',\mathfrak D')/k)$. The family of modules with integrable connection
$$
\{(H_\alpha,\nabla_\alpha):=((F^*E)|_{U_\alpha},\nabla_\mathrm{can}+(\mathrm{id}\otimes\zeta_\alpha)F^*\theta)\}_{\alpha\in I_{\mathcal L}}
$$
can be glued together by a family of transition morphisms
$$
\{G_{\alpha\beta}:=\exp((\mathrm{id}\otimes h_{\alpha\beta})F^*\theta)=\sum_{i=0}^{p-1}\frac{((\mathrm{id}\otimes h_{\alpha\beta})F^*\theta)^i}{i!}\}_{\alpha,\beta\in I_{\mathcal L}},
$$
In other words, we have a cocycle condition
$$
G_{\alpha\beta}G_{\beta\gamma}=G_{\alpha\gamma},~\forall\alpha,\beta,\gamma
$$
and a commutative diagram 
$$
\xymatrix{
H_\beta|_{\mathfrak U_{\alpha\beta}}\ar[r]^-{\nabla_\beta}\ar[d]^-{G_{\alpha\beta}}&H_\beta|_{\mathfrak U_{\alpha\beta}}\otimes\Omega^1_{\mathfrak U_{\alpha\beta}/k}\ar[d]^-{G_{\alpha\beta}\otimes\mathrm{id}}\\
H_\alpha|_{\mathfrak U_{\alpha\beta}}\ar[r]^-{\nabla_\alpha}&H_\alpha|_{\mathfrak U_{\alpha\beta}}\otimes\Omega^1_{\mathfrak U_{\alpha\beta}/k}.
}
$$
\end{lemma}
\begin{proof}
The cocycle condition above follows immediately from \eqref{nilpotent h}. By using the loc. cit. again, one has
$$
\begin{array}{rcl}
\nabla_\alpha G_{\alpha\beta}(F^*e)&=&(G_{\alpha\beta}\otimes\mathrm{id})(\mathrm{id}\otimes dh_{\alpha\beta})(F^*(\theta e))+(G_{\alpha\beta}\otimes\zeta_\alpha)(F^*(\theta e))\\
&=&(G_{\alpha\beta}\otimes dh_{\alpha\beta})(F^*(\theta e))+(G_{\alpha\beta}\otimes\zeta_\alpha)(F^*(\theta e))\\
&=&(G_{\alpha\beta}\otimes\zeta_\beta)(F^*(\theta e))\\
&=&(G_{\alpha\beta}\otimes\mathrm{id})\nabla_\beta(F^*e),
\end{array}
$$
from which the commutative diagram above follows. This completes the proof.
\end{proof}
\begin{definition}
The lemma above gives rise to a functor
\begin{eqnarray}\label{inverse Cartier non-nilpotent}
C^{-1}_\mathcal{L}:\mathrm{HIG}_{(\mathcal L,h)\leq p-1}((\mathfrak X',\mathfrak D')/k)\to\mathrm{MIC}((\mathfrak X,\mathfrak D)/k).
\end{eqnarray}
For simplicity, we drop the lower decoration of $C^{-1}_{\mathcal L}$ and abbreviate it to $C^{-1}$.
\end{definition}
\subsection{The $\mathcal L$-indexed $\infty$-homotopy}
\cite[Theorem 3.5]{SZ2} showed that there is an $F_*\mathcal L_{\mathcal X/\mathcal S}$-indexed $\infty$-homotopy
\begin{eqnarray}\label{F_*L indexed homotopy}
\qquad\mathrm{Ho}:\Delta_*(F_*\mathcal L_{\mathcal X/\mathcal S})\to\mathcal Hom^*_{\mathcal O_{X'}}(\tau_{<p-\ell}\Omega^*(E,\theta),\tau_{<p-\ell}F_*\Omega^*C^{-1}(E,\theta))
\end{eqnarray}
for any $(E,\theta)\in\mathrm{HIG}_{\leq\ell}((X',D')/k)$ ($0\leq\ell<p$). Moreover, $\mathrm{Ho}^0(\tilde F)$ ($\tilde F\in\mathcal L_{\mathcal X/\mathcal S}$) is induced by the morphism of complexes
\begin{eqnarray}\label{F_alpha}
\begin{array}{c}
\varphi_\alpha:\Omega^*(E|_{U'},\theta|_{U'})\to F_*\Omega^*(H_\alpha,\nabla_\alpha),\\
e\mapsto e\otimes 1,\\
e\otimes\omega_1\wedge\cdots\wedge\omega_s\mapsto e\otimes\zeta_\alpha(\omega_1)\wedge\cdots\wedge\zeta_\alpha(\omega_s).
\end{array}
\end{eqnarray}
Here $\alpha=(U',\tilde F)\in I_{F_*\mathcal L_{\mathcal X/\mathcal S}}$ and we have used a natural identification 
\begin{eqnarray}\label{identification}
F_*\Omega^\bullet(H_\alpha,\nabla_\alpha)\cong E|_{U_\alpha'}\otimes F_*\Omega^\bullet_{U_\alpha/k}.
\end{eqnarray}

Set $\hat{\mathcal L}$ to be the pullback of $\mathcal L$ to $\mathfrak X'$. In this subsection, we aim to extend \eqref{F_*L indexed homotopy} to the following 
\begin{theorem}\label{infty homotopy non-nilpotent}
There are $\hat{\mathcal L}$-indexed $\infty$-homotopies
\begin{eqnarray}\label{L indexed homotopy}
\mathrm{Ho}:\Delta_*(\hat{\mathcal L})\to\mathcal Hom^*_{\mathcal O_{\mathfrak X'}}(\tau_{<p-\ell}\Omega^*(E,\theta),\tau_{<p-\ell}F_*\Omega^*C^{-1}(E,\theta))
\end{eqnarray}
and
\begin{eqnarray}\label{L indexed homotopy int}
\mathrm{Ho}:\Delta_*(\hat{\mathcal L})\to\mathcal Hom^*_{\mathcal O_{\mathfrak X'}}(\tau_{<p-\ell}\Omega_\mathrm{int}^*(E,\theta),\tau_{<p-\ell}F_*\Omega_\mathrm{int}^*C^{-1}(E,\theta)),
\end{eqnarray}
for $(E,\theta)\in\mathrm{HIG}_{(\mathcal L,h)\leq\ell-1}((\mathfrak X',\mathfrak D')/k)~(0\leq\ell<p)$. Here \eqref{L indexed homotopy int} is induced by \eqref{L indexed homotopy} and $\mathrm{Ho}^0(\tilde F_\alpha)$ is induced by \eqref{F_alpha} for $\alpha\in I_{\mathcal L}$.
 \end{theorem}
\begin{proof}
Firstly, we construct \eqref{L indexed homotopy} as a morphism of graded shaves of abelian groups. For any $r>0$, we define the $r$th component of $\mathrm{Ho}$ by
$$
\mathrm{Ho}^r(\tilde F_0,\cdots,\tilde F_r):=\{\varphi(r,s)_{\underline\alpha}:\Omega^{r+s}(E|_{\mathfrak U'},\theta|_{\mathfrak U'})\to F_*\Omega^sC^{-1}(E,\theta)|_{\mathfrak U'}\}_{s\geq0},
$$
where $\tilde F_0,\cdots,\tilde F_r\in\Gamma(U',\mathcal L)$, $\alpha_i:=(U,\tilde F_i),~\underline\alpha:=(\alpha_0,\cdots,\alpha_r)$. Note that $\tilde F_{\alpha_0}$ induces an identification 
\begin{eqnarray}\label{iota}
\iota_{\alpha_0}:(H_{\alpha_0},\nabla_{\alpha_0})\cong C^{-1}(E,\theta)|_{\mathfrak U}.
\end{eqnarray}
Define
\begin{eqnarray}\label{construction varphi}
\varphi(r,s)_{\underline\alpha}=\sum a(\overline j,\underline s)(\iota_{\alpha_0}\otimes\mathrm{id})(\sigma_{\underline\alpha}^{[\overline j]}\otimes\zeta_{\underline\alpha,\underline s}h_{\underline\alpha})F^*(\mathrm{id}\otimes\delta_{r+s}),
\end{eqnarray}
where the summation takes over 
$$\overline j=(j^1,\cdots,j^r)\in\mathbb N^r,~\underline s=(s_0,\cdots,s_r)\in\mathbb N^{r+1},~s_0+\cdots+s_r=s$$
 such that $j^1+\cdots+j^r+r+s<p$ and
\begin{itemize}
\item $F^*(\mathrm{id}\otimes\delta_{r+s}):E|_{\mathfrak U'}\otimes\Omega^{r+s}_{\mathfrak U'/k}\to (F^*E)|_{\mathfrak U}\otimes(F^*\Omega^1_{\mathfrak U'/k})^{\otimes(r+s)}$,
\item $\zeta_{\underline\alpha,\underline s}h_{\underline\alpha}:(F^*\Omega^1_{\mathfrak U'/k})^{\otimes(r+s)}\to\Omega^s_{\mathfrak U/k}$,
\item $\sigma^{[\overline j]}_{\underline\alpha}=\prod_{k=1}^r\frac{[(\mathrm{id}\otimes h_{\alpha_{k-1}\alpha_k})F^*\theta]^{j^k}}{j^k!}$,
\item $a(\overline j,\underline s)\in\mathbb F_p$.
\end{itemize}

To check that the family $\{\mathrm{Ho}^r\}$ forms a morphism of complexes, it suffices to verify 
$$
d_{\mathcal Hom^*_{\mathcal O_{X'}}}\mathrm{Ho}^r(\tilde F_0,\cdots,\tilde F_r)=\mathrm{Ho}^{r-1}\partial(\tilde F_0,\cdots,\tilde F_r).$$
Equivalently, for any $r>0,s\geq0,r+s<p-\ell$, the expression
\begin{eqnarray}\label{morphism of complexes}
\qquad(-1)^s\nabla\varphi(r,s)_{\underline\alpha}+(-1)^{s+1}\varphi(r,s)_{\underline\alpha}\theta-\sum_{k=0}^r(-1)^k\varphi(r-1,s+1)_{\hat{\underline\alpha}(k)}
\end{eqnarray}
vanishes, where $\hat{\underline\alpha}(k)$ is obtained from $\underline\alpha=(\alpha_0,\cdots,\alpha_k,\cdots,\alpha_r)$ by deleting $\alpha_k$. According to the proof of SZ, $(\iota_{\alpha_0}^{-1}\otimes\mathrm{id})\eqref{morphism of complexes}$ equals to
\begin{eqnarray}\label{difference}
\begin{array}{cl}
&(\mathrm{id}\otimes\zeta_{\alpha_0})F^*\theta\sum a(\overline j,\underline s)(\sigma_{\underline\alpha}^{[\overline j]}\otimes\zeta_{\underline\alpha,\underline s}h_{\underline\alpha})F^*(\mathrm{id}\otimes\delta_{r+s})\\
-&\sum a(\overline j,\underline s)(\sigma_{\underline\alpha}^{[\overline j]}\otimes\zeta_{\hat{\underline\alpha}(0),\underline s}h_{\hat{\underline\alpha}(0)})F^*(\mathrm{id}\otimes\delta_{r+s}),
\end{array}
\end{eqnarray}
where the first summation takes over 
$$\overline j=(j^1,\cdots,j^r)\in\mathbb N^r,~\underline s=(s_0,\cdots,s_r)\in\mathbb N^{r+1},~s_0+\cdots+s_r=s$$
 such that $j^1+\cdots+j^r+r+s+1=p$ and the second summation takes over
 $$\overline j=(j^1,\cdots,j^r)\in\mathbb N^r,~\underline s=(s_1,\cdots,s_r)\in\mathbb N^r,~s_1+\cdots+s_r=s$$
 such that $j^1+\cdots+j^r+r+s\geq p$. Since $(E,\theta)$ is $(\mathcal L,h)$-nilpotent of level $\leq\ell-1$, the factor $\sigma_{\underline\alpha}^{[\overline j]}$ in both summations of \eqref{difference} vanishes. It follows that the equality \eqref{morphism of complexes} holds and hence \eqref{L indexed homotopy}, \eqref{L indexed homotopy int} follows. This completes the proof.
\end{proof}
\subsection{The formal Cartier isomorphism}
Suppose that  $(X,D)/k$ admits a system of coordinates $(t_1,\cdots,t_n)$ and let $\tilde F$ be a global  lifting of $F$ such that $\tilde F(\tilde t'_i)=\tilde t_i^p$ for any $1\leq i\leq n$. Let $(E,\theta)$ be a Higgs module on $(\mathfrak X',\mathfrak D')/k$ and set $(H,\nabla)=C^{-1}_{\tilde F}(E,\theta)$. Define a morphism of complexes of $\mathcal O_{\mathfrak X'}$-modules as \eqref{F_alpha}:
$$
\varphi:\Omega^*(E,\theta)\to F_*\Omega^*(H,\nabla).
$$
Clearly, it induces a morphism of subcomplexes
$$
\varphi_\mathrm{int}:\Omega_\mathrm{int}^*(E,\theta)\to F_*\Omega_\mathrm{int}^*(H,\nabla).
$$

Using similar arguments as in the proof of \cite[Proposition 3.8]{SZ2}, we have
\begin{proposition}\label{formal Cartier iso}
$\varphi,\varphi_\mathrm{int}$ are quasi-isomorphism if $\theta$ is topologically nilpotent.
 \end{proposition}
In fact, the theorem holds for any global lifting $\tilde F$ of $F$ which will be proved in my next paper.
\subsection{Comparison with non-isolated critical set}
In this subsection, we shall attack Question \ref{question} under the case B. This case is very subtle, we can only tackle a few situations. To simplify the explanation, we only consider a typical example: $X=Y\times\mathbb A^1$, $0\neq f\in k[t]$ ($t$ is the standard coordinate on $\mathbb A^1$), $D=\pi_Y^*D_1+\pi^*_{\mathbb A^1}D_2$ ($D_1$ is a SNCD on $Y$ and $D_2$ is divisor on $\mathbb A^1$ ) and $(X,D)/k$ has a $W_2(k)$-lifting $(\tilde X,\tilde D)/W_2(k)$.
\iffalse
the following situations:
\begin{itemize}
\item[(1)\ ] $(E,\theta)=(\mathcal O_{X'},0)$ and $f\in\Gamma(X,\mathcal O_X)$ admits a lifting $\tilde f\in\Gamma(\tilde X,\mathcal O_{\tilde X})$ such that $\tilde f:\tilde X\to\mathbb A^1_{W_2(k)}$ is a semistable morphism of $W_2(k)$-schemes.
\item[(2)\ ]
\end{itemize}
 Suppose that $f(D)$ is a finite set. Define $\mathcal L\subset F_*\mathcal L_{\mathcal X/\mathcal S}$ to be the subsheaf consisting of $\tilde F\in\mathcal L_{\mathcal X/\mathcal S}$ such that $\tilde F^*(\tilde f')=\tilde f^p$. Let $Z$ be the reduced closed subscheme of $X$ with support $\mathrm{CS}(\mathcal O_{X'},-df'\wedge)$. 

Using the results established in previous subsections, we have  
\begin{theorem}
Notation and assumption as above. Then
$$
\tau_{<p-1}F_*\Omega^*(\mathcal O_X,d+df\wedge)\cong\tau_{<p-1}\Omega^*(\mathcal O_{X'},-df'\wedge)~\mathrm{in}~D(X').
$$
\end{theorem}
\fi




\iffalse
Let $0<\ell\leq p$ and let $(E,\theta)$ be a coherent Higgs module on $(X',D')/k$ of level $\leq\ell-1$. To avoid some technique problem, we make the following assumptions:

Note that locally on $\mathfrak X'$, there is an isomorphism of complexes between the two complexes above. We prefer to believe that {\itshape the assumptions are already true}, but we do not know how to check it so far. Nevertheless, there are some examples justify the truth of \eqref{assumption on Higgs}.
\begin{lemma}
If $(E,\theta)=(\mathcal O_{X'},0)$ or $f:X\to\mathbb A^1$ is a $Y$-bundle over $\mathbb A^1$ for some smooth $k$-variety $Y$, then \eqref{assumption on Higgs} holds.
\end{lemma}
\begin{proof}
If $(E,\theta)=(\mathcal O_{X'},0)$, then  \eqref{assumption on Higgs} obviously holds. Assume $f:X\to\mathbb A^1$ is a $Y$-bundle over $\mathbb A^1$ for some smooth $k$-variety $Y$. This assumption implies that the short exact sequence
$$
0\to f'^*\Omega^1_{\mathbb A_{\log}^1/k}\to\Omega^1_{X_{\log}'/k}\to\mathcal Q\to0
$$
is locally split over $\mathbb A^1$. For simplicity, we may assume that $E=E'=[0]$. Choose a splitting of the sequence above around $f'^{-1}(0)$, then it gives rise to a decomposition
$$
\Omega^1_{X_{\log}'/k}\cong f'^*\Omega^1_{\mathbb A_{\log}^1/k}\bigoplus\mathcal Q.
$$
Note that
$$
\Omega^i_{X_{\log}'/k}\cong\Omega^i_1\bigoplus\Omega^i_2,~\Omega^i_1:=(f'^*\Omega^1_{\mathbb A_{\log}^1/k}\bigotimes\bigwedge^{i-1}\mathcal Q),~\Omega^i_2:=\bigwedge^i\mathcal Q.
$$
Set $\theta_1\otimes d\log f'$ to be the composite of
$$
E\stackrel{\theta}{\rightarrow}E\otimes\Omega^1_{X_{\log}'/k}\to E\otimes f'^*\Omega^1_{\mathbb A_{\log}^1/k}.
$$

Now we can define an isomorphism of complexes
$$
\begin{array}{c}
\Omega^*(\hat E,\hat\theta-df'\wedge)\longrightarrow\Omega^*(\hat E,\hat\theta-\sum_{i=0}^\infty f'^{p^i-1}df'\wedge),\\
\hat e_1\otimes\omega^i_1+\hat e_2\otimes^i_2\mapsto\frac{\hat\theta_1-\sum_{i=0}^\infty f'^{p^i}}{\hat\theta_1-f'}(\hat e_1)\otimes\omega^i_1+\hat e_2\otimes\omega^i_2,\\
\hat e_1,\hat e_2\in\hat E,~\omega^i_1\in\Omega^i_1,~\omega^i_2\in\Omega^i_2.
\end{array}
$$
One can check that the isomorphism above induces an isomorphism of subcomplexes
$$
\Omega_\mathrm{int}^*(\hat E,\hat\theta-df'\wedge)\longrightarrow\Omega_\mathrm{int}^*(\hat E,\hat\theta-\sum_{i=0}^\infty f'^{p^i-1}df'\wedge).
$$
This completes the proof.
\end{proof}
\fi
Applying the results established in previous subsections,  we have
\begin{theorem}\label{twisted de Rham-Higgs thm}
Let $(E,\theta)$ be a coherent Higgs module on $(X',D')/k$ of level $\leq\ell-1$. Set $(H,\nabla):=C^{-1}(E,\theta)$. Then
$$
\tau_{<p-\ell}F_*\Omega^*(H,\nabla+df\wedge)\cong\tau_{<p-\ell}\Omega^*(E,\theta-df'\wedge)~\mathrm{in}~D(X')
$$
and
$$
\tau_{<p-\ell}F_*\Omega_\mathrm{int}^*(H,\nabla+df\wedge)\cong\tau_{<p-\ell}\Omega_\mathrm{int}^*(E,\theta-df'\wedge)~\mathrm{in}~D(X').
$$
\end{theorem}
\begin{proof}
 Let  $\mathcal L$ be the  subsheaf of $\mathcal L_{\mathcal X/\mathcal S}$ consisting of local liftings $\tilde F$ such that $\tilde F(\tilde t')=\tilde t^p$. By the construction of $X$ and the assumption on $f$, there exist $c_1,\cdots,c_m\in k$ such that 
 $$\mathrm{CS}_\mathrm{int}(E,\theta-df'\wedge)\subset\cup_{i=1}^mf'^{-1}(c_i).$$
 For any $c\in k$, Let $Z_c$ be the reduced closed subscheme of $X$ with support $ f^{-1}(c)$. Let $\mathfrak X_c,\mathfrak X_c'$ be formal completions of $X,X'$ along $F^*Z'_c,Z'_c$, respectively.
Set 
$$\hat{\mathcal L}_c:=\mathcal L|_{\mathfrak X_c'},~(\hat E_c,\hat\theta_c):=(E,\theta)_{/Z_c'},~(\hat H_c,\hat\nabla_c):=(H,\nabla)_{/F^*Z_c'}.$$

To prove the theorem, it suffices to show that for any $c\in k$, we have
\begin{eqnarray}\label{c}
\tau_{<p-\ell}\Omega^*(\hat E_c,\hat\theta_c-df'\wedge)\cong\tau_{<p-\ell}F_*\Omega^*(\hat H_c,\hat\nabla_c+df\wedge)~\mathrm{in}~D(\mathfrak X_c')
\end{eqnarray}
and
\begin{eqnarray}\label{c int}
\quad \tau_{<p-\ell}\Omega_\mathrm{int}^*(\hat E_c,\hat\theta_c-df'\wedge)\cong\tau_{<p-\ell}F_*\Omega_\mathrm{int}^*(\hat H_c,\hat\nabla_c+df\wedge)~\mathrm{in}~D(\mathfrak X_c').
\end{eqnarray}
 Without loss of generality, we may assume that 
$$
f=a_i(t-c)^i+\cdots+a_{i+m}(t-c)^{i+m},~a_i\neq0,~p\nmid i,~i>0.
$$
It is harmless to drop the subscript $c$, set $c:=0$ and $f_j:=a_{i+j}t^{i+j}$ for $0\leq j\leq m$. The poof contains three steps. 

{\itshape Step 1.} This step aims to show that
\begin{eqnarray}\label{twisted de Rham-Higgs}
\Omega^*(\hat E,\hat\theta-df'\wedge)\cong\Omega^*(\hat E,\hat\theta-\sum_{i=0}^m\sum_{j=0}^\infty f'^{p^j-1}_idf_i'\wedge)
~\mathrm{in}~D(\mathfrak X')
\end{eqnarray}
and
\begin{eqnarray}\label{twisted de Rham-Higgs int}
\Omega_\mathrm{int}^*(\hat E,\hat\theta-df'\wedge)\cong\Omega_\mathrm{int}^*(\hat E,\hat\theta-\sum_{i=0}^m\sum_{j=0}^\infty f'^{p^j-1}_idf_i'\wedge)
~\mathrm{in}~D(\mathfrak X').
\end{eqnarray}
Since there is a natural decomposition
$$
\Omega^1_{X'/k}(\log D')=\pi_{Y'}^*\Omega^1_{Y'/k}(\log D'_1)\bigoplus\pi^*_{\mathbb A^1}\Omega^1_{\mathbb A^1/k}(\log D'_2),
$$
coupled with {\itshape Step 1} in the proof of Theorem \ref{formal de Rham-Higgs isolated}, \eqref{twisted de Rham-Higgs} and \eqref{twisted de Rham-Higgs int} follow.

{\itshape Step 2.} Combing Theorem \ref{infty homotopy non-nilpotent} and Proposition \ref{formal Cartier iso}, we have
$$
\tau_{<p-\ell}F_*\Omega^*(\hat H,\hat\nabla-\sum_{i=0}^m\sum_{j=1}^\infty f^{p^j-1}_idf_i\wedge)\cong\tau_{<p-\ell}\Omega^*(\hat E,\hat\theta-\sum_{i=0}^m\sum_{j=0}^\infty f'^{p^j-1}_idf'_i\wedge)
$$
and
$$
\tau_{<p-\ell}F_*\Omega_\mathrm{int}^*(\hat H,\hat\nabla-\sum_{i=0}^m\sum_{j=1}^\infty f^{p^j-1}_idf_i\wedge)\cong\tau_{<p-\ell}\Omega_\mathrm{int}^*(\hat E,\hat\theta-\sum_{i=0}^m\sum_{j=0}^\infty f'^{p^j-1}_idf'_i\wedge)
$$
in $D(X')$.

{\itshape Step 3.} Set $g:=\prod_{i=0}^m\mathrm{AH}(f_i)$. One can check that multiplication by $g$ induces an isomorphism of modules with integrable connection on $(\mathfrak X,\mathfrak D)/k$:
$$
(\hat H,\hat\nabla+df\wedge)\cong(\hat H,\hat\nabla-\sum_{i=0}^m\sum_{j=1}^\infty f^{p^j-1}_idf_i\wedge).
$$

Combing the three steps above, \eqref{c} and \eqref{c int} follow. This completes the proof.
\end{proof}
\begin{remark}
When $D=\emptyset$ and $(E,\theta)=(\mathcal O_{X'},0)$, this theorem is a special case of \cite[Teorem 4.23 (2)]{OV}. In the cited theorem,  we need a technical condition, namely $p>\mathrm{max}\{rs,s+n\}$. Here $r$ is the codimension of $Z$, $s$ is a positive number such that
$$
\Omega^*(\mathcal O_{X'},-df'\wedge)\to\Omega^*(\mathcal O_{X'},-df'\wedge)\otimes\mathcal O_{X'}/\mathcal I'^s
$$
admits a section ($\mathcal I$ is the the defining ideal of $Z$).  By \cite[Corollary 4.32]{OV}, there always exists a positive number $s$ such that the above projection admits a section. However, it may not easy to guarantee that $s$ is \emph{small}  relative to $p$. Our theorem shows that for some special $f$, the technical condition above may be replaced by an easier one, namely $n<p-1$. 
\end{remark}

\subsection{Vanishing cycle complexes in positive characteristic}
Let $X$ be a smooth quasi-projective variety over $\mathbb C$ and let $f:X\to\mathbb A^1$ be a projective morphism of algebraic varieties. Let $(M,F,K^*)$ be a mixed Hodge module on $X$ (see \cite[\S4]{Saito}). 
\begin{theorem}[{\cite[Theorem 2]{Sah1}}]
If the restriction of $f$ to the support of $M$ is projective in a neighborhood of $f^{-1}(0)$, then for any $i\in\mathbb Z$ one has
$$
\mathrm{dim}~\mathbb H^{i-1}(f^{-1}(0)^\mathrm{an},\phi_fK^*)=\mathbb H^{i+\mathrm{dim}(X)}(f^{-1}(0)^\mathrm{an}, \Omega^*(\mathrm{Gr}^FM,\mathrm{Gr}^F\nabla-df\wedge)^\mathrm{an}).
$$
\end{theorem}
The theorem relates vanishing cycle complexes to twisted Higgs complexes. 
Combing it with the de Rham-Higgs comparisons established before this subsection, we make the following
\begin{definition}
Let $k$ be an algebraically closed field of characteristic $p>0$ and let $X$ be a smooth variety over $k$ endowed with a SNCD $D$. For a coherent Higgs module $(E,\theta)$ on $(X,D)/k$, we define the vanishing cycle complex $\phi_f\Omega^*(E,\theta)$ (resp. $\phi_f\Omega_\mathrm{int}^*(E,\theta)$) by
$$
\Omega^*(\hat E,\hat\theta-df\wedge)~(\mathrm{resp.}~\Omega_\mathrm{int}^*(\hat E,\hat\theta-df\wedge)),~(\hat E,\hat\theta):=(E,\theta)_{/(f=0)}.
$$

Similarly, for a coherent module with integrable connection $(H,\nabla)$ on $(X,D)/k$, we define the vanishing cycle complex $\phi_f\Omega^*(H,\nabla)$ (resp. $\phi_f\Omega_\mathrm{int}^*(H,\nabla)$) by
$$
\Omega^*(\hat H,\hat\nabla+df\wedge)~(\mathrm{resp.}~\Omega^*_{\mathrm{int}}(\hat H,\hat\nabla+df\wedge)),~(\hat H,\hat\nabla):=(H,\nabla)_{/(f^p=0)}.
$$
\end{definition}
Using this definition, Theorem \ref{formal de Rham-Higgs isolated}, Corollary \ref{C2} and Theorem \ref{twisted de Rham-Higgs thm} can be reformulated as the following de Rham-Higgs comparison for vanishing cycle complexes.
\begin{theorem}\label{vanishing cycle}
Let $(X,D)/k$ be as in the definition above. Suppose that $(X,D)/k$ has a $W_2(k)$-lifting $(\tilde X,\tilde D)/W_2(k)$. Let $(E,\theta)$ be a coherent Higgs module on $(X',D')/k$.
\begin{itemize}
\item[(i)\ ] Assume that $(E,\theta)$ is nilpotent of level $\leq p-1$ and set $(H,\nabla):=C^{-1}(E.\theta)$. Let $f\in\Gamma(X,\mathcal O_X)$ such that either $\mathrm{CS}(E,\theta-df'\wedge)$ (resp. $\mathrm{CS}_{\mathrm{int}}(E,\theta-df'\wedge)$) is a set of isolated points around $f'^{-1}(0)$ or it as in Corollary \ref{C2}. Then 
$$
\begin{array}{c}
F_*\phi_f\Omega^*(H,\nabla)\cong\phi_{f'}\Omega^*(E,\theta)~\mathrm{in}~D(X'),\\
(\mathrm{resp.}~F_*\phi_f\Omega_\mathrm{int}^*(H,\nabla)\cong\phi_{f'}\Omega_\mathrm{int}^*(E,\theta)~\mathrm{in}~D(X').)
\end{array}
$$
\item[(ii)\ ]Let $(E,\theta)$ be a coherent Higgs module on $(X',D')/k$ of level $\leq\ell-1$. Set $(H,\nabla):=C^{-1}(E,\theta)$. Then
$$
\tau_{<p-\ell}F_*\phi_f\Omega^*(H,\nabla)\cong\tau_{<p-\ell}\phi_{f'}\Omega^*(E,\theta)~\mathrm{in}~D(\mathfrak X')
$$
and
$$
\tau_{<p-\ell}F_*\phi_f\Omega_\mathrm{int}^*(H,\nabla)\cong\tau_{<p-\ell}\phi_{f'}\Omega_\mathrm{int}^*(E,\theta)~\mathrm{in}~D(\mathfrak X').
$$
\end{itemize}
\end{theorem}
Let $(X_i,D_i)/k$, $i=1,2$ be as in the above definition. Set $X:=X_1\times X_2$ and $D:=\pi_1^*D_1+\pi_2^*D_2$, where $\pi_i$ is the natural projection $X\to X_i$. Let $f_i:X_i\to\mathbb A^1$ be a morphism of $k$-varieties for $i=1,2$. We state a Thom-Sebastiani result for vanishing cycle complexes in positive characteristic.
\begin{proposition}
Let $(H_i,\nabla_i)$ be a bundle with integrable connection on $(X_i,D_i)/k$. Then
$$
\begin{array}{rcl}
\phi_{f_1\boxtimes f_2}\Omega^*((H_1,\nabla_1)\boxtimes(H_2,\nabla_2))&\cong&\phi_{f_1}\Omega^*(H_1,\nabla_1)\boxtimes\phi_{f_2}\Omega^*(H_2,\nabla_2),\\
\phi_{f_1\boxtimes f_2}\Omega_\mathrm{int}^*((H_1,\nabla_1)\boxtimes(H_2,\nabla_2))&\cong&\phi_{f_1}\Omega_\mathrm{int}^*(H_1,\nabla_1)\boxtimes\phi_{f_2}\Omega_\mathrm{int}^*(H_2,\nabla_2)
\end{array}
$$
\end{proposition}
The proposition may be regarded as a positive characteristic analogy of \cite{TS}, \cite{Massy} (see also \cite[corollary 2.17, Theorem 4.25]{Max}). Among other things, the $\ell$-adic analogy of loc.cit. had been established by \cite{Fu} and \cite{I17}.

Next we consider the setting: $k$ is an algebraically closed field of characteristic  $p$ and $X$ is a smooth scheme over $k$. Let $(H,\nabla)\in\mathrm{MIC}(X/k)$ such that it has nilpotent $p$-curvature $\psi$ and $H$ is locally free of finite rank over $\mathcal O_X$. Suppose that $f$ has only isolated singularities around $f^{-1}(0)$.
\begin{proposition}\label{amplitude n}
The vanishing cycle complex $\phi_f\Omega^*(H,\nabla)$ has cohomological amplitude in $\{n\}$.
\end{proposition}
\begin{proof}
Without loss of generality, we may assume that $x\in f^{-1}(0)$ is the only critical point of $f$ on $X$. Let $(t_1,\cdots,t_n)$ be a system of coordinates around $x$. By Theorem \ref{Cartier Ogus}, it suffices to show that the $F$-Higgs complex $\Omega^*(H,\psi+F^*df\wedge)$ has cohomological amplitude in $\{n\}$. Write $\psi=\sum_{i=1}^n\psi_i\otimes F^*dt$. Clearly, $\Omega^*(H,\psi+F^*df\wedge)$ is isomorphic to the Koszul complex 
$$K^*(H;\psi_1-(\partial_1f)^p,\cdots,\psi_n-(\partial_nf)^p),~\partial_if:=\frac{\partial f}{\partial t_i}.$$
Thanks to \cite[Corollary 2.2.3]{Ogus04}, to prove this Koszul complex has cohomological amplitude in $\{n\}$, we need to check that $\psi_i+(\partial_i f)^p$ acts injectively on the cokernel of
$$
(\psi_1+(\partial_1 f)^p,\cdots,\psi_{i-1}+(\partial_{i-1} f)^p):H_x^{i-1}\to H_x
$$
for any $i$. Here $H^0:=0$.

Regard $H$ as a coherent sheaf on $Y:=V(F^*\Omega^1_{X/k})=\mathrm{Spec}_X\mathrm{Sym}^\bullet F^*T_{X/k}$. Using the nilpotency of $\psi$, one finds that $H_x$ is a Cohen-Macaulay module over $\mathcal O_{Y,x}$ of dimension $n$ and which is supported on $X$. Combing \cite[Lemma 10.103.3]{SP} and $x$ is an isolated singularity of $f$, one obtains an $H_x$-regular sequence $F^*\partial_1+(\partial_1f)^p,\cdots,F^*\partial_n+(\partial_nf)^p$. Consequently, the desired injectivity follows. This completes the proof.
\iffalse
Without loss of generality again, we may assume that $E=E_0\bigoplus E_1$. For any $i$, consider the diagram
\begin{eqnarray}\label{diagram cokernel}
\xymatrix{
\frac{\hat E}{\sum_{j=1}^{i-1}(\hat\theta_j-\partial f/\partial t_j)\hat E}\ar[r]^-{\hat\theta_i-\partial f/\partial t_i}\ar[d]&\frac{\hat E}{\sum_{j=1}^{i-1}(\hat\theta_j-\partial f/\partial t_j)}\ar[d]\\
\frac{\hat E_2}{\sum_{j=1}^{i-1}\partial f/\partial t_j\hat E_2}\ar[r]^-{-\partial f/\partial t_i}&\frac{\hat E_2}{\sum_{j=1}^{i-1}\partial f/\partial t_j\hat E_2}.}
\end{eqnarray}
By \cite[Theorem 2.1]{Fa} and \cite[Theorem 4.17]{OV}, $E,E_i$ are locally free $\mathcal O_X$-modules of finite rank. Since $\hat{\mathcal O}_{X,x}$ is a regular local ring, coupled with the locally freeness of $E,E_i$, we know that $\hat E,\hat E_i$ are Cohen-Macaulay $\hat{\mathcal O}_{X,x}$-modules. Combing \cite[Lemma 10.103.3]{SP} and $x$ is an isolated critical point of $f$, one gets an $E_i$-regular sequence $\partial f/\partial t_1,\cdots,\partial f/\partial t_n$. It follows that the second row of \eqref{diagram cokernel} is injective. In particular, if 
$$(\hat\theta_i-\partial f/\partial t_i)(\hat e_1,\hat e_2)\in\sum_{j=1}^{i-1}(\hat\theta_j-\partial f/\partial t_j)\hat E,$$
then $\hat e_2\in\sum_{j=1}^{i-1}(\partial f/\partial t_j)E_2$. Clearly, for some $\hat e_1'\in\hat E_1$ we have
$$
(\hat e_1,\hat e_2)\equiv(\hat e_1',0)~\mod~\sum_{j=1}^{i-1}(\hat\theta_j-\partial f/\partial t_j)\hat E.$$
It is easy to see that $\partial f/\partial t_i\hat e_1'\in\sum_{j=1}^{i-1}\partial f/\partial t_jE_1$. Using the injectivity of 
$$
\partial f/\partial t_i:\frac{\hat E_1}{\sum_{j=1}^{i-1}(\partial f/\partial t_j)E_1}\to\frac{\hat E_1}{\sum_{j=1}^{i-1}(\partial f/\partial t_j)E_1},$$
one has $\hat e_1'\in\sum_{j=1}^{i-1}(\partial f/\partial t_j)E_1$. Consequently, one gets
$$
(\hat e_1,\hat e_2)\in\sum_{j=1}^{i-1}(\hat\theta_j-\partial f/\partial t_j)\hat E,
$$
from which the injectivity of \eqref{cokernel} follows. This completes the proof.\fi
\end{proof}
If $(H,\nabla)=(\mathcal O_X,d)$, then it is easy to see that
$$
\mathrm{dim}_k~\mathcal H^n(\phi_f\Omega^*(\mathcal O_X,d))_x=\mu(f,x),~x\in f^{-1}(0).
$$
Here $\mu(f,x):=\mathrm{dim}_k~\mathcal O_{X,x}/J_{f,x}$ is the Milnor number of $f$ at $x$ and $J_f$ is the Jacobian ideal of $f$. The proposition may be regarded as a positive characteristic analogy of \cite[Proposition 4.24 (i)]{Max}.
For the $\ell$-adic analogy of loc.cit., the reader may refer to \cite[\uppercase\expandafter{\romannumeral1} 4.6]{IL03} and \cite[2.10]{SGA}.

Over $\mathbb C$, the monodromy operator of the cohomology of vanishing cycle complexes had been studied in \cite{SS} by using twisted de Rham complexes. Based on \cite[Theorem 2]{SS} as a blueprint, we prove a similar result in positive characteristic. 
To this end, we need the following notion which comes from microlocal analysis, see loc.cit..
\begin{definition}
Let $(H,\nabla)$ be a module with integrable connection on $X/k$. The formal microlocalization of $(H,\nabla)$ with respect to $f$ is defined by 
$$
(H,\nabla)((u))^f:=(H((u)),\nabla-\frac{df}{u}\wedge).
$$
\end{definition}
One can check that $\nabla_u:=u\partial_u+f/u$ is an endomorphism of $(H,\nabla)((u))^f$. In particular, it induces an endomorphism of the complex 
$\Omega^*(H,\nabla)((u))^f$ and then of $\mathcal H^i\Omega^*(H,\nabla)((u))^f$ for any $i$. In the results below, to avoid some technical problem, we take $(H,\nabla):=C^{-1}(E,\theta)$, where $(E=\bigoplus_{i=0}^{p-1}E_i,\theta)$ is a Hodge bundle on $X'/k$, namely $E_i$ is a locally free $\mathcal O_X$-module of finite rank and $\theta(E_i)\subset E_{i-1}$. 
Firstly, we have the following simple facts.
\begin{lemma}\label{equality dim}
$\Omega^*(H,\nabla)((u))^f$ has amplitude in $\{n\}$ and $\mathcal H^n\Omega^*(H,\nabla)((u))^f$ is supported on $\{x\}$. Moreover, we have an equality
\begin{eqnarray}\label{formal equality}
\dim_{k((u))}[\mathcal H^n\Omega^*(H,\nabla)((u))^f]_x=\dim_k[\mathcal H^n\phi_f\Omega^*(H,\nabla)]_x.
\end{eqnarray}
\end{lemma}
\begin{proof}
Set $\mathcal K^*:=\Omega^*(H,\nabla)((u))^f$. Using a similar argument as in the proof of Proposition \ref{amplitude n}, one sees that $\mathcal K^*$ has amplitude $\{n\}$. Obviously, $\mathcal H^n(\mathcal K^*)$ supported on $\{x\}$. It remains to show \eqref{formal equality}.

Let $(t_1,\cdots,t_n)$ be a system of coordinates around $x$ and write $\theta=\theta_i\otimes dt_i'$. Near $x$, we may take
$$
(H,\nabla)=(F^*E,\nabla_\mathrm{can}+t_i^{p-1}F^*\theta_i\otimes dt_i).
$$
Note that $\nabla-df/u\wedge$ is $k((u))$-linear.  According to Theorems \ref{Cartier Ogus}, one has an isomorphisms of $k((u))$-vector spaces
$$ 
\mathcal H^n(\mathcal K^*)\cong(\frac{H((u))}{\sum_{i=1}^n(F^*\theta_i+(\frac{\partial_i f}{u})^p))H((u))})^{\nabla-df/u\wedge}.
$$
Here and next we will use the convention: if a sheaf $\mathcal F$ on $X$ is supported on $\{x\}$, then we identify $\mathcal F_x$ with $\mathcal F$.
By direct computation, there are an isomorphism of $k((u))$-vector spaces
$$
\frac{H((u))}{\sum_{i=1}^n(F^*\theta_i+(\frac{\partial_i f}{u})^p))H((u))}\cong\frac{H((u))}{\sum_{i=1}^n(\frac{\partial_i f}{u})^pH((u))}\cong(F^*\frac{E}{\sum_{i=1}^n\partial_i fE})((u)).
$$
Using Cartier descent theorem \cite[Theorem 5.1]{Katz}, one gets
$$
p^n\dim_{k((u))}\mathcal H^n(\mathcal K^*)=\dim_{k((u))}(F^*\frac{E}{\sum_{i=1}^n\partial_i fE})((u))=p^n\dim_k\frac{E}{\sum_{i=1}^n\partial_i fE}.$$
On the other hand, Theorem \ref{vanishing cycle} shows that
$$
\dim_k\mathcal H^n\phi_f\Omega^*(H,\nabla)=\dim_k\frac{E}{\sum_{i=1}^n(\theta_i+\partial_i f)E}.
$$
Consequently, \eqref{formal equality} follows. This completes the proof.
\end{proof}
 Assume that $x\in X$ is a Brieskorn-Pham isolated singularity of $f$, namely there exists a system of coordinates $(t_1,\cdots,t_n)$ around $x$ such that
$$
f=t_1^{a_1}+\cdots+t_n^{a_n},~p\nmid a_1,\cdots,p\nmid a_n.
$$



\begin{theorem}\label{log monodromy}
For the variable $u$, there is an isomorphism of formal meromorphic connections
\begin{eqnarray}\label{Jordan decomposition}
(\mathcal H^n\Omega^*(H,\nabla)((u))_x^f,\nabla_u)\cong(\mathcal H^n\phi_f\Omega^*(H,\nabla)_x((u)),u\partial_u+M).
\end{eqnarray}
Moreover, $M$ is induced by an endomorphism of $\mathcal H^n\phi_f\Omega^*(H,\nabla)$ and the Jordan decomposition $M=M_s+M_n$ of $M$ can be characterized as follows:
\begin{itemize}
\item The  eigenvalues of  $M_s$  together with multiplicities are given by
$$
\sum_{1\leq i\leq n}-\frac{b_i}{a_i},~1\leq b_i\leq a_i-1;
$$
\item $M_n$ involves the Higgs field $\theta$ and $M_n=0$ for $\theta=0$.
\end{itemize}
\end{theorem}
\begin{proof}
Keep the notations and convention in the lemma above. For simplicity, we consider the case of $n=1$ and assume that $1<a_1<p$. Given a $k$-linear section $s:E(x)\to E_x$ of the natural projection $E_x\to E(x)$. We drop the lower decorations of $t_1,a_1$. 
\iffalse
Set
$$
V_{ij}:=E_i(x),~0\leq i\leq p-1,~0\leq j\leq a-2.
$$
Using the given section $s$, there is  isomorphism of $k$-vector spaces 
\begin{eqnarray}\label{iso V E}
\quad h:\bigoplus_{ij} V_{ij}\cong E/\partial f E,~\bigoplus_{ij}v_{ij}\mapsto\sum_{ij}s(v_{ij})t^j~\mathrm{mod}~\partial f E,~v_{ij}\in V_{ij}.
\end{eqnarray}
\fi
Set
$$
V^\tau_{ij}:=E_i(x),~0\leq i\leq p-1,~0\leq j\leq a-2,~p-1-a\leq\tau\leq p-1
$$
and $V:=\bigoplus_{ij\tau}V_{ij}^\tau$.
According to the proof of the lemma above, one has
\begin{eqnarray}\label{iso to V}
\begin{array}{c}
V((u))\cong\frac{H((u))}{(F^*\theta_1+(\frac{\partial f}{u})^p))H((u))},\\
\bigoplus_{ij\tau}L_{ij}^\tau\mapsto\sum_{ij\tau}F^*(s(L_{ij}^\tau))t^{pj+\tau}~\mathrm{mod}~(F^*\theta_1+(\frac{\partial f}{u})^p))H((u)),\\
L_{ij}^\tau\in V_{ij}^\tau((u)).
\end{array}
\end{eqnarray}
Here $s:V((u))\to E_x((u))$ is induced by $s:V\to E_x$ and
$$
F^*(\sum_i e_iu^i):=\sum_i(F^*e_i)u^i,~\sum_i e_iu^i\in E_x((u)).
$$
Under the isomorphism above, $\partial, t, F^*\theta_1$ induce $k((u))$-linear endomorphisms of $V((u))$ which are denoted by $\partial^V, t^V,\theta^V$, respectively. Obviously, one has
$$
\partial^V(\bigoplus_{ij\tau}x_{ij}^\tau)=\bigoplus_{ij\tau}(\tau+1)x_{ij}^{(\tau+1)},~px^p_{ij}:=0.
$$
Let $t^V(x_{i'j'}^{\tau'})=\bigoplus_{ij\tau}y_{ij}^\tau$. By direct computation, one has
\begin{itemize}
\item if $\tau'<p-1$, then
$$
y_{ij}^\tau=\left\{
\begin{matrix}
x_{i'j'}^{\tau'},&(i,j,\tau)=(i',j',\tau'+1),\\
0,&\mathrm{otherwise};
\end{matrix}
\right.
$$
\item if $\tau'=p-1,j'<a-2$, then
$$
y_{ij}^\tau=\left\{\begin{matrix}
x_{i'j'}^{\tau'},&(i,j,\tau)=(i',j'+1,0),\\
0,&\mathrm{otherwise};
\end{matrix}\right.
$$
\item if $\tau'=p-1,j'<a-2$, then $y_{ij}^\tau=0,~i\leq i'$.
\end{itemize}
Let $\theta^V(x_{i'j'}^{\tau'})=\bigoplus_{ij\tau}y_{ij}^\tau$. Clearly, one has $y_{ij}^\tau=0$ for $i\leq i'$ or $\tau\neq\tau'$.

Put $\nabla^V:=\partial^V+(t^V)^{p-1}\theta^V-a(t^V)^{a-1}/u$. It is a $k((u))$-linear endomorphism of $V((u))$
and $\mathcal H^n(\mathcal K^*)$ is isomorphic to $\mathrm{ker}(\nabla^V)$. It follows that \eqref{iso to V} gives rise to an isomorphism of formal meromorphic connections
$$
(\mathcal H^n(\mathcal K^*)_x,\nabla_u)\cong(\mathrm{ker}(\nabla^V),u\partial_u+(t^V)^a/u).
$$
Meanwhile, by construction one has
$$
(t^V)^a=\frac{t^V\partial^V}{a}+\frac{(t^V)^p\theta^V}{a}~\mathrm{on}~\mathrm{ker}(\nabla^V).
$$
Note that $(t^V)^p\theta^V$ commutes with $u\partial_u, t^V\partial^V$, respectively and $\mathrm{ker}(\nabla^V)$ is a $(t^V)^p\theta^V$-invariant subspace. 
On the other hand, the lemma below shows that the projection 
\begin{eqnarray}\label{projection}
\pi:\mathrm{ker}(\nabla^V)\to W((u)),~W:=\bigoplus_{0\leq i\leq p-1,p-a+1\leq\tau\leq p-1}V_{i(a-2)}^\tau
\end{eqnarray}
is an isomorphism. Combing the discussions above, one has an isomorphism of formal meromorphic connections
$$
(\mathcal H^n(\mathcal K^*)_x,\nabla_u)\cong(W((u)),u\partial_u+M').
$$
Here $M'$ has Jordan decomposition $M'=M_s+M'_n$, where 
\begin{itemize}
\item $M'_s$ acts on $V_{i(a-2)}^\tau((u))$ as $\tau\mathrm{id}$, and
\item $M'_n$ is uniquely determined by the projection
$$
(\mathrm{ker}(\nabla^V),\frac{1}{a}(t^V)^p\theta^V)\cong(W((u)),M'_n).
$$
\end{itemize}
Since $[\partial_u,M']=0$ and the eigenvalues of $M'$ belong to $\mathbb F_p$, one can find a basis $(L_i)$ of $W((u))$ over $k((u))$ such that $\partial_u(L_i)=0$ and the matrix representations of $M'$ with respect to $(L_i)$ is a Jordan matrix. It follows that 
$$
(W((u)),u\partial_u+M')\cong(k((u))^{\oplus\dim_kW},u\partial_u+M),
$$
where the matrix representation of $M$ with respect to the standard basis is a Jordan matrix.
Lemma \ref{equality dim} shows that 
 $$
 \dim_k[\mathcal H^n\phi_f\Omega^*(H,\nabla)]_x=\dim_k W,$$
 from which the theorem follows.
\end{proof}
\begin{lemma}
The projection \eqref{projection} is an isomorphism.
\end{lemma}
\begin{proof}
Put a decreasing order $>_a$ on $\mathbb F_p$ as follows:
$$
p-a>_ap-2a>_a\cdots>_ap-pa.
$$
Using the order, set 
$$\{i_1,\cdots,i_{a-1}\}:=\{p-1,\cdots,p-a+1\}\subset\mathbb F_p,~i_1>_a\cdots>_a i_{a-1}.$$
Put a decreasing order $>$ on $\mathbb F_p\times\mathbb F_p$ as follows: $(\tau_1,i_1)>(\tau_2,j_2)$ if $\tau_1>_a\tau_2$ or $\tau_1=\tau_2$ and $i_1>_1i_2$.  Set $J:=\{p-a,\cdots,p-1\}\subset\mathbb F_p$.

According to Lemma \ref{equality dim}, it suffices to show that \eqref{projection} is an isomorphism. Let
$\nabla^V(\bigoplus_{ij\tau}L^\tau_{ij})=0$. Spelling out the first equality, one has
\begin{itemize}
\item[(1)\ ] $\theta^V(\bigoplus_{ij}L_{ij}^0)=(a/u)\bigoplus_{ij}L^{p-a}_{ij}$,
\item[(2)\ ] $(\tau+a+(t^V)^p\theta^V)(\bigoplus_{ij}L_{ij}^{\tau+a})=(a/u)\bigoplus_{ij}L^\tau_{ij}$ for $\tau\notin J$, and
\item[(3)\ ] $(\tau+a+(t^V)^p\theta^V)(\bigoplus_{ij}L_{ij}^{\tau+a})=(a/u)(t^V)^p(\bigoplus_{ij}L^\tau_{ij})$ for $\tau\in J-\{p-a\}$.
\end{itemize}
There are some consequences:
\begin{itemize}
\item (1) implies that if $\bigoplus_{i\geq i_0,j}L^0_{ij}=0$, then $\bigoplus_{i\geq i_0-1,j}L^{p-a}_{ij}=0$;
\item (2) implies that for $\tau\notin J$, if $\bigoplus_{i\geq i_0,j}L^{\tau+a}_{ij}=0$, then $\bigoplus_{i\geq i_0,j}L^\tau_{ij}=0$;
\item (3) implies that for $\tau\in J-\{p-a\}$, if 
$\bigoplus_{i\geq i_0,j}L^{\tau+a}_{ij}=0$, $\bigoplus_{i\geq i_0+1,j}L^{\tau}_{ij}=0$ and $L^\tau_{i_0(a-2)}=0$, then $\bigoplus_{i\geq i_0,j}L^{\tau}_{ij}=0$.
\end{itemize}
It follows that if $\pi(\bigoplus_{ij\tau}L^\tau_{ij})=0$ for $\bigoplus_{ij\tau}L^\tau_{ij}\in\mathrm{ker}(\nabla^V)$, then $\bigoplus_{ij\tau}L^\tau_{ij}=0$. Equivalently, we have $\bigoplus_jL^\tau_{ij}=0$ holds for any $(\tau,i)$.
This conclusion can be obtained by induction on $(\tau,i)\in\mathbb F_p\times\mathbb F_p$ with respect to the decreasing order $>$. This completes the proof.
\end{proof}
\begin{remark}
It is hard to define the monodromy operator $T$ of $\mathcal H^n\phi_f\Omega^*(H,\nabla)_x$. Nevertheless, \cite[Introduction]{SS} and Theorem \ref{log monodromy} show that $-\frac{\log T}{2\pi\sqrt{-1}}$ may be defined, which is nothing but $M$. Similar to \cite[Theorem 1]{SS} which is a special case of the classical Levelt-Turrittin theorem, Theorem \ref{log monodromy} says that the left hand side of \eqref{Jordan decomposition} admits Jordan decomposition. Note that in positive characteristic, the existence of the Levelt-Turrittin theorem fails in general but the uniqueness is true. The interested reader may check that \cite[\S4]{Lev} is available in positive characteristic.
\end{remark}
Due to some technical problems, Theorem \ref{log monodromy} can only be stated for the Brieskorn-Pham isolated singularities of $f$ and $(E,\theta)$ is a Hodge bundle. However, we think it may hold in a more general situation.
\begin{question}
Let $(H,\nabla)\in\mathrm{MIC}(X/k)$ with $H$ is locally free of finite rank. Let $0\neq f\in\Gamma(X,\mathcal O_X)$. Is there an endomorphism $M$ of $\phi_f\Omega^*(H,\nabla)$ such that we have an isomorphism of formal meromorphic connections
\begin{eqnarray}\label{derived Saito-Sabbah}
(\mathcal H^n\Omega^*(\hat H,\hat \nabla)((u))^f,\nabla_u)\cong(\mathcal H^n\phi_f\Omega^*(H,\nabla)((u)),u\partial_u+M).\end{eqnarray}
Here $(\hat H,\hat\nabla):=(H,\nabla)_{/f^*[0]}$ and $[0]$ is the reduced divisor on $\mathbb A^1$ corresponding to $0\in\mathbb A^1$.
\end{question}
\eqref{derived Saito-Sabbah} may be regarded as a derived version of \cite[Theorem 1]{SS} in positive characteristic.

\iffalse
 Consider the diagram 
$$
\xymatrix{
F_*\Omega^*(H,\nabla+df\wedge)\ar[r]&F_*\Omega^*(\hat H,\hat\nabla+df\wedge)\ar[r]\ar@{-->}[d]^-{\hat\phi}&F_*\Omega^*(H_s,\nabla+df\wedge)\ar[d]^-{\phi}\\
\Omega^*(E,\theta-df'\wedge)\ar[r]&\Omega^*(\hat E,\hat\theta-df')\ar[r]&\Omega^*(E_s,\theta-df'\wedge)\ar@/^1pc/[l]^-{\gamma}.}
$$

\begin{lemma}
$C^{-1}_{\mathcal X/\mathcal S}(E_s,\theta-df'\wedge)\cong (H_s,\nabla+df\wedge)$.
\end{lemma}
\begin{proof}
 Clearly, there are decompositions
$
E_s=\bigoplus_{i=1}^r E_{s,i},H_s=\bigoplus_{i=1}^r H_{s,i},
$
where $E_{s,i},F_*H_{s,i}$ are supported on $Z_i$. By construction, the Higgs field $\theta$ on $E_s$ induces a Higgs field on $E_{s,i}$ and the integrable connection $\nabla$ on $H_s$ induces an integrable connection on $H_{s,i}$. For simplicity, we denote them by the same symbols $\theta,\nabla$, respectively.
Consequently, it suffices to show that for any $i$, one has
\begin{eqnarray}\label{inverse Cartier i}
C^{-1}_{\mathcal X/\mathcal S}(E_{s,i},\theta-df'\wedge)\cong (H_{s,i},\nabla+df\wedge).
\end{eqnarray}
Without loss of generality, we assume $r=1$ and $(f-c)|_Z=0$.

Suppose $f$ {\itshape is of type \uppercase\expandafter{\romannumeral1}}. Choose a family of open subsets of $X$ $\{U_\alpha\}$ satisfying the following conditions:
\begin{itemize}
\item it cover $Z$;
\item there is a system of coordinates $(t_{\alpha1},\cdots,t_{\alpha n})$ on $U_\alpha/S$ such that 
$$
(f-c)|_U=t^{i_1}_1\cdots t_n^{i_n},~D_U=(t^{j_1}_1\cdots t^{j_n}_n=0),$$
where $0\leq i_1,\cdots,i_n<p$ and $j_1,\cdots,j_n\in\{0,1\}$.
\item there is a family of local liftings of $F$ $\{\tilde F_\alpha:\tilde U_\alpha\to\tilde U'_\alpha\}$ such that
$$
\tilde F^*_{ij}(\tilde t'_{\alpha i})=\tilde t^p_{\alpha i}.
$$
\end{itemize}
By construction we know that the left hand side of \eqref{inverse Cartier i} is the gluing of a family
$$
\{(F^*(E_s|_{U_\alpha}),\nabla_{\mathrm{can}}+(\mathrm{id}\otimes\zeta_{\tilde F_\alpha})F^*\theta-(f-c)^{p-1}d(f-c)\wedge)\}
$$
via the family of transition morphisms $\{G_{\alpha\beta}=\exp((\mathrm{id}\otimes h_{\tilde F_\alpha\tilde F_\beta})F^*\theta)\}$. On the other hand,  the right hand side of \eqref{inverse Cartier i} is the gluing of a family
$$
\{(F^*(E_s|_{U_\alpha}),\nabla_{\mathrm{can}}+(\mathrm{id}\otimes\zeta_{\tilde F_\alpha})F^*\theta+(f-c)d(f-c)\wedge)\}
$$
via $\{G_{\alpha\beta}\}$. Now consider the Artin-Hasse exponential of $f-c$, which is given formally by
$$
\mathrm{AH}(f-c)=\exp((f-c)+(f-c)^p/p+(f-c)^{p^2}/p^2+\cdots).
$$
One can check that the automorphism of $F^*E_s$ given by multiplication by $\mathrm{AH}(-f+c)$ induces an isomorphism between the two sides of \eqref{inverse Cartier i}. 

Suppose $f$ {\itshape is of type \uppercase\expandafter{\romannumeral2}}. We may further assume that $Z$ is single point. By definition, there is an open affine subset $U\subset X$ containing $Z$ and a coordinate $t$ around $Z$ such that
$$
(f-c)|_U=gt^i,~g\in\mathcal O_X(U)^*,~g|_Z=1,~t|_Z=0.
$$
If $i=1$, then $f-c$ is coordinate around $Z$. This situation has been discussed above. Let us assume $1<i<p$. Consider the Taylor expansion
$$
g^{1/i}=(1+(1-g))^{1/i}=\sum_{j=0}^\infty\left(\begin{matrix}
1/i\\
j
\end{matrix}
\right)(1-g)^j\in\mathcal O_{X_{/Z}},
$$
where
$$
\left(\begin{matrix}
1/i\\
j
\end{matrix}
\right)=\frac{\frac{1}{i}\cdots(\frac{1}{i}-j+1)}{j!}.
$$
One can check that this number is a $p$-adic integer. Set
$$
m:=\sum_{j=0}^p\left(\begin{matrix}
1/i\\
j
\end{matrix}
\right)(1-g)^jt.$$
Clearly, $m$ is a coordinate around $Z$.
Choose a local lifting $\tilde F$ of $F$ around $Z$ such that $\zeta_{\tilde F}(dm')=m^{p-1}dm$. obviously, one has
$$
\zeta_{\tilde F}(df')\equiv f^{p-1}df\mod \mathcal I^s\Omega^1_{X/k}(\mathrm{log}~D).
$$
 It follows that \eqref{inverse Cartier i} holds. This completes the proof.
\end{proof}
\begin{proposition}
Zariski locally, $\phi$ can be lifted to an isomorphism (in a suitable derived category)
$$
\hat\phi:\Omega^*(\hat H,\hat\nabla+df\wedge)\dashrightarrow\Omega^*(\hat E,\hat\theta-df'\wedge).
$$
\end{proposition}
\begin{proof}
Suppose {\itshape $f$ is of type \uppercase\expandafter{\romannumeral1}}. Without loss of generality, we may assume that 
\begin{itemize}
\item[(i)\ ] $X/S$ admits a system of coordinates $(t_1,\cdots,t_n)$,
\item[(ii)\ ] there is a global lifting of relative Frobenius $\tilde F:\tilde X\to\tilde X'$ such that $\tilde F(\tilde t'_i)=\tilde t_i^p$, and
\item[(iii)\ ] $f=\prod_{i=1}^nt_i^{j_i}$.
\end{itemize}
Under these assumptions, $\phi^{-1}$ becomes
$$\mathrm{AH}(f)\varphi_{\tilde F}:\Omega^*(E_s,\theta-df'\wedge)\to F_*\Omega^*(F^*E_s,\nabla_\mathrm{can}+(\mathrm{id}\otimes\zeta_{\tilde F})F^*\theta+df\wedge).$$
It suffices to lift $\varphi_{\tilde F}$ to a quasi-isomorphism
$$
\hat\varphi_{\tilde F}:\Omega^*(\hat E,\theta-df'\wedge)\to F_*\Omega^*(F^*\hat E,\nabla_\mathrm{can}+(\mathrm{id}\otimes\zeta_{\tilde F})F^*\theta-\sum_{i=1}^\infty f^{p^i-1}df\wedge).
$$

Set $\hat\varphi_{\tilde F}$ to be the composite of an isomorphism of Koszul complexes 
$$
\begin{array}{c}
\hat\varphi_1:\Omega^*(\hat E,\hat\theta-df'\wedge)\to\Omega^*(\hat E,\hat\theta-\sum_{i=0}^\infty f'^{p^i-1}df'\wedge)\\
\sum_I e_I\otimes\omega'_I\mapsto\sum_I\prod_{i\in I}\frac{\hat\theta_i-j_if'-j_if'\sum_{i=1}^\infty f'^{p^i-1}}{\theta_i-j_if'}(e_I)\otimes\omega'_I
\end{array}
$$
and a quasi-isomorphism of complexes
$$
\begin{array}{c}
\hat\varphi_2:\Omega^*(\hat E,\hat\theta-\sum_{i=0}^\infty f'^{p^i-1}df'\wedge)\to F_*\Omega^*(F^*\hat E,\nabla_\mathrm{can}+(\mathrm{id}\otimes\zeta_{\tilde F})F^*\hat\theta-\sum_{i=1}^\infty f^{p^i-1}df\wedge)\\
\sum_Ie_I\otimes\omega_I'\mapsto\sum_IF^*e_I\otimes\omega_I.
\end{array}
$$
Here if $j_i\neq0$, set
$$
\frac{\hat\theta_i-j_if'-j_if'\sum_{i=1}^\infty f'^{p^i-1}}{\hat\theta_i-j_if'}=1+\frac{\sum_{i=1}^\infty f'^{p^i-1}}{1-\frac{\hat\theta_i}{j_if'}}=1+(\sum_{i=0}^\infty f'^{p^i-1})^p\sum_{k=0}^\ell f^{p-1-k}\hat\theta_i^k;
$$
if $j_i=0$, set
$$
\frac{\hat\theta_i-j_if'-j_if'\sum_{i=1}^\infty f'^{p^i-1}}{\theta_i-j_if'}=1.$$

Suppose {\itshape $f$ is of type \uppercase\expandafter{\romannumeral2}}. Clearly, we can lift $\zeta_{\tilde F}$ to a morphism
$$
\hat\zeta_{\tilde F}:\Omega_{X'/k}^1(\mathrm{log}~D')_{/Z'}\to F_*\Omega^1_{X/k}(\mathrm{log}~D)_{/Z}
$$
which sends $d(g'^{1/i}t')$ to $(g^{1/i}t)^{p-1}d(g^{1/i}t)$. In particular, we have $\hat\zeta_{\tilde F}(df')=f^{p-1}df$. On the other hand, one can check that
$$
F_*\mathcal O_{X_{/Z}}\cong\mathcal O_{X'_{/Z'}}\{1,g^{1/i}t,\cdots,(g^{1/i}t)^{p-1}\}.
$$
This completes the proof.
\end{proof}
\fi
\section{De Rham-Higgs comparison for nearby cycles}

\subsection{Kontsevich complexes}
Let $X$ be a smooth variety over $k$ of dimension $n$, $D$ a simple normal crossing divisor on $X$, and let $f:X\to\mathbb P^1_k$ be morphism of $k$-varieties such that $D_\infty:=f^{-1}(\infty)$ is a subdivisor of $D$.
\begin{definition}\label{Kontsevich complex}
The Kontsevich complex associated to $f$, namely $\Omega^*_f$, is a subcomplex of $\Omega_{X_{\log}/k}^*$ which is defined by
$$
\Omega^i_f:=\{\beta\in\Omega_{X_{\log}/k}^i|\beta\wedge df\in\Omega_{X_{\log}/k}^{i+1}\}.
$$
\end{definition}
Clearly, one has $\Omega^0_f=\mathcal O_X(-D_\infty),~\Omega^n_f=\Omega^n_{X_{\log}/k}$.
Let $(E,\nabla)$ be a module with integrable connection (resp. a Higgs module) on $(X,D)/k$.
Set
$$
\Omega^i_f(E,\nabla):=E\otimes \Omega^i_f.
$$
Since $(\Omega^\bullet_f,\wedge)$ is a two-sided ideal of $(\Omega^\bullet_{X_{\log}/k},\wedge)$, one can check that $\nabla$ induces a $k$-linear (resp. an $\mathcal O_X$-linear) morphism
$$
\nabla:E\otimes\Omega^i_f\to E\otimes\Omega^{i+1}_f.
$$
\begin{definition}
Define the Kontsevich complex of $(E,\nabla)$ associated to $f$ by
$$
\Omega_f^*(E,\nabla):=[E\otimes\Omega^0_f\stackrel{\nabla}{\longrightarrow}E\otimes\Omega^1_f\stackrel{\nabla}{\longrightarrow}\cdots\stackrel{\nabla}{\longrightarrow}E\otimes\Omega^n_f].
$$
\end{definition}
\begin{remark}Note that if $E$ is not flat over $\mathcal O_X$, then $E\otimes\Omega^*_f$ may not a subcomplex of $\Omega^*(E,\nabla)$. 
\end{remark}
Assume that $E$ is a locally free $\mathcal O_X$-module of finite type. Let $(E^\vee,\nabla^\vee)$ be the dual of $(E,\nabla)$. Clearly, there is a morphism of complexes
$$
\begin{array}{c}
\Omega_f^*(E,\nabla)\otimes_k\Omega_f^*(E^\vee,\nabla^\vee)(D_\infty-D)\to\Omega^*(\mathcal O_X,d)\\
(\mathrm{resp}.~\Omega_f^*(E,\nabla)\otimes_{\mathcal O_X}\Omega_f^*(E^\vee,\nabla^\vee)(D_\infty-D)\to\Omega^*(\mathcal O_X,0)).
\end{array}
$$
In particular, we have a pairing 
$$
\begin{array}{c}
\mathbb H^{n-i}(X,\Omega_f^*(E,\nabla))\otimes_k\mathbb H^{n+i}(X,\Omega_f^*(E^\vee,\nabla^\vee)(D_\infty-D))\to k.
\end{array}
$$
\begin{proposition}
The pairing above is perfect.
\end{proposition}

\subsection{The Cartier isomorphism of Kontsevich type}
In this subsection, we aim to generalize the Cartier isomorphism of ESP to a version with coefficients. Assume that there is a global Frobenius lifting $\tilde F:\tilde X\to\tilde X'$ of $F:X\to X'$ such that $\tilde F^*\tilde D'=p\tilde D$ and $\tilde F^*(\tilde f')=\tilde f^p$. Here $\tilde f:\tilde X\to\mathbb P^1_{W_2(k)}$ is a $W_2(k)$-morphism which has closed fibre $f$. Set
$$
(H,\nabla):=(F^*E,\nabla_\mathrm{can}+(\mathrm{id}\otimes\zeta_{\tilde F})F^*\theta).
$$
Using the natural identification
\begin{eqnarray}\label{identification f}
F_*\Omega^\bullet_f(H,\nabla)\cong E\otimes F_*\Omega^\bullet_f,
\end{eqnarray}
we construct a morphism of complexes
$$
\varphi_{\tilde F,f}:\Omega_{f'}^*(E,\theta)\to F_*\Omega^*_f(H,\nabla)
$$
as \eqref{F_alpha}.
\begin{proposition}\label{Cartier f}
$\varphi_{\tilde F,f}$ is a quasi-isomorphism.
\end{proposition}
\begin{proof}
It suffices to prove this lemma around $D_\infty$, hence without loss of generality, we may assume that there exists a system of coordinates $(t_1,\cdots,t_n)$ on $(X,D)/k$ such that $f=(t_1\cdots t_r)^{-1}$. Clearly, 
$$\Omega^1_{X_{\log}/k}\cong\mathcal O_X\{d\log f,d\log t_2,\cdots,d\log t_n\}.$$
Set 
$$\theta=\theta_1\otimes d\log f'+\sum_{i=2}^n\theta_i\otimes d\log t'_i.$$
Obviously, one has
$$
\nabla=\nabla_\mathrm{can}+F^*\theta_1\otimes d\log f+\sum_{i=2}^nF^*\theta_i\otimes d\log t_i.
$$
Recall that
$$
F_*\mathcal O_X\cong\mathcal O_{X'}\{t^{i_1}_1\cdots t^{i_n}_n:0\leq i_1,\cdots,i_n<p\}
$$
and the $\mathcal O_X$-module $\Omega_f^i$ is generated by 
\begin{eqnarray}\label{Kontsevich basis 1}
d\log f\wedge(\wedge^{i-1}\{d\log t_2,\cdots,d\log t_n\})
\end{eqnarray}
and
\begin{eqnarray}\label{Kontsevich basis 2}
f^{-1}\wedge^i\{d\log t_2,\cdots,d\log t_n\}.
\end{eqnarray}
Consequently, as an $\mathcal O_{X'}$-module, $F_*\Omega^i_f$ is generated by 
\begin{eqnarray}\label{basis for F_*Omega}
\begin{array}{c}
t_1^{w_1}\cdots t_n^{w_n}d\log f\wedge(\wedge^{w-1}\{d\log t_2,\cdots,d\log t_n\}),\\
t_1^{w_1}\cdots t_n^{w_n}f^{-1}\wedge^w\{d\log t_2,\cdots,d\log t_n\},\\
0\leq w_1,\cdots,w_n<p.
\end{array}
\end{eqnarray}

For any $w=(w_1,\cdots,w_n)\in\{0,\cdots,p-1\}^n$,  define a Higgs field on $E$ by setting
 $$\theta^w:=(\prod_{1\leq i\leq r,w_i=0}t_i')\theta_1\otimes d\log f'+\sum_{i=2}^n\theta_i\otimes d\log t_i'+\sum_{i=1}^nw_id\log t'_i.$$
Construct morphisms of complexes of $\mathcal O_{X'}$-modules
$$
\begin{array}{c}
\varphi^w_{\tilde F,f}:\Omega^*(E,\theta^w)\to F_*\Omega^*(H,\nabla),\\
e\otimes\beta'\mapsto e\otimes t^w\beta,~\beta\in\eqref{Kontsevich basis 1},\\
e\otimes\beta'\mapsto e\otimes(\prod_{1\leq i\leq r,w_i=0}t_i^p)t^wf(f^{-1}\beta),~f^{-1}\beta\in\eqref{Kontsevich basis 2}.
\end{array}
$$
By using \eqref{basis for F_*Omega}, one can check that
$$
\bigoplus_w\varphi^w_{\tilde F,f}:\bigoplus_w\Omega^*(E,\theta^w)\to F_*\Omega^*(H,\nabla)
$$
gives rise to an isomorphism of complexes and $\varphi^0_{\tilde F,f}=\varphi_{\tilde F,f}$. On the other hand, as the complex $\Omega^*(E,\theta^w)$ can be reinterpreted in a Koszul complex, one easily sees that it  is acyclic for $w\neq0$. This completes the proof.
 \end{proof}
\subsection{De Rham-Higgs comparison of Kontsevich type}




\iffalse
\begin{proposition}[ESP]\label{prop ESP}
The Cartier isomorphism 
$$
C^{-1}:\bigoplus_i\Omega^i_{X'/k}(\mathrm{log}~D')\to\bigoplus_i\mathcal H^i(F_*(\Omega^\bullet(\mathrm{log}~D),d))
$$
of $\mathcal O_{X'}$-algebras induces an isomorphism
\begin{eqnarray}\label{Cartier iso ESP}
C^{-1}:\bigoplus_i\Omega^i_{f'}\to\bigoplus_i\mathcal H^i(F_*(\Omega^\bullet_f,d))
\end{eqnarray}
of $\mathcal O_{X'}$-algebras.
\end{proposition}






Let $\mathcal Q^\bullet_f,\mathcal N^\bullet_f$ be the $\mathcal O_{X'}$-submodules of $F_*\Omega^\bullet_f$ generated by $Q^\bullet_f,N^\bullet_f:=B^\bullet_f-Q^\bullet_f$, respectively. Here
 $$
 \begin{array}{c}
 Q^\bullet_f:=\{f^{-p},d\mathrm{log}~f\}\wedge(\wedge^\bullet\{d\mathrm{log} ~t_2,\cdots,d\mathrm{log}~t_n\}),\\ 
B^\bullet_f:=\bigcup_{0\leq\alpha_1,\cdots,\alpha_n\leq p-1} t_1^{\alpha_1}\cdots t_n^{\alpha_n}\{f^{-1},d\mathrm{log}~f\}\wedge(\wedge^\bullet\{d\mathrm{log} ~t_2,\cdots,d\mathrm{log}~t_n\}).\\
\end{array}
 $$
 Clearly, $\mathcal N^\bullet_f,\mathcal Q^\bullet_f$ are closed under the differential $d$ and we have a decomposition of complexes
 $$
(F_*\Omega^\bullet_f,d)=(\mathcal Q^\bullet_f,0)\bigoplus(\mathcal N_f^\bullet,d).
 $$
One can check that the isomorphism \eqref{Cartier iso ESP} is induced by a morphism of complexes of $\mathcal O_{X'}$-modules
$$
\varphi_{\tilde F,0,f}=\wedge^\bullet\zeta_{\tilde F}:(\Omega^\bullet_{f'},0)\to F_*(\Omega^\bullet_f,d).
$$
Denote by $\varphi^{\mathcal Q}_{\tilde F,0,f}$ the projection of $\varphi_{\tilde F,0,f}$ to $(\mathcal Q^\bullet_f,0)$, which is an isomorphism. In particular, we have
\begin{lemma}
$(\mathcal N_f^\bullet,d)$ is acyclic.
\end{lemma}
Let $V^\bullet_{N_f}$ be the $\mathbb F_p$-vector space freely spanned by $N_f$. Obviously, 
$$
(V^\bullet_{N_f},d)\subset\Gamma(X',(\mathcal N_f^\bullet,d))
$$
is a subcomplex and 
$$
(\mathcal N_f^\bullet,d)\cong(V^\bullet_{N_f},d)\otimes_{\mathbb F_p}\mathcal O_{X'}.$$
By basic linear algebra and the lemma above, we get
\begin{lemma}\label{N decomp}
There have decompositions
$$
N_f^\bullet=N^\bullet_{f,i1}\cup N^\bullet_{f,i2},~N^\bullet_{f,i1}\cap N^\bullet_{f,i2}=\emptyset,~i=1,2,3$$
such that 
$$
V^\bullet_{N_f}\stackrel{d}{\longrightarrow}V^\bullet_{N_f}\stackrel{d}{\longrightarrow}V^\bullet_{N_f}
$$
becomes
\begin{eqnarray}\label{decomposable exact}
~~~~V^{N,\bullet}_{f,11}\bigoplus V^{N,\bullet}_{f,12}\xrightarrow{ \left(\begin{matrix}\mu^N_{11}&\mu^N_{12}\\
\mu^N_{21}&\mu^N_{22}\end{matrix}\right)}V^{N,\bullet}_{f,21}\bigoplus V^{N,\bullet}_{f,22}\xrightarrow{ \left(\begin{matrix}\nu^N_{11}&\nu^N_{12}\\
\nu^N_{21}&\nu^N_{22}\end{matrix}\right)}V^{N,\bullet}_{f,31}\bigoplus V^{N,\bullet}_{f,32}	\end{eqnarray}	
with $\mu^N_{22},\nu^N_{11}$	are isomorphisms and $V^{N,\bullet}_{f,ij}$ is generated by $N^\bullet_{f,ij}$. Here
$$
\mu^N_{ij}:V^{N,\bullet}_{f,1j}\to V^{N,\bullet}_{f,2i},~\nu^N_{ij}:V^{N,\bullet}_{f,2j}\to V^{N,\bullet}_{f,3i},~i,j=1,2.
$$	
\end{lemma}
Set 
$$\mathrm{Cone}(\varphi_{\tilde F,0,f})=(\mathcal V_f^\bullet,d_{\tilde F,0,f}),~\mathrm{Cone}(\varphi^\mathcal{Q}_{\tilde F,0,f})=(\mathcal V_f^{\mathcal Q,\bullet},d^{\mathcal Q}_{\tilde F,0,f}).
$$
Now Proposition \ref{prop ESP} can be reformulated as
\begin{proposition}\label{ESP decomp}
 There have decompositions
$$
\mathcal V_f^\bullet=\mathcal V^\bullet_{f,i1}\bigoplus \mathcal V^\bullet_{f,i2},~i=1,2,3$$
such that 
$$
\mathcal V_f^\bullet\xrightarrow{d_{\tilde F,0,f}}\mathcal V_f^\bullet\xrightarrow{d_{\tilde F,0,f}}\mathcal V_f^\bullet
$$
becomes
\begin{eqnarray}\label{decomposable exact}
\mathcal V^\bullet_{f,11}\bigoplus \mathcal V^\bullet_{f,12}\xrightarrow{ \left(\begin{matrix}\mu_{11}&\mu_{12}\\
\mu_{21}&\mu_{22}\end{matrix}\right)}\mathcal V^\bullet_{f,21}\bigoplus \mathcal V^\bullet_{f,22}\xrightarrow{ \left(\begin{matrix}\nu_{11}&\nu_{12}\\
\nu_{21}&\nu_{22}\end{matrix}\right)}\mathcal V^\bullet_{f,31}\bigoplus \mathcal V^\bullet_{f,32}	\end{eqnarray}	
with $\mu_{22},\nu_{11}$	are isomorphisms. Here
$$
\mu_{ij}:\mathcal V^\bullet_{f,1j}\to \mathcal V^\bullet_{f,2i},~\nu_{ij}:\mathcal V^\bullet_{f,2j}\to \mathcal V^\bullet_{f,3i},~i,j=1,2.
$$	
\end{proposition}
\begin{proof}
Set
$$
\begin{array}{c}
\mathcal V^\bullet_{f,11}=(V^{N,\bullet}_{f,11}\otimes_{\mathbb F_p}\mathcal O_{X'})\bigoplus\mathcal Q^\bullet_f,~\mathcal V^\bullet_{f,12}=(V^{N,\bullet}_{f,12}\otimes_{\mathbb F_p}\mathcal O_{X'})\bigoplus\Omega^\bullet_{f'},\\
\mathcal V^\bullet_{f,21}=(V^{N,\bullet}_{f,21}\otimes_{\mathbb F_p}\mathcal O_{X'})\bigoplus\Omega^\bullet_{f'},~\mathcal V^\bullet_{f,22}=(V^{N,\bullet}_{f,22}\otimes_{\mathbb F_p}\mathcal O_{X'})\bigoplus\mathcal Q^\bullet_f,\\
\mathcal V^\bullet_{f,31}=(V^{N,\bullet}_{f,31}\otimes_{\mathbb F_p}\mathcal O_{X'})\bigoplus\mathcal Q^\bullet_f,~\mathcal V^\bullet_{f,32}=(V^{N,\bullet}_{f,32}\otimes_{\mathbb F_p}\mathcal O_{X'})\bigoplus\Omega^\bullet_{f'}.
\end{array}
$$
By combing the exact sequence of complexes
$$
0\to(\mathcal N^\bullet_f,d)\to(\mathcal V_f^\bullet,d_{\tilde F,0,f})\to(\mathcal V_f^{\mathcal Q,\bullet},d^{\mathcal Q}_{\tilde F,0,f})\to0,$$
the decomposition $\mathcal V_f^\bullet=\mathcal N^\bullet_f\bigoplus\mathcal V_f^{\mathcal Q,\bullet}$ and $\varphi^\mathcal{Q}_{\tilde F,0,f}$ is an isomorphism of complexes, it follows that $\mu_{22},\nu_{11}$ are isomorphisms.
\end{proof}
Set $\mathrm{Cone}(\varphi_{\tilde F,0})=(\mathcal V^\bullet,d_{\tilde F,0})$, where
$$
\varphi_{\tilde F,0}=\wedge^\bullet\zeta_{\tilde F}:(\Omega^\bullet_{X'/S}(\mathrm{log}~D'),0)\to F_*(\Omega^\bullet_{X/S}(\mathrm{log}~D),d).
$$
By similar arguments as above, we have 
\begin{corollary}[SZ]\label{log Cartier}
For any lifting $\tilde F:\tilde X\to\tilde X'$ such that $\tilde F^*(\tilde D')=p\tilde D$ (the condition $\tilde F^*(\tilde f')=\tilde f^p$ is not required), there have decompositions
$$
\mathcal V^\bullet=\mathcal V^\bullet_{i1}\bigoplus \mathcal V^\bullet_{i2},~i=1,2,3$$
such that 
$$
\mathcal V^\bullet\xrightarrow{d_{\tilde F,0}}\mathcal V^\bullet\xrightarrow{d_{\tilde F,0}}\mathcal V^\bullet
$$
becomes
\begin{eqnarray}\label{decomposable exact}
\mathcal V^\bullet_{11}\bigoplus \mathcal V^\bullet_{12}\xrightarrow{ \left(\begin{matrix}\mu_{11}&\mu_{12}\\
\mu_{21}&\mu_{22}\end{matrix}\right)}\mathcal V^\bullet_{21}\bigoplus \mathcal V^\bullet_{22}\xrightarrow{ \left(\begin{matrix}\nu_{11}&\nu_{12}\\
\nu_{21}&\nu_{22}\end{matrix}\right)}\mathcal V^\bullet_{31}\bigoplus \mathcal V^\bullet_{32}	\end{eqnarray}	
with $\mu_{22},\nu_{11}$	are isomorphisms. Here
$$
\mu_{ij}:\mathcal V^\bullet_{1j}\to \mathcal V^\bullet_{2i},~\nu_{ij}:\mathcal V^\bullet_{2j}\to \mathcal V^\bullet_{3i},~i,j=1,2.
$$	
\end{corollary}
This corollary is slightly stronger than SZ which is stated only for the standard lifting, namely $\zeta_{\tilde F}(d\mathrm{log}~t'_i)=d\mathrm{log}~t_i$ with $1\leq i\leq n$, but they share the same idea. 
\subsubsection{Cartier isomorphism with non-trivial coefficients}

To generalize Proposition \ref{Cartier iso ESP} to non-trivial coefficients, we need to recall the notion of {\itshape scalarization} of nilpotent Higgs sheaves on $(X',D')/S$, which is introduced in SZ. Let $(E,\theta)$ be a nilpotent Higgs sheaves on $(X',D')/S$ of level $\leq \ell$. Let $T'$ be the dual of $\Omega^1_{X'/S}(\mathrm{log}~D')$. Then the Higgs structure on  $E$ gives rise to a morphism of $\mathcal O_{X'}$-algebras:
$$
\rho_\theta:S^\bullet T'\to\mathcal End_{\mathcal O_{X'}}(E),
$$
so that $E$ is an $S^\bullet T'$-module. Let $S^\bullet T'\oplus E$ be the trivial infinitesimal extension of $S^\bullet T'$ by $E$. Let $\ker(\rho^+_{\theta})$ be the kernel of $\rho_{\theta}|_{S^+ T'}$. It is an ideal of $S^\bullet T'\oplus E$. We put
$$
R(E,\theta):=S^\bullet T'\oplus E/\ker(\rho^+_{\theta})\cong \mathcal O_{X'}\oplus\rho_\theta(S^+ T')\oplus E,
$$
which inherits an $\mathcal O_{X'}$-algebra structure from $S^\bullet T'\oplus E$. More importantly, there is a tautological Higgs field over $R(E,\theta)$: Note that the Higgs field $\theta$ is a global section of
$$
\rho_\theta(S^+ T') \otimes_{\mathcal O_{X'}}\Omega^1_{X'/S}(\mathrm{log}~D') \subset\mathcal End_{\mathcal O_{X'}}(E)\otimes_{\mathcal O_{X'}}\Omega^1_{X'/S}(\mathrm{log}~D').
$$
So we obtain a global section $\theta$ of $R(E,\theta)\otimes_{\mathcal O_{X'}}\Omega^1_{X'/S}(\mathrm{log}~D')$.
Thus, we obtain the following Higgs structure on $R(E,\theta)$:
$$
R(E,\theta)\to R(E,\theta)\otimes_{\mathcal O_{X'}}\Omega^1_{X'/S}(\mathrm{log}~D'),~r\mapsto r\theta,~r\in R(E,\theta),
$$
where $R(E,\theta)\otimes_{\mathcal O_{X'}}\Omega^1_{X'/S}(\mathrm{log}~D')$ endowed with a natural $R(E,\theta)$-module structure.  By abuse of notation, we denote this Higgs field again by $\theta$. 
\begin{definition}
The Higgs bundle $(R(E,\theta),\theta)$ is said to be the scalarization of $(E,\theta)$.
\end{definition}
Clearly, $(R(E,\theta),\theta)$ has the following properties:
\begin{itemize}
	\item[(i)\ ] $R(E,\theta)$ is an $\mathcal O_{X'}$-algebra, $\mathcal O_{X'}\subset R(E,\theta)$ as an $\mathcal O_{X'}$-subalgebra and there is a projection of $\mathcal O_{X'}$-algebras $R(E,\theta)\to \mathcal O_{X'}$,
		\item[(ii)\ ] the projection above induces a morphism of Higgs sheaves $(E,\theta)\to (\mathcal O_{X'},0)$ and $(E,\theta)$ is a direct summand of $(R(E,\theta),\theta)$,
	\item[(iii)\ ] any section lies in the kernel of the projection above is nilpotent.\end{itemize}

Set $(H,\nabla):=C^{-1}_{\mathcal X/\mathcal S}(E,\theta)$. Assume that there is a global lifting $\tilde F:\tilde X\to\tilde X'$ such that $\tilde F^*(\tilde D')=p\tilde D$. Under the identifications
$$
(H,\nabla)\cong(F^*E,\nabla_\mathrm{can}+(\mathrm{id}_{F^*E}\otimes\zeta_{\tilde F})F^*\theta),~F_*\Omega^i(H,\nabla)\cong E\otimes F_*\Omega^i_{X/S}(\mathrm{\log}~D),
$$
we construct a quasi-isomorphism between complexes of $\mathcal O_{X'}$-modules
$$
\begin{array}{c}
\varphi_{\tilde F,\theta}:(\Omega^\bullet(E,\theta),\theta\wedge)\to (F_*\Omega^\bullet(H,\nabla),\nabla),\\
e\mapsto e\otimes 1,
\\e\otimes\beta_1\wedge\cdots\wedge\beta_i\mapsto e\otimes\zeta_{\tilde F}(\beta_1)\wedge\cdots\wedge\zeta_{\tilde F}(\beta_i).
\end{array}
$$
Clearly, $\varphi_{\tilde F}$ induces a morphism between subcomplexes
\begin{eqnarray}\label{Cartier iso Kontesvich}
\varphi_{\tilde F,\theta,f}:(\Omega^\bullet_f(E,\theta),\theta\wedge)\to (F_*\Omega_f^\bullet(H,\nabla),\nabla).
\end{eqnarray}
Moreover, we have
\begin{theorem}\label{Cartier iso f}
The morphism $\varphi_{\tilde F,\theta,f}$
is a quasi-isomorphism.
\end{theorem}
\begin{proof}
 The problem is local, hence we may assume that there is a system of coordinates $(t_1,\cdots,t_n)$ on $(X,D)/S$ such that $f=(t_1\cdots t_r)^{-1}$. Meanwhile, it is harmless to replace $(E,\theta)$ by its scalarization. Set $R:=R(E,\theta)$ and 
$$\mathrm{Cone}(\varphi_{\tilde F,0,f})=(R\otimes\mathcal V_f^\bullet,d_{\tilde F,\theta,f}).$$
By Proposition \ref{ESP decomp}, the complex
\begin{eqnarray}\label{Higgs R}
R\otimes\mathcal V_f^\bullet\xrightarrow{d_{\tilde F,\theta,f}}\mathcal V_f^\bullet\xrightarrow{d_{\tilde F,\theta,f}}R\otimes\mathcal V_f^\bullet
\end{eqnarray}
can be rewritten as
$$
\mathcal V^{R,\bullet}_{f,11}\bigoplus \mathcal V^{R,\bullet}_{f,12}\xrightarrow{ \left(\begin{matrix}\mu^\theta_{11}&\mu^\theta_{12}\\
\mu^\theta_{21}&\mu^\theta_{22}\end{matrix}\right)}\mathcal V^{R,\bullet}_{f,21}\bigoplus \mathcal V^{R,\bullet}_{f,22}\xrightarrow{ \left(\begin{matrix}\nu^\theta_{11}&\nu^\theta_{12}\\
\nu^\theta_{21}&\nu^\theta_{22}\end{matrix}\right)}\mathcal V^{R,\bullet}_{f,31}\bigoplus \mathcal V^{R,\bullet}_{f,32},
$$
where $\mathcal V^{R,\bullet}_{f,ij}=R\otimes\mathcal V^\bullet_{ij}$ and $\mu^\theta_{ij},\nu^\theta_{ij}$ are $R$-linear morphisms. Clearly, the images of 
$\mu^\theta_{ij}-\mu_{ij},~\nu^\theta_{ij}-\nu_{ij}$ are contained in $\rho_\theta(S^+T')\mathcal V^{R,\bullet}_{f,2i},~\rho_\theta(S^+T')\mathcal V^{R,\bullet}_{f,3i}$, respectively. Since $\rho_\theta(S^+T')\subset R$ is a nilpotent ideal, coupled with $\mu_{22},\nu_{11}$ are isomorphisms by Proposition \ref{ESP decomp} again, we conclude that $\mu^\theta_{22},\nu_{11}^\theta$ are isomorphisms. Consequently, the complex \eqref{Higgs R} is exact at the middle place, from which this theorem follows.
 \end{proof}
By a similar argument, we have
\begin{corollary}
With assumption in Corollary \ref{log Cartier}, $\varphi_{\tilde F,\theta}$ is a quasi-isomorphism.
\end{corollary}
\subsection{De Rham-Higgs comparison}
Choose an open covering $\mathcal U=\{U_\alpha\}$ of $X$ together with a family of liftings $\{\tilde F_\alpha:\tilde U_\alpha\to\tilde U'_\alpha\}$ such that $\tilde F_\alpha(\tilde f'_\alpha)=\tilde f_\alpha$ if $D_\infty\cap U_\alpha\neq\emptyset$. Recall that in SZ the authors constructed a morphism of graded $\mathcal O_{X'}$-modules
$$
\varphi:\Omega^*(E,\theta)\to\check{\mathcal C}(\mathcal U,F_*\Omega^*(H,\nabla)),$$
whose $(r,s,\underline\alpha)$-component, namely a morphism in
$$
\mathrm{Hom}_{\mathcal O_{X'}}(\Omega^{r+s}(E,\theta),i_{\underline\alpha*} F_*\Omega^s(H,\nabla)|_{U'_{\underline\alpha}}),
$$
 is given by
\begin{eqnarray}\label{construction varphi}
\varphi(r,s)_{\underline\alpha}=\sum a(\overline j,\underline s)(\iota_{\alpha_0}\otimes\mathrm{id})(\sigma_{\underline\alpha}^{[\overline j]}\otimes\zeta_{\underline\alpha,\underline s}h_{\underline\alpha})F^*(\mathrm{id}\otimes\delta_{r+s}).
\end{eqnarray}
Here 
\begin{itemize}
\item $r,s\geq0$ and $\underline\alpha=(\alpha_0,\cdots,\alpha_r)$,
\item $i_{\underline\alpha}:U'_{\underline\alpha}=U'_{\alpha_0}\cap\cdots\cap U'_{\alpha_r}\hookrightarrow X'$ is the natural inclusion,
\item for $r>0$, the summation takes over $\overline j=(j^1,\cdots,j^r)\in\mathbb N^r,\underline s=(s_0,\cdots,s_r)\in\mathbb N^{r+1}$ subject to $0<j^1+\cdots+j^r+r+s<p$ and $s_0+\cdots+s_r=s$; for $r=0$, the summation becomes
$$\varphi(r,s)_{\alpha_0}(e\otimes\beta_1\wedge\cdots\wedge\beta_{r+s})=\iota_{\alpha_0}(F^*e)\otimes\zeta_{\tilde F_{\alpha_0}}(\beta_1)\wedge\cdots\wedge\zeta_{\tilde F_{\alpha_0}}(\beta_r),$$
\item $a(\overline j,\underline s)\in\mathbb F_p$,
\item $\sigma^{[\overline j]}_{\underline\alpha}=\prod_{k=1}^r\frac{[(\mathrm{id}\otimes h_{\alpha_{k-1}\alpha_k})F^*\theta]^{j^k}}{j^k!}$ for $r>0$ and $\mathrm{id}_{F^*E}$ for $r=0$,
\item $\zeta_{\underline\alpha,\underline s}h_{\underline\alpha}$ is the composite of 
$$\zeta^{\otimes s_0}_{\tilde F_{\alpha_0}}\otimes\cdots\otimes\zeta^{\otimes s_r}_{\tilde F_{\alpha_r}}\otimes h_{\alpha_0\alpha_1}\otimes\cdots\otimes h_{\alpha_{r-1}\alpha_r}$$
with the natural projection
$$
(\Omega^1_{X/S}(\mathrm{log}~D))^{\otimes s}\to\Omega^s_{X/S}(\mathrm{log}~D),
$$ 
\item $F^*:E\otimes(\Omega^1_{X'/S}(\mathrm{log}~D'))^{\otimes(r+s)}\to F_*[F^*E\otimes (F^*\Omega^1_{X'/S}(\mathrm{log}~D'))^{\otimes(r+s)}]$ is the natural morphism,
\item $\delta$ is the natural section of the projection $$(\Omega^1_{X'/S}(\mathrm{log}~D'))^{\otimes(r+s)}\to\Omega^{r+s}_{X'/S}(\mathrm{log}~D').$$
\end{itemize}

Note that by construction, for any $\alpha_0,\alpha_1$, we have
$$
\zeta_{\tilde F_{\alpha_0}}:\Omega^1_{f'}|_{U'_{\alpha_0}}\to F_*\Omega^1_f|_{U_{\alpha_0}},~h_{\alpha_0\alpha_1}:\Omega^1_{f'}|_{U'_{\alpha_0,\alpha_1}}\to F_*\Omega^0_f|_{U_{\alpha_0\alpha_1}}.$$
It follows that the formula \eqref{construction varphi} applies to Kontsevich complexes, namely one has
$$
\varphi:\Omega_f^\bullet(E,\theta)\to\check{\mathcal C}(\mathcal U,F_*\Omega_f^\bullet(H,\nabla))^\bullet.
$$
\begin{lemma}
The morphism above induces a morphism of complexes of $\mathcal O_{X'}$-modules
\begin{eqnarray}\label{Kontsevich varphi}
\varphi:\tau_{<p-\ell}(\Omega_f^\bullet(E,\theta),\theta\wedge)\to\tau_{<p-\ell}(\check{\mathcal C}(\mathcal U,F_*\Omega_f^\bullet(H,\nabla))^\bullet,\check\nabla).
\end{eqnarray}
\end{lemma}
\begin{proof}
We firstly assume that $E$ is a locally free $\mathcal O_{X'}$-module. In this situation we have inclusions of complexes
\begin{eqnarray}\label{inclusion 1}
\tau_{<p-\ell}(\Omega_f^\bullet(E,\theta),\theta\wedge)\subset\tau_{<p-\ell}(\Omega^\bullet(E,\theta),\theta\wedge)\end{eqnarray}
and
\begin{eqnarray}\label{inclusion 2}
\tau_{<p-\ell}(\check{\mathcal C}(\mathcal U,F_*\Omega_f^\bullet(H,\nabla))^\bullet,\check\nabla)\subset\tau_{<p-\ell}(\check{\mathcal C}(\mathcal U,F_*\Omega^\bullet(H,\nabla))^\bullet,\check\nabla).
\end{eqnarray}
On the other hand, SZ shows that
\begin{eqnarray}\label{varphi morphism}
\varphi:\tau_{<p-\ell}(\Omega^\bullet(E,\theta),\theta\wedge)\to\tau_{<p-\ell}(\check{\mathcal C}(\mathcal U,F_*\Omega^\bullet(H,\nabla))^\bullet,\check\nabla)
\end{eqnarray}
is a morphism of complexes of $\mathcal O_{X'}$-modules. Consequently, this lemma follows.

The general case is more involved. We need a more subtle version of \eqref{varphi morphism}, i.e. for any $s<p-\ell$, we have the following commutative diagram 
\begin{eqnarray}\label{commutative s}
\xymatrix{
\Omega^s(E,\theta)\ar[r]^-{\theta\wedge}\ar[d]^-{\varphi}&\Omega^{s+1}(E,\theta)\ar[d]^-{\varphi}\\
\check{\mathcal C}(\mathcal U,F_*\Omega^\bullet(H,\nabla))^s\ar[r]^-{\check\nabla}&\check{\mathcal C}(\mathcal U,F_*\Omega^\bullet(H,\nabla))^{s+1}.
}
\end{eqnarray}
It suffices to show the diagram
\begin{eqnarray}\label{commutative s f}
\xymatrix{
\Omega_f^s(E,\theta)\ar[r]^-{\theta\wedge}\ar[d]^-{\varphi}&\Omega_f^{s+1}(E,\theta)\ar[d]^-{\varphi}\\
\check{\mathcal C}(\mathcal U,F_*\Omega_f^\bullet(H,\nabla))^s\ar[r]^-{\check\nabla}&\check{\mathcal C}(\mathcal U,F_*\Omega_f^\bullet(H,\nabla))^{s+1}
}
\end{eqnarray}
is commutative. Since we do not have inclusions \eqref{inclusion 1} and \eqref{inclusion 2} in general, the commutativity of \eqref{commutative s f} can  not be derived from the commutativity of \eqref{commutative s} directly. Alternatively, we use the fact that \eqref{commutative s f} is functorial with respect to $(E,\theta)$. 

Set $E':=S^\bullet T'/(S^+T')^{\ell+1}$ which is a locally free $\mathcal O_{X'}$-module. Denote by $\theta'$ the natural Higgs field on $E'$. Obviously, for any $e\in\Gamma(U',E)$, there is a natural morphism $(E',\theta')|_{U'}\to(E,\theta)_{U'}$ sending $1$ to $e$. By combing the functoriality of \eqref{commutative s f} with the result obtained in the first paragraph, the commutativity of \eqref{commutative s f} follows. This completes the proof.
\end{proof}
\fi

Let $\mathcal L_f\subset F_*\mathcal L_{\mathcal X/\mathcal S}$ be the subsheaf consisting of $\tilde F\in\Gamma(U,\mathcal L_{\mathcal X/\mathcal S})$ such that either $U\cap D_\infty=\emptyset$ or $\tilde F(\tilde f')=\tilde f^p$. Clearly, $\mathcal L_{f,x'}\neq\emptyset$ for any $x'\in X'$.
\begin{theorem}\label{infty homotopy f}
Let $(E,\theta)$ be a nilpotent Higgs module on $(X,D)/k$ of level $\leq\ell$. Set $(H,\nabla):=C^{-1}_{\mathcal L_f}(E,\theta)$. Then there is an $\mathcal L$-indexed homotopy 
\begin{eqnarray}\label{F_*L indexed homotopy f}
\mathrm{Ho}:\Delta_*(\mathcal L_f)\to\mathcal Hom^*_{\mathcal O_{X'}}(\tau_{<p-\ell}\Omega_{f'}^*(E,\theta),\tau_{<p-\ell}F_*\Omega_f^*(H,\nabla)).
\end{eqnarray}
\end{theorem}
\begin{proof}
Construct \eqref{F_*L indexed homotopy f} as a morphism of sheaves of abelian groups by a variant of \eqref{construction varphi}, namely
\begin{eqnarray}\label{construction varphi f}
\varphi(r,s)_{\underline\alpha}=\sum a(\overline j,\underline s)(\iota_{\alpha_0}\circ\sigma_{\underline\alpha}^{[\overline j]}\circ F^*)\otimes(\zeta_{\underline\alpha,\underline s}h_{\underline\alpha}\circ F^*\circ\delta_{r+s})|_{\Omega_{f'}^{r+s}}.
\end{eqnarray}
Here $\iota_{\alpha_0}\circ\sigma_{\underline\alpha}^{[\overline j]}\circ F^*$ is the composite of 
$$
E\stackrel{F^*}{\longrightarrow}F_*F^*E\stackrel{\sigma_{\underline\alpha}^{[\overline j]}}{\longrightarrow}F_*F^*E\stackrel{\iota_{\alpha_0}}{\longrightarrow}F_*H_{\alpha_0}
$$
and $F_*\Omega^s_f\hookrightarrow F_*\Omega^s_{X'_{\log}/k}$ precomposite with
$$(\zeta_{\underline\alpha,\underline s}h_{\underline\alpha}\circ F^*\circ\delta_{r+s})|_{\Omega_{f'}^{r+s}}:\Omega^{r+s}_{f'}\to F_*\Omega^s_f$$
coincides with the composite of
$$
\Omega^{r+s}_{f'}\hookrightarrow\Omega^{r+s}_{X'_{\log}/k}\stackrel{\delta_{r+s}}{\longrightarrow}(\Omega^1_{X'_{\log}/k})^{\otimes(r+s)}\stackrel{F^*}{\longrightarrow}F_*(F^*\Omega^1_{X'_{\log}/k})^{\otimes(r+s)}\stackrel{\zeta_{\underline\alpha,\underline s}h_{\underline\alpha}}{\longrightarrow}F_*\Omega_{X_{\log}/k}^s.
$$

If $E$ is locally free over $\mathcal O_{X'}$, then
$$
\Omega_{f'}^*(E,\theta)\subset\Omega^*(E,\theta),~\Omega_f^*(H,\nabla)\subset\Omega^*(H,\nabla).
$$
It follows that \eqref{F_*L indexed homotopy f} is a morphism of complexes. For a general $\mathcal O_{X'}$-module $E$, to obtain the same conclusion, we need to show the vanishing of $\eqref{morphism of complexes}(e)$ for any $e\in\Gamma(U',E)$. Define a morphism of Higgs fields
\begin{eqnarray}\label{lift e to 1}
(S^\bullet T_{U'_{\log}/k},\theta)\to(E,\theta)|_{U'},~P\mapsto Pe,
\end{eqnarray}
where the source is the tautological Higgs module. Here we regard $E$ as an $S^\bullet T_{X'_{\log}/k}$-module via $\theta$. Since $S^\bullet T_{U'_{\log}/k}$ is locally free over $\mathcal O_{X'}$, we get the vanishing of $\eqref{morphism of complexes}(1)$ for $1\in S^\bullet T_{X'_{\log}/k}$. It immediately implies the vanishing of $\eqref{morphism of complexes}(e)$, thanks to the fact that \eqref{construction varphi f} is functorial with respect to $(E,\theta)$ and \eqref{lift e to 1}. This completes the proof.
\end{proof}
By combing SZ, Proposition \ref{Cartier f} and the theorem above, we have
\begin{corollary}
Notation and assumption as in Theorem \ref{infty homotopy f}. Then
$$
\tau_{<p-\ell}F_*\Omega_f^*(H,\nabla)\cong\tau_{<p-\ell}\Omega_{f'}^*(E,\theta)~\mathrm{in}~D(X').
$$
\end{corollary}
The corollary can be extended to an intersection version.
\begin{definition}
 Let $(E,\nabla)$ be a module with integrable connection or a Higgs module on $(X,D)/k$. Define $$\Omega^*_{\mathrm{int},f}(E,\nabla)
:=\Omega^*_\mathrm{int}(E,\nabla)\cap\Omega^*_f(E,\nabla).
$$
\end{definition}
By combing SZ and the corollary above, we have
\begin{corollary}\label{nearby de Rham-Higgs}
Notation and assumption as in Theorem \ref{infty homotopy f}. Then
$$
\tau_{<p-\ell} F_*\Omega^*_{\mathrm{int},f}(H,\nabla)\cong\tau_{<p-\ell}\Omega^*_{\mathrm{int},f'}(E,\theta)~\mathrm{in}~D(X').
$$
\end{corollary}
\subsection{$E_1$-degeneration}
 Recall that
\begin{definition}
A Fontaine module on $\mathcal X/\mathcal S$ is a quadruple $(H,\nabla,\mathrm{Fil},\psi)$ satisfying the following conditions:
\begin{itemize}
\item $H$ is coherent $\mathcal O_X$-module, $\mathrm{Fil}$ is a finite decreasing filtration on $H$;
\item $\nabla$ is an integrable connection on $H$ with pole along $D$ satisfying the Griffiths transversality, namely
$$
\nabla(\mathrm{Fil}^iH)\subset\mathrm{Fil}^{i-1}H;
$$
\item $\psi:C^{-1}_{\mathcal X/\mathcal S}(E',\theta')\cong(H,\nabla)$, where $(E,\theta):=\mathrm{Gr}_\mathrm{Fil}(H,\nabla)$.
\end{itemize}
\end{definition}
Clearly, the filtration $\mathrm{Fil}$ induces filtrations on $\Omega_f^*(H,\nabla)$ and $\Omega^*_{\mathrm{int},f}(H,\nabla)$ if $(D-D_\infty)\cap D_\infty=\emptyset$, respectively. For simplicity, the resulting filtrations are defined by the same symbol $\mathrm{Fil}$.
\begin{theorem}\label{E_1 MF}
 Let $(H,\nabla,\mathrm{Fil},\psi)$ be a Fontaine module on $\mathcal X/\mathcal S$. Assume that  $X$ is proper over $k$ and $n+\mathrm{rank}(H)\leq p$.
Then the Hodge to de Rham spectral sequences
$$
E_1^{i,j}=\mathbb H^i(X,\mathrm{Gr}^i_\mathrm{Fil}\Omega^*_f(H,\nabla))\Rightarrow \mathbb H^{i+j}(X,\Omega^*_f(H,\nabla))
$$
and
$$
E_1^{i,j}=\mathbb H^i(X,\mathrm{Gr}^i_\mathrm{Fil}\Omega^*_{\mathrm{int},f}(H,\nabla))\Rightarrow \mathbb H^{i+j}(X,\Omega^*_{\mathrm{int},f}(H,\nabla))
$$
degenerates at $E_1$, respectively.
\end{theorem}
\begin{proof}
The proof is standard, see DI.
\end{proof}
\begin{theorem}\label{vanishing f}
Notation and assumption as in the theorem above. We further assume that $X$ admits an ample line bundle $L$. Set $(E,\theta):=\mathrm{Gr}_\mathrm{Fil}(H,\nabla)$. Then 
$$
\mathbb H^i(X,\Omega_f^*(E,\theta)\otimes L)=0,~
\mathbb H^i(X,\Omega_{\mathrm{int},f}^*(E,\theta)\otimes L)=0
$$
for any $i>n$.
\end{theorem}
\begin{proof}
The proof is standard, see DI.
\end{proof}
\subsection{Nearby cycle complexes in positive characteristic}
Let $X$ be a smooth projective algebraic variety of dimension $n$ over $\mathbb C$, $D\subset X$ a simple normal crossing divisor and $f:X\to\mathbb P^1_\mathbb{C}$ is a morphism of complex varieties such that $D_\infty:=f^*([\infty])\subset D$. Set $U:=X-D$. Kontsevich conjectured in \cite{Kon} and then confirmed by \cite[Theorem 1.3.2]{ESY} that the spectral sequence 
\begin{eqnarray}\label{E_1 char 0}
E_1^{p,q}=H^q(X,\Omega_f^p)\Rightarrow \mathbb H^{p+q}(X,\Omega^*_f)
\end{eqnarray}
degenerates at $E_1$. There are two proofs for the above $E_1$-degeneration given in loc.cit., they are \cite[Appendix D]{ESY} and \cite[Appendix E]{ESY}. In this paper, we generalize the first approach to twisted coefficients and twisted connections. However, the second approach told us more Hodge-theoretic informations about the Kontsevich complexes. We briefly recall them below.
 
 Let $\mathrm{DR}_X$ be the filtered de Rham functor, $g=f^{-1}$, 
$$\Xi_g:\mathrm{MHM}(X-D_\infty)\to\mathrm{MHM}(X)$$
  Beilinson's maximal extension functor (see e.g. \cite{Bei}, \cite[(E.1.5.2)]{ESY}), $j':U\to X-D_\infty$ and $j:U\to X$ the canonical  inclusions, and $\mathbb Q_{h,U}[n]$ the pure Hodge module whose underlying perverse  sheaf is $\mathbb Q_U[n]$. In loc.cit., Saito showed that there is a short exact sequence of mixed Hodge modules 
 \begin{eqnarray}\label{ses MHM}
  0\to\psi_gj_*\mathbb Q_{h,U}[n-1]\to\Xi_g(j'_*\mathbb Q_{h,U}[n])\to j_*\mathbb Q_{h,U}[n]\to0
  \end{eqnarray}
  whose image under $\mathrm{DR}_X$ is a distinguished triangle
$$
(\overline\Omega^*_{X/\mathbb P^1}(\log D),F)[n-1]\to(\Omega^*_f,F)[n]\to(\Omega^*_X(\log D),F)[n]\stackrel{+1}{\longrightarrow}.
$$
Up to a shift of triangles, it becomes a short exact sequence
\begin{eqnarray}\label{ses DR}
0\to(\Omega^*_f,F)\to(\Omega^*_X(\log D),F)\to(\overline\Omega^*_{X/\mathbb P^1}(\log D),F)\to0.
\end{eqnarray}
\emph{We believe that the discuss above can be extended to general mixed Hodge modules on $U$.}
 
 Theorem \ref{E_1 MF} may be regarded as a positive characteristic analogy of \eqref{E_1 char 0} with coefficient. On the other hand, Theorem \ref{vanishing f} may be regarded as a positive characteristic analogy of Kodaira-Saito vanishing \cite[Proposition 2.33]{Saito}.
 Consequently, similar to \cite[\S5]{SZ1}, by spreading-out technique, one obtains some $E_1$-degenerations and vanishings in characteristic $0$. For examples, we have
\begin{theorem}
 Let $\mathbb V$ be a variation of Hodge structure of geometric origin on $X$. Let $(H,\nabla,F)$ be the corresponding de Rham bundle on $(X,D)/\mathbb C$. Then the Hodge to de Rham spectral sequence
$$
E^{p,q}=\mathbb H^{p+q}(X,\mathrm{Gr}^p_F\Omega_f^*(H,\nabla))\Rightarrow\mathbb H^{p+q}(X,\Omega_f^*(H,\nabla))
$$
degenerates at $E_1$.  On the other hand, let $L$ be an ample line bundle on $X$ and set $(E,\theta):=\mathrm{Gr}_F(H,\nabla)$, then
$$
\mathbb H^i(X,\Omega_f^*(E,\theta)\otimes L)=0,~i>n.
$$
\end{theorem}
\begin{remark}
We believe that 
$$\mathrm{DR}_X(\Xi_g(j'_*\mathbb V_U[n])\cong(\Omega_f^*(H,\nabla),F).$$
 It follows that the theorem above can be proved by Saito's theory of mixed Hodge modules \cite{Saito}.
\end{remark}

The discussion above shows that the Kontsevich complexes is a valuable notion in positive characteristic. Thanks to \eqref{ses MHM} and \eqref{ses DR}, the Kontsevich complexes can be used to define the nearby cycle sheaves in characteristic $0$. It is reasonable to extend this observation to positive characteristic.
\begin{definition}
Notations as in the very beginning of \S3.1. Set $g=f^{-1}$. Let $(E,\nabla)$ be a bundle with nilpotent integrable connection or a nilpotent Higgs bundle on $(X,D)/k$. Define the nearby cycle sheaf $\psi_g\Omega^*(E,\nabla)$ by the quotient of complexes
$$
\Omega^*_f(E,\nabla)\to\Omega^*(E,\nabla).
$$
As an immediate consequence of the definition, we have
\end{definition}
\begin{theorem}
Let $(E,\theta)$ be a nilpotent Higgs bundle on $(X',D')/k$ of level $\leq\ell$. Set $(H,\nabla):=C^{-1}_{\mathcal L_{\tilde f}}(E,\theta)$. Then
$$
\tau_{<p-l}F_*\psi_g\Omega^*(H,\nabla)\cong\tau_{<p-l}\psi_g\Omega^*(E,\theta)~\mathrm{in}~D(X').
$$
\end{theorem}
Using the theorem above, we obtain the $E_1$-degeneration and vanishing for nearby cycles shaves in positive characteristic.
\begin{corollary}
Let $(H,\nabla,F)$ be a Fontaine module on $\mathcal X/\mathcal S$. Suppose that $\dim(X)+n<p$.Then the Hodge to de Rham spectral sequence
$$
E_1^{i,j}=\mathbb H^{i+j}(X,\mathrm{Gr}_F^i\psi_g\Omega^*(H,\nabla))\Rightarrow\mathbb H^{i+j}(X,\psi_g\Omega^*(H,\nabla))
$$
degenerates at $E_1$. Let $L$ be an ample line bundle. Set $(E,\theta):=\mathrm{Gr}_F(H,\nabla)$. Then
$$
\mathbb H^i(X,\psi_g\Omega^*(E,\theta)\otimes L)=0,~i>n-1.
$$
\end{corollary}
Finally, we study the monodromy operator of $\psi_g$. Firstly, we recall some basic facts concern the nearby cycle complexes over $\mathbb C$. Let $\Delta$ be a unit disk and let $g:(X,X_0)\to(\Delta,[0])$ be a semistable analytic morphism of complex manifolds. Let $\mathbb V$ be a local system on $X$ and let $(H,\nabla)$ be the corresponding bundle with integrable connection. It is well-known that
$$
\psi_g\mathbb V\cong\Omega^*_{X/\Delta}(H,\nabla)\otimes\mathcal O_{X_0}~\mathrm{in}~D^+(\mathbb C_{X_0}),
$$
where $\Omega^*_{X/\Delta}(H,\nabla):=(H\otimes\Omega^\bullet_{X/\Delta}(\log D/[0]),\nabla)$.
Let $T$ be the monodromy operator of $\psi_g\mathbb V$ which is unipotent by the assumption on $g$. In particular, if $X_0$ is a smooth fiber, then $T=\mathrm{id}$.
It is well-known again that $-\frac{\log T}{2\pi\sqrt{-1}}$ coincides with the extension class of
$$
0\to\Omega_{X/\Delta}^*(H,\nabla)\otimes\mathcal O_{X_0}[-1]\to\Omega^*(H,\nabla)\otimes\mathcal O_{X_0}\to\Omega_{X/\Delta}^*(H,\nabla)\otimes\mathcal O_{X_0}\to0.
$$

Return to the setting of positive characteristic, namely $X/k$ as before, $f:(X,X_0)\to(\mathbb A^1,[0])$ is a semistable morphism and they are $W_2(k)$-liftable. Let $(E,\theta)$ be a nilpotent Higgs bundle of level $\leq p-n$ on $(X',X_0')/k$ and let $(H,\nabla):=C^{-1}(E,\theta)$. In positive characteristic, it is hard to define the monodromy operator $T$ of $\psi_g$. However, we can define its logarithm, namely $-\frac{\log T}{2\pi\sqrt{-1}}$. Note that by definition one has
$$
\psi_f\Omega^*(H,\nabla)=\Omega_{X/\mathbb A^1}^*(H,\nabla)\otimes\mathcal O_{X_0}.
$$
\begin{definition}\label{monodromy nearby cycle}
Define the log monodromy $N_\mathrm{dR}$ of $\psi_f\Omega^*(H,\nabla)$ by the extension class of 
$$
0\to\Omega_{X/\mathbb A^1}^*(H,\nabla)\otimes\mathcal O_{X_0}[-1]\to\Omega^*(H,\nabla)\otimes\mathcal O_{X_0}\to\Omega_{X/\mathbb A^1}^*(H,\nabla)\otimes\mathcal O_{X_0}\to0.
$$
Define the log monodromy $N_\mathrm{Hig}$ of $\psi_f\Omega^*(E,\theta)$ similarly.
\end{definition}
\begin{theorem}
The log monodromy $N_\mathrm{dR},N_\mathrm{Hig}$ are nilpotent. Moreover, if 
$(E,\theta)\in\mathrm{MIC}(X'/k)$ 
and $X_0$ is smooth, then $N_\mathrm{dR}=0,N_\mathrm{Hig}=0$.
\end{theorem}
\begin{proof}
Firstly, we reduce the nilpotence and vanishing of $N_\mathrm{dR}$ to the  nilpotence and vanishing of $N_\mathrm{Hig}$, respectively.
The de Rham-Higgs comparison \cite[Corollary 5.7]{Schepler} gives rise to a morphism of distinguished triangles
\begin{eqnarray}\label{distinguished triangle}
\xymatrix{\Omega_{X/\mathbb A^1}^*(E,\theta)[-1]\ar[r]\ar[d]&\Omega^*(E,\theta)\ar[r]\ar[d]&\Omega_{X/\mathbb A^1}^*(E,\theta)\ar[r]^-{+1}\ar[d]&\\
\Omega_{X/\mathbb A^1}^*(H,\nabla)[-1]\ar[r]&\Omega^*(H,\nabla)\ar[r]&\Omega_{X/\mathbb A^1}^*(H,\nabla)\ar[r]^-{+1}&.
}
\end{eqnarray}
In particular, one gets a commutative diagram
$$
\xymatrix{
\Omega_{X/\mathbb A^1}^*(E,\theta)\ar[d]\ar[r]^-{\delta_\mathrm{Hig}}&\Omega_{X/\mathbb A^1}^*(E,\theta)\ar[d]\\
\Omega_{X/\mathbb A^1}^*(H,\nabla)\ar[r]^-{\delta_\mathrm{dR}}&\Omega_{X/\mathbb A^1}^*(H,\nabla).
}
$$
Here the horizontal morphisms are connecting morphisms and the vertical morphisms are given by the de Rham-Higgs comparison (see loc.cit.). Note that $\delta_\mathrm{Hig}$ is $\mathcal O_X$-linear and $\delta_\mathrm{dR}$ is $k$-linear. One can check that the diagram above induces 
\begin{eqnarray}\label{derived higher direct image}
\xymatrix{
\Omega_{X/\mathbb A^1}^*(E,\theta)\otimes_{\mathcal O_\mathbb A^1,F_{\mathbb A^1}^*}\mathcal O_\mathbb A^1\ar[d]^-{\cong}\ar[rr]^-{\mathrm{id}\otimes t\partial_t+\delta_\mathrm{Hig}\otimes\mathrm{id}}&&\Omega_{X/\mathbb A^1}^*(E,\theta)\otimes_{\mathcal O_\mathbb A^1,F_{\mathbb A^1}^*}\mathcal O_\mathbb A^1\ar[d]^-{\cong}\\
\Omega_{X/\mathbb A^1}^*(H,\nabla)\ar[rr]^-{\delta_\mathrm{dR}}&&\Omega_{X/\mathbb A^1}^*(H,\nabla),
}
\end{eqnarray}
where $t$ is the standard coordinate on $\mathbb A^1$. Taking the closed fiber of \eqref{derived higher direct image} at $0\in\mathbb A^1$, one gets
$$
\xymatrix{
\psi_f\Omega^*(E,\theta)\otimes_{k,\sigma}k\ar[d]^-{\cong}\ar[r]^-{N_\mathrm{Hig}\otimes\mathrm{id}}&\psi_f\Omega^*(E,\theta)\otimes_{k,\sigma}k\ar[d]^-{\cong}\\
\psi_f\Omega^*(H,\nabla)\ar[r]^-{N_\mathrm{dR}}&\psi_f\Omega^*(H,\nabla).
}
$$
Here $\sigma:k\to k$ is the Frobenius endomorphism of $k$ and the vertical isomorphism is ensured by Corollary \ref{nearby de Rham-Higgs}. From which the reduction follows.

Next we show the nilpotence of $N_\mathrm{Hig}$. It suffices to show the nilpotence of $\delta_\mathrm{Hig}$.
Assume that there is a system of coordinates $(t_1,\cdots,t_n)$ on $X/k$ such that $f=t_1\cdots t_n$. Write 
$$
\theta=\theta_f\otimes d\log f+\theta_2\otimes d\log t_2+\cdots+\theta_n\otimes d\log t_n.
$$
Clearly, under this assumption, the first row of \eqref{distinguished triangle} splits as graded $\mathcal O_X$-modules. More precisely, one has
\begin{eqnarray}\label{splitting}
\begin{array}{c}
\Omega^i(E,\theta)\cong\Omega^{i-1}_{X/\mathbb A^1}(E,\theta)\bigoplus\Omega^i_{X/\mathbb A^1}(E,\theta),~\forall i,\\
d\log f\wedge\omega_1+\omega_2\mapsto(\omega_1,\omega_2),\\
\omega_1\in\wedge^{i-1}\{d\log t_2,\cdots,d\log t_n\},~\omega_1\in\wedge^i\{d\log t_2,\cdots,d\log t_n\}.
\end{array}
\end{eqnarray}
It follows that $\delta_\mathrm{Hig}=\theta_f$. In general, this assumption may fail around $X_0$, but it always true around any point of $X$. Choosing an open covering    $\mathcal U:=\{U_\alpha\}$ of $X$ such that the above assumption is true on each $U_\alpha$. Combing \eqref{splitting} and $\mathcal U$, $\delta_\mathrm{Hig}$ can be represented by a morphism of complexes
$$
\rho:\Omega^*_{X/\mathbb A^1}(E,\theta)\to\check{\mathcal C}(\mathcal U,\Omega^*_{X/\mathbb A^1}(E,\theta)),
$$
where its component $\Omega^*_{X/\mathbb A^1}(E,\theta)\to\mathcal C^0(\mathcal U,\Omega^*_{X/\mathbb A^1}(E,\theta))$ is given by a family of endomorphisms of complexes
$$
\{
\theta_{f,\alpha}:\Omega^*_{X/\mathbb A^1}(E,\theta)|_{U_\alpha}\to\Omega^*_{X/\mathbb A^1}(E,\theta)|_{U_\alpha}
\}.
$$
Clearly, 
\begin{eqnarray}\label{nilpotent derived}
\prod_{i=1}^{p-n}\theta_{f,\alpha_i}=0,~\forall \alpha_1,\cdots,\alpha_{p-n}.
\end{eqnarray}

We claim that $\delta_\mathrm{Hig}^p=0$. Let us prove this claim. Note that $\delta^p_\mathrm{Hig}$ is represented by a morphism 
$$
\rho^p:\Omega^*_{X/\mathbb A^1}(E,\theta)\to\mathcal K^*:=\underbrace{\check{\mathcal C}(\mathcal U,\cdots,\check{\mathcal C}(\mathcal U,}_{p~\mathrm{times}}\Omega^*_{X/\mathbb A^1}(E,\theta))\cdots),$$
where 
$$\rho^p:=\underbrace{\check{\mathcal C}(\mathcal U,\cdots,\check{\mathcal C}(\mathcal U,}_{(p-1)~\mathrm{times}}\rho)\cdots)\circ\cdots\circ\rho.$$
By definition, one has
$$
\mathcal K^m=\bigoplus_{r_1+\cdots+r_p+s=m}\underbrace{\check{\mathcal C}^{r_1}(\mathcal U,\cdots,\check{\mathcal C}^{r_p}(\mathcal U,}_{p~\mathrm{times}}\Omega^s_{X/\mathbb A^1}(E,\theta))\cdots).$$
Observe that if $r_1+\cdots+r_p+s\leq n$ and $r_1,\cdots,r_p,s\geq0$, then at least $p-n$ of $r_1,\cdots,r_p$ are $0$. Combing this observation with \eqref{nilpotent derived}, one concludes that $\rho^p=0$. 

We finally show that if $(E,\theta)\in\mathrm{MIC}(X'/k)$ and $X_0$ is smooth, then $N_\mathrm{Hig}=0$. It suffices to construct a splitting for the exact sequence in Definition \ref{monodromy nearby cycle}. Let $s_0$ be any section of the graded $\mathcal O_X$-linear morphism 
$$E\otimes\Omega^\bullet_{X/k}\to\Omega_{X/\mathbb A^1}^\bullet(E,\theta)$$
and let $s:=(i\circ s_0)\otimes\mathcal O_{X_0}$. Here $i:E\otimes\Omega^\bullet_{X/k}\hookrightarrow\Omega^\bullet(E,\theta)$ is the natural inclusion. One can check that $s$ is independent of the choice of $s_0$ and it is a morphism of complexes. It turns out that the aforementioned exact sequence splits, or equivalently $N_\mathrm{Hig}=0$. This completes the proof.
\end{proof}
We believe that the arguments in \cite[\S3.1-3.2]{OV} may imply the theorem above, but  we prefer to give its proof in our favorite way. \eqref{derived higher direct image} may be regarded as a special case of \cite[Theorem 3.8]{OV}.





\fi

\begin{thebibliography}{X-X00}
 	\addcontentsline{toc}{chapter}{Bibliography}
 
 \bibitem[AS]{AS}
P. Achinger, J. Suh, Some refinements of the Deligne-Illusie theorem, Algebra Number Theory17(2023), no.2, 465–496.

\bibitem[AWZ]{AWZ}
P. Achinger, J. Witaszek, M. Zdanowicz, 
Global Frobenius liftability I,
J. Eur. Math. Soc. (JEMS) 23 (2021), no. 8, 2601–2648.

\bibitem[Ar]{Arapura}
 	D. Arapura, Kodaira-Saito vanishing theorem via Higgs bundles in positive characteristic, Journal f\"{u}r angewandte und reine Mathematik 755 (2019), 293-312.
	
\bibitem[ACH]{ACH}
D. Arinkin, A. C\u{a}ld\u{a}raru, M. Hablicsek, Derived intersections and the Hodge theorem, Algebr. Geom. 4 (2017), no. 4, 394–423.
	
 	\iffalse
 	\bibitem[AK]{Alt}
 	A. Altman, S. Kleiman, Introduction to Grothendieck duality theory. Lecture Notes in Mathematics 146, Springer-Verlag, Berlin-Heidelberg-New York, 1970.\fi
 	\iffalse
 	\bibitem[AM]{Atiyah}
 	M. Atiyah, I. Macdonald, Introduction to commutative algebra, Addison-Weiley Series in Mathematics, Addison-Weiley Publishing Cimpany, 1969.
 	\fi
\bibitem[Bea]{Bea}
A. Beauville, Complex manifolds with split tangent bundle, In Complex analysis
 and algebraic geometry, pages 61–70, de Gruyter, Berlin, 2000.
    
\bibitem[BL]{BL}
B. Bhatt J. Lurie, The prismatization of $p$-adic formal schemes, 2022, arXiv:2201.06124.

\bibitem[Bei]{Bei}
	A. Beilinson, How to glue perverse sheaves, in K-theory, arithmetic and geometry (Moscow, 1984-1986), Lect. Notes in Math., vol. 1289, Springer, 1987, p. 42–51.	
	
	\bibitem[Ber]{Ber}
 	 P. Berthelot, Cohomologie cristalline des sch\'emas de caract\'eristique $p>0$. (French) Lecture Notes in Mathematics, Vol. 407. Springer-Verlag, Berlin-New York, 1974. 604 pp.

     \bibitem[BPT]{BPT}
M. Brunella, J. Pereira, F. Touzet. K\"ahler manifolds with split
 tangent bundle. preprint, to appear in Bull. Soc. Math. France, 2004.
 
     \bibitem[CP]{CP}
 	F. Campana and T. Peternell. Projective manifolds with splitting tangent
 bundle, I,  Math. Z., 241(3):613–637, 2002.
 	\bibitem[CEGT]{CEGT}
 	E. Cattani, F. El Zein, P. Griffiths, L.-D. Tr\'{a}ng, Hodge theory, Mathematical Notes 49, Princeton University Press, 2014.
	
 	\bibitem[CKS]{CKS}
 	E. Cattani, A. Kaplan, W. Schmid, $L^2$ and intersection cohomologies for a polarizable variation of Hodge structure. Invent. Math. 87, 217-252, 1987,
	
	
	
	
	\bibitem[Del]{Del}
	P. Deligne, Equations diff\'erentielles \`a points singuliers r\'eguliers, Lect. Notes in Math. 163, Springer, Berlin, 1970.  	

 	\bibitem[DI]{DI}
 	P. Deligne, L. Illusie, Rel\`{e}vements modulo $p^2$ et decomposition du complexe de de Rham, Invent. Math. 89 (1987), 247-270.

\bibitem[Dru]{Dru}
     S. Druel, Variétés algébriques dont le fibré tangent est totalement décomposé,
 J. Reine Angew. Math., 522:161–171, 2000.
    
 	\iffalse
 	\bibitem[Ei]{Eisenbud}
 	D. Eisenbud, Commutative algebra with a wiew towards algebraic geometry, GTM 150, New York, Springer-Verlag, 1999.
 	\fi
	 \bibitem[ESY]{ESY} 
	 H. Esnault, Y. Sabbah, J. Yu, $E_1$-degeneration of the irregular Hodge filtration. With an appendix by Morihiko Saito, J. Reine Angew. Math. 729 (2017), 171-227. 	
	
 	\bibitem[EV]{EV}
 	H. Esnault, E. Viehweg, Lectures on vanishing theorems. DMV Seminar, 20. Birkh\"{a}user Verlag, Basel, 1992.
 	
 	\bibitem[Fa]{Fa}
 	G. Faltings, Crystalline cohomology and $p$-adic Galois-representations, Algebraic analysis, geometry, and number theory (Baltimore, MD, 1988), 25-80, Johns Hopkins Univ. Press,
 	Baltimore, MD, 1989.
 	\iffalse
 	\bibitem[Fu]{Fu}
 	L.  Fu, \'Etale cohomology theory, Nankai Tracts in Mathematics, Vol.13, World Scientific, 2011.
 	\fi
 	\bibitem[FL]{FL}
 	J.-M. Fontaine, G. Laffaille, Construction de repr\'{e}sentation
 	$p$-adiques, Ann. Sci. Ec. Norm. Sup. 15 (1982), 547-608.
    \iffalse
	\bibitem[Fu]{Fu}
	 L. Fu, A Thom-Sebastiani theorem in characteristic $p$, Math. Res. Lett. 21 (2014), no. 1, 101-119.\fi	
     
     \bibitem[J]{J}
K. Joshi, On varieties with trivial tangent bundle in characteristic $p>0$, Nagoya Math. J. 242 (2021), 35–51.
    
 	\bibitem[H]{H} 
 	A. Höring, Uniruled varieties with split tangent bundle, Math. Z. 256 (2007), no. 3, 465–479.
 	\iffalse
 	\bibitem[HL]{HL}
 	D. Huybrechts, M. Lehn,  The geometry of moduli spaces of sheaves. Second Edition, Cambridge University Press.
 	\fi
 	\bibitem[I90]{IL90}
 	L. Illusie, R\'eduction semistable et d\'ecomposition de complexes de de Rham
 	\`a\ coefficients, Duke Math. J., vol. 60, 1990, no. 1, 139-185.
 	
 	
 	\bibitem[I03]{IL03}
 	L. Illusie, Perversit\'e et variation, Manuscripta math. 112 (2003), 271-295.
	\bibitem[I17]{I17}
	L. Illusie, Around the Thom-Sebastiani theorem, with an appendix by Weizhe Zheng, 
Manuscripta Math.152(2017), no.1-2, 61-125.	
 	\iffalse
 	\bibitem[JYZ]{JYZ}
 	J. Jost, Y.-H. Yang and K. Zuo, The cohomology of a variation of
 	Hodge structure over a quasi-compact K\"{a}hler manifold, Journal of
 	Algebraic Geometry 16 (2007), No. 3, 401-434.
 	\fi
 	\bibitem[KK86]{KK86}
 	M. Kashiwara, T. Kawai, Hodge structure and holonomic systems, Proc. Japan Acad. Ser. A Math. Sci., Volume 62, Number 1 (1986), 1-4.
 	
 	\bibitem[KK]{KK}
 	M. Kashiwara, T. Kawai, The Poincar\'e lemma for a variation of Hodge structure. Publ. RIMS, Kyoto University 23, 345-407, 1987.
 	
 	\bibitem[Ka]{Kato}
 	K. Kato, Logarithmic structures of Fontaine-Illusie, Algebraic
 	analysis, geometry, and number theory (Baltimore, MD, 1988),
 	191-224, Johns Hopkins Univ. Press, Baltimore, MD, 1989.
	\bibitem[Kat]{Katz}
	N. Katz, Nilpotent connections and the monodromy theorem: Applications of a result of Turrittin, Inst. Hautes Études Sci. Publ. Math.(1970), no.39, 175–232.	
 \bibitem[KKP]{KKP}L. Katzarkov, M. Kontsevich, T. Pantev, Bogomolov–Tian–Todorov theorems for Landau
Ginzburg models, J. Differential Geom. 105 (2017), 55–117.

    
    \iffalse
	\bibitem[Kon1]{Kon1}
	Kontsevich, Letters to L. Katzarkov$\And$T. Pantev, March 14$\And$31, 2012. 
	
	\bibitem[Kon2]{Kon2}
	Kontsevich, Letters to H. Esnault $\And$ J.-D. Yu, May 3 $\And$ 7, 2012.	\fi
 	\bibitem[LSYZ]{LSYZ}
 	G.-T. Lan, M. Sheng, Y.-H. Yang, K. Zuo, Uniformization of $p$-adic curves via Higgs-de Rham flows. Journal f\"{u}r angewandte und reine Mathematik, Vol.747, pp 63-108, 2019.
 	
 	\bibitem[LSZ14]{LSZ}
 	G.-T. Lan, M. Sheng, K. Zuo, Nonabelian Hodge theory in positive characteristic via exponential twisting, Mathematical Research Letter, Vol. 22, Pages 859-879, 2015.
 	
 	\bibitem[LSZ]{LSZ1}
 	G.-T. Lan, M. Sheng, K. Zuo, Semistable Higgs bundles,
 	periodic Higgs bundles and representations of algebraic fundamental
 	groups, Journal of European Mathematical Society, Vol.21, pp 3053-3112, 2019.

    \iffalse
 	\bibitem[La]{La}
 	 A. Langer,  Nearby cycles and semipositivity in positive characteristic, J. Eur. Math. Soc. (JEMS) 24 (2022), no. 11, 3829–3872.
	\bibitem[Lev]{Lev}
	 A. Levelt, Jordan decomposition for a class of singular differential operators.Ark. Mat.13(1975), 1–27.\fi
     \bibitem[Li]{Li}
K. Li, Actions of group schemes. Compositio Math., 80:55–74, 1991.
     
	
\bibitem[LM]{LM}
S. Li, S. Mondal, On endomorphisms of the de Rham cohomology functor.(English summary)Geom. Topol.28(2024), no.2, 759–802.
\bibitem[LZ]{LZ}
 	S. Li, D. Zhang, Exponentially twisted de Rham cohomology and rigid cohomology, arXiv:2111.05689.

\bibitem[Mum]{Mumford}
D. Mumford, Abelian varieties, Tata Institute of Fundamental Research
Studies in Mathematics, No. 5, Oxford University Press, London 1970.

       \bibitem[MS]{MS}
    V. Mehta, V. Srinivas, Varieties in positive characteristic with trivial tangent bundle,
 Comp. Math., 64:191–212, 1987.

    \iffalse
	\bibitem[Mas]{Massy}
	D. B. Massey, The Sebastiani-Thom isomorphism in the derived category, Compos. Math. 125 (2001),
353–362.	 
	\bibitem[Max]{Max}
	 L. Maxim, Notes on vanishing cycles and applications,
J. Aust. Math. Soc. 109 (2020), no. 3, 371-415.	
\bibitem[Mil]{Milnor}
J. W. Milnor, Singular points of complex hypersurfaces, Annals of Mathematics Studies, No. 61, 1968.
	\bibitem[Mo]{Mo}	
 	T. Mochizuki, A twistor approach to the Kontsevich complexes. Manuscripta Math. 157 (2018), no. 1-2, 193–231.
 	\fi
\bibitem[Oda]{Oda}

     T. Oda, The first de Rham cohomology group and Dieudonné modules, Ann. scient.
 Éc. Norm. Sup. (4), 2, (1969), 63–135.
 	\bibitem[O94]{Ogus}
 	A. Ogus, F-crystals, Griffiths transversality, and the Hodge decomposition, Ast\'{e}risque 221, 1994.
 	
 	\bibitem[O04]{Ogus04}
 	A. Ogus, Higgs cohomology, $p$-curvature, and the Cartier isomorphism. Compositio Math. 140 (2004), 145-164.
	
	\bibitem[O18]{O18}
	 A. Ogus, Lectures on logarithmic algebraic geometry, Cambridge Studies in Advanced Mathematics, 178. Cambridge University Press, Cambridge, 2018. xviii+539 pp.  	
       \bibitem[O24]{O24}
       A. Ogus, Crystalline prisms: Reflections and diffractions, present and past, arXiv:2204.06621.

     
 	\bibitem[OV]{OV}
 	A. Ogus, V. Vologodsky, Nonabelian Hodge theory in characteristic $p$,
 	Publ. Math. Inst. Hautes \'{e}tudes Sci. 106 (2007), 1-138.
    \bibitem[P1]{P1}
A. Petrov. Non-decomposability of the de Rham complex and non-semisimplicity of
the Sen operator, arXiv:2302.11389.
\bibitem[P2]{P2}
A. Petrov, Decomposition of de Rham complex for quasi-F-split varieties, arXiv:2502.13356.
 	\iffalse
 	\bibitem[Oh]{Oh}
 	S. Ohkawa, On logarithmic nonabelian Hodge theory of higher level in characteristic $p$, Rendiconti del Seminario Matematico della Universit di Padova, Volume 134 (2015), 47-92.
 	\fi
 		 
	\bibitem[Sah1]{Sah1} 
	 C. Sabbah, On a twisted de Rham complex, Tohoku Math. J. (2) 51 (1999), no. 1, 125–140.
	
	\bibitem[Sah2]{Sah2}
	C. Sabbah, Duality for Landau-Ginzburg models,  arXiv:2212.07745. 
	\bibitem[SS]{SS}
	C. Sabbah,  M. Saito,
Kontsevich's conjecture on an algebraic formula for vanishing cycles of local systems,
Algebr. Geom.1(2014), no.1, 107–130.
	
 	
	\bibitem[Sai]{Saito} M. Saito, Mixed Hodge modules, Publ. Res. Inst. Math. Sci. 26 (1990), no. 2, 221-333.
	\bibitem[TS]{TS}
	M. Sebastiani, R. Thom, Unr\'esultat sur la monodromie, Invent. Math. 13 (1971), 90-96.	
 	
 	\bibitem[Sche]{Schepler}
 	D. Schepler,  Logarithmic nonabelian Hodge theory in characteristic $p$. arXiv:0802.1977.
	\bibitem[SGA7]{SGA}
	SGA 7, Groupes de monodromie en g\'eom\'etrie alg\'ebrique, S\'eminaire de G\'eom\'etrie Alg\'ebrique du Bois-Marie 1967-1969, I dirig\'e par A. Grothendieck, II par P. Deligne et N. Katz, Lecture Notes in Math. 288, 340, Springer-Verlag 1972, 1973.
       \bibitem[SZ1]{SZ1}
        M. Sheng, Z. Zhang, An explicit infinite homotopy in nonabelian Hodge theory in positive characteristic, arXiv:2302.09231.
       
       \bibitem[SZ2]{SZ2}
       M. Sheng, Z. Zhang, Intersection de Rham complexes in positive characteristic, J. Differential Geom. 127(2): 551-602 (June 2024). DOI: 10.4310/jdg/1717772421.

       \bibitem[Shi]{Shi}
       A. Shiho, Notes on generalizations of local Ogus-Vologodsky correspondence. J. Math. Sci. Univ. Tokyo, 22:798–875, 2015.
  
 	\bibitem[Sim]{Sim}
 	C. Simpson, Higgs bundles and local systems, Inst. Hautes \'{E}tudes Sci. Publ.
 	Math. No. 75 (1992), 5-95.
	\bibitem[SP]{SP}
	The Stacks Project Authors Stacks Project. https://stacks.math.columbia.edu.
	
	\bibitem[Tsu]{Tsu}
	T. Tsuji, Poincaré duality for logarithmic crystalline cohomology, Compositio Math. 118 (1999), no. 1, 11–41. 

       \bibitem[W]{W}
Y. Wang, Prismatic crystals for smooth schemes in characteristic $p$ with Frobenius lifting mod $p^2$, arXiv:2402.02109.
        
       \bibitem[X]{X}
 	D. Xu, Lifting the Cartier transform of Ogus-Vologodsky modulo $p^n$, Mém. Soc. Math. Fr. (N.S.) No. 163 (2019), 133 pp.
\bibitem[Y]{Y}
F. Yobuko, Quasi-F-split and Hodge-Witt, arXiv:2312.00682.

\bibitem[Z]{Z}
Z. Zhang, A note on the filtered decomposition theorem, Commun. Math. Stat.11(2023), no.3, 519–539.

  
	\iffalse
 	\bibitem[Zuc]{Zuc}
 	S. Zucker, Hodge theory with degenerating coefficients: $L_2$ cohomology in the Poincar\'{e} metric, Ann. of Math. 109, 1979, 415-476.
 	
 	\bibitem[EGA IV]{EGA IV}
 	A. Grothendieck and J. Dieudonn\'{e}, \'{E}tude locale des sch\'{e}mas et des morphisms de sch\'{e}mas, Publ. Math. IHES, 20(1964), 24(1965), 28(1966), 32(1967).
 	
 	\bibitem[SGA 1]{SGA 1}
 	A. Grothendieck et al., Rev\^{e}tements \'{e}tales et groupe fondamental, Lecture notes in Mathematics 224, Springer-Verlag 1971.
	\fi
 \end{thebibliography}
	















\end{document}
